\section{Tipos de gráficos avanzados}

\subsection{Gráfico de torta (\textit{.pie()})}

Este gráfico representa datos categóricos como "rebanadas" de un círculo proporcional a su magnitud. Es efectivo cuando se tienen pocas categorías (máximo 5 o 6) y las diferencias entre ellas son fáciles de distinguir visualmente.

Para generar un gráfico profesional y legible, se utilizan los siguientes parámetros clave:

\begin{itemize}
\item \texttt{labels}: Lista de nombres para cada sector.
\item \texttt{autopct}: Formatea automáticamente los valores numéricos en porcentajes dentro de la torta. \texttt{'\%1.1f\%\%'} muestra un decimal (ej: 25.4\%).
\item \texttt{startangle}: Define dónde empieza la primera rebanada. Usar \texttt{90} hace que el gráfico comience en la parte superior ("las 12 en punto"), lo cual es un estándar visual.
\item \texttt{explode}: Permite "despegar" o resaltar una porción específica del resto del círculo.
\item \texttt{wedgeprops}: Permite añadir bordes a los sectores (usualmente blancos) para que los colores no se toquen directamente, mejorando la nitidez.
\end{itemize}

\begin{pycode}{Gráfico de torta}
import matplotlib.pyplot as plt

fig, ax = plt.subplots()

# Datos 
paises = ['Brasil', 'China', 'EE.UU.', 'Otros']
valores = [35, 30, 20, 15]
desplazamiento = [0.1, 0, 0, 0] # Resaltamos a Brasil

ax.pie(
    valores,
    labels=paises,
    explode=desplazamiento,
    autopct="%1.1f%%",
    startangle=140,
    shadow=True, # Agrega una sombra sutil para dar profundidad
    wedgeprops={"linewidth": 1.5, "edgecolor": "white"}
)

ax.set_aspect('equal') # Para que sea circular
ax.set_title("Participación de países en el mercado")

plt.show()
\end{pycode}

\begin{nota}
Por defecto, si el lienzo de la figura es rectangular, la torta se verá ovalada. Es obligatorio usar \texttt{ax.set\_aspect('equal')} para forzar una relación de aspecto 1:1 y que el círculo sea perfecto.
\end{nota}

\begin{figure}[ht]
    \centering
    %% Creator: Matplotlib, PGF backend
%%
%% To include the figure in your LaTeX document, write
%%   \input{<filename>.pgf}
%%
%% Make sure the required packages are loaded in your preamble
%%   \usepackage{pgf}
%%
%% Also ensure that all the required font packages are loaded; for instance,
%% the lmodern package is sometimes necessary when using math font.
%%   \usepackage{lmodern}
%%
%% Figures using additional raster images can only be included by \input if
%% they are in the same directory as the main LaTeX file. For loading figures
%% from other directories you can use the `import` package
%%   \usepackage{import}
%%
%% and then include the figures with
%%   \import{<path to file>}{<filename>.pgf}
%%
%% Matplotlib used the following preamble
%%   \def\mathdefault#1{#1}
%%   \everymath=\expandafter{\the\everymath\displaystyle}
%%   \IfFileExists{scrextend.sty}{
%%     \usepackage[fontsize=10.000000pt]{scrextend}
%%   }{
%%     \renewcommand{\normalsize}{\fontsize{10.000000}{12.000000}\selectfont}
%%     \normalsize
%%   }
%%   \usepackage{fontspec}
%%   \setmainfont{CMU Serif}
%%   \makeatletter\@ifpackageloaded{underscore}{}{\usepackage[strings]{underscore}}\makeatother
%%
\begingroup%
\makeatletter%
\begin{pgfpicture}%
\pgfpathrectangle{\pgfpointorigin}{\pgfqpoint{4.179843in}{4.268278in}}%
\pgfusepath{use as bounding box, clip}%
\begin{pgfscope}%
\pgfsetbuttcap%
\pgfsetmiterjoin%
\definecolor{currentfill}{rgb}{1.000000,1.000000,1.000000}%
\pgfsetfillcolor{currentfill}%
\pgfsetlinewidth{0.000000pt}%
\definecolor{currentstroke}{rgb}{1.000000,1.000000,1.000000}%
\pgfsetstrokecolor{currentstroke}%
\pgfsetdash{}{0pt}%
\pgfpathmoveto{\pgfqpoint{0.000000in}{0.000000in}}%
\pgfpathlineto{\pgfqpoint{4.179843in}{0.000000in}}%
\pgfpathlineto{\pgfqpoint{4.179843in}{4.268278in}}%
\pgfpathlineto{\pgfqpoint{0.000000in}{4.268278in}}%
\pgfpathlineto{\pgfqpoint{0.000000in}{0.000000in}}%
\pgfpathclose%
\pgfusepath{fill}%
\end{pgfscope}%
\begin{pgfscope}%
\pgfsetbuttcap%
\pgfsetmiterjoin%
\definecolor{currentfill}{rgb}{0.036471,0.140000,0.211765}%
\pgfsetfillcolor{currentfill}%
\pgfsetfillopacity{0.500000}%
\pgfsetlinewidth{1.505625pt}%
\definecolor{currentstroke}{rgb}{0.036471,0.140000,0.211765}%
\pgfsetstrokecolor{currentstroke}%
\pgfsetstrokeopacity{0.500000}%
\pgfsetdash{}{0pt}%
\pgfpathmoveto{\pgfqpoint{0.790599in}{2.911943in}}%
\pgfpathcurveto{\pgfqpoint{0.608057in}{2.694398in}}{\pgfqpoint{0.489198in}{2.430541in}}{\pgfqpoint{0.447223in}{2.149676in}}%
\pgfpathcurveto{\pgfqpoint{0.405247in}{1.868811in}}{\pgfqpoint{0.441768in}{1.581733in}}{\pgfqpoint{0.552730in}{1.320324in}}%
\pgfpathcurveto{\pgfqpoint{0.663691in}{1.058914in}}{\pgfqpoint{0.844828in}{0.833222in}}{\pgfqpoint{1.076025in}{0.668312in}}%
\pgfpathcurveto{\pgfqpoint{1.307221in}{0.503401in}}{\pgfqpoint{1.579589in}{0.405611in}}{\pgfqpoint{1.862882in}{0.385801in}}%
\pgfpathlineto{\pgfqpoint{1.970307in}{1.922050in}}%
\pgfpathlineto{\pgfqpoint{0.790599in}{2.911943in}}%
\pgfpathclose%
\pgfusepath{stroke,fill}%
\end{pgfscope}%
\begin{pgfscope}%
\pgfsetbuttcap%
\pgfsetmiterjoin%
\definecolor{currentfill}{rgb}{0.300000,0.149412,0.016471}%
\pgfsetfillcolor{currentfill}%
\pgfsetfillopacity{0.500000}%
\pgfsetlinewidth{1.505625pt}%
\definecolor{currentstroke}{rgb}{0.300000,0.149412,0.016471}%
\pgfsetstrokecolor{currentstroke}%
\pgfsetstrokeopacity{0.500000}%
\pgfsetdash{}{0pt}%
\pgfpathmoveto{\pgfqpoint{2.004640in}{0.445974in}}%
\pgfpathcurveto{\pgfqpoint{2.247064in}{0.429022in}}{\pgfqpoint{2.490093in}{0.469691in}}{\pgfqpoint{2.713790in}{0.564645in}}%
\pgfpathcurveto{\pgfqpoint{2.937488in}{0.659598in}}{\pgfqpoint{3.135566in}{0.806168in}}{\pgfqpoint{3.291773in}{0.992329in}}%
\pgfpathcurveto{\pgfqpoint{3.447981in}{1.178490in}}{\pgfqpoint{3.557928in}{1.399010in}}{\pgfqpoint{3.612595in}{1.635798in}}%
\pgfpathcurveto{\pgfqpoint{3.667262in}{1.872585in}}{\pgfqpoint{3.665112in}{2.118984in}}{\pgfqpoint{3.606321in}{2.354782in}}%
\pgfpathlineto{\pgfqpoint{2.112065in}{1.982222in}}%
\pgfpathlineto{\pgfqpoint{2.004640in}{0.445974in}}%
\pgfpathclose%
\pgfusepath{stroke,fill}%
\end{pgfscope}%
\begin{pgfscope}%
\pgfsetbuttcap%
\pgfsetmiterjoin%
\definecolor{currentfill}{rgb}{0.051765,0.188235,0.051765}%
\pgfsetfillcolor{currentfill}%
\pgfsetfillopacity{0.500000}%
\pgfsetlinewidth{1.505625pt}%
\definecolor{currentstroke}{rgb}{0.051765,0.188235,0.051765}%
\pgfsetstrokecolor{currentstroke}%
\pgfsetstrokeopacity{0.500000}%
\pgfsetdash{}{0pt}%
\pgfpathmoveto{\pgfqpoint{3.606321in}{2.354782in}}%
\pgfpathcurveto{\pgfqpoint{3.527656in}{2.670288in}}{\pgfqpoint{3.351049in}{2.952918in}}{\pgfqpoint{3.101958in}{3.161931in}}%
\pgfpathcurveto{\pgfqpoint{2.852867in}{3.370943in}}{\pgfqpoint{2.543863in}{3.495789in}}{\pgfqpoint{2.219490in}{3.518471in}}%
\pgfpathlineto{\pgfqpoint{2.112065in}{1.982222in}}%
\pgfpathlineto{\pgfqpoint{3.606321in}{2.354782in}}%
\pgfpathclose%
\pgfusepath{stroke,fill}%
\end{pgfscope}%
\begin{pgfscope}%
\pgfsetbuttcap%
\pgfsetmiterjoin%
\definecolor{currentfill}{rgb}{0.251765,0.045882,0.047059}%
\pgfsetfillcolor{currentfill}%
\pgfsetfillopacity{0.500000}%
\pgfsetlinewidth{1.505625pt}%
\definecolor{currentstroke}{rgb}{0.251765,0.045882,0.047059}%
\pgfsetstrokecolor{currentstroke}%
\pgfsetstrokeopacity{0.500000}%
\pgfsetdash{}{0pt}%
\pgfpathmoveto{\pgfqpoint{2.219490in}{3.518471in}}%
\pgfpathcurveto{\pgfqpoint{1.977065in}{3.535423in}}{\pgfqpoint{1.734037in}{3.494754in}}{\pgfqpoint{1.510339in}{3.399800in}}%
\pgfpathcurveto{\pgfqpoint{1.286642in}{3.304846in}}{\pgfqpoint{1.088564in}{3.158276in}}{\pgfqpoint{0.932356in}{2.972115in}}%
\pgfpathlineto{\pgfqpoint{2.112065in}{1.982222in}}%
\pgfpathlineto{\pgfqpoint{2.219490in}{3.518471in}}%
\pgfpathclose%
\pgfusepath{stroke,fill}%
\end{pgfscope}%
\begin{pgfscope}%
\pgfsetbuttcap%
\pgfsetmiterjoin%
\definecolor{currentfill}{rgb}{0.121569,0.466667,0.705882}%
\pgfsetfillcolor{currentfill}%
\pgfsetlinewidth{1.505625pt}%
\definecolor{currentstroke}{rgb}{1.000000,1.000000,1.000000}%
\pgfsetstrokecolor{currentstroke}%
\pgfsetdash{}{0pt}%
\pgfpathmoveto{\pgfqpoint{0.833377in}{2.954720in}}%
\pgfpathcurveto{\pgfqpoint{0.650835in}{2.737175in}}{\pgfqpoint{0.531976in}{2.473319in}}{\pgfqpoint{0.490001in}{2.192454in}}%
\pgfpathcurveto{\pgfqpoint{0.448025in}{1.911589in}}{\pgfqpoint{0.484546in}{1.624511in}}{\pgfqpoint{0.595508in}{1.363101in}}%
\pgfpathcurveto{\pgfqpoint{0.706469in}{1.101692in}}{\pgfqpoint{0.887606in}{0.876000in}}{\pgfqpoint{1.118803in}{0.711089in}}%
\pgfpathcurveto{\pgfqpoint{1.349999in}{0.546179in}}{\pgfqpoint{1.622367in}{0.448389in}}{\pgfqpoint{1.905660in}{0.428579in}}%
\pgfpathlineto{\pgfqpoint{2.013085in}{1.964827in}}%
\pgfpathlineto{\pgfqpoint{0.833377in}{2.954720in}}%
\pgfpathclose%
\pgfusepath{stroke,fill}%
\end{pgfscope}%
\begin{pgfscope}%
\pgfsetbuttcap%
\pgfsetmiterjoin%
\definecolor{currentfill}{rgb}{1.000000,0.498039,0.054902}%
\pgfsetfillcolor{currentfill}%
\pgfsetlinewidth{1.505625pt}%
\definecolor{currentstroke}{rgb}{1.000000,1.000000,1.000000}%
\pgfsetstrokecolor{currentstroke}%
\pgfsetdash{}{0pt}%
\pgfpathmoveto{\pgfqpoint{2.047418in}{0.488751in}}%
\pgfpathcurveto{\pgfqpoint{2.289842in}{0.471800in}}{\pgfqpoint{2.532871in}{0.512468in}}{\pgfqpoint{2.756568in}{0.607422in}}%
\pgfpathcurveto{\pgfqpoint{2.980266in}{0.702376in}}{\pgfqpoint{3.178343in}{0.848946in}}{\pgfqpoint{3.334551in}{1.035107in}}%
\pgfpathcurveto{\pgfqpoint{3.490759in}{1.221268in}}{\pgfqpoint{3.600706in}{1.441788in}}{\pgfqpoint{3.655373in}{1.678575in}}%
\pgfpathcurveto{\pgfqpoint{3.710039in}{1.915363in}}{\pgfqpoint{3.707889in}{2.161762in}}{\pgfqpoint{3.649098in}{2.397559in}}%
\pgfpathlineto{\pgfqpoint{2.154843in}{2.025000in}}%
\pgfpathlineto{\pgfqpoint{2.047418in}{0.488751in}}%
\pgfpathclose%
\pgfusepath{stroke,fill}%
\end{pgfscope}%
\begin{pgfscope}%
\pgfsetbuttcap%
\pgfsetmiterjoin%
\definecolor{currentfill}{rgb}{0.172549,0.627451,0.172549}%
\pgfsetfillcolor{currentfill}%
\pgfsetlinewidth{1.505625pt}%
\definecolor{currentstroke}{rgb}{1.000000,1.000000,1.000000}%
\pgfsetstrokecolor{currentstroke}%
\pgfsetdash{}{0pt}%
\pgfpathmoveto{\pgfqpoint{3.649098in}{2.397559in}}%
\pgfpathcurveto{\pgfqpoint{3.570434in}{2.713066in}}{\pgfqpoint{3.393827in}{2.995696in}}{\pgfqpoint{3.144736in}{3.204708in}}%
\pgfpathcurveto{\pgfqpoint{2.895644in}{3.413721in}}{\pgfqpoint{2.586641in}{3.538566in}}{\pgfqpoint{2.262267in}{3.561249in}}%
\pgfpathlineto{\pgfqpoint{2.154843in}{2.025000in}}%
\pgfpathlineto{\pgfqpoint{3.649098in}{2.397559in}}%
\pgfpathclose%
\pgfusepath{stroke,fill}%
\end{pgfscope}%
\begin{pgfscope}%
\pgfsetbuttcap%
\pgfsetmiterjoin%
\definecolor{currentfill}{rgb}{0.839216,0.152941,0.156863}%
\pgfsetfillcolor{currentfill}%
\pgfsetlinewidth{1.505625pt}%
\definecolor{currentstroke}{rgb}{1.000000,1.000000,1.000000}%
\pgfsetstrokecolor{currentstroke}%
\pgfsetdash{}{0pt}%
\pgfpathmoveto{\pgfqpoint{2.262267in}{3.561249in}}%
\pgfpathcurveto{\pgfqpoint{2.019843in}{3.578200in}}{\pgfqpoint{1.776815in}{3.537531in}}{\pgfqpoint{1.553117in}{3.442577in}}%
\pgfpathcurveto{\pgfqpoint{1.329419in}{3.347623in}}{\pgfqpoint{1.131342in}{3.201054in}}{\pgfqpoint{0.975134in}{3.014893in}}%
\pgfpathlineto{\pgfqpoint{2.154843in}{2.025000in}}%
\pgfpathlineto{\pgfqpoint{2.262267in}{3.561249in}}%
\pgfpathclose%
\pgfusepath{stroke,fill}%
\end{pgfscope}%
\begin{pgfscope}%
\definecolor{textcolor}{rgb}{0.000000,0.000000,0.000000}%
\pgfsetstrokecolor{textcolor}%
\pgfsetfillcolor{textcolor}%
\pgftext[x=0.453750in,y=1.302929in,right,]{\color{textcolor}{\rmfamily\fontsize{10.000000}{12.000000}\selectfont\catcode`\^=\active\def^{\ifmmode\sp\else\^{}\fi}\catcode`\%=\active\def%{\%}Brasil}}%
\end{pgfscope}%
\begin{pgfscope}%
\definecolor{textcolor}{rgb}{0.000000,0.000000,0.000000}%
\pgfsetstrokecolor{textcolor}%
\pgfsetfillcolor{textcolor}%
\pgftext[x=1.162539in,y=1.603792in,,]{\color{textcolor}{\rmfamily\fontsize{10.000000}{12.000000}\selectfont\catcode`\^=\active\def^{\ifmmode\sp\else\^{}\fi}\catcode`\%=\active\def%{\%}35.0\%}}%
\end{pgfscope}%
\begin{pgfscope}%
\definecolor{textcolor}{rgb}{0.000000,0.000000,0.000000}%
\pgfsetstrokecolor{textcolor}%
\pgfsetfillcolor{textcolor}%
\pgftext[x=3.452522in,y=0.936118in,left,]{\color{textcolor}{\rmfamily\fontsize{10.000000}{12.000000}\selectfont\catcode`\^=\active\def^{\ifmmode\sp\else\^{}\fi}\catcode`\%=\active\def%{\%}China}}%
\end{pgfscope}%
\begin{pgfscope}%
\definecolor{textcolor}{rgb}{0.000000,0.000000,0.000000}%
\pgfsetstrokecolor{textcolor}%
\pgfsetfillcolor{textcolor}%
\pgftext[x=2.862668in,y=1.431064in,,]{\color{textcolor}{\rmfamily\fontsize{10.000000}{12.000000}\selectfont\catcode`\^=\active\def^{\ifmmode\sp\else\^{}\fi}\catcode`\%=\active\def%{\%}30.0\%}}%
\end{pgfscope}%
\begin{pgfscope}%
\definecolor{textcolor}{rgb}{0.000000,0.000000,0.000000}%
\pgfsetstrokecolor{textcolor}%
\pgfsetfillcolor{textcolor}%
\pgftext[x=3.243724in,y=3.322680in,left,]{\color{textcolor}{\rmfamily\fontsize{10.000000}{12.000000}\selectfont\catcode`\^=\active\def^{\ifmmode\sp\else\^{}\fi}\catcode`\%=\active\def%{\%}EE.UU.}}%
\end{pgfscope}%
\begin{pgfscope}%
\definecolor{textcolor}{rgb}{0.000000,0.000000,0.000000}%
\pgfsetstrokecolor{textcolor}%
\pgfsetfillcolor{textcolor}%
\pgftext[x=2.748778in,y=2.732825in,,]{\color{textcolor}{\rmfamily\fontsize{10.000000}{12.000000}\selectfont\catcode`\^=\active\def^{\ifmmode\sp\else\^{}\fi}\catcode`\%=\active\def%{\%}20.0\%}}%
\end{pgfscope}%
\begin{pgfscope}%
\definecolor{textcolor}{rgb}{0.000000,0.000000,0.000000}%
\pgfsetstrokecolor{textcolor}%
\pgfsetfillcolor{textcolor}%
\pgftext[x=1.492943in,y=3.584335in,right,]{\color{textcolor}{\rmfamily\fontsize{10.000000}{12.000000}\selectfont\catcode`\^=\active\def^{\ifmmode\sp\else\^{}\fi}\catcode`\%=\active\def%{\%}Otros}}%
\end{pgfscope}%
\begin{pgfscope}%
\definecolor{textcolor}{rgb}{0.000000,0.000000,0.000000}%
\pgfsetstrokecolor{textcolor}%
\pgfsetfillcolor{textcolor}%
\pgftext[x=1.793806in,y=2.875546in,,]{\color{textcolor}{\rmfamily\fontsize{10.000000}{12.000000}\selectfont\catcode`\^=\active\def^{\ifmmode\sp\else\^{}\fi}\catcode`\%=\active\def%{\%}15.0\%}}%
\end{pgfscope}%
\begin{pgfscope}%
\definecolor{textcolor}{rgb}{0.000000,0.000000,0.000000}%
\pgfsetstrokecolor{textcolor}%
\pgfsetfillcolor{textcolor}%
\pgftext[x=2.154843in,y=4.033333in,,base]{\color{textcolor}{\rmfamily\fontsize{14.000000}{16.800000}\selectfont\catcode`\^=\active\def^{\ifmmode\sp\else\^{}\fi}\catcode`\%=\active\def%{\%}Participación de países en el mercado}}%
\end{pgfscope}%
\end{pgfpicture}%
\makeatother%
\endgroup%
 
    \caption{Gráfico de torta}
    \label{fig:grafico_torta_avanzado}
\end{figure}



\subsection{Gráficos de área (\texttt{.fill\_between()})}

La función \texttt{.fill\_between()} permite colorear el espacio comprendido entre dos líneas horizontales o entre una curva y un valor de referencia (generalmente el eje 0).

Esta técnica es útil para:
\begin{itemize}
\item Resaltar volumen: Da peso visual a la magnitud total acumulada a lo largo del tiempo.
\item Mostrar desviaciones: Sombrear el rango entre un valor proyectado y uno real.
\item Diferenciar condiciones: Usar el parámetro \texttt{where} para cambiar el color del área según qué curva sea superior.
\end{itemize}

\paragraph{Sintaxis básica}
\begin{verbatim}
ax.fill_between(x, y1, y2, where=condicion, alpha=transparencia)
\end{verbatim}

\begin{pycode}{Gráfico de área bajo la curva}
import numpy as np
import matplotlib.pyplot as plt

# Datos
x = np.linspace(0, 10, 200)
y = np.sin(x) + 1.5

fig, ax = plt.subplots()
ax.plot(x, y, label="Señal")

# Área bajo la curva
ax.fill_between(x, y, 0, alpha=0.3) # Relleno hasta el eje X

ax.set_title("Área bajo la curva")
ax.set_xlabel("Tiempo")
ax.set_ylabel("Amplitud")
ax.legend()
plt.show()
\end{pycode}

\begin{pycode}{Área entre curvas}
import numpy as np
import matplotlib.pyplot as plt

# Datos
x = np.linspace(0, 2 * np.pi, 300)
y1 = np.sin(x) + 1
y2 = np.cos(x) + 1

fig, ax = plt.subplots()
ax.plot(x, y1, label="Seno")
ax.plot(x, y2, label="Coseno")

# Área bajo la curva
ax.fill_between(x, y1, y2, where=(y1 >= y2), alpha=0.25, label="Área (y1 >=y2)")
ax.fill_between(x, y1, y2, where=(y1 < y2), alpha=0.25, label="Área (y1 < y2)")

ax.set_title("Área entre dos curvas")
ax.set_xlabel("Ángulo")
ax.set_ylabel("Valor")
ax.legend()
plt.show()    
\end{pycode}

\begin{pycode}{Ambos gráficos en uno}
import numpy as np
import matplotlib.pyplot as plt

fig, ax = plt.subplots(1,2, figsize=(5, 5))

# ------------------ #
#       Ax[0]        #
# ------------------ #
x = np.linspace(0, 10, 200)
y = np.sin(x) + 1.5

ax[0].plot(x,y, label='señal')
ax[0].fill_between(x, y, 0, alpha=0.3) # 0 es para que llegue al eje x, alpha es la transparencia.
ax[0].set_title("Área bajo la curva")
ax[0].set_xlabel("Tiempo")
ax[0].set_ylabel("Amplitud")
ax[0].legend()

# ------------------ #
#       Ax[1]        #
# ------------------ #
x_1 = np.linspace(0, 2 * np.pi, 300)
y_1 = np.sin(x_1) + 1
y_2 = np.cos(x_1) + 1

ax[1].plot(x_1, y_1, label="Seno")
ax[1].plot(x_1, y_2, label="Coseno")
ax[1].fill_between(x_1, y_1, y_2, where=(y_1 >= y_2), alpha=0.25, label="Área (y1 >=y2)")
ax[1].fill_between(x_1, y_1, y_2, where=(y_1 < y_2), alpha=0.25, label="Área (y1 < y2)")
ax[1].set_title("Área entre dos curvas")
ax[1].set_xlabel("Ángulo")
ax[1].set_ylabel("Valor")
ax[1].legend()

# Título global
fig.suptitle('Gráfico de Áreas', fontsize=14)  
\end{pycode}

La convinación de ambos graficos nos deja la figura \ref{fig:grafico_area_curva}

\begin{figure}[ht]
    \centering
    %% Creator: Matplotlib, PGF backend
%%
%% To include the figure in your LaTeX document, write
%%   \input{<filename>.pgf}
%%
%% Make sure the required packages are loaded in your preamble
%%   \usepackage{pgf}
%%
%% Also ensure that all the required font packages are loaded; for instance,
%% the lmodern package is sometimes necessary when using math font.
%%   \usepackage{lmodern}
%%
%% Figures using additional raster images can only be included by \input if
%% they are in the same directory as the main LaTeX file. For loading figures
%% from other directories you can use the `import` package
%%   \usepackage{import}
%%
%% and then include the figures with
%%   \import{<path to file>}{<filename>.pgf}
%%
%% Matplotlib used the following preamble
%%   \def\mathdefault#1{#1}
%%   \everymath=\expandafter{\the\everymath\displaystyle}
%%   \IfFileExists{scrextend.sty}{
%%     \usepackage[fontsize=10.000000pt]{scrextend}
%%   }{
%%     \renewcommand{\normalsize}{\fontsize{10.000000}{12.000000}\selectfont}
%%     \normalsize
%%   }
%%   \usepackage{fontspec}
%%   \setmainfont{CMU Serif}
%%   \makeatletter\@ifpackageloaded{underscore}{}{\usepackage[strings]{underscore}}\makeatother
%%
\begingroup%
\makeatletter%
\begin{pgfpicture}%
\pgfpathrectangle{\pgfpointorigin}{\pgfqpoint{4.648501in}{4.812921in}}%
\pgfusepath{use as bounding box, clip}%
\begin{pgfscope}%
\pgfsetbuttcap%
\pgfsetmiterjoin%
\definecolor{currentfill}{rgb}{1.000000,1.000000,1.000000}%
\pgfsetfillcolor{currentfill}%
\pgfsetlinewidth{0.000000pt}%
\definecolor{currentstroke}{rgb}{1.000000,1.000000,1.000000}%
\pgfsetstrokecolor{currentstroke}%
\pgfsetdash{}{0pt}%
\pgfpathmoveto{\pgfqpoint{0.000000in}{0.000000in}}%
\pgfpathlineto{\pgfqpoint{4.648501in}{0.000000in}}%
\pgfpathlineto{\pgfqpoint{4.648501in}{4.812921in}}%
\pgfpathlineto{\pgfqpoint{0.000000in}{4.812921in}}%
\pgfpathlineto{\pgfqpoint{0.000000in}{0.000000in}}%
\pgfpathclose%
\pgfusepath{fill}%
\end{pgfscope}%
\begin{pgfscope}%
\pgfsetbuttcap%
\pgfsetmiterjoin%
\definecolor{currentfill}{rgb}{1.000000,1.000000,1.000000}%
\pgfsetfillcolor{currentfill}%
\pgfsetlinewidth{0.000000pt}%
\definecolor{currentstroke}{rgb}{0.000000,0.000000,0.000000}%
\pgfsetstrokecolor{currentstroke}%
\pgfsetstrokeopacity{0.000000}%
\pgfsetdash{}{0pt}%
\pgfpathmoveto{\pgfqpoint{0.556636in}{0.524444in}}%
\pgfpathlineto{\pgfqpoint{2.152797in}{0.524444in}}%
\pgfpathlineto{\pgfqpoint{2.152797in}{3.833258in}}%
\pgfpathlineto{\pgfqpoint{0.556636in}{3.833258in}}%
\pgfpathlineto{\pgfqpoint{0.556636in}{0.524444in}}%
\pgfpathclose%
\pgfusepath{fill}%
\end{pgfscope}%
\begin{pgfscope}%
\pgfpathrectangle{\pgfqpoint{0.556636in}{0.524444in}}{\pgfqpoint{1.596161in}{3.308814in}}%
\pgfusepath{clip}%
\pgfsetbuttcap%
\pgfsetroundjoin%
\definecolor{currentfill}{rgb}{0.121569,0.466667,0.705882}%
\pgfsetfillcolor{currentfill}%
\pgfsetfillopacity{0.300000}%
\pgfsetlinewidth{0.000000pt}%
\definecolor{currentstroke}{rgb}{0.000000,0.000000,0.000000}%
\pgfsetstrokecolor{currentstroke}%
\pgfsetdash{}{0pt}%
\pgfpathmoveto{\pgfqpoint{0.629189in}{0.674845in}}%
\pgfpathlineto{\pgfqpoint{0.629189in}{2.479713in}}%
\pgfpathlineto{\pgfqpoint{0.636481in}{2.540153in}}%
\pgfpathlineto{\pgfqpoint{0.643772in}{2.600439in}}%
\pgfpathlineto{\pgfqpoint{0.651064in}{2.660421in}}%
\pgfpathlineto{\pgfqpoint{0.658356in}{2.719947in}}%
\pgfpathlineto{\pgfqpoint{0.665648in}{2.778866in}}%
\pgfpathlineto{\pgfqpoint{0.672939in}{2.837029in}}%
\pgfpathlineto{\pgfqpoint{0.680231in}{2.894291in}}%
\pgfpathlineto{\pgfqpoint{0.687523in}{2.950506in}}%
\pgfpathlineto{\pgfqpoint{0.694815in}{3.005533in}}%
\pgfpathlineto{\pgfqpoint{0.702106in}{3.059231in}}%
\pgfpathlineto{\pgfqpoint{0.709398in}{3.111467in}}%
\pgfpathlineto{\pgfqpoint{0.716690in}{3.162108in}}%
\pgfpathlineto{\pgfqpoint{0.723982in}{3.211026in}}%
\pgfpathlineto{\pgfqpoint{0.731273in}{3.258098in}}%
\pgfpathlineto{\pgfqpoint{0.738565in}{3.303205in}}%
\pgfpathlineto{\pgfqpoint{0.745857in}{3.346232in}}%
\pgfpathlineto{\pgfqpoint{0.753149in}{3.387072in}}%
\pgfpathlineto{\pgfqpoint{0.760440in}{3.425621in}}%
\pgfpathlineto{\pgfqpoint{0.767732in}{3.461782in}}%
\pgfpathlineto{\pgfqpoint{0.775024in}{3.495464in}}%
\pgfpathlineto{\pgfqpoint{0.782315in}{3.526581in}}%
\pgfpathlineto{\pgfqpoint{0.789607in}{3.555055in}}%
\pgfpathlineto{\pgfqpoint{0.796899in}{3.580815in}}%
\pgfpathlineto{\pgfqpoint{0.804191in}{3.603794in}}%
\pgfpathlineto{\pgfqpoint{0.811482in}{3.623936in}}%
\pgfpathlineto{\pgfqpoint{0.818774in}{3.641189in}}%
\pgfpathlineto{\pgfqpoint{0.826066in}{3.655509in}}%
\pgfpathlineto{\pgfqpoint{0.833358in}{3.666861in}}%
\pgfpathlineto{\pgfqpoint{0.840649in}{3.675216in}}%
\pgfpathlineto{\pgfqpoint{0.847941in}{3.680553in}}%
\pgfpathlineto{\pgfqpoint{0.855233in}{3.682857in}}%
\pgfpathlineto{\pgfqpoint{0.862525in}{3.682125in}}%
\pgfpathlineto{\pgfqpoint{0.869816in}{3.678357in}}%
\pgfpathlineto{\pgfqpoint{0.877108in}{3.671562in}}%
\pgfpathlineto{\pgfqpoint{0.884400in}{3.661759in}}%
\pgfpathlineto{\pgfqpoint{0.891691in}{3.648971in}}%
\pgfpathlineto{\pgfqpoint{0.898983in}{3.633231in}}%
\pgfpathlineto{\pgfqpoint{0.906275in}{3.614579in}}%
\pgfpathlineto{\pgfqpoint{0.913567in}{3.593062in}}%
\pgfpathlineto{\pgfqpoint{0.920858in}{3.568734in}}%
\pgfpathlineto{\pgfqpoint{0.928150in}{3.541657in}}%
\pgfpathlineto{\pgfqpoint{0.935442in}{3.511898in}}%
\pgfpathlineto{\pgfqpoint{0.942734in}{3.479534in}}%
\pgfpathlineto{\pgfqpoint{0.950025in}{3.444646in}}%
\pgfpathlineto{\pgfqpoint{0.957317in}{3.407321in}}%
\pgfpathlineto{\pgfqpoint{0.964609in}{3.367655in}}%
\pgfpathlineto{\pgfqpoint{0.971901in}{3.325746in}}%
\pgfpathlineto{\pgfqpoint{0.979192in}{3.281702in}}%
\pgfpathlineto{\pgfqpoint{0.986484in}{3.235633in}}%
\pgfpathlineto{\pgfqpoint{0.993776in}{3.187656in}}%
\pgfpathlineto{\pgfqpoint{1.001067in}{3.137891in}}%
\pgfpathlineto{\pgfqpoint{1.008359in}{3.086465in}}%
\pgfpathlineto{\pgfqpoint{1.015651in}{3.033507in}}%
\pgfpathlineto{\pgfqpoint{1.022943in}{2.979150in}}%
\pgfpathlineto{\pgfqpoint{1.030234in}{2.923533in}}%
\pgfpathlineto{\pgfqpoint{1.037526in}{2.866795in}}%
\pgfpathlineto{\pgfqpoint{1.044818in}{2.809080in}}%
\pgfpathlineto{\pgfqpoint{1.052110in}{2.750534in}}%
\pgfpathlineto{\pgfqpoint{1.059401in}{2.691303in}}%
\pgfpathlineto{\pgfqpoint{1.066693in}{2.631539in}}%
\pgfpathlineto{\pgfqpoint{1.073985in}{2.571391in}}%
\pgfpathlineto{\pgfqpoint{1.081277in}{2.511012in}}%
\pgfpathlineto{\pgfqpoint{1.088568in}{2.450554in}}%
\pgfpathlineto{\pgfqpoint{1.095860in}{2.390169in}}%
\pgfpathlineto{\pgfqpoint{1.103152in}{2.330011in}}%
\pgfpathlineto{\pgfqpoint{1.110444in}{2.270230in}}%
\pgfpathlineto{\pgfqpoint{1.117735in}{2.210978in}}%
\pgfpathlineto{\pgfqpoint{1.125027in}{2.152405in}}%
\pgfpathlineto{\pgfqpoint{1.132319in}{2.094658in}}%
\pgfpathlineto{\pgfqpoint{1.139610in}{2.037883in}}%
\pgfpathlineto{\pgfqpoint{1.146902in}{1.982224in}}%
\pgfpathlineto{\pgfqpoint{1.154194in}{1.927821in}}%
\pgfpathlineto{\pgfqpoint{1.161486in}{1.874811in}}%
\pgfpathlineto{\pgfqpoint{1.168777in}{1.823328in}}%
\pgfpathlineto{\pgfqpoint{1.176069in}{1.773502in}}%
\pgfpathlineto{\pgfqpoint{1.183361in}{1.725460in}}%
\pgfpathlineto{\pgfqpoint{1.190653in}{1.679321in}}%
\pgfpathlineto{\pgfqpoint{1.197944in}{1.635203in}}%
\pgfpathlineto{\pgfqpoint{1.205236in}{1.593218in}}%
\pgfpathlineto{\pgfqpoint{1.212528in}{1.553470in}}%
\pgfpathlineto{\pgfqpoint{1.219820in}{1.516061in}}%
\pgfpathlineto{\pgfqpoint{1.227111in}{1.481085in}}%
\pgfpathlineto{\pgfqpoint{1.234403in}{1.448630in}}%
\pgfpathlineto{\pgfqpoint{1.241695in}{1.418778in}}%
\pgfpathlineto{\pgfqpoint{1.248986in}{1.391604in}}%
\pgfpathlineto{\pgfqpoint{1.256278in}{1.367178in}}%
\pgfpathlineto{\pgfqpoint{1.263570in}{1.345560in}}%
\pgfpathlineto{\pgfqpoint{1.270862in}{1.326806in}}%
\pgfpathlineto{\pgfqpoint{1.278153in}{1.310963in}}%
\pgfpathlineto{\pgfqpoint{1.285445in}{1.298070in}}%
\pgfpathlineto{\pgfqpoint{1.292737in}{1.288160in}}%
\pgfpathlineto{\pgfqpoint{1.300029in}{1.281259in}}%
\pgfpathlineto{\pgfqpoint{1.307320in}{1.277384in}}%
\pgfpathlineto{\pgfqpoint{1.314612in}{1.276543in}}%
\pgfpathlineto{\pgfqpoint{1.321904in}{1.278741in}}%
\pgfpathlineto{\pgfqpoint{1.329196in}{1.283970in}}%
\pgfpathlineto{\pgfqpoint{1.336487in}{1.292219in}}%
\pgfpathlineto{\pgfqpoint{1.343779in}{1.303465in}}%
\pgfpathlineto{\pgfqpoint{1.351071in}{1.317681in}}%
\pgfpathlineto{\pgfqpoint{1.358363in}{1.334831in}}%
\pgfpathlineto{\pgfqpoint{1.365654in}{1.354871in}}%
\pgfpathlineto{\pgfqpoint{1.372946in}{1.377751in}}%
\pgfpathlineto{\pgfqpoint{1.380238in}{1.403413in}}%
\pgfpathlineto{\pgfqpoint{1.387529in}{1.431792in}}%
\pgfpathlineto{\pgfqpoint{1.394821in}{1.462817in}}%
\pgfpathlineto{\pgfqpoint{1.402113in}{1.496410in}}%
\pgfpathlineto{\pgfqpoint{1.409405in}{1.532485in}}%
\pgfpathlineto{\pgfqpoint{1.416696in}{1.570951in}}%
\pgfpathlineto{\pgfqpoint{1.423988in}{1.611711in}}%
\pgfpathlineto{\pgfqpoint{1.431280in}{1.654663in}}%
\pgfpathlineto{\pgfqpoint{1.438572in}{1.699698in}}%
\pgfpathlineto{\pgfqpoint{1.445863in}{1.746703in}}%
\pgfpathlineto{\pgfqpoint{1.453155in}{1.795557in}}%
\pgfpathlineto{\pgfqpoint{1.460447in}{1.846139in}}%
\pgfpathlineto{\pgfqpoint{1.467739in}{1.898321in}}%
\pgfpathlineto{\pgfqpoint{1.475030in}{1.951971in}}%
\pgfpathlineto{\pgfqpoint{1.482322in}{2.006952in}}%
\pgfpathlineto{\pgfqpoint{1.489614in}{2.063128in}}%
\pgfpathlineto{\pgfqpoint{1.496905in}{2.120355in}}%
\pgfpathlineto{\pgfqpoint{1.504197in}{2.178489in}}%
\pgfpathlineto{\pgfqpoint{1.511489in}{2.237384in}}%
\pgfpathlineto{\pgfqpoint{1.518781in}{2.296891in}}%
\pgfpathlineto{\pgfqpoint{1.526072in}{2.356859in}}%
\pgfpathlineto{\pgfqpoint{1.533364in}{2.417137in}}%
\pgfpathlineto{\pgfqpoint{1.540656in}{2.477574in}}%
\pgfpathlineto{\pgfqpoint{1.547948in}{2.538015in}}%
\pgfpathlineto{\pgfqpoint{1.555239in}{2.598310in}}%
\pgfpathlineto{\pgfqpoint{1.562531in}{2.658305in}}%
\pgfpathlineto{\pgfqpoint{1.569823in}{2.717850in}}%
\pgfpathlineto{\pgfqpoint{1.577115in}{2.776793in}}%
\pgfpathlineto{\pgfqpoint{1.584406in}{2.834986in}}%
\pgfpathlineto{\pgfqpoint{1.591698in}{2.892282in}}%
\pgfpathlineto{\pgfqpoint{1.598990in}{2.948536in}}%
\pgfpathlineto{\pgfqpoint{1.606282in}{3.003607in}}%
\pgfpathlineto{\pgfqpoint{1.613573in}{3.057355in}}%
\pgfpathlineto{\pgfqpoint{1.620865in}{3.109645in}}%
\pgfpathlineto{\pgfqpoint{1.628157in}{3.160345in}}%
\pgfpathlineto{\pgfqpoint{1.635448in}{3.209326in}}%
\pgfpathlineto{\pgfqpoint{1.642740in}{3.256465in}}%
\pgfpathlineto{\pgfqpoint{1.650032in}{3.301643in}}%
\pgfpathlineto{\pgfqpoint{1.657324in}{3.344746in}}%
\pgfpathlineto{\pgfqpoint{1.664615in}{3.385665in}}%
\pgfpathlineto{\pgfqpoint{1.671907in}{3.424297in}}%
\pgfpathlineto{\pgfqpoint{1.679199in}{3.460544in}}%
\pgfpathlineto{\pgfqpoint{1.686491in}{3.494315in}}%
\pgfpathlineto{\pgfqpoint{1.693782in}{3.525525in}}%
\pgfpathlineto{\pgfqpoint{1.701074in}{3.554094in}}%
\pgfpathlineto{\pgfqpoint{1.708366in}{3.579950in}}%
\pgfpathlineto{\pgfqpoint{1.715658in}{3.603029in}}%
\pgfpathlineto{\pgfqpoint{1.722949in}{3.623272in}}%
\pgfpathlineto{\pgfqpoint{1.730241in}{3.640628in}}%
\pgfpathlineto{\pgfqpoint{1.737533in}{3.655053in}}%
\pgfpathlineto{\pgfqpoint{1.744824in}{3.666510in}}%
\pgfpathlineto{\pgfqpoint{1.752116in}{3.674972in}}%
\pgfpathlineto{\pgfqpoint{1.759408in}{3.680415in}}%
\pgfpathlineto{\pgfqpoint{1.766700in}{3.682828in}}%
\pgfpathlineto{\pgfqpoint{1.773991in}{3.682203in}}%
\pgfpathlineto{\pgfqpoint{1.781283in}{3.678542in}}%
\pgfpathlineto{\pgfqpoint{1.788575in}{3.671854in}}%
\pgfpathlineto{\pgfqpoint{1.795867in}{3.662157in}}%
\pgfpathlineto{\pgfqpoint{1.803158in}{3.649474in}}%
\pgfpathlineto{\pgfqpoint{1.810450in}{3.633838in}}%
\pgfpathlineto{\pgfqpoint{1.817742in}{3.615288in}}%
\pgfpathlineto{\pgfqpoint{1.825034in}{3.593872in}}%
\pgfpathlineto{\pgfqpoint{1.832325in}{3.569642in}}%
\pgfpathlineto{\pgfqpoint{1.839617in}{3.542661in}}%
\pgfpathlineto{\pgfqpoint{1.846909in}{3.512996in}}%
\pgfpathlineto{\pgfqpoint{1.854200in}{3.480723in}}%
\pgfpathlineto{\pgfqpoint{1.861492in}{3.445922in}}%
\pgfpathlineto{\pgfqpoint{1.868784in}{3.408682in}}%
\pgfpathlineto{\pgfqpoint{1.876076in}{3.369097in}}%
\pgfpathlineto{\pgfqpoint{1.883367in}{3.327266in}}%
\pgfpathlineto{\pgfqpoint{1.890659in}{3.283296in}}%
\pgfpathlineto{\pgfqpoint{1.897951in}{3.237297in}}%
\pgfpathlineto{\pgfqpoint{1.905243in}{3.189385in}}%
\pgfpathlineto{\pgfqpoint{1.912534in}{3.139681in}}%
\pgfpathlineto{\pgfqpoint{1.919826in}{3.088312in}}%
\pgfpathlineto{\pgfqpoint{1.927118in}{3.035405in}}%
\pgfpathlineto{\pgfqpoint{1.934410in}{2.981096in}}%
\pgfpathlineto{\pgfqpoint{1.941701in}{2.925521in}}%
\pgfpathlineto{\pgfqpoint{1.948993in}{2.868821in}}%
\pgfpathlineto{\pgfqpoint{1.956285in}{2.811138in}}%
\pgfpathlineto{\pgfqpoint{1.963577in}{2.752618in}}%
\pgfpathlineto{\pgfqpoint{1.970868in}{2.693409in}}%
\pgfpathlineto{\pgfqpoint{1.978160in}{2.633661in}}%
\pgfpathlineto{\pgfqpoint{1.985452in}{2.573525in}}%
\pgfpathlineto{\pgfqpoint{1.992743in}{2.513151in}}%
\pgfpathlineto{\pgfqpoint{2.000035in}{2.452693in}}%
\pgfpathlineto{\pgfqpoint{2.007327in}{2.392303in}}%
\pgfpathlineto{\pgfqpoint{2.014619in}{2.332134in}}%
\pgfpathlineto{\pgfqpoint{2.021910in}{2.272337in}}%
\pgfpathlineto{\pgfqpoint{2.029202in}{2.213064in}}%
\pgfpathlineto{\pgfqpoint{2.036494in}{2.154465in}}%
\pgfpathlineto{\pgfqpoint{2.043786in}{2.096686in}}%
\pgfpathlineto{\pgfqpoint{2.051077in}{2.039874in}}%
\pgfpathlineto{\pgfqpoint{2.058369in}{1.984173in}}%
\pgfpathlineto{\pgfqpoint{2.065661in}{1.929723in}}%
\pgfpathlineto{\pgfqpoint{2.072953in}{1.876661in}}%
\pgfpathlineto{\pgfqpoint{2.080244in}{1.825122in}}%
\pgfpathlineto{\pgfqpoint{2.080244in}{0.674845in}}%
\pgfpathlineto{\pgfqpoint{2.080244in}{0.674845in}}%
\pgfpathlineto{\pgfqpoint{2.072953in}{0.674845in}}%
\pgfpathlineto{\pgfqpoint{2.065661in}{0.674845in}}%
\pgfpathlineto{\pgfqpoint{2.058369in}{0.674845in}}%
\pgfpathlineto{\pgfqpoint{2.051077in}{0.674845in}}%
\pgfpathlineto{\pgfqpoint{2.043786in}{0.674845in}}%
\pgfpathlineto{\pgfqpoint{2.036494in}{0.674845in}}%
\pgfpathlineto{\pgfqpoint{2.029202in}{0.674845in}}%
\pgfpathlineto{\pgfqpoint{2.021910in}{0.674845in}}%
\pgfpathlineto{\pgfqpoint{2.014619in}{0.674845in}}%
\pgfpathlineto{\pgfqpoint{2.007327in}{0.674845in}}%
\pgfpathlineto{\pgfqpoint{2.000035in}{0.674845in}}%
\pgfpathlineto{\pgfqpoint{1.992743in}{0.674845in}}%
\pgfpathlineto{\pgfqpoint{1.985452in}{0.674845in}}%
\pgfpathlineto{\pgfqpoint{1.978160in}{0.674845in}}%
\pgfpathlineto{\pgfqpoint{1.970868in}{0.674845in}}%
\pgfpathlineto{\pgfqpoint{1.963577in}{0.674845in}}%
\pgfpathlineto{\pgfqpoint{1.956285in}{0.674845in}}%
\pgfpathlineto{\pgfqpoint{1.948993in}{0.674845in}}%
\pgfpathlineto{\pgfqpoint{1.941701in}{0.674845in}}%
\pgfpathlineto{\pgfqpoint{1.934410in}{0.674845in}}%
\pgfpathlineto{\pgfqpoint{1.927118in}{0.674845in}}%
\pgfpathlineto{\pgfqpoint{1.919826in}{0.674845in}}%
\pgfpathlineto{\pgfqpoint{1.912534in}{0.674845in}}%
\pgfpathlineto{\pgfqpoint{1.905243in}{0.674845in}}%
\pgfpathlineto{\pgfqpoint{1.897951in}{0.674845in}}%
\pgfpathlineto{\pgfqpoint{1.890659in}{0.674845in}}%
\pgfpathlineto{\pgfqpoint{1.883367in}{0.674845in}}%
\pgfpathlineto{\pgfqpoint{1.876076in}{0.674845in}}%
\pgfpathlineto{\pgfqpoint{1.868784in}{0.674845in}}%
\pgfpathlineto{\pgfqpoint{1.861492in}{0.674845in}}%
\pgfpathlineto{\pgfqpoint{1.854200in}{0.674845in}}%
\pgfpathlineto{\pgfqpoint{1.846909in}{0.674845in}}%
\pgfpathlineto{\pgfqpoint{1.839617in}{0.674845in}}%
\pgfpathlineto{\pgfqpoint{1.832325in}{0.674845in}}%
\pgfpathlineto{\pgfqpoint{1.825034in}{0.674845in}}%
\pgfpathlineto{\pgfqpoint{1.817742in}{0.674845in}}%
\pgfpathlineto{\pgfqpoint{1.810450in}{0.674845in}}%
\pgfpathlineto{\pgfqpoint{1.803158in}{0.674845in}}%
\pgfpathlineto{\pgfqpoint{1.795867in}{0.674845in}}%
\pgfpathlineto{\pgfqpoint{1.788575in}{0.674845in}}%
\pgfpathlineto{\pgfqpoint{1.781283in}{0.674845in}}%
\pgfpathlineto{\pgfqpoint{1.773991in}{0.674845in}}%
\pgfpathlineto{\pgfqpoint{1.766700in}{0.674845in}}%
\pgfpathlineto{\pgfqpoint{1.759408in}{0.674845in}}%
\pgfpathlineto{\pgfqpoint{1.752116in}{0.674845in}}%
\pgfpathlineto{\pgfqpoint{1.744824in}{0.674845in}}%
\pgfpathlineto{\pgfqpoint{1.737533in}{0.674845in}}%
\pgfpathlineto{\pgfqpoint{1.730241in}{0.674845in}}%
\pgfpathlineto{\pgfqpoint{1.722949in}{0.674845in}}%
\pgfpathlineto{\pgfqpoint{1.715658in}{0.674845in}}%
\pgfpathlineto{\pgfqpoint{1.708366in}{0.674845in}}%
\pgfpathlineto{\pgfqpoint{1.701074in}{0.674845in}}%
\pgfpathlineto{\pgfqpoint{1.693782in}{0.674845in}}%
\pgfpathlineto{\pgfqpoint{1.686491in}{0.674845in}}%
\pgfpathlineto{\pgfqpoint{1.679199in}{0.674845in}}%
\pgfpathlineto{\pgfqpoint{1.671907in}{0.674845in}}%
\pgfpathlineto{\pgfqpoint{1.664615in}{0.674845in}}%
\pgfpathlineto{\pgfqpoint{1.657324in}{0.674845in}}%
\pgfpathlineto{\pgfqpoint{1.650032in}{0.674845in}}%
\pgfpathlineto{\pgfqpoint{1.642740in}{0.674845in}}%
\pgfpathlineto{\pgfqpoint{1.635448in}{0.674845in}}%
\pgfpathlineto{\pgfqpoint{1.628157in}{0.674845in}}%
\pgfpathlineto{\pgfqpoint{1.620865in}{0.674845in}}%
\pgfpathlineto{\pgfqpoint{1.613573in}{0.674845in}}%
\pgfpathlineto{\pgfqpoint{1.606282in}{0.674845in}}%
\pgfpathlineto{\pgfqpoint{1.598990in}{0.674845in}}%
\pgfpathlineto{\pgfqpoint{1.591698in}{0.674845in}}%
\pgfpathlineto{\pgfqpoint{1.584406in}{0.674845in}}%
\pgfpathlineto{\pgfqpoint{1.577115in}{0.674845in}}%
\pgfpathlineto{\pgfqpoint{1.569823in}{0.674845in}}%
\pgfpathlineto{\pgfqpoint{1.562531in}{0.674845in}}%
\pgfpathlineto{\pgfqpoint{1.555239in}{0.674845in}}%
\pgfpathlineto{\pgfqpoint{1.547948in}{0.674845in}}%
\pgfpathlineto{\pgfqpoint{1.540656in}{0.674845in}}%
\pgfpathlineto{\pgfqpoint{1.533364in}{0.674845in}}%
\pgfpathlineto{\pgfqpoint{1.526072in}{0.674845in}}%
\pgfpathlineto{\pgfqpoint{1.518781in}{0.674845in}}%
\pgfpathlineto{\pgfqpoint{1.511489in}{0.674845in}}%
\pgfpathlineto{\pgfqpoint{1.504197in}{0.674845in}}%
\pgfpathlineto{\pgfqpoint{1.496905in}{0.674845in}}%
\pgfpathlineto{\pgfqpoint{1.489614in}{0.674845in}}%
\pgfpathlineto{\pgfqpoint{1.482322in}{0.674845in}}%
\pgfpathlineto{\pgfqpoint{1.475030in}{0.674845in}}%
\pgfpathlineto{\pgfqpoint{1.467739in}{0.674845in}}%
\pgfpathlineto{\pgfqpoint{1.460447in}{0.674845in}}%
\pgfpathlineto{\pgfqpoint{1.453155in}{0.674845in}}%
\pgfpathlineto{\pgfqpoint{1.445863in}{0.674845in}}%
\pgfpathlineto{\pgfqpoint{1.438572in}{0.674845in}}%
\pgfpathlineto{\pgfqpoint{1.431280in}{0.674845in}}%
\pgfpathlineto{\pgfqpoint{1.423988in}{0.674845in}}%
\pgfpathlineto{\pgfqpoint{1.416696in}{0.674845in}}%
\pgfpathlineto{\pgfqpoint{1.409405in}{0.674845in}}%
\pgfpathlineto{\pgfqpoint{1.402113in}{0.674845in}}%
\pgfpathlineto{\pgfqpoint{1.394821in}{0.674845in}}%
\pgfpathlineto{\pgfqpoint{1.387529in}{0.674845in}}%
\pgfpathlineto{\pgfqpoint{1.380238in}{0.674845in}}%
\pgfpathlineto{\pgfqpoint{1.372946in}{0.674845in}}%
\pgfpathlineto{\pgfqpoint{1.365654in}{0.674845in}}%
\pgfpathlineto{\pgfqpoint{1.358363in}{0.674845in}}%
\pgfpathlineto{\pgfqpoint{1.351071in}{0.674845in}}%
\pgfpathlineto{\pgfqpoint{1.343779in}{0.674845in}}%
\pgfpathlineto{\pgfqpoint{1.336487in}{0.674845in}}%
\pgfpathlineto{\pgfqpoint{1.329196in}{0.674845in}}%
\pgfpathlineto{\pgfqpoint{1.321904in}{0.674845in}}%
\pgfpathlineto{\pgfqpoint{1.314612in}{0.674845in}}%
\pgfpathlineto{\pgfqpoint{1.307320in}{0.674845in}}%
\pgfpathlineto{\pgfqpoint{1.300029in}{0.674845in}}%
\pgfpathlineto{\pgfqpoint{1.292737in}{0.674845in}}%
\pgfpathlineto{\pgfqpoint{1.285445in}{0.674845in}}%
\pgfpathlineto{\pgfqpoint{1.278153in}{0.674845in}}%
\pgfpathlineto{\pgfqpoint{1.270862in}{0.674845in}}%
\pgfpathlineto{\pgfqpoint{1.263570in}{0.674845in}}%
\pgfpathlineto{\pgfqpoint{1.256278in}{0.674845in}}%
\pgfpathlineto{\pgfqpoint{1.248986in}{0.674845in}}%
\pgfpathlineto{\pgfqpoint{1.241695in}{0.674845in}}%
\pgfpathlineto{\pgfqpoint{1.234403in}{0.674845in}}%
\pgfpathlineto{\pgfqpoint{1.227111in}{0.674845in}}%
\pgfpathlineto{\pgfqpoint{1.219820in}{0.674845in}}%
\pgfpathlineto{\pgfqpoint{1.212528in}{0.674845in}}%
\pgfpathlineto{\pgfqpoint{1.205236in}{0.674845in}}%
\pgfpathlineto{\pgfqpoint{1.197944in}{0.674845in}}%
\pgfpathlineto{\pgfqpoint{1.190653in}{0.674845in}}%
\pgfpathlineto{\pgfqpoint{1.183361in}{0.674845in}}%
\pgfpathlineto{\pgfqpoint{1.176069in}{0.674845in}}%
\pgfpathlineto{\pgfqpoint{1.168777in}{0.674845in}}%
\pgfpathlineto{\pgfqpoint{1.161486in}{0.674845in}}%
\pgfpathlineto{\pgfqpoint{1.154194in}{0.674845in}}%
\pgfpathlineto{\pgfqpoint{1.146902in}{0.674845in}}%
\pgfpathlineto{\pgfqpoint{1.139610in}{0.674845in}}%
\pgfpathlineto{\pgfqpoint{1.132319in}{0.674845in}}%
\pgfpathlineto{\pgfqpoint{1.125027in}{0.674845in}}%
\pgfpathlineto{\pgfqpoint{1.117735in}{0.674845in}}%
\pgfpathlineto{\pgfqpoint{1.110444in}{0.674845in}}%
\pgfpathlineto{\pgfqpoint{1.103152in}{0.674845in}}%
\pgfpathlineto{\pgfqpoint{1.095860in}{0.674845in}}%
\pgfpathlineto{\pgfqpoint{1.088568in}{0.674845in}}%
\pgfpathlineto{\pgfqpoint{1.081277in}{0.674845in}}%
\pgfpathlineto{\pgfqpoint{1.073985in}{0.674845in}}%
\pgfpathlineto{\pgfqpoint{1.066693in}{0.674845in}}%
\pgfpathlineto{\pgfqpoint{1.059401in}{0.674845in}}%
\pgfpathlineto{\pgfqpoint{1.052110in}{0.674845in}}%
\pgfpathlineto{\pgfqpoint{1.044818in}{0.674845in}}%
\pgfpathlineto{\pgfqpoint{1.037526in}{0.674845in}}%
\pgfpathlineto{\pgfqpoint{1.030234in}{0.674845in}}%
\pgfpathlineto{\pgfqpoint{1.022943in}{0.674845in}}%
\pgfpathlineto{\pgfqpoint{1.015651in}{0.674845in}}%
\pgfpathlineto{\pgfqpoint{1.008359in}{0.674845in}}%
\pgfpathlineto{\pgfqpoint{1.001067in}{0.674845in}}%
\pgfpathlineto{\pgfqpoint{0.993776in}{0.674845in}}%
\pgfpathlineto{\pgfqpoint{0.986484in}{0.674845in}}%
\pgfpathlineto{\pgfqpoint{0.979192in}{0.674845in}}%
\pgfpathlineto{\pgfqpoint{0.971901in}{0.674845in}}%
\pgfpathlineto{\pgfqpoint{0.964609in}{0.674845in}}%
\pgfpathlineto{\pgfqpoint{0.957317in}{0.674845in}}%
\pgfpathlineto{\pgfqpoint{0.950025in}{0.674845in}}%
\pgfpathlineto{\pgfqpoint{0.942734in}{0.674845in}}%
\pgfpathlineto{\pgfqpoint{0.935442in}{0.674845in}}%
\pgfpathlineto{\pgfqpoint{0.928150in}{0.674845in}}%
\pgfpathlineto{\pgfqpoint{0.920858in}{0.674845in}}%
\pgfpathlineto{\pgfqpoint{0.913567in}{0.674845in}}%
\pgfpathlineto{\pgfqpoint{0.906275in}{0.674845in}}%
\pgfpathlineto{\pgfqpoint{0.898983in}{0.674845in}}%
\pgfpathlineto{\pgfqpoint{0.891691in}{0.674845in}}%
\pgfpathlineto{\pgfqpoint{0.884400in}{0.674845in}}%
\pgfpathlineto{\pgfqpoint{0.877108in}{0.674845in}}%
\pgfpathlineto{\pgfqpoint{0.869816in}{0.674845in}}%
\pgfpathlineto{\pgfqpoint{0.862525in}{0.674845in}}%
\pgfpathlineto{\pgfqpoint{0.855233in}{0.674845in}}%
\pgfpathlineto{\pgfqpoint{0.847941in}{0.674845in}}%
\pgfpathlineto{\pgfqpoint{0.840649in}{0.674845in}}%
\pgfpathlineto{\pgfqpoint{0.833358in}{0.674845in}}%
\pgfpathlineto{\pgfqpoint{0.826066in}{0.674845in}}%
\pgfpathlineto{\pgfqpoint{0.818774in}{0.674845in}}%
\pgfpathlineto{\pgfqpoint{0.811482in}{0.674845in}}%
\pgfpathlineto{\pgfqpoint{0.804191in}{0.674845in}}%
\pgfpathlineto{\pgfqpoint{0.796899in}{0.674845in}}%
\pgfpathlineto{\pgfqpoint{0.789607in}{0.674845in}}%
\pgfpathlineto{\pgfqpoint{0.782315in}{0.674845in}}%
\pgfpathlineto{\pgfqpoint{0.775024in}{0.674845in}}%
\pgfpathlineto{\pgfqpoint{0.767732in}{0.674845in}}%
\pgfpathlineto{\pgfqpoint{0.760440in}{0.674845in}}%
\pgfpathlineto{\pgfqpoint{0.753149in}{0.674845in}}%
\pgfpathlineto{\pgfqpoint{0.745857in}{0.674845in}}%
\pgfpathlineto{\pgfqpoint{0.738565in}{0.674845in}}%
\pgfpathlineto{\pgfqpoint{0.731273in}{0.674845in}}%
\pgfpathlineto{\pgfqpoint{0.723982in}{0.674845in}}%
\pgfpathlineto{\pgfqpoint{0.716690in}{0.674845in}}%
\pgfpathlineto{\pgfqpoint{0.709398in}{0.674845in}}%
\pgfpathlineto{\pgfqpoint{0.702106in}{0.674845in}}%
\pgfpathlineto{\pgfqpoint{0.694815in}{0.674845in}}%
\pgfpathlineto{\pgfqpoint{0.687523in}{0.674845in}}%
\pgfpathlineto{\pgfqpoint{0.680231in}{0.674845in}}%
\pgfpathlineto{\pgfqpoint{0.672939in}{0.674845in}}%
\pgfpathlineto{\pgfqpoint{0.665648in}{0.674845in}}%
\pgfpathlineto{\pgfqpoint{0.658356in}{0.674845in}}%
\pgfpathlineto{\pgfqpoint{0.651064in}{0.674845in}}%
\pgfpathlineto{\pgfqpoint{0.643772in}{0.674845in}}%
\pgfpathlineto{\pgfqpoint{0.636481in}{0.674845in}}%
\pgfpathlineto{\pgfqpoint{0.629189in}{0.674845in}}%
\pgfpathlineto{\pgfqpoint{0.629189in}{0.674845in}}%
\pgfpathclose%
\pgfusepath{fill}%
\end{pgfscope}%
\begin{pgfscope}%
\pgfpathrectangle{\pgfqpoint{0.556636in}{0.524444in}}{\pgfqpoint{1.596161in}{3.308814in}}%
\pgfusepath{clip}%
\pgfsetbuttcap%
\pgfsetroundjoin%
\pgfsetlinewidth{0.501875pt}%
\definecolor{currentstroke}{rgb}{0.501961,0.501961,0.501961}%
\pgfsetstrokecolor{currentstroke}%
\pgfsetstrokeopacity{0.500000}%
\pgfsetdash{{1.850000pt}{0.800000pt}}{0.000000pt}%
\pgfpathmoveto{\pgfqpoint{0.629189in}{0.524444in}}%
\pgfpathlineto{\pgfqpoint{0.629189in}{3.833258in}}%
\pgfusepath{stroke}%
\end{pgfscope}%
\begin{pgfscope}%
\pgfsetbuttcap%
\pgfsetroundjoin%
\definecolor{currentfill}{rgb}{0.000000,0.000000,0.000000}%
\pgfsetfillcolor{currentfill}%
\pgfsetlinewidth{0.803000pt}%
\definecolor{currentstroke}{rgb}{0.000000,0.000000,0.000000}%
\pgfsetstrokecolor{currentstroke}%
\pgfsetdash{}{0pt}%
\pgfsys@defobject{currentmarker}{\pgfqpoint{0.000000in}{-0.048611in}}{\pgfqpoint{0.000000in}{0.000000in}}{%
\pgfpathmoveto{\pgfqpoint{0.000000in}{0.000000in}}%
\pgfpathlineto{\pgfqpoint{0.000000in}{-0.048611in}}%
\pgfusepath{stroke,fill}%
}%
\begin{pgfscope}%
\pgfsys@transformshift{0.629189in}{0.524444in}%
\pgfsys@useobject{currentmarker}{}%
\end{pgfscope}%
\end{pgfscope}%
\begin{pgfscope}%
\definecolor{textcolor}{rgb}{0.000000,0.000000,0.000000}%
\pgfsetstrokecolor{textcolor}%
\pgfsetfillcolor{textcolor}%
\pgftext[x=0.629189in,y=0.427222in,,top]{\color{textcolor}{\rmfamily\fontsize{10.000000}{12.000000}\selectfont\catcode`\^=\active\def^{\ifmmode\sp\else\^{}\fi}\catcode`\%=\active\def%{\%}$\mathdefault{0}$}}%
\end{pgfscope}%
\begin{pgfscope}%
\pgfpathrectangle{\pgfqpoint{0.556636in}{0.524444in}}{\pgfqpoint{1.596161in}{3.308814in}}%
\pgfusepath{clip}%
\pgfsetbuttcap%
\pgfsetroundjoin%
\pgfsetlinewidth{0.501875pt}%
\definecolor{currentstroke}{rgb}{0.501961,0.501961,0.501961}%
\pgfsetstrokecolor{currentstroke}%
\pgfsetstrokeopacity{0.500000}%
\pgfsetdash{{1.850000pt}{0.800000pt}}{0.000000pt}%
\pgfpathmoveto{\pgfqpoint{1.354717in}{0.524444in}}%
\pgfpathlineto{\pgfqpoint{1.354717in}{3.833258in}}%
\pgfusepath{stroke}%
\end{pgfscope}%
\begin{pgfscope}%
\pgfsetbuttcap%
\pgfsetroundjoin%
\definecolor{currentfill}{rgb}{0.000000,0.000000,0.000000}%
\pgfsetfillcolor{currentfill}%
\pgfsetlinewidth{0.803000pt}%
\definecolor{currentstroke}{rgb}{0.000000,0.000000,0.000000}%
\pgfsetstrokecolor{currentstroke}%
\pgfsetdash{}{0pt}%
\pgfsys@defobject{currentmarker}{\pgfqpoint{0.000000in}{-0.048611in}}{\pgfqpoint{0.000000in}{0.000000in}}{%
\pgfpathmoveto{\pgfqpoint{0.000000in}{0.000000in}}%
\pgfpathlineto{\pgfqpoint{0.000000in}{-0.048611in}}%
\pgfusepath{stroke,fill}%
}%
\begin{pgfscope}%
\pgfsys@transformshift{1.354717in}{0.524444in}%
\pgfsys@useobject{currentmarker}{}%
\end{pgfscope}%
\end{pgfscope}%
\begin{pgfscope}%
\definecolor{textcolor}{rgb}{0.000000,0.000000,0.000000}%
\pgfsetstrokecolor{textcolor}%
\pgfsetfillcolor{textcolor}%
\pgftext[x=1.354717in,y=0.427222in,,top]{\color{textcolor}{\rmfamily\fontsize{10.000000}{12.000000}\selectfont\catcode`\^=\active\def^{\ifmmode\sp\else\^{}\fi}\catcode`\%=\active\def%{\%}$\mathdefault{5}$}}%
\end{pgfscope}%
\begin{pgfscope}%
\pgfpathrectangle{\pgfqpoint{0.556636in}{0.524444in}}{\pgfqpoint{1.596161in}{3.308814in}}%
\pgfusepath{clip}%
\pgfsetbuttcap%
\pgfsetroundjoin%
\pgfsetlinewidth{0.501875pt}%
\definecolor{currentstroke}{rgb}{0.501961,0.501961,0.501961}%
\pgfsetstrokecolor{currentstroke}%
\pgfsetstrokeopacity{0.500000}%
\pgfsetdash{{1.850000pt}{0.800000pt}}{0.000000pt}%
\pgfpathmoveto{\pgfqpoint{2.080244in}{0.524444in}}%
\pgfpathlineto{\pgfqpoint{2.080244in}{3.833258in}}%
\pgfusepath{stroke}%
\end{pgfscope}%
\begin{pgfscope}%
\pgfsetbuttcap%
\pgfsetroundjoin%
\definecolor{currentfill}{rgb}{0.000000,0.000000,0.000000}%
\pgfsetfillcolor{currentfill}%
\pgfsetlinewidth{0.803000pt}%
\definecolor{currentstroke}{rgb}{0.000000,0.000000,0.000000}%
\pgfsetstrokecolor{currentstroke}%
\pgfsetdash{}{0pt}%
\pgfsys@defobject{currentmarker}{\pgfqpoint{0.000000in}{-0.048611in}}{\pgfqpoint{0.000000in}{0.000000in}}{%
\pgfpathmoveto{\pgfqpoint{0.000000in}{0.000000in}}%
\pgfpathlineto{\pgfqpoint{0.000000in}{-0.048611in}}%
\pgfusepath{stroke,fill}%
}%
\begin{pgfscope}%
\pgfsys@transformshift{2.080244in}{0.524444in}%
\pgfsys@useobject{currentmarker}{}%
\end{pgfscope}%
\end{pgfscope}%
\begin{pgfscope}%
\definecolor{textcolor}{rgb}{0.000000,0.000000,0.000000}%
\pgfsetstrokecolor{textcolor}%
\pgfsetfillcolor{textcolor}%
\pgftext[x=2.080244in,y=0.427222in,,top]{\color{textcolor}{\rmfamily\fontsize{10.000000}{12.000000}\selectfont\catcode`\^=\active\def^{\ifmmode\sp\else\^{}\fi}\catcode`\%=\active\def%{\%}$\mathdefault{10}$}}%
\end{pgfscope}%
\begin{pgfscope}%
\definecolor{textcolor}{rgb}{0.000000,0.000000,0.000000}%
\pgfsetstrokecolor{textcolor}%
\pgfsetfillcolor{textcolor}%
\pgftext[x=1.354717in,y=0.248333in,,top]{\color{textcolor}{\rmfamily\fontsize{10.000000}{12.000000}\selectfont\catcode`\^=\active\def^{\ifmmode\sp\else\^{}\fi}\catcode`\%=\active\def%{\%}Tiempo}}%
\end{pgfscope}%
\begin{pgfscope}%
\pgfpathrectangle{\pgfqpoint{0.556636in}{0.524444in}}{\pgfqpoint{1.596161in}{3.308814in}}%
\pgfusepath{clip}%
\pgfsetbuttcap%
\pgfsetroundjoin%
\pgfsetlinewidth{0.501875pt}%
\definecolor{currentstroke}{rgb}{0.501961,0.501961,0.501961}%
\pgfsetstrokecolor{currentstroke}%
\pgfsetstrokeopacity{0.500000}%
\pgfsetdash{{1.850000pt}{0.800000pt}}{0.000000pt}%
\pgfpathmoveto{\pgfqpoint{0.556636in}{0.674845in}}%
\pgfpathlineto{\pgfqpoint{2.152797in}{0.674845in}}%
\pgfusepath{stroke}%
\end{pgfscope}%
\begin{pgfscope}%
\pgfsetbuttcap%
\pgfsetroundjoin%
\definecolor{currentfill}{rgb}{0.000000,0.000000,0.000000}%
\pgfsetfillcolor{currentfill}%
\pgfsetlinewidth{0.803000pt}%
\definecolor{currentstroke}{rgb}{0.000000,0.000000,0.000000}%
\pgfsetstrokecolor{currentstroke}%
\pgfsetdash{}{0pt}%
\pgfsys@defobject{currentmarker}{\pgfqpoint{-0.048611in}{0.000000in}}{\pgfqpoint{-0.000000in}{0.000000in}}{%
\pgfpathmoveto{\pgfqpoint{-0.000000in}{0.000000in}}%
\pgfpathlineto{\pgfqpoint{-0.048611in}{0.000000in}}%
\pgfusepath{stroke,fill}%
}%
\begin{pgfscope}%
\pgfsys@transformshift{0.556636in}{0.674845in}%
\pgfsys@useobject{currentmarker}{}%
\end{pgfscope}%
\end{pgfscope}%
\begin{pgfscope}%
\definecolor{textcolor}{rgb}{0.000000,0.000000,0.000000}%
\pgfsetstrokecolor{textcolor}%
\pgfsetfillcolor{textcolor}%
\pgftext[x=0.281944in, y=0.626650in, left, base]{\color{textcolor}{\rmfamily\fontsize{10.000000}{12.000000}\selectfont\catcode`\^=\active\def^{\ifmmode\sp\else\^{}\fi}\catcode`\%=\active\def%{\%}$\mathdefault{0.0}$}}%
\end{pgfscope}%
\begin{pgfscope}%
\pgfpathrectangle{\pgfqpoint{0.556636in}{0.524444in}}{\pgfqpoint{1.596161in}{3.308814in}}%
\pgfusepath{clip}%
\pgfsetbuttcap%
\pgfsetroundjoin%
\pgfsetlinewidth{0.501875pt}%
\definecolor{currentstroke}{rgb}{0.501961,0.501961,0.501961}%
\pgfsetstrokecolor{currentstroke}%
\pgfsetstrokeopacity{0.500000}%
\pgfsetdash{{1.850000pt}{0.800000pt}}{0.000000pt}%
\pgfpathmoveto{\pgfqpoint{0.556636in}{1.276468in}}%
\pgfpathlineto{\pgfqpoint{2.152797in}{1.276468in}}%
\pgfusepath{stroke}%
\end{pgfscope}%
\begin{pgfscope}%
\pgfsetbuttcap%
\pgfsetroundjoin%
\definecolor{currentfill}{rgb}{0.000000,0.000000,0.000000}%
\pgfsetfillcolor{currentfill}%
\pgfsetlinewidth{0.803000pt}%
\definecolor{currentstroke}{rgb}{0.000000,0.000000,0.000000}%
\pgfsetstrokecolor{currentstroke}%
\pgfsetdash{}{0pt}%
\pgfsys@defobject{currentmarker}{\pgfqpoint{-0.048611in}{0.000000in}}{\pgfqpoint{-0.000000in}{0.000000in}}{%
\pgfpathmoveto{\pgfqpoint{-0.000000in}{0.000000in}}%
\pgfpathlineto{\pgfqpoint{-0.048611in}{0.000000in}}%
\pgfusepath{stroke,fill}%
}%
\begin{pgfscope}%
\pgfsys@transformshift{0.556636in}{1.276468in}%
\pgfsys@useobject{currentmarker}{}%
\end{pgfscope}%
\end{pgfscope}%
\begin{pgfscope}%
\definecolor{textcolor}{rgb}{0.000000,0.000000,0.000000}%
\pgfsetstrokecolor{textcolor}%
\pgfsetfillcolor{textcolor}%
\pgftext[x=0.281944in, y=1.228273in, left, base]{\color{textcolor}{\rmfamily\fontsize{10.000000}{12.000000}\selectfont\catcode`\^=\active\def^{\ifmmode\sp\else\^{}\fi}\catcode`\%=\active\def%{\%}$\mathdefault{0.5}$}}%
\end{pgfscope}%
\begin{pgfscope}%
\pgfpathrectangle{\pgfqpoint{0.556636in}{0.524444in}}{\pgfqpoint{1.596161in}{3.308814in}}%
\pgfusepath{clip}%
\pgfsetbuttcap%
\pgfsetroundjoin%
\pgfsetlinewidth{0.501875pt}%
\definecolor{currentstroke}{rgb}{0.501961,0.501961,0.501961}%
\pgfsetstrokecolor{currentstroke}%
\pgfsetstrokeopacity{0.500000}%
\pgfsetdash{{1.850000pt}{0.800000pt}}{0.000000pt}%
\pgfpathmoveto{\pgfqpoint{0.556636in}{1.878090in}}%
\pgfpathlineto{\pgfqpoint{2.152797in}{1.878090in}}%
\pgfusepath{stroke}%
\end{pgfscope}%
\begin{pgfscope}%
\pgfsetbuttcap%
\pgfsetroundjoin%
\definecolor{currentfill}{rgb}{0.000000,0.000000,0.000000}%
\pgfsetfillcolor{currentfill}%
\pgfsetlinewidth{0.803000pt}%
\definecolor{currentstroke}{rgb}{0.000000,0.000000,0.000000}%
\pgfsetstrokecolor{currentstroke}%
\pgfsetdash{}{0pt}%
\pgfsys@defobject{currentmarker}{\pgfqpoint{-0.048611in}{0.000000in}}{\pgfqpoint{-0.000000in}{0.000000in}}{%
\pgfpathmoveto{\pgfqpoint{-0.000000in}{0.000000in}}%
\pgfpathlineto{\pgfqpoint{-0.048611in}{0.000000in}}%
\pgfusepath{stroke,fill}%
}%
\begin{pgfscope}%
\pgfsys@transformshift{0.556636in}{1.878090in}%
\pgfsys@useobject{currentmarker}{}%
\end{pgfscope}%
\end{pgfscope}%
\begin{pgfscope}%
\definecolor{textcolor}{rgb}{0.000000,0.000000,0.000000}%
\pgfsetstrokecolor{textcolor}%
\pgfsetfillcolor{textcolor}%
\pgftext[x=0.281944in, y=1.829896in, left, base]{\color{textcolor}{\rmfamily\fontsize{10.000000}{12.000000}\selectfont\catcode`\^=\active\def^{\ifmmode\sp\else\^{}\fi}\catcode`\%=\active\def%{\%}$\mathdefault{1.0}$}}%
\end{pgfscope}%
\begin{pgfscope}%
\pgfpathrectangle{\pgfqpoint{0.556636in}{0.524444in}}{\pgfqpoint{1.596161in}{3.308814in}}%
\pgfusepath{clip}%
\pgfsetbuttcap%
\pgfsetroundjoin%
\pgfsetlinewidth{0.501875pt}%
\definecolor{currentstroke}{rgb}{0.501961,0.501961,0.501961}%
\pgfsetstrokecolor{currentstroke}%
\pgfsetstrokeopacity{0.500000}%
\pgfsetdash{{1.850000pt}{0.800000pt}}{0.000000pt}%
\pgfpathmoveto{\pgfqpoint{0.556636in}{2.479713in}}%
\pgfpathlineto{\pgfqpoint{2.152797in}{2.479713in}}%
\pgfusepath{stroke}%
\end{pgfscope}%
\begin{pgfscope}%
\pgfsetbuttcap%
\pgfsetroundjoin%
\definecolor{currentfill}{rgb}{0.000000,0.000000,0.000000}%
\pgfsetfillcolor{currentfill}%
\pgfsetlinewidth{0.803000pt}%
\definecolor{currentstroke}{rgb}{0.000000,0.000000,0.000000}%
\pgfsetstrokecolor{currentstroke}%
\pgfsetdash{}{0pt}%
\pgfsys@defobject{currentmarker}{\pgfqpoint{-0.048611in}{0.000000in}}{\pgfqpoint{-0.000000in}{0.000000in}}{%
\pgfpathmoveto{\pgfqpoint{-0.000000in}{0.000000in}}%
\pgfpathlineto{\pgfqpoint{-0.048611in}{0.000000in}}%
\pgfusepath{stroke,fill}%
}%
\begin{pgfscope}%
\pgfsys@transformshift{0.556636in}{2.479713in}%
\pgfsys@useobject{currentmarker}{}%
\end{pgfscope}%
\end{pgfscope}%
\begin{pgfscope}%
\definecolor{textcolor}{rgb}{0.000000,0.000000,0.000000}%
\pgfsetstrokecolor{textcolor}%
\pgfsetfillcolor{textcolor}%
\pgftext[x=0.281944in, y=2.431519in, left, base]{\color{textcolor}{\rmfamily\fontsize{10.000000}{12.000000}\selectfont\catcode`\^=\active\def^{\ifmmode\sp\else\^{}\fi}\catcode`\%=\active\def%{\%}$\mathdefault{1.5}$}}%
\end{pgfscope}%
\begin{pgfscope}%
\pgfpathrectangle{\pgfqpoint{0.556636in}{0.524444in}}{\pgfqpoint{1.596161in}{3.308814in}}%
\pgfusepath{clip}%
\pgfsetbuttcap%
\pgfsetroundjoin%
\pgfsetlinewidth{0.501875pt}%
\definecolor{currentstroke}{rgb}{0.501961,0.501961,0.501961}%
\pgfsetstrokecolor{currentstroke}%
\pgfsetstrokeopacity{0.500000}%
\pgfsetdash{{1.850000pt}{0.800000pt}}{0.000000pt}%
\pgfpathmoveto{\pgfqpoint{0.556636in}{3.081336in}}%
\pgfpathlineto{\pgfqpoint{2.152797in}{3.081336in}}%
\pgfusepath{stroke}%
\end{pgfscope}%
\begin{pgfscope}%
\pgfsetbuttcap%
\pgfsetroundjoin%
\definecolor{currentfill}{rgb}{0.000000,0.000000,0.000000}%
\pgfsetfillcolor{currentfill}%
\pgfsetlinewidth{0.803000pt}%
\definecolor{currentstroke}{rgb}{0.000000,0.000000,0.000000}%
\pgfsetstrokecolor{currentstroke}%
\pgfsetdash{}{0pt}%
\pgfsys@defobject{currentmarker}{\pgfqpoint{-0.048611in}{0.000000in}}{\pgfqpoint{-0.000000in}{0.000000in}}{%
\pgfpathmoveto{\pgfqpoint{-0.000000in}{0.000000in}}%
\pgfpathlineto{\pgfqpoint{-0.048611in}{0.000000in}}%
\pgfusepath{stroke,fill}%
}%
\begin{pgfscope}%
\pgfsys@transformshift{0.556636in}{3.081336in}%
\pgfsys@useobject{currentmarker}{}%
\end{pgfscope}%
\end{pgfscope}%
\begin{pgfscope}%
\definecolor{textcolor}{rgb}{0.000000,0.000000,0.000000}%
\pgfsetstrokecolor{textcolor}%
\pgfsetfillcolor{textcolor}%
\pgftext[x=0.281944in, y=3.033142in, left, base]{\color{textcolor}{\rmfamily\fontsize{10.000000}{12.000000}\selectfont\catcode`\^=\active\def^{\ifmmode\sp\else\^{}\fi}\catcode`\%=\active\def%{\%}$\mathdefault{2.0}$}}%
\end{pgfscope}%
\begin{pgfscope}%
\pgfpathrectangle{\pgfqpoint{0.556636in}{0.524444in}}{\pgfqpoint{1.596161in}{3.308814in}}%
\pgfusepath{clip}%
\pgfsetbuttcap%
\pgfsetroundjoin%
\pgfsetlinewidth{0.501875pt}%
\definecolor{currentstroke}{rgb}{0.501961,0.501961,0.501961}%
\pgfsetstrokecolor{currentstroke}%
\pgfsetstrokeopacity{0.500000}%
\pgfsetdash{{1.850000pt}{0.800000pt}}{0.000000pt}%
\pgfpathmoveto{\pgfqpoint{0.556636in}{3.682959in}}%
\pgfpathlineto{\pgfqpoint{2.152797in}{3.682959in}}%
\pgfusepath{stroke}%
\end{pgfscope}%
\begin{pgfscope}%
\pgfsetbuttcap%
\pgfsetroundjoin%
\definecolor{currentfill}{rgb}{0.000000,0.000000,0.000000}%
\pgfsetfillcolor{currentfill}%
\pgfsetlinewidth{0.803000pt}%
\definecolor{currentstroke}{rgb}{0.000000,0.000000,0.000000}%
\pgfsetstrokecolor{currentstroke}%
\pgfsetdash{}{0pt}%
\pgfsys@defobject{currentmarker}{\pgfqpoint{-0.048611in}{0.000000in}}{\pgfqpoint{-0.000000in}{0.000000in}}{%
\pgfpathmoveto{\pgfqpoint{-0.000000in}{0.000000in}}%
\pgfpathlineto{\pgfqpoint{-0.048611in}{0.000000in}}%
\pgfusepath{stroke,fill}%
}%
\begin{pgfscope}%
\pgfsys@transformshift{0.556636in}{3.682959in}%
\pgfsys@useobject{currentmarker}{}%
\end{pgfscope}%
\end{pgfscope}%
\begin{pgfscope}%
\definecolor{textcolor}{rgb}{0.000000,0.000000,0.000000}%
\pgfsetstrokecolor{textcolor}%
\pgfsetfillcolor{textcolor}%
\pgftext[x=0.281944in, y=3.634765in, left, base]{\color{textcolor}{\rmfamily\fontsize{10.000000}{12.000000}\selectfont\catcode`\^=\active\def^{\ifmmode\sp\else\^{}\fi}\catcode`\%=\active\def%{\%}$\mathdefault{2.5}$}}%
\end{pgfscope}%
\begin{pgfscope}%
\definecolor{textcolor}{rgb}{0.000000,0.000000,0.000000}%
\pgfsetstrokecolor{textcolor}%
\pgfsetfillcolor{textcolor}%
\pgftext[x=0.226389in,y=2.178851in,,bottom,rotate=90.000000]{\color{textcolor}{\rmfamily\fontsize{10.000000}{12.000000}\selectfont\catcode`\^=\active\def^{\ifmmode\sp\else\^{}\fi}\catcode`\%=\active\def%{\%}Amplitud}}%
\end{pgfscope}%
\begin{pgfscope}%
\pgfpathrectangle{\pgfqpoint{0.556636in}{0.524444in}}{\pgfqpoint{1.596161in}{3.308814in}}%
\pgfusepath{clip}%
\pgfsetrectcap%
\pgfsetroundjoin%
\pgfsetlinewidth{1.505625pt}%
\definecolor{currentstroke}{rgb}{0.121569,0.466667,0.705882}%
\pgfsetstrokecolor{currentstroke}%
\pgfsetdash{}{0pt}%
\pgfpathmoveto{\pgfqpoint{0.629189in}{2.479713in}}%
\pgfpathlineto{\pgfqpoint{0.680231in}{2.894291in}}%
\pgfpathlineto{\pgfqpoint{0.709398in}{3.111467in}}%
\pgfpathlineto{\pgfqpoint{0.731273in}{3.258098in}}%
\pgfpathlineto{\pgfqpoint{0.753149in}{3.387072in}}%
\pgfpathlineto{\pgfqpoint{0.767732in}{3.461782in}}%
\pgfpathlineto{\pgfqpoint{0.782315in}{3.526581in}}%
\pgfpathlineto{\pgfqpoint{0.796899in}{3.580815in}}%
\pgfpathlineto{\pgfqpoint{0.811482in}{3.623936in}}%
\pgfpathlineto{\pgfqpoint{0.818774in}{3.641189in}}%
\pgfpathlineto{\pgfqpoint{0.826066in}{3.655509in}}%
\pgfpathlineto{\pgfqpoint{0.833358in}{3.666861in}}%
\pgfpathlineto{\pgfqpoint{0.840649in}{3.675216in}}%
\pgfpathlineto{\pgfqpoint{0.847941in}{3.680553in}}%
\pgfpathlineto{\pgfqpoint{0.855233in}{3.682857in}}%
\pgfpathlineto{\pgfqpoint{0.862525in}{3.682125in}}%
\pgfpathlineto{\pgfqpoint{0.869816in}{3.678357in}}%
\pgfpathlineto{\pgfqpoint{0.877108in}{3.671562in}}%
\pgfpathlineto{\pgfqpoint{0.884400in}{3.661759in}}%
\pgfpathlineto{\pgfqpoint{0.891691in}{3.648971in}}%
\pgfpathlineto{\pgfqpoint{0.898983in}{3.633231in}}%
\pgfpathlineto{\pgfqpoint{0.913567in}{3.593062in}}%
\pgfpathlineto{\pgfqpoint{0.928150in}{3.541657in}}%
\pgfpathlineto{\pgfqpoint{0.942734in}{3.479534in}}%
\pgfpathlineto{\pgfqpoint{0.957317in}{3.407321in}}%
\pgfpathlineto{\pgfqpoint{0.971901in}{3.325746in}}%
\pgfpathlineto{\pgfqpoint{0.993776in}{3.187656in}}%
\pgfpathlineto{\pgfqpoint{1.015651in}{3.033507in}}%
\pgfpathlineto{\pgfqpoint{1.044818in}{2.809080in}}%
\pgfpathlineto{\pgfqpoint{1.088568in}{2.450554in}}%
\pgfpathlineto{\pgfqpoint{1.132319in}{2.094658in}}%
\pgfpathlineto{\pgfqpoint{1.161486in}{1.874811in}}%
\pgfpathlineto{\pgfqpoint{1.183361in}{1.725460in}}%
\pgfpathlineto{\pgfqpoint{1.205236in}{1.593218in}}%
\pgfpathlineto{\pgfqpoint{1.219820in}{1.516061in}}%
\pgfpathlineto{\pgfqpoint{1.234403in}{1.448630in}}%
\pgfpathlineto{\pgfqpoint{1.248986in}{1.391604in}}%
\pgfpathlineto{\pgfqpoint{1.263570in}{1.345560in}}%
\pgfpathlineto{\pgfqpoint{1.278153in}{1.310963in}}%
\pgfpathlineto{\pgfqpoint{1.285445in}{1.298070in}}%
\pgfpathlineto{\pgfqpoint{1.292737in}{1.288160in}}%
\pgfpathlineto{\pgfqpoint{1.300029in}{1.281259in}}%
\pgfpathlineto{\pgfqpoint{1.307320in}{1.277384in}}%
\pgfpathlineto{\pgfqpoint{1.314612in}{1.276543in}}%
\pgfpathlineto{\pgfqpoint{1.321904in}{1.278741in}}%
\pgfpathlineto{\pgfqpoint{1.329196in}{1.283970in}}%
\pgfpathlineto{\pgfqpoint{1.336487in}{1.292219in}}%
\pgfpathlineto{\pgfqpoint{1.343779in}{1.303465in}}%
\pgfpathlineto{\pgfqpoint{1.351071in}{1.317681in}}%
\pgfpathlineto{\pgfqpoint{1.358363in}{1.334831in}}%
\pgfpathlineto{\pgfqpoint{1.372946in}{1.377751in}}%
\pgfpathlineto{\pgfqpoint{1.387529in}{1.431792in}}%
\pgfpathlineto{\pgfqpoint{1.402113in}{1.496410in}}%
\pgfpathlineto{\pgfqpoint{1.416696in}{1.570951in}}%
\pgfpathlineto{\pgfqpoint{1.431280in}{1.654663in}}%
\pgfpathlineto{\pgfqpoint{1.453155in}{1.795557in}}%
\pgfpathlineto{\pgfqpoint{1.475030in}{1.951971in}}%
\pgfpathlineto{\pgfqpoint{1.504197in}{2.178489in}}%
\pgfpathlineto{\pgfqpoint{1.562531in}{2.658305in}}%
\pgfpathlineto{\pgfqpoint{1.598990in}{2.948536in}}%
\pgfpathlineto{\pgfqpoint{1.620865in}{3.109645in}}%
\pgfpathlineto{\pgfqpoint{1.642740in}{3.256465in}}%
\pgfpathlineto{\pgfqpoint{1.664615in}{3.385665in}}%
\pgfpathlineto{\pgfqpoint{1.679199in}{3.460544in}}%
\pgfpathlineto{\pgfqpoint{1.693782in}{3.525525in}}%
\pgfpathlineto{\pgfqpoint{1.708366in}{3.579950in}}%
\pgfpathlineto{\pgfqpoint{1.722949in}{3.623272in}}%
\pgfpathlineto{\pgfqpoint{1.730241in}{3.640628in}}%
\pgfpathlineto{\pgfqpoint{1.737533in}{3.655053in}}%
\pgfpathlineto{\pgfqpoint{1.744824in}{3.666510in}}%
\pgfpathlineto{\pgfqpoint{1.752116in}{3.674972in}}%
\pgfpathlineto{\pgfqpoint{1.759408in}{3.680415in}}%
\pgfpathlineto{\pgfqpoint{1.766700in}{3.682828in}}%
\pgfpathlineto{\pgfqpoint{1.773991in}{3.682203in}}%
\pgfpathlineto{\pgfqpoint{1.781283in}{3.678542in}}%
\pgfpathlineto{\pgfqpoint{1.788575in}{3.671854in}}%
\pgfpathlineto{\pgfqpoint{1.795867in}{3.662157in}}%
\pgfpathlineto{\pgfqpoint{1.803158in}{3.649474in}}%
\pgfpathlineto{\pgfqpoint{1.810450in}{3.633838in}}%
\pgfpathlineto{\pgfqpoint{1.825034in}{3.593872in}}%
\pgfpathlineto{\pgfqpoint{1.839617in}{3.542661in}}%
\pgfpathlineto{\pgfqpoint{1.854200in}{3.480723in}}%
\pgfpathlineto{\pgfqpoint{1.868784in}{3.408682in}}%
\pgfpathlineto{\pgfqpoint{1.883367in}{3.327266in}}%
\pgfpathlineto{\pgfqpoint{1.905243in}{3.189385in}}%
\pgfpathlineto{\pgfqpoint{1.927118in}{3.035405in}}%
\pgfpathlineto{\pgfqpoint{1.956285in}{2.811138in}}%
\pgfpathlineto{\pgfqpoint{2.000035in}{2.452693in}}%
\pgfpathlineto{\pgfqpoint{2.043786in}{2.096686in}}%
\pgfpathlineto{\pgfqpoint{2.072953in}{1.876661in}}%
\pgfpathlineto{\pgfqpoint{2.080244in}{1.825122in}}%
\pgfpathlineto{\pgfqpoint{2.080244in}{1.825122in}}%
\pgfusepath{stroke}%
\end{pgfscope}%
\begin{pgfscope}%
\pgfsetrectcap%
\pgfsetmiterjoin%
\pgfsetlinewidth{0.803000pt}%
\definecolor{currentstroke}{rgb}{0.000000,0.000000,0.000000}%
\pgfsetstrokecolor{currentstroke}%
\pgfsetdash{}{0pt}%
\pgfpathmoveto{\pgfqpoint{0.556636in}{0.524444in}}%
\pgfpathlineto{\pgfqpoint{0.556636in}{3.833258in}}%
\pgfusepath{stroke}%
\end{pgfscope}%
\begin{pgfscope}%
\pgfsetrectcap%
\pgfsetmiterjoin%
\pgfsetlinewidth{0.803000pt}%
\definecolor{currentstroke}{rgb}{0.000000,0.000000,0.000000}%
\pgfsetstrokecolor{currentstroke}%
\pgfsetdash{}{0pt}%
\pgfpathmoveto{\pgfqpoint{2.152797in}{0.524444in}}%
\pgfpathlineto{\pgfqpoint{2.152797in}{3.833258in}}%
\pgfusepath{stroke}%
\end{pgfscope}%
\begin{pgfscope}%
\pgfsetrectcap%
\pgfsetmiterjoin%
\pgfsetlinewidth{0.803000pt}%
\definecolor{currentstroke}{rgb}{0.000000,0.000000,0.000000}%
\pgfsetstrokecolor{currentstroke}%
\pgfsetdash{}{0pt}%
\pgfpathmoveto{\pgfqpoint{0.556636in}{0.524444in}}%
\pgfpathlineto{\pgfqpoint{2.152797in}{0.524444in}}%
\pgfusepath{stroke}%
\end{pgfscope}%
\begin{pgfscope}%
\pgfsetrectcap%
\pgfsetmiterjoin%
\pgfsetlinewidth{0.803000pt}%
\definecolor{currentstroke}{rgb}{0.000000,0.000000,0.000000}%
\pgfsetstrokecolor{currentstroke}%
\pgfsetdash{}{0pt}%
\pgfpathmoveto{\pgfqpoint{0.556636in}{3.833258in}}%
\pgfpathlineto{\pgfqpoint{2.152797in}{3.833258in}}%
\pgfusepath{stroke}%
\end{pgfscope}%
\begin{pgfscope}%
\definecolor{textcolor}{rgb}{0.000000,0.000000,0.000000}%
\pgfsetstrokecolor{textcolor}%
\pgfsetfillcolor{textcolor}%
\pgftext[x=1.354717in,y=3.916591in,,base]{\color{textcolor}{\rmfamily\fontsize{12.000000}{14.400000}\selectfont\catcode`\^=\active\def^{\ifmmode\sp\else\^{}\fi}\catcode`\%=\active\def%{\%}Área bajo la curva}}%
\end{pgfscope}%
\begin{pgfscope}%
\pgfsetbuttcap%
\pgfsetmiterjoin%
\definecolor{currentfill}{rgb}{1.000000,1.000000,1.000000}%
\pgfsetfillcolor{currentfill}%
\pgfsetfillopacity{0.800000}%
\pgfsetlinewidth{1.003750pt}%
\definecolor{currentstroke}{rgb}{0.800000,0.800000,0.800000}%
\pgfsetstrokecolor{currentstroke}%
\pgfsetstrokeopacity{0.800000}%
\pgfsetdash{}{0pt}%
\pgfpathmoveto{\pgfqpoint{0.653858in}{2.061212in}}%
\pgfpathlineto{\pgfqpoint{1.399692in}{2.061212in}}%
\pgfpathquadraticcurveto{\pgfqpoint{1.427470in}{2.061212in}}{\pgfqpoint{1.427470in}{2.088990in}}%
\pgfpathlineto{\pgfqpoint{1.427470in}{2.268712in}}%
\pgfpathquadraticcurveto{\pgfqpoint{1.427470in}{2.296490in}}{\pgfqpoint{1.399692in}{2.296490in}}%
\pgfpathlineto{\pgfqpoint{0.653858in}{2.296490in}}%
\pgfpathquadraticcurveto{\pgfqpoint{0.626081in}{2.296490in}}{\pgfqpoint{0.626081in}{2.268712in}}%
\pgfpathlineto{\pgfqpoint{0.626081in}{2.088990in}}%
\pgfpathquadraticcurveto{\pgfqpoint{0.626081in}{2.061212in}}{\pgfqpoint{0.653858in}{2.061212in}}%
\pgfpathlineto{\pgfqpoint{0.653858in}{2.061212in}}%
\pgfpathclose%
\pgfusepath{stroke,fill}%
\end{pgfscope}%
\begin{pgfscope}%
\pgfsetrectcap%
\pgfsetroundjoin%
\pgfsetlinewidth{1.505625pt}%
\definecolor{currentstroke}{rgb}{0.121569,0.466667,0.705882}%
\pgfsetstrokecolor{currentstroke}%
\pgfsetdash{}{0pt}%
\pgfpathmoveto{\pgfqpoint{0.681636in}{2.192323in}}%
\pgfpathlineto{\pgfqpoint{0.820525in}{2.192323in}}%
\pgfpathlineto{\pgfqpoint{0.959414in}{2.192323in}}%
\pgfusepath{stroke}%
\end{pgfscope}%
\begin{pgfscope}%
\definecolor{textcolor}{rgb}{0.000000,0.000000,0.000000}%
\pgfsetstrokecolor{textcolor}%
\pgfsetfillcolor{textcolor}%
\pgftext[x=1.070525in,y=2.143712in,left,base]{\color{textcolor}{\rmfamily\fontsize{10.000000}{12.000000}\selectfont\catcode`\^=\active\def^{\ifmmode\sp\else\^{}\fi}\catcode`\%=\active\def%{\%}señal}}%
\end{pgfscope}%
\begin{pgfscope}%
\pgfsetbuttcap%
\pgfsetmiterjoin%
\definecolor{currentfill}{rgb}{1.000000,1.000000,1.000000}%
\pgfsetfillcolor{currentfill}%
\pgfsetlinewidth{0.000000pt}%
\definecolor{currentstroke}{rgb}{0.000000,0.000000,0.000000}%
\pgfsetstrokecolor{currentstroke}%
\pgfsetstrokeopacity{0.000000}%
\pgfsetdash{}{0pt}%
\pgfpathmoveto{\pgfqpoint{2.952340in}{0.524444in}}%
\pgfpathlineto{\pgfqpoint{4.548501in}{0.524444in}}%
\pgfpathlineto{\pgfqpoint{4.548501in}{3.833258in}}%
\pgfpathlineto{\pgfqpoint{2.952340in}{3.833258in}}%
\pgfpathlineto{\pgfqpoint{2.952340in}{0.524444in}}%
\pgfpathclose%
\pgfusepath{fill}%
\end{pgfscope}%
\begin{pgfscope}%
\pgfpathrectangle{\pgfqpoint{2.952340in}{0.524444in}}{\pgfqpoint{1.596161in}{3.308814in}}%
\pgfusepath{clip}%
\pgfsetbuttcap%
\pgfsetroundjoin%
\definecolor{currentfill}{rgb}{0.121569,0.466667,0.705882}%
\pgfsetfillcolor{currentfill}%
\pgfsetfillopacity{0.250000}%
\pgfsetlinewidth{0.000000pt}%
\definecolor{currentstroke}{rgb}{0.000000,0.000000,0.000000}%
\pgfsetstrokecolor{currentstroke}%
\pgfsetdash{}{0pt}%
\pgfpathmoveto{\pgfqpoint{3.209308in}{3.228282in}}%
\pgfpathlineto{\pgfqpoint{3.209308in}{3.256217in}}%
\pgfpathlineto{\pgfqpoint{3.214161in}{3.278030in}}%
\pgfpathlineto{\pgfqpoint{3.219014in}{3.299358in}}%
\pgfpathlineto{\pgfqpoint{3.223867in}{3.320192in}}%
\pgfpathlineto{\pgfqpoint{3.228720in}{3.340521in}}%
\pgfpathlineto{\pgfqpoint{3.233573in}{3.360337in}}%
\pgfpathlineto{\pgfqpoint{3.238426in}{3.379632in}}%
\pgfpathlineto{\pgfqpoint{3.243279in}{3.398396in}}%
\pgfpathlineto{\pgfqpoint{3.248132in}{3.416622in}}%
\pgfpathlineto{\pgfqpoint{3.252985in}{3.434302in}}%
\pgfpathlineto{\pgfqpoint{3.257838in}{3.451426in}}%
\pgfpathlineto{\pgfqpoint{3.262691in}{3.467989in}}%
\pgfpathlineto{\pgfqpoint{3.267544in}{3.483983in}}%
\pgfpathlineto{\pgfqpoint{3.272397in}{3.499401in}}%
\pgfpathlineto{\pgfqpoint{3.277250in}{3.514235in}}%
\pgfpathlineto{\pgfqpoint{3.282103in}{3.528480in}}%
\pgfpathlineto{\pgfqpoint{3.286956in}{3.542128in}}%
\pgfpathlineto{\pgfqpoint{3.291809in}{3.555175in}}%
\pgfpathlineto{\pgfqpoint{3.296662in}{3.567614in}}%
\pgfpathlineto{\pgfqpoint{3.301515in}{3.579440in}}%
\pgfpathlineto{\pgfqpoint{3.306368in}{3.590647in}}%
\pgfpathlineto{\pgfqpoint{3.311221in}{3.601231in}}%
\pgfpathlineto{\pgfqpoint{3.316074in}{3.611186in}}%
\pgfpathlineto{\pgfqpoint{3.320927in}{3.620510in}}%
\pgfpathlineto{\pgfqpoint{3.325780in}{3.629196in}}%
\pgfpathlineto{\pgfqpoint{3.330633in}{3.637242in}}%
\pgfpathlineto{\pgfqpoint{3.335486in}{3.644645in}}%
\pgfpathlineto{\pgfqpoint{3.340340in}{3.651400in}}%
\pgfpathlineto{\pgfqpoint{3.345193in}{3.657504in}}%
\pgfpathlineto{\pgfqpoint{3.350046in}{3.662956in}}%
\pgfpathlineto{\pgfqpoint{3.354899in}{3.667752in}}%
\pgfpathlineto{\pgfqpoint{3.359752in}{3.671891in}}%
\pgfpathlineto{\pgfqpoint{3.364605in}{3.675371in}}%
\pgfpathlineto{\pgfqpoint{3.369458in}{3.678190in}}%
\pgfpathlineto{\pgfqpoint{3.374311in}{3.680347in}}%
\pgfpathlineto{\pgfqpoint{3.379164in}{3.681841in}}%
\pgfpathlineto{\pgfqpoint{3.384017in}{3.682671in}}%
\pgfpathlineto{\pgfqpoint{3.388870in}{3.682837in}}%
\pgfpathlineto{\pgfqpoint{3.393723in}{3.682339in}}%
\pgfpathlineto{\pgfqpoint{3.398576in}{3.681177in}}%
\pgfpathlineto{\pgfqpoint{3.403429in}{3.679351in}}%
\pgfpathlineto{\pgfqpoint{3.408282in}{3.676863in}}%
\pgfpathlineto{\pgfqpoint{3.413135in}{3.673714in}}%
\pgfpathlineto{\pgfqpoint{3.417988in}{3.669904in}}%
\pgfpathlineto{\pgfqpoint{3.422841in}{3.665436in}}%
\pgfpathlineto{\pgfqpoint{3.427694in}{3.660312in}}%
\pgfpathlineto{\pgfqpoint{3.432547in}{3.654533in}}%
\pgfpathlineto{\pgfqpoint{3.437400in}{3.648103in}}%
\pgfpathlineto{\pgfqpoint{3.442253in}{3.641024in}}%
\pgfpathlineto{\pgfqpoint{3.447106in}{3.633300in}}%
\pgfpathlineto{\pgfqpoint{3.451959in}{3.624933in}}%
\pgfpathlineto{\pgfqpoint{3.456812in}{3.615927in}}%
\pgfpathlineto{\pgfqpoint{3.461665in}{3.606287in}}%
\pgfpathlineto{\pgfqpoint{3.466518in}{3.596017in}}%
\pgfpathlineto{\pgfqpoint{3.471371in}{3.585121in}}%
\pgfpathlineto{\pgfqpoint{3.476224in}{3.573604in}}%
\pgfpathlineto{\pgfqpoint{3.481077in}{3.561471in}}%
\pgfpathlineto{\pgfqpoint{3.485930in}{3.548727in}}%
\pgfpathlineto{\pgfqpoint{3.490783in}{3.535379in}}%
\pgfpathlineto{\pgfqpoint{3.495636in}{3.521431in}}%
\pgfpathlineto{\pgfqpoint{3.500489in}{3.506891in}}%
\pgfpathlineto{\pgfqpoint{3.505342in}{3.491764in}}%
\pgfpathlineto{\pgfqpoint{3.510195in}{3.476058in}}%
\pgfpathlineto{\pgfqpoint{3.515049in}{3.459779in}}%
\pgfpathlineto{\pgfqpoint{3.519902in}{3.442934in}}%
\pgfpathlineto{\pgfqpoint{3.524755in}{3.425531in}}%
\pgfpathlineto{\pgfqpoint{3.529608in}{3.407577in}}%
\pgfpathlineto{\pgfqpoint{3.534461in}{3.389081in}}%
\pgfpathlineto{\pgfqpoint{3.539314in}{3.370050in}}%
\pgfpathlineto{\pgfqpoint{3.544167in}{3.350494in}}%
\pgfpathlineto{\pgfqpoint{3.549020in}{3.330420in}}%
\pgfpathlineto{\pgfqpoint{3.553873in}{3.309837in}}%
\pgfpathlineto{\pgfqpoint{3.558726in}{3.288756in}}%
\pgfpathlineto{\pgfqpoint{3.563579in}{3.267184in}}%
\pgfpathlineto{\pgfqpoint{3.568432in}{3.245131in}}%
\pgfpathlineto{\pgfqpoint{3.573285in}{3.222608in}}%
\pgfpathlineto{\pgfqpoint{3.578138in}{3.199623in}}%
\pgfpathlineto{\pgfqpoint{3.582991in}{3.176188in}}%
\pgfpathlineto{\pgfqpoint{3.587844in}{3.152313in}}%
\pgfpathlineto{\pgfqpoint{3.592697in}{3.128008in}}%
\pgfpathlineto{\pgfqpoint{3.597550in}{3.103283in}}%
\pgfpathlineto{\pgfqpoint{3.602403in}{3.078150in}}%
\pgfpathlineto{\pgfqpoint{3.607256in}{3.052621in}}%
\pgfpathlineto{\pgfqpoint{3.612109in}{3.026705in}}%
\pgfpathlineto{\pgfqpoint{3.616962in}{3.000415in}}%
\pgfpathlineto{\pgfqpoint{3.621815in}{2.973762in}}%
\pgfpathlineto{\pgfqpoint{3.626668in}{2.946759in}}%
\pgfpathlineto{\pgfqpoint{3.631521in}{2.919416in}}%
\pgfpathlineto{\pgfqpoint{3.636374in}{2.891746in}}%
\pgfpathlineto{\pgfqpoint{3.641227in}{2.863761in}}%
\pgfpathlineto{\pgfqpoint{3.646080in}{2.835474in}}%
\pgfpathlineto{\pgfqpoint{3.650933in}{2.806897in}}%
\pgfpathlineto{\pgfqpoint{3.655786in}{2.778043in}}%
\pgfpathlineto{\pgfqpoint{3.660639in}{2.748924in}}%
\pgfpathlineto{\pgfqpoint{3.665492in}{2.719553in}}%
\pgfpathlineto{\pgfqpoint{3.670345in}{2.689943in}}%
\pgfpathlineto{\pgfqpoint{3.675198in}{2.660108in}}%
\pgfpathlineto{\pgfqpoint{3.680051in}{2.630061in}}%
\pgfpathlineto{\pgfqpoint{3.684904in}{2.599814in}}%
\pgfpathlineto{\pgfqpoint{3.689758in}{2.569381in}}%
\pgfpathlineto{\pgfqpoint{3.694611in}{2.538775in}}%
\pgfpathlineto{\pgfqpoint{3.699464in}{2.508011in}}%
\pgfpathlineto{\pgfqpoint{3.704317in}{2.477102in}}%
\pgfpathlineto{\pgfqpoint{3.709170in}{2.446060in}}%
\pgfpathlineto{\pgfqpoint{3.714023in}{2.414901in}}%
\pgfpathlineto{\pgfqpoint{3.718876in}{2.383638in}}%
\pgfpathlineto{\pgfqpoint{3.723729in}{2.352284in}}%
\pgfpathlineto{\pgfqpoint{3.728582in}{2.320853in}}%
\pgfpathlineto{\pgfqpoint{3.733435in}{2.289360in}}%
\pgfpathlineto{\pgfqpoint{3.738288in}{2.257818in}}%
\pgfpathlineto{\pgfqpoint{3.743141in}{2.226241in}}%
\pgfpathlineto{\pgfqpoint{3.747994in}{2.194643in}}%
\pgfpathlineto{\pgfqpoint{3.752847in}{2.163038in}}%
\pgfpathlineto{\pgfqpoint{3.757700in}{2.131440in}}%
\pgfpathlineto{\pgfqpoint{3.762553in}{2.099864in}}%
\pgfpathlineto{\pgfqpoint{3.767406in}{2.068321in}}%
\pgfpathlineto{\pgfqpoint{3.772259in}{2.036828in}}%
\pgfpathlineto{\pgfqpoint{3.777112in}{2.005398in}}%
\pgfpathlineto{\pgfqpoint{3.781965in}{1.974044in}}%
\pgfpathlineto{\pgfqpoint{3.786818in}{1.942780in}}%
\pgfpathlineto{\pgfqpoint{3.791671in}{1.911621in}}%
\pgfpathlineto{\pgfqpoint{3.796524in}{1.880580in}}%
\pgfpathlineto{\pgfqpoint{3.801377in}{1.849670in}}%
\pgfpathlineto{\pgfqpoint{3.806230in}{1.818906in}}%
\pgfpathlineto{\pgfqpoint{3.811083in}{1.788301in}}%
\pgfpathlineto{\pgfqpoint{3.815936in}{1.757868in}}%
\pgfpathlineto{\pgfqpoint{3.820789in}{1.727621in}}%
\pgfpathlineto{\pgfqpoint{3.825642in}{1.697573in}}%
\pgfpathlineto{\pgfqpoint{3.830495in}{1.667738in}}%
\pgfpathlineto{\pgfqpoint{3.835348in}{1.638128in}}%
\pgfpathlineto{\pgfqpoint{3.840201in}{1.608758in}}%
\pgfpathlineto{\pgfqpoint{3.845054in}{1.579639in}}%
\pgfpathlineto{\pgfqpoint{3.849907in}{1.550784in}}%
\pgfpathlineto{\pgfqpoint{3.854760in}{1.522207in}}%
\pgfpathlineto{\pgfqpoint{3.859613in}{1.493920in}}%
\pgfpathlineto{\pgfqpoint{3.864467in}{1.465935in}}%
\pgfpathlineto{\pgfqpoint{3.869320in}{1.438266in}}%
\pgfpathlineto{\pgfqpoint{3.874173in}{1.410923in}}%
\pgfpathlineto{\pgfqpoint{3.879026in}{1.383919in}}%
\pgfpathlineto{\pgfqpoint{3.883879in}{1.357266in}}%
\pgfpathlineto{\pgfqpoint{3.888732in}{1.330976in}}%
\pgfpathlineto{\pgfqpoint{3.893585in}{1.305061in}}%
\pgfpathlineto{\pgfqpoint{3.898438in}{1.279531in}}%
\pgfpathlineto{\pgfqpoint{3.903291in}{1.254398in}}%
\pgfpathlineto{\pgfqpoint{3.908144in}{1.229674in}}%
\pgfpathlineto{\pgfqpoint{3.912997in}{1.205369in}}%
\pgfpathlineto{\pgfqpoint{3.917850in}{1.181493in}}%
\pgfpathlineto{\pgfqpoint{3.922703in}{1.158058in}}%
\pgfpathlineto{\pgfqpoint{3.927556in}{1.135074in}}%
\pgfpathlineto{\pgfqpoint{3.927556in}{1.095966in}}%
\pgfpathlineto{\pgfqpoint{3.927556in}{1.095966in}}%
\pgfpathlineto{\pgfqpoint{3.922703in}{1.074273in}}%
\pgfpathlineto{\pgfqpoint{3.917850in}{1.053068in}}%
\pgfpathlineto{\pgfqpoint{3.912997in}{1.032360in}}%
\pgfpathlineto{\pgfqpoint{3.908144in}{1.012158in}}%
\pgfpathlineto{\pgfqpoint{3.903291in}{0.992471in}}%
\pgfpathlineto{\pgfqpoint{3.898438in}{0.973308in}}%
\pgfpathlineto{\pgfqpoint{3.893585in}{0.954678in}}%
\pgfpathlineto{\pgfqpoint{3.888732in}{0.936588in}}%
\pgfpathlineto{\pgfqpoint{3.883879in}{0.919046in}}%
\pgfpathlineto{\pgfqpoint{3.879026in}{0.902061in}}%
\pgfpathlineto{\pgfqpoint{3.874173in}{0.885640in}}%
\pgfpathlineto{\pgfqpoint{3.869320in}{0.869790in}}%
\pgfpathlineto{\pgfqpoint{3.864467in}{0.854517in}}%
\pgfpathlineto{\pgfqpoint{3.859613in}{0.839830in}}%
\pgfpathlineto{\pgfqpoint{3.854760in}{0.825734in}}%
\pgfpathlineto{\pgfqpoint{3.849907in}{0.812235in}}%
\pgfpathlineto{\pgfqpoint{3.845054in}{0.799340in}}%
\pgfpathlineto{\pgfqpoint{3.840201in}{0.787053in}}%
\pgfpathlineto{\pgfqpoint{3.835348in}{0.775382in}}%
\pgfpathlineto{\pgfqpoint{3.830495in}{0.764330in}}%
\pgfpathlineto{\pgfqpoint{3.825642in}{0.753903in}}%
\pgfpathlineto{\pgfqpoint{3.820789in}{0.744105in}}%
\pgfpathlineto{\pgfqpoint{3.815936in}{0.734940in}}%
\pgfpathlineto{\pgfqpoint{3.811083in}{0.726413in}}%
\pgfpathlineto{\pgfqpoint{3.806230in}{0.718528in}}%
\pgfpathlineto{\pgfqpoint{3.801377in}{0.711287in}}%
\pgfpathlineto{\pgfqpoint{3.796524in}{0.704695in}}%
\pgfpathlineto{\pgfqpoint{3.791671in}{0.698753in}}%
\pgfpathlineto{\pgfqpoint{3.786818in}{0.693465in}}%
\pgfpathlineto{\pgfqpoint{3.781965in}{0.688832in}}%
\pgfpathlineto{\pgfqpoint{3.777112in}{0.684858in}}%
\pgfpathlineto{\pgfqpoint{3.772259in}{0.681543in}}%
\pgfpathlineto{\pgfqpoint{3.767406in}{0.678890in}}%
\pgfpathlineto{\pgfqpoint{3.762553in}{0.676899in}}%
\pgfpathlineto{\pgfqpoint{3.757700in}{0.675571in}}%
\pgfpathlineto{\pgfqpoint{3.752847in}{0.674907in}}%
\pgfpathlineto{\pgfqpoint{3.747994in}{0.674907in}}%
\pgfpathlineto{\pgfqpoint{3.743141in}{0.675571in}}%
\pgfpathlineto{\pgfqpoint{3.738288in}{0.676899in}}%
\pgfpathlineto{\pgfqpoint{3.733435in}{0.678890in}}%
\pgfpathlineto{\pgfqpoint{3.728582in}{0.681543in}}%
\pgfpathlineto{\pgfqpoint{3.723729in}{0.684858in}}%
\pgfpathlineto{\pgfqpoint{3.718876in}{0.688832in}}%
\pgfpathlineto{\pgfqpoint{3.714023in}{0.693465in}}%
\pgfpathlineto{\pgfqpoint{3.709170in}{0.698753in}}%
\pgfpathlineto{\pgfqpoint{3.704317in}{0.704695in}}%
\pgfpathlineto{\pgfqpoint{3.699464in}{0.711287in}}%
\pgfpathlineto{\pgfqpoint{3.694611in}{0.718528in}}%
\pgfpathlineto{\pgfqpoint{3.689758in}{0.726413in}}%
\pgfpathlineto{\pgfqpoint{3.684904in}{0.734940in}}%
\pgfpathlineto{\pgfqpoint{3.680051in}{0.744105in}}%
\pgfpathlineto{\pgfqpoint{3.675198in}{0.753903in}}%
\pgfpathlineto{\pgfqpoint{3.670345in}{0.764330in}}%
\pgfpathlineto{\pgfqpoint{3.665492in}{0.775382in}}%
\pgfpathlineto{\pgfqpoint{3.660639in}{0.787053in}}%
\pgfpathlineto{\pgfqpoint{3.655786in}{0.799340in}}%
\pgfpathlineto{\pgfqpoint{3.650933in}{0.812235in}}%
\pgfpathlineto{\pgfqpoint{3.646080in}{0.825734in}}%
\pgfpathlineto{\pgfqpoint{3.641227in}{0.839830in}}%
\pgfpathlineto{\pgfqpoint{3.636374in}{0.854517in}}%
\pgfpathlineto{\pgfqpoint{3.631521in}{0.869790in}}%
\pgfpathlineto{\pgfqpoint{3.626668in}{0.885640in}}%
\pgfpathlineto{\pgfqpoint{3.621815in}{0.902061in}}%
\pgfpathlineto{\pgfqpoint{3.616962in}{0.919046in}}%
\pgfpathlineto{\pgfqpoint{3.612109in}{0.936588in}}%
\pgfpathlineto{\pgfqpoint{3.607256in}{0.954678in}}%
\pgfpathlineto{\pgfqpoint{3.602403in}{0.973308in}}%
\pgfpathlineto{\pgfqpoint{3.597550in}{0.992471in}}%
\pgfpathlineto{\pgfqpoint{3.592697in}{1.012158in}}%
\pgfpathlineto{\pgfqpoint{3.587844in}{1.032360in}}%
\pgfpathlineto{\pgfqpoint{3.582991in}{1.053068in}}%
\pgfpathlineto{\pgfqpoint{3.578138in}{1.074273in}}%
\pgfpathlineto{\pgfqpoint{3.573285in}{1.095966in}}%
\pgfpathlineto{\pgfqpoint{3.568432in}{1.118137in}}%
\pgfpathlineto{\pgfqpoint{3.563579in}{1.140777in}}%
\pgfpathlineto{\pgfqpoint{3.558726in}{1.163875in}}%
\pgfpathlineto{\pgfqpoint{3.553873in}{1.187421in}}%
\pgfpathlineto{\pgfqpoint{3.549020in}{1.211405in}}%
\pgfpathlineto{\pgfqpoint{3.544167in}{1.235816in}}%
\pgfpathlineto{\pgfqpoint{3.539314in}{1.260644in}}%
\pgfpathlineto{\pgfqpoint{3.534461in}{1.285877in}}%
\pgfpathlineto{\pgfqpoint{3.529608in}{1.311504in}}%
\pgfpathlineto{\pgfqpoint{3.524755in}{1.337514in}}%
\pgfpathlineto{\pgfqpoint{3.519902in}{1.363896in}}%
\pgfpathlineto{\pgfqpoint{3.515049in}{1.390637in}}%
\pgfpathlineto{\pgfqpoint{3.510195in}{1.417727in}}%
\pgfpathlineto{\pgfqpoint{3.505342in}{1.445153in}}%
\pgfpathlineto{\pgfqpoint{3.500489in}{1.472903in}}%
\pgfpathlineto{\pgfqpoint{3.495636in}{1.500964in}}%
\pgfpathlineto{\pgfqpoint{3.490783in}{1.529325in}}%
\pgfpathlineto{\pgfqpoint{3.485930in}{1.557972in}}%
\pgfpathlineto{\pgfqpoint{3.481077in}{1.586894in}}%
\pgfpathlineto{\pgfqpoint{3.476224in}{1.616077in}}%
\pgfpathlineto{\pgfqpoint{3.471371in}{1.645509in}}%
\pgfpathlineto{\pgfqpoint{3.466518in}{1.675176in}}%
\pgfpathlineto{\pgfqpoint{3.461665in}{1.705066in}}%
\pgfpathlineto{\pgfqpoint{3.456812in}{1.735164in}}%
\pgfpathlineto{\pgfqpoint{3.451959in}{1.765459in}}%
\pgfpathlineto{\pgfqpoint{3.447106in}{1.795936in}}%
\pgfpathlineto{\pgfqpoint{3.442253in}{1.826583in}}%
\pgfpathlineto{\pgfqpoint{3.437400in}{1.857384in}}%
\pgfpathlineto{\pgfqpoint{3.432547in}{1.888328in}}%
\pgfpathlineto{\pgfqpoint{3.427694in}{1.919400in}}%
\pgfpathlineto{\pgfqpoint{3.422841in}{1.950587in}}%
\pgfpathlineto{\pgfqpoint{3.417988in}{1.981874in}}%
\pgfpathlineto{\pgfqpoint{3.413135in}{2.013249in}}%
\pgfpathlineto{\pgfqpoint{3.408282in}{2.044696in}}%
\pgfpathlineto{\pgfqpoint{3.403429in}{2.076203in}}%
\pgfpathlineto{\pgfqpoint{3.398576in}{2.107755in}}%
\pgfpathlineto{\pgfqpoint{3.393723in}{2.139338in}}%
\pgfpathlineto{\pgfqpoint{3.388870in}{2.170939in}}%
\pgfpathlineto{\pgfqpoint{3.384017in}{2.202544in}}%
\pgfpathlineto{\pgfqpoint{3.379164in}{2.234138in}}%
\pgfpathlineto{\pgfqpoint{3.374311in}{2.265707in}}%
\pgfpathlineto{\pgfqpoint{3.369458in}{2.297238in}}%
\pgfpathlineto{\pgfqpoint{3.364605in}{2.328717in}}%
\pgfpathlineto{\pgfqpoint{3.359752in}{2.360130in}}%
\pgfpathlineto{\pgfqpoint{3.354899in}{2.391462in}}%
\pgfpathlineto{\pgfqpoint{3.350046in}{2.422701in}}%
\pgfpathlineto{\pgfqpoint{3.345193in}{2.453832in}}%
\pgfpathlineto{\pgfqpoint{3.340340in}{2.484842in}}%
\pgfpathlineto{\pgfqpoint{3.335486in}{2.515716in}}%
\pgfpathlineto{\pgfqpoint{3.330633in}{2.546442in}}%
\pgfpathlineto{\pgfqpoint{3.325780in}{2.577006in}}%
\pgfpathlineto{\pgfqpoint{3.320927in}{2.607393in}}%
\pgfpathlineto{\pgfqpoint{3.316074in}{2.637592in}}%
\pgfpathlineto{\pgfqpoint{3.311221in}{2.667587in}}%
\pgfpathlineto{\pgfqpoint{3.306368in}{2.697367in}}%
\pgfpathlineto{\pgfqpoint{3.301515in}{2.726918in}}%
\pgfpathlineto{\pgfqpoint{3.296662in}{2.756227in}}%
\pgfpathlineto{\pgfqpoint{3.291809in}{2.785282in}}%
\pgfpathlineto{\pgfqpoint{3.286956in}{2.814068in}}%
\pgfpathlineto{\pgfqpoint{3.282103in}{2.842574in}}%
\pgfpathlineto{\pgfqpoint{3.277250in}{2.870786in}}%
\pgfpathlineto{\pgfqpoint{3.272397in}{2.898693in}}%
\pgfpathlineto{\pgfqpoint{3.267544in}{2.926283in}}%
\pgfpathlineto{\pgfqpoint{3.262691in}{2.953542in}}%
\pgfpathlineto{\pgfqpoint{3.257838in}{2.980459in}}%
\pgfpathlineto{\pgfqpoint{3.252985in}{3.007022in}}%
\pgfpathlineto{\pgfqpoint{3.248132in}{3.033220in}}%
\pgfpathlineto{\pgfqpoint{3.243279in}{3.059040in}}%
\pgfpathlineto{\pgfqpoint{3.238426in}{3.084471in}}%
\pgfpathlineto{\pgfqpoint{3.233573in}{3.109503in}}%
\pgfpathlineto{\pgfqpoint{3.228720in}{3.134124in}}%
\pgfpathlineto{\pgfqpoint{3.223867in}{3.158322in}}%
\pgfpathlineto{\pgfqpoint{3.219014in}{3.182089in}}%
\pgfpathlineto{\pgfqpoint{3.214161in}{3.205412in}}%
\pgfpathlineto{\pgfqpoint{3.209308in}{3.228282in}}%
\pgfpathlineto{\pgfqpoint{3.209308in}{3.228282in}}%
\pgfpathclose%
\pgfusepath{fill}%
\end{pgfscope}%
\begin{pgfscope}%
\pgfpathrectangle{\pgfqpoint{2.952340in}{0.524444in}}{\pgfqpoint{1.596161in}{3.308814in}}%
\pgfusepath{clip}%
\pgfsetbuttcap%
\pgfsetroundjoin%
\definecolor{currentfill}{rgb}{1.000000,0.498039,0.054902}%
\pgfsetfillcolor{currentfill}%
\pgfsetfillopacity{0.250000}%
\pgfsetlinewidth{0.000000pt}%
\definecolor{currentstroke}{rgb}{0.000000,0.000000,0.000000}%
\pgfsetstrokecolor{currentstroke}%
\pgfsetdash{}{0pt}%
\pgfpathmoveto{\pgfqpoint{3.024893in}{3.682857in}}%
\pgfpathlineto{\pgfqpoint{3.024893in}{2.178841in}}%
\pgfpathlineto{\pgfqpoint{3.029746in}{2.210444in}}%
\pgfpathlineto{\pgfqpoint{3.034599in}{2.242033in}}%
\pgfpathlineto{\pgfqpoint{3.039452in}{2.273594in}}%
\pgfpathlineto{\pgfqpoint{3.044305in}{2.305113in}}%
\pgfpathlineto{\pgfqpoint{3.049158in}{2.336577in}}%
\pgfpathlineto{\pgfqpoint{3.054011in}{2.367971in}}%
\pgfpathlineto{\pgfqpoint{3.058864in}{2.399282in}}%
\pgfpathlineto{\pgfqpoint{3.063717in}{2.430495in}}%
\pgfpathlineto{\pgfqpoint{3.068570in}{2.461597in}}%
\pgfpathlineto{\pgfqpoint{3.073423in}{2.492574in}}%
\pgfpathlineto{\pgfqpoint{3.078276in}{2.523412in}}%
\pgfpathlineto{\pgfqpoint{3.083129in}{2.554099in}}%
\pgfpathlineto{\pgfqpoint{3.087982in}{2.584620in}}%
\pgfpathlineto{\pgfqpoint{3.092835in}{2.614961in}}%
\pgfpathlineto{\pgfqpoint{3.097688in}{2.645110in}}%
\pgfpathlineto{\pgfqpoint{3.102541in}{2.675053in}}%
\pgfpathlineto{\pgfqpoint{3.107394in}{2.704777in}}%
\pgfpathlineto{\pgfqpoint{3.112247in}{2.734269in}}%
\pgfpathlineto{\pgfqpoint{3.117100in}{2.763515in}}%
\pgfpathlineto{\pgfqpoint{3.121953in}{2.792504in}}%
\pgfpathlineto{\pgfqpoint{3.126806in}{2.821221in}}%
\pgfpathlineto{\pgfqpoint{3.131659in}{2.849655in}}%
\pgfpathlineto{\pgfqpoint{3.136512in}{2.877792in}}%
\pgfpathlineto{\pgfqpoint{3.141365in}{2.905621in}}%
\pgfpathlineto{\pgfqpoint{3.146218in}{2.933129in}}%
\pgfpathlineto{\pgfqpoint{3.151071in}{2.960304in}}%
\pgfpathlineto{\pgfqpoint{3.155924in}{2.987133in}}%
\pgfpathlineto{\pgfqpoint{3.160778in}{3.013606in}}%
\pgfpathlineto{\pgfqpoint{3.165631in}{3.039710in}}%
\pgfpathlineto{\pgfqpoint{3.170484in}{3.065435in}}%
\pgfpathlineto{\pgfqpoint{3.175337in}{3.090767in}}%
\pgfpathlineto{\pgfqpoint{3.180190in}{3.115697in}}%
\pgfpathlineto{\pgfqpoint{3.185043in}{3.140213in}}%
\pgfpathlineto{\pgfqpoint{3.189896in}{3.164305in}}%
\pgfpathlineto{\pgfqpoint{3.194749in}{3.187961in}}%
\pgfpathlineto{\pgfqpoint{3.199602in}{3.211172in}}%
\pgfpathlineto{\pgfqpoint{3.204455in}{3.233928in}}%
\pgfpathlineto{\pgfqpoint{3.204455in}{3.250689in}}%
\pgfpathlineto{\pgfqpoint{3.204455in}{3.250689in}}%
\pgfpathlineto{\pgfqpoint{3.199602in}{3.272622in}}%
\pgfpathlineto{\pgfqpoint{3.194749in}{3.294072in}}%
\pgfpathlineto{\pgfqpoint{3.189896in}{3.315030in}}%
\pgfpathlineto{\pgfqpoint{3.185043in}{3.335486in}}%
\pgfpathlineto{\pgfqpoint{3.180190in}{3.355432in}}%
\pgfpathlineto{\pgfqpoint{3.175337in}{3.374858in}}%
\pgfpathlineto{\pgfqpoint{3.170484in}{3.393755in}}%
\pgfpathlineto{\pgfqpoint{3.165631in}{3.412117in}}%
\pgfpathlineto{\pgfqpoint{3.160778in}{3.429933in}}%
\pgfpathlineto{\pgfqpoint{3.155924in}{3.447198in}}%
\pgfpathlineto{\pgfqpoint{3.151071in}{3.463902in}}%
\pgfpathlineto{\pgfqpoint{3.146218in}{3.480039in}}%
\pgfpathlineto{\pgfqpoint{3.141365in}{3.495601in}}%
\pgfpathlineto{\pgfqpoint{3.136512in}{3.510581in}}%
\pgfpathlineto{\pgfqpoint{3.131659in}{3.524974in}}%
\pgfpathlineto{\pgfqpoint{3.126806in}{3.538772in}}%
\pgfpathlineto{\pgfqpoint{3.121953in}{3.551970in}}%
\pgfpathlineto{\pgfqpoint{3.117100in}{3.564561in}}%
\pgfpathlineto{\pgfqpoint{3.112247in}{3.576541in}}%
\pgfpathlineto{\pgfqpoint{3.107394in}{3.587903in}}%
\pgfpathlineto{\pgfqpoint{3.102541in}{3.598643in}}%
\pgfpathlineto{\pgfqpoint{3.097688in}{3.608757in}}%
\pgfpathlineto{\pgfqpoint{3.092835in}{3.618238in}}%
\pgfpathlineto{\pgfqpoint{3.087982in}{3.627084in}}%
\pgfpathlineto{\pgfqpoint{3.083129in}{3.635291in}}%
\pgfpathlineto{\pgfqpoint{3.078276in}{3.642855in}}%
\pgfpathlineto{\pgfqpoint{3.073423in}{3.649772in}}%
\pgfpathlineto{\pgfqpoint{3.068570in}{3.656039in}}%
\pgfpathlineto{\pgfqpoint{3.063717in}{3.661654in}}%
\pgfpathlineto{\pgfqpoint{3.058864in}{3.666615in}}%
\pgfpathlineto{\pgfqpoint{3.054011in}{3.670918in}}%
\pgfpathlineto{\pgfqpoint{3.049158in}{3.674563in}}%
\pgfpathlineto{\pgfqpoint{3.044305in}{3.677547in}}%
\pgfpathlineto{\pgfqpoint{3.039452in}{3.679870in}}%
\pgfpathlineto{\pgfqpoint{3.034599in}{3.681529in}}%
\pgfpathlineto{\pgfqpoint{3.029746in}{3.682525in}}%
\pgfpathlineto{\pgfqpoint{3.024893in}{3.682857in}}%
\pgfpathlineto{\pgfqpoint{3.024893in}{3.682857in}}%
\pgfpathclose%
\pgfusepath{fill}%
\end{pgfscope}%
\begin{pgfscope}%
\pgfpathrectangle{\pgfqpoint{2.952340in}{0.524444in}}{\pgfqpoint{1.596161in}{3.308814in}}%
\pgfusepath{clip}%
\pgfsetbuttcap%
\pgfsetroundjoin%
\definecolor{currentfill}{rgb}{1.000000,0.498039,0.054902}%
\pgfsetfillcolor{currentfill}%
\pgfsetfillopacity{0.250000}%
\pgfsetlinewidth{0.000000pt}%
\definecolor{currentstroke}{rgb}{0.000000,0.000000,0.000000}%
\pgfsetstrokecolor{currentstroke}%
\pgfsetdash{}{0pt}%
\pgfpathmoveto{\pgfqpoint{3.932409in}{1.118137in}}%
\pgfpathlineto{\pgfqpoint{3.932409in}{1.112550in}}%
\pgfpathlineto{\pgfqpoint{3.937262in}{1.090498in}}%
\pgfpathlineto{\pgfqpoint{3.942115in}{1.068926in}}%
\pgfpathlineto{\pgfqpoint{3.946968in}{1.047844in}}%
\pgfpathlineto{\pgfqpoint{3.951821in}{1.027262in}}%
\pgfpathlineto{\pgfqpoint{3.956674in}{1.007188in}}%
\pgfpathlineto{\pgfqpoint{3.961527in}{0.987631in}}%
\pgfpathlineto{\pgfqpoint{3.966380in}{0.968600in}}%
\pgfpathlineto{\pgfqpoint{3.971233in}{0.950104in}}%
\pgfpathlineto{\pgfqpoint{3.976086in}{0.932151in}}%
\pgfpathlineto{\pgfqpoint{3.980939in}{0.914748in}}%
\pgfpathlineto{\pgfqpoint{3.985792in}{0.897903in}}%
\pgfpathlineto{\pgfqpoint{3.990645in}{0.881623in}}%
\pgfpathlineto{\pgfqpoint{3.995498in}{0.865917in}}%
\pgfpathlineto{\pgfqpoint{4.000351in}{0.850790in}}%
\pgfpathlineto{\pgfqpoint{4.005204in}{0.836250in}}%
\pgfpathlineto{\pgfqpoint{4.010057in}{0.822303in}}%
\pgfpathlineto{\pgfqpoint{4.014910in}{0.808954in}}%
\pgfpathlineto{\pgfqpoint{4.019763in}{0.796211in}}%
\pgfpathlineto{\pgfqpoint{4.024616in}{0.784078in}}%
\pgfpathlineto{\pgfqpoint{4.029469in}{0.772561in}}%
\pgfpathlineto{\pgfqpoint{4.034322in}{0.761664in}}%
\pgfpathlineto{\pgfqpoint{4.039176in}{0.751394in}}%
\pgfpathlineto{\pgfqpoint{4.044029in}{0.741754in}}%
\pgfpathlineto{\pgfqpoint{4.048882in}{0.732749in}}%
\pgfpathlineto{\pgfqpoint{4.053735in}{0.724382in}}%
\pgfpathlineto{\pgfqpoint{4.058588in}{0.716657in}}%
\pgfpathlineto{\pgfqpoint{4.063441in}{0.709578in}}%
\pgfpathlineto{\pgfqpoint{4.068294in}{0.703148in}}%
\pgfpathlineto{\pgfqpoint{4.073147in}{0.697369in}}%
\pgfpathlineto{\pgfqpoint{4.078000in}{0.692245in}}%
\pgfpathlineto{\pgfqpoint{4.082853in}{0.687777in}}%
\pgfpathlineto{\pgfqpoint{4.087706in}{0.683968in}}%
\pgfpathlineto{\pgfqpoint{4.092559in}{0.680818in}}%
\pgfpathlineto{\pgfqpoint{4.097412in}{0.678330in}}%
\pgfpathlineto{\pgfqpoint{4.102265in}{0.676505in}}%
\pgfpathlineto{\pgfqpoint{4.107118in}{0.675343in}}%
\pgfpathlineto{\pgfqpoint{4.111971in}{0.674845in}}%
\pgfpathlineto{\pgfqpoint{4.116824in}{0.675011in}}%
\pgfpathlineto{\pgfqpoint{4.121677in}{0.675841in}}%
\pgfpathlineto{\pgfqpoint{4.126530in}{0.677335in}}%
\pgfpathlineto{\pgfqpoint{4.131383in}{0.679491in}}%
\pgfpathlineto{\pgfqpoint{4.136236in}{0.682310in}}%
\pgfpathlineto{\pgfqpoint{4.141089in}{0.685790in}}%
\pgfpathlineto{\pgfqpoint{4.145942in}{0.689929in}}%
\pgfpathlineto{\pgfqpoint{4.150795in}{0.694725in}}%
\pgfpathlineto{\pgfqpoint{4.155648in}{0.700177in}}%
\pgfpathlineto{\pgfqpoint{4.160501in}{0.706282in}}%
\pgfpathlineto{\pgfqpoint{4.165354in}{0.713037in}}%
\pgfpathlineto{\pgfqpoint{4.170207in}{0.720439in}}%
\pgfpathlineto{\pgfqpoint{4.175060in}{0.728485in}}%
\pgfpathlineto{\pgfqpoint{4.179913in}{0.737172in}}%
\pgfpathlineto{\pgfqpoint{4.184766in}{0.746495in}}%
\pgfpathlineto{\pgfqpoint{4.189619in}{0.756451in}}%
\pgfpathlineto{\pgfqpoint{4.194472in}{0.767035in}}%
\pgfpathlineto{\pgfqpoint{4.199325in}{0.778242in}}%
\pgfpathlineto{\pgfqpoint{4.204178in}{0.790067in}}%
\pgfpathlineto{\pgfqpoint{4.209031in}{0.802506in}}%
\pgfpathlineto{\pgfqpoint{4.213884in}{0.815553in}}%
\pgfpathlineto{\pgfqpoint{4.218738in}{0.829202in}}%
\pgfpathlineto{\pgfqpoint{4.223591in}{0.843446in}}%
\pgfpathlineto{\pgfqpoint{4.228444in}{0.858281in}}%
\pgfpathlineto{\pgfqpoint{4.233297in}{0.873698in}}%
\pgfpathlineto{\pgfqpoint{4.238150in}{0.889692in}}%
\pgfpathlineto{\pgfqpoint{4.243003in}{0.906255in}}%
\pgfpathlineto{\pgfqpoint{4.247856in}{0.923380in}}%
\pgfpathlineto{\pgfqpoint{4.252709in}{0.941059in}}%
\pgfpathlineto{\pgfqpoint{4.257562in}{0.959285in}}%
\pgfpathlineto{\pgfqpoint{4.262415in}{0.978049in}}%
\pgfpathlineto{\pgfqpoint{4.267268in}{0.997344in}}%
\pgfpathlineto{\pgfqpoint{4.272121in}{1.017160in}}%
\pgfpathlineto{\pgfqpoint{4.276974in}{1.037490in}}%
\pgfpathlineto{\pgfqpoint{4.281827in}{1.058323in}}%
\pgfpathlineto{\pgfqpoint{4.286680in}{1.079651in}}%
\pgfpathlineto{\pgfqpoint{4.291533in}{1.101465in}}%
\pgfpathlineto{\pgfqpoint{4.296386in}{1.123754in}}%
\pgfpathlineto{\pgfqpoint{4.301239in}{1.146509in}}%
\pgfpathlineto{\pgfqpoint{4.306092in}{1.169720in}}%
\pgfpathlineto{\pgfqpoint{4.310945in}{1.193376in}}%
\pgfpathlineto{\pgfqpoint{4.315798in}{1.217468in}}%
\pgfpathlineto{\pgfqpoint{4.320651in}{1.241984in}}%
\pgfpathlineto{\pgfqpoint{4.325504in}{1.266914in}}%
\pgfpathlineto{\pgfqpoint{4.330357in}{1.292247in}}%
\pgfpathlineto{\pgfqpoint{4.335210in}{1.317971in}}%
\pgfpathlineto{\pgfqpoint{4.340063in}{1.344075in}}%
\pgfpathlineto{\pgfqpoint{4.344916in}{1.370548in}}%
\pgfpathlineto{\pgfqpoint{4.349769in}{1.397378in}}%
\pgfpathlineto{\pgfqpoint{4.354622in}{1.424553in}}%
\pgfpathlineto{\pgfqpoint{4.359475in}{1.452060in}}%
\pgfpathlineto{\pgfqpoint{4.364328in}{1.479889in}}%
\pgfpathlineto{\pgfqpoint{4.369181in}{1.508027in}}%
\pgfpathlineto{\pgfqpoint{4.374034in}{1.536460in}}%
\pgfpathlineto{\pgfqpoint{4.378887in}{1.565178in}}%
\pgfpathlineto{\pgfqpoint{4.383740in}{1.594166in}}%
\pgfpathlineto{\pgfqpoint{4.388593in}{1.623412in}}%
\pgfpathlineto{\pgfqpoint{4.393447in}{1.652904in}}%
\pgfpathlineto{\pgfqpoint{4.398300in}{1.682628in}}%
\pgfpathlineto{\pgfqpoint{4.403153in}{1.712571in}}%
\pgfpathlineto{\pgfqpoint{4.408006in}{1.742720in}}%
\pgfpathlineto{\pgfqpoint{4.412859in}{1.773062in}}%
\pgfpathlineto{\pgfqpoint{4.417712in}{1.803583in}}%
\pgfpathlineto{\pgfqpoint{4.422565in}{1.834269in}}%
\pgfpathlineto{\pgfqpoint{4.427418in}{1.865108in}}%
\pgfpathlineto{\pgfqpoint{4.432271in}{1.896085in}}%
\pgfpathlineto{\pgfqpoint{4.437124in}{1.927187in}}%
\pgfpathlineto{\pgfqpoint{4.441977in}{1.958400in}}%
\pgfpathlineto{\pgfqpoint{4.446830in}{1.989710in}}%
\pgfpathlineto{\pgfqpoint{4.451683in}{2.021104in}}%
\pgfpathlineto{\pgfqpoint{4.456536in}{2.052568in}}%
\pgfpathlineto{\pgfqpoint{4.461389in}{2.084087in}}%
\pgfpathlineto{\pgfqpoint{4.466242in}{2.115648in}}%
\pgfpathlineto{\pgfqpoint{4.471095in}{2.147238in}}%
\pgfpathlineto{\pgfqpoint{4.475948in}{2.178841in}}%
\pgfpathlineto{\pgfqpoint{4.475948in}{3.682857in}}%
\pgfpathlineto{\pgfqpoint{4.475948in}{3.682857in}}%
\pgfpathlineto{\pgfqpoint{4.471095in}{3.682525in}}%
\pgfpathlineto{\pgfqpoint{4.466242in}{3.681529in}}%
\pgfpathlineto{\pgfqpoint{4.461389in}{3.679870in}}%
\pgfpathlineto{\pgfqpoint{4.456536in}{3.677547in}}%
\pgfpathlineto{\pgfqpoint{4.451683in}{3.674563in}}%
\pgfpathlineto{\pgfqpoint{4.446830in}{3.670918in}}%
\pgfpathlineto{\pgfqpoint{4.441977in}{3.666615in}}%
\pgfpathlineto{\pgfqpoint{4.437124in}{3.661654in}}%
\pgfpathlineto{\pgfqpoint{4.432271in}{3.656039in}}%
\pgfpathlineto{\pgfqpoint{4.427418in}{3.649772in}}%
\pgfpathlineto{\pgfqpoint{4.422565in}{3.642855in}}%
\pgfpathlineto{\pgfqpoint{4.417712in}{3.635291in}}%
\pgfpathlineto{\pgfqpoint{4.412859in}{3.627084in}}%
\pgfpathlineto{\pgfqpoint{4.408006in}{3.618238in}}%
\pgfpathlineto{\pgfqpoint{4.403153in}{3.608757in}}%
\pgfpathlineto{\pgfqpoint{4.398300in}{3.598643in}}%
\pgfpathlineto{\pgfqpoint{4.393447in}{3.587903in}}%
\pgfpathlineto{\pgfqpoint{4.388593in}{3.576541in}}%
\pgfpathlineto{\pgfqpoint{4.383740in}{3.564561in}}%
\pgfpathlineto{\pgfqpoint{4.378887in}{3.551970in}}%
\pgfpathlineto{\pgfqpoint{4.374034in}{3.538772in}}%
\pgfpathlineto{\pgfqpoint{4.369181in}{3.524974in}}%
\pgfpathlineto{\pgfqpoint{4.364328in}{3.510581in}}%
\pgfpathlineto{\pgfqpoint{4.359475in}{3.495601in}}%
\pgfpathlineto{\pgfqpoint{4.354622in}{3.480039in}}%
\pgfpathlineto{\pgfqpoint{4.349769in}{3.463902in}}%
\pgfpathlineto{\pgfqpoint{4.344916in}{3.447198in}}%
\pgfpathlineto{\pgfqpoint{4.340063in}{3.429933in}}%
\pgfpathlineto{\pgfqpoint{4.335210in}{3.412117in}}%
\pgfpathlineto{\pgfqpoint{4.330357in}{3.393755in}}%
\pgfpathlineto{\pgfqpoint{4.325504in}{3.374858in}}%
\pgfpathlineto{\pgfqpoint{4.320651in}{3.355432in}}%
\pgfpathlineto{\pgfqpoint{4.315798in}{3.335486in}}%
\pgfpathlineto{\pgfqpoint{4.310945in}{3.315030in}}%
\pgfpathlineto{\pgfqpoint{4.306092in}{3.294072in}}%
\pgfpathlineto{\pgfqpoint{4.301239in}{3.272622in}}%
\pgfpathlineto{\pgfqpoint{4.296386in}{3.250689in}}%
\pgfpathlineto{\pgfqpoint{4.291533in}{3.228282in}}%
\pgfpathlineto{\pgfqpoint{4.286680in}{3.205412in}}%
\pgfpathlineto{\pgfqpoint{4.281827in}{3.182089in}}%
\pgfpathlineto{\pgfqpoint{4.276974in}{3.158322in}}%
\pgfpathlineto{\pgfqpoint{4.272121in}{3.134124in}}%
\pgfpathlineto{\pgfqpoint{4.267268in}{3.109503in}}%
\pgfpathlineto{\pgfqpoint{4.262415in}{3.084471in}}%
\pgfpathlineto{\pgfqpoint{4.257562in}{3.059040in}}%
\pgfpathlineto{\pgfqpoint{4.252709in}{3.033220in}}%
\pgfpathlineto{\pgfqpoint{4.247856in}{3.007022in}}%
\pgfpathlineto{\pgfqpoint{4.243003in}{2.980459in}}%
\pgfpathlineto{\pgfqpoint{4.238150in}{2.953542in}}%
\pgfpathlineto{\pgfqpoint{4.233297in}{2.926283in}}%
\pgfpathlineto{\pgfqpoint{4.228444in}{2.898693in}}%
\pgfpathlineto{\pgfqpoint{4.223591in}{2.870786in}}%
\pgfpathlineto{\pgfqpoint{4.218738in}{2.842574in}}%
\pgfpathlineto{\pgfqpoint{4.213884in}{2.814068in}}%
\pgfpathlineto{\pgfqpoint{4.209031in}{2.785282in}}%
\pgfpathlineto{\pgfqpoint{4.204178in}{2.756227in}}%
\pgfpathlineto{\pgfqpoint{4.199325in}{2.726918in}}%
\pgfpathlineto{\pgfqpoint{4.194472in}{2.697367in}}%
\pgfpathlineto{\pgfqpoint{4.189619in}{2.667587in}}%
\pgfpathlineto{\pgfqpoint{4.184766in}{2.637592in}}%
\pgfpathlineto{\pgfqpoint{4.179913in}{2.607393in}}%
\pgfpathlineto{\pgfqpoint{4.175060in}{2.577006in}}%
\pgfpathlineto{\pgfqpoint{4.170207in}{2.546442in}}%
\pgfpathlineto{\pgfqpoint{4.165354in}{2.515716in}}%
\pgfpathlineto{\pgfqpoint{4.160501in}{2.484842in}}%
\pgfpathlineto{\pgfqpoint{4.155648in}{2.453832in}}%
\pgfpathlineto{\pgfqpoint{4.150795in}{2.422701in}}%
\pgfpathlineto{\pgfqpoint{4.145942in}{2.391462in}}%
\pgfpathlineto{\pgfqpoint{4.141089in}{2.360130in}}%
\pgfpathlineto{\pgfqpoint{4.136236in}{2.328717in}}%
\pgfpathlineto{\pgfqpoint{4.131383in}{2.297238in}}%
\pgfpathlineto{\pgfqpoint{4.126530in}{2.265707in}}%
\pgfpathlineto{\pgfqpoint{4.121677in}{2.234138in}}%
\pgfpathlineto{\pgfqpoint{4.116824in}{2.202544in}}%
\pgfpathlineto{\pgfqpoint{4.111971in}{2.170939in}}%
\pgfpathlineto{\pgfqpoint{4.107118in}{2.139338in}}%
\pgfpathlineto{\pgfqpoint{4.102265in}{2.107755in}}%
\pgfpathlineto{\pgfqpoint{4.097412in}{2.076203in}}%
\pgfpathlineto{\pgfqpoint{4.092559in}{2.044696in}}%
\pgfpathlineto{\pgfqpoint{4.087706in}{2.013249in}}%
\pgfpathlineto{\pgfqpoint{4.082853in}{1.981874in}}%
\pgfpathlineto{\pgfqpoint{4.078000in}{1.950587in}}%
\pgfpathlineto{\pgfqpoint{4.073147in}{1.919400in}}%
\pgfpathlineto{\pgfqpoint{4.068294in}{1.888328in}}%
\pgfpathlineto{\pgfqpoint{4.063441in}{1.857384in}}%
\pgfpathlineto{\pgfqpoint{4.058588in}{1.826583in}}%
\pgfpathlineto{\pgfqpoint{4.053735in}{1.795936in}}%
\pgfpathlineto{\pgfqpoint{4.048882in}{1.765459in}}%
\pgfpathlineto{\pgfqpoint{4.044029in}{1.735164in}}%
\pgfpathlineto{\pgfqpoint{4.039176in}{1.705066in}}%
\pgfpathlineto{\pgfqpoint{4.034322in}{1.675176in}}%
\pgfpathlineto{\pgfqpoint{4.029469in}{1.645509in}}%
\pgfpathlineto{\pgfqpoint{4.024616in}{1.616077in}}%
\pgfpathlineto{\pgfqpoint{4.019763in}{1.586894in}}%
\pgfpathlineto{\pgfqpoint{4.014910in}{1.557972in}}%
\pgfpathlineto{\pgfqpoint{4.010057in}{1.529325in}}%
\pgfpathlineto{\pgfqpoint{4.005204in}{1.500964in}}%
\pgfpathlineto{\pgfqpoint{4.000351in}{1.472903in}}%
\pgfpathlineto{\pgfqpoint{3.995498in}{1.445153in}}%
\pgfpathlineto{\pgfqpoint{3.990645in}{1.417727in}}%
\pgfpathlineto{\pgfqpoint{3.985792in}{1.390637in}}%
\pgfpathlineto{\pgfqpoint{3.980939in}{1.363896in}}%
\pgfpathlineto{\pgfqpoint{3.976086in}{1.337514in}}%
\pgfpathlineto{\pgfqpoint{3.971233in}{1.311504in}}%
\pgfpathlineto{\pgfqpoint{3.966380in}{1.285877in}}%
\pgfpathlineto{\pgfqpoint{3.961527in}{1.260644in}}%
\pgfpathlineto{\pgfqpoint{3.956674in}{1.235816in}}%
\pgfpathlineto{\pgfqpoint{3.951821in}{1.211405in}}%
\pgfpathlineto{\pgfqpoint{3.946968in}{1.187421in}}%
\pgfpathlineto{\pgfqpoint{3.942115in}{1.163875in}}%
\pgfpathlineto{\pgfqpoint{3.937262in}{1.140777in}}%
\pgfpathlineto{\pgfqpoint{3.932409in}{1.118137in}}%
\pgfpathlineto{\pgfqpoint{3.932409in}{1.118137in}}%
\pgfpathclose%
\pgfusepath{fill}%
\end{pgfscope}%
\begin{pgfscope}%
\pgfpathrectangle{\pgfqpoint{2.952340in}{0.524444in}}{\pgfqpoint{1.596161in}{3.308814in}}%
\pgfusepath{clip}%
\pgfsetbuttcap%
\pgfsetroundjoin%
\pgfsetlinewidth{0.501875pt}%
\definecolor{currentstroke}{rgb}{0.501961,0.501961,0.501961}%
\pgfsetstrokecolor{currentstroke}%
\pgfsetstrokeopacity{0.500000}%
\pgfsetdash{{1.850000pt}{0.800000pt}}{0.000000pt}%
\pgfpathmoveto{\pgfqpoint{3.024893in}{0.524444in}}%
\pgfpathlineto{\pgfqpoint{3.024893in}{3.833258in}}%
\pgfusepath{stroke}%
\end{pgfscope}%
\begin{pgfscope}%
\pgfsetbuttcap%
\pgfsetroundjoin%
\definecolor{currentfill}{rgb}{0.000000,0.000000,0.000000}%
\pgfsetfillcolor{currentfill}%
\pgfsetlinewidth{0.803000pt}%
\definecolor{currentstroke}{rgb}{0.000000,0.000000,0.000000}%
\pgfsetstrokecolor{currentstroke}%
\pgfsetdash{}{0pt}%
\pgfsys@defobject{currentmarker}{\pgfqpoint{0.000000in}{-0.048611in}}{\pgfqpoint{0.000000in}{0.000000in}}{%
\pgfpathmoveto{\pgfqpoint{0.000000in}{0.000000in}}%
\pgfpathlineto{\pgfqpoint{0.000000in}{-0.048611in}}%
\pgfusepath{stroke,fill}%
}%
\begin{pgfscope}%
\pgfsys@transformshift{3.024893in}{0.524444in}%
\pgfsys@useobject{currentmarker}{}%
\end{pgfscope}%
\end{pgfscope}%
\begin{pgfscope}%
\definecolor{textcolor}{rgb}{0.000000,0.000000,0.000000}%
\pgfsetstrokecolor{textcolor}%
\pgfsetfillcolor{textcolor}%
\pgftext[x=3.024893in,y=0.427222in,,top]{\color{textcolor}{\rmfamily\fontsize{10.000000}{12.000000}\selectfont\catcode`\^=\active\def^{\ifmmode\sp\else\^{}\fi}\catcode`\%=\active\def%{\%}$\mathdefault{0.0}$}}%
\end{pgfscope}%
\begin{pgfscope}%
\pgfpathrectangle{\pgfqpoint{2.952340in}{0.524444in}}{\pgfqpoint{1.596161in}{3.308814in}}%
\pgfusepath{clip}%
\pgfsetbuttcap%
\pgfsetroundjoin%
\pgfsetlinewidth{0.501875pt}%
\definecolor{currentstroke}{rgb}{0.501961,0.501961,0.501961}%
\pgfsetstrokecolor{currentstroke}%
\pgfsetstrokeopacity{0.500000}%
\pgfsetdash{{1.850000pt}{0.800000pt}}{0.000000pt}%
\pgfpathmoveto{\pgfqpoint{3.602249in}{0.524444in}}%
\pgfpathlineto{\pgfqpoint{3.602249in}{3.833258in}}%
\pgfusepath{stroke}%
\end{pgfscope}%
\begin{pgfscope}%
\pgfsetbuttcap%
\pgfsetroundjoin%
\definecolor{currentfill}{rgb}{0.000000,0.000000,0.000000}%
\pgfsetfillcolor{currentfill}%
\pgfsetlinewidth{0.803000pt}%
\definecolor{currentstroke}{rgb}{0.000000,0.000000,0.000000}%
\pgfsetstrokecolor{currentstroke}%
\pgfsetdash{}{0pt}%
\pgfsys@defobject{currentmarker}{\pgfqpoint{0.000000in}{-0.048611in}}{\pgfqpoint{0.000000in}{0.000000in}}{%
\pgfpathmoveto{\pgfqpoint{0.000000in}{0.000000in}}%
\pgfpathlineto{\pgfqpoint{0.000000in}{-0.048611in}}%
\pgfusepath{stroke,fill}%
}%
\begin{pgfscope}%
\pgfsys@transformshift{3.602249in}{0.524444in}%
\pgfsys@useobject{currentmarker}{}%
\end{pgfscope}%
\end{pgfscope}%
\begin{pgfscope}%
\definecolor{textcolor}{rgb}{0.000000,0.000000,0.000000}%
\pgfsetstrokecolor{textcolor}%
\pgfsetfillcolor{textcolor}%
\pgftext[x=3.602249in,y=0.427222in,,top]{\color{textcolor}{\rmfamily\fontsize{10.000000}{12.000000}\selectfont\catcode`\^=\active\def^{\ifmmode\sp\else\^{}\fi}\catcode`\%=\active\def%{\%}$\mathdefault{2.5}$}}%
\end{pgfscope}%
\begin{pgfscope}%
\pgfpathrectangle{\pgfqpoint{2.952340in}{0.524444in}}{\pgfqpoint{1.596161in}{3.308814in}}%
\pgfusepath{clip}%
\pgfsetbuttcap%
\pgfsetroundjoin%
\pgfsetlinewidth{0.501875pt}%
\definecolor{currentstroke}{rgb}{0.501961,0.501961,0.501961}%
\pgfsetstrokecolor{currentstroke}%
\pgfsetstrokeopacity{0.500000}%
\pgfsetdash{{1.850000pt}{0.800000pt}}{0.000000pt}%
\pgfpathmoveto{\pgfqpoint{4.179606in}{0.524444in}}%
\pgfpathlineto{\pgfqpoint{4.179606in}{3.833258in}}%
\pgfusepath{stroke}%
\end{pgfscope}%
\begin{pgfscope}%
\pgfsetbuttcap%
\pgfsetroundjoin%
\definecolor{currentfill}{rgb}{0.000000,0.000000,0.000000}%
\pgfsetfillcolor{currentfill}%
\pgfsetlinewidth{0.803000pt}%
\definecolor{currentstroke}{rgb}{0.000000,0.000000,0.000000}%
\pgfsetstrokecolor{currentstroke}%
\pgfsetdash{}{0pt}%
\pgfsys@defobject{currentmarker}{\pgfqpoint{0.000000in}{-0.048611in}}{\pgfqpoint{0.000000in}{0.000000in}}{%
\pgfpathmoveto{\pgfqpoint{0.000000in}{0.000000in}}%
\pgfpathlineto{\pgfqpoint{0.000000in}{-0.048611in}}%
\pgfusepath{stroke,fill}%
}%
\begin{pgfscope}%
\pgfsys@transformshift{4.179606in}{0.524444in}%
\pgfsys@useobject{currentmarker}{}%
\end{pgfscope}%
\end{pgfscope}%
\begin{pgfscope}%
\definecolor{textcolor}{rgb}{0.000000,0.000000,0.000000}%
\pgfsetstrokecolor{textcolor}%
\pgfsetfillcolor{textcolor}%
\pgftext[x=4.179606in,y=0.427222in,,top]{\color{textcolor}{\rmfamily\fontsize{10.000000}{12.000000}\selectfont\catcode`\^=\active\def^{\ifmmode\sp\else\^{}\fi}\catcode`\%=\active\def%{\%}$\mathdefault{5.0}$}}%
\end{pgfscope}%
\begin{pgfscope}%
\definecolor{textcolor}{rgb}{0.000000,0.000000,0.000000}%
\pgfsetstrokecolor{textcolor}%
\pgfsetfillcolor{textcolor}%
\pgftext[x=3.750420in,y=0.248333in,,top]{\color{textcolor}{\rmfamily\fontsize{10.000000}{12.000000}\selectfont\catcode`\^=\active\def^{\ifmmode\sp\else\^{}\fi}\catcode`\%=\active\def%{\%}Ángulo}}%
\end{pgfscope}%
\begin{pgfscope}%
\pgfpathrectangle{\pgfqpoint{2.952340in}{0.524444in}}{\pgfqpoint{1.596161in}{3.308814in}}%
\pgfusepath{clip}%
\pgfsetbuttcap%
\pgfsetroundjoin%
\pgfsetlinewidth{0.501875pt}%
\definecolor{currentstroke}{rgb}{0.501961,0.501961,0.501961}%
\pgfsetstrokecolor{currentstroke}%
\pgfsetstrokeopacity{0.500000}%
\pgfsetdash{{1.850000pt}{0.800000pt}}{0.000000pt}%
\pgfpathmoveto{\pgfqpoint{2.952340in}{0.674824in}}%
\pgfpathlineto{\pgfqpoint{4.548501in}{0.674824in}}%
\pgfusepath{stroke}%
\end{pgfscope}%
\begin{pgfscope}%
\pgfsetbuttcap%
\pgfsetroundjoin%
\definecolor{currentfill}{rgb}{0.000000,0.000000,0.000000}%
\pgfsetfillcolor{currentfill}%
\pgfsetlinewidth{0.803000pt}%
\definecolor{currentstroke}{rgb}{0.000000,0.000000,0.000000}%
\pgfsetstrokecolor{currentstroke}%
\pgfsetdash{}{0pt}%
\pgfsys@defobject{currentmarker}{\pgfqpoint{-0.048611in}{0.000000in}}{\pgfqpoint{-0.000000in}{0.000000in}}{%
\pgfpathmoveto{\pgfqpoint{-0.000000in}{0.000000in}}%
\pgfpathlineto{\pgfqpoint{-0.048611in}{0.000000in}}%
\pgfusepath{stroke,fill}%
}%
\begin{pgfscope}%
\pgfsys@transformshift{2.952340in}{0.674824in}%
\pgfsys@useobject{currentmarker}{}%
\end{pgfscope}%
\end{pgfscope}%
\begin{pgfscope}%
\definecolor{textcolor}{rgb}{0.000000,0.000000,0.000000}%
\pgfsetstrokecolor{textcolor}%
\pgfsetfillcolor{textcolor}%
\pgftext[x=2.608203in, y=0.626629in, left, base]{\color{textcolor}{\rmfamily\fontsize{10.000000}{12.000000}\selectfont\catcode`\^=\active\def^{\ifmmode\sp\else\^{}\fi}\catcode`\%=\active\def%{\%}$\mathdefault{0.00}$}}%
\end{pgfscope}%
\begin{pgfscope}%
\pgfpathrectangle{\pgfqpoint{2.952340in}{0.524444in}}{\pgfqpoint{1.596161in}{3.308814in}}%
\pgfusepath{clip}%
\pgfsetbuttcap%
\pgfsetroundjoin%
\pgfsetlinewidth{0.501875pt}%
\definecolor{currentstroke}{rgb}{0.501961,0.501961,0.501961}%
\pgfsetstrokecolor{currentstroke}%
\pgfsetstrokeopacity{0.500000}%
\pgfsetdash{{1.850000pt}{0.800000pt}}{0.000000pt}%
\pgfpathmoveto{\pgfqpoint{2.952340in}{1.050828in}}%
\pgfpathlineto{\pgfqpoint{4.548501in}{1.050828in}}%
\pgfusepath{stroke}%
\end{pgfscope}%
\begin{pgfscope}%
\pgfsetbuttcap%
\pgfsetroundjoin%
\definecolor{currentfill}{rgb}{0.000000,0.000000,0.000000}%
\pgfsetfillcolor{currentfill}%
\pgfsetlinewidth{0.803000pt}%
\definecolor{currentstroke}{rgb}{0.000000,0.000000,0.000000}%
\pgfsetstrokecolor{currentstroke}%
\pgfsetdash{}{0pt}%
\pgfsys@defobject{currentmarker}{\pgfqpoint{-0.048611in}{0.000000in}}{\pgfqpoint{-0.000000in}{0.000000in}}{%
\pgfpathmoveto{\pgfqpoint{-0.000000in}{0.000000in}}%
\pgfpathlineto{\pgfqpoint{-0.048611in}{0.000000in}}%
\pgfusepath{stroke,fill}%
}%
\begin{pgfscope}%
\pgfsys@transformshift{2.952340in}{1.050828in}%
\pgfsys@useobject{currentmarker}{}%
\end{pgfscope}%
\end{pgfscope}%
\begin{pgfscope}%
\definecolor{textcolor}{rgb}{0.000000,0.000000,0.000000}%
\pgfsetstrokecolor{textcolor}%
\pgfsetfillcolor{textcolor}%
\pgftext[x=2.608203in, y=1.002634in, left, base]{\color{textcolor}{\rmfamily\fontsize{10.000000}{12.000000}\selectfont\catcode`\^=\active\def^{\ifmmode\sp\else\^{}\fi}\catcode`\%=\active\def%{\%}$\mathdefault{0.25}$}}%
\end{pgfscope}%
\begin{pgfscope}%
\pgfpathrectangle{\pgfqpoint{2.952340in}{0.524444in}}{\pgfqpoint{1.596161in}{3.308814in}}%
\pgfusepath{clip}%
\pgfsetbuttcap%
\pgfsetroundjoin%
\pgfsetlinewidth{0.501875pt}%
\definecolor{currentstroke}{rgb}{0.501961,0.501961,0.501961}%
\pgfsetstrokecolor{currentstroke}%
\pgfsetstrokeopacity{0.500000}%
\pgfsetdash{{1.850000pt}{0.800000pt}}{0.000000pt}%
\pgfpathmoveto{\pgfqpoint{2.952340in}{1.426832in}}%
\pgfpathlineto{\pgfqpoint{4.548501in}{1.426832in}}%
\pgfusepath{stroke}%
\end{pgfscope}%
\begin{pgfscope}%
\pgfsetbuttcap%
\pgfsetroundjoin%
\definecolor{currentfill}{rgb}{0.000000,0.000000,0.000000}%
\pgfsetfillcolor{currentfill}%
\pgfsetlinewidth{0.803000pt}%
\definecolor{currentstroke}{rgb}{0.000000,0.000000,0.000000}%
\pgfsetstrokecolor{currentstroke}%
\pgfsetdash{}{0pt}%
\pgfsys@defobject{currentmarker}{\pgfqpoint{-0.048611in}{0.000000in}}{\pgfqpoint{-0.000000in}{0.000000in}}{%
\pgfpathmoveto{\pgfqpoint{-0.000000in}{0.000000in}}%
\pgfpathlineto{\pgfqpoint{-0.048611in}{0.000000in}}%
\pgfusepath{stroke,fill}%
}%
\begin{pgfscope}%
\pgfsys@transformshift{2.952340in}{1.426832in}%
\pgfsys@useobject{currentmarker}{}%
\end{pgfscope}%
\end{pgfscope}%
\begin{pgfscope}%
\definecolor{textcolor}{rgb}{0.000000,0.000000,0.000000}%
\pgfsetstrokecolor{textcolor}%
\pgfsetfillcolor{textcolor}%
\pgftext[x=2.608203in, y=1.378638in, left, base]{\color{textcolor}{\rmfamily\fontsize{10.000000}{12.000000}\selectfont\catcode`\^=\active\def^{\ifmmode\sp\else\^{}\fi}\catcode`\%=\active\def%{\%}$\mathdefault{0.50}$}}%
\end{pgfscope}%
\begin{pgfscope}%
\pgfpathrectangle{\pgfqpoint{2.952340in}{0.524444in}}{\pgfqpoint{1.596161in}{3.308814in}}%
\pgfusepath{clip}%
\pgfsetbuttcap%
\pgfsetroundjoin%
\pgfsetlinewidth{0.501875pt}%
\definecolor{currentstroke}{rgb}{0.501961,0.501961,0.501961}%
\pgfsetstrokecolor{currentstroke}%
\pgfsetstrokeopacity{0.500000}%
\pgfsetdash{{1.850000pt}{0.800000pt}}{0.000000pt}%
\pgfpathmoveto{\pgfqpoint{2.952340in}{1.802836in}}%
\pgfpathlineto{\pgfqpoint{4.548501in}{1.802836in}}%
\pgfusepath{stroke}%
\end{pgfscope}%
\begin{pgfscope}%
\pgfsetbuttcap%
\pgfsetroundjoin%
\definecolor{currentfill}{rgb}{0.000000,0.000000,0.000000}%
\pgfsetfillcolor{currentfill}%
\pgfsetlinewidth{0.803000pt}%
\definecolor{currentstroke}{rgb}{0.000000,0.000000,0.000000}%
\pgfsetstrokecolor{currentstroke}%
\pgfsetdash{}{0pt}%
\pgfsys@defobject{currentmarker}{\pgfqpoint{-0.048611in}{0.000000in}}{\pgfqpoint{-0.000000in}{0.000000in}}{%
\pgfpathmoveto{\pgfqpoint{-0.000000in}{0.000000in}}%
\pgfpathlineto{\pgfqpoint{-0.048611in}{0.000000in}}%
\pgfusepath{stroke,fill}%
}%
\begin{pgfscope}%
\pgfsys@transformshift{2.952340in}{1.802836in}%
\pgfsys@useobject{currentmarker}{}%
\end{pgfscope}%
\end{pgfscope}%
\begin{pgfscope}%
\definecolor{textcolor}{rgb}{0.000000,0.000000,0.000000}%
\pgfsetstrokecolor{textcolor}%
\pgfsetfillcolor{textcolor}%
\pgftext[x=2.608203in, y=1.754642in, left, base]{\color{textcolor}{\rmfamily\fontsize{10.000000}{12.000000}\selectfont\catcode`\^=\active\def^{\ifmmode\sp\else\^{}\fi}\catcode`\%=\active\def%{\%}$\mathdefault{0.75}$}}%
\end{pgfscope}%
\begin{pgfscope}%
\pgfpathrectangle{\pgfqpoint{2.952340in}{0.524444in}}{\pgfqpoint{1.596161in}{3.308814in}}%
\pgfusepath{clip}%
\pgfsetbuttcap%
\pgfsetroundjoin%
\pgfsetlinewidth{0.501875pt}%
\definecolor{currentstroke}{rgb}{0.501961,0.501961,0.501961}%
\pgfsetstrokecolor{currentstroke}%
\pgfsetstrokeopacity{0.500000}%
\pgfsetdash{{1.850000pt}{0.800000pt}}{0.000000pt}%
\pgfpathmoveto{\pgfqpoint{2.952340in}{2.178841in}}%
\pgfpathlineto{\pgfqpoint{4.548501in}{2.178841in}}%
\pgfusepath{stroke}%
\end{pgfscope}%
\begin{pgfscope}%
\pgfsetbuttcap%
\pgfsetroundjoin%
\definecolor{currentfill}{rgb}{0.000000,0.000000,0.000000}%
\pgfsetfillcolor{currentfill}%
\pgfsetlinewidth{0.803000pt}%
\definecolor{currentstroke}{rgb}{0.000000,0.000000,0.000000}%
\pgfsetstrokecolor{currentstroke}%
\pgfsetdash{}{0pt}%
\pgfsys@defobject{currentmarker}{\pgfqpoint{-0.048611in}{0.000000in}}{\pgfqpoint{-0.000000in}{0.000000in}}{%
\pgfpathmoveto{\pgfqpoint{-0.000000in}{0.000000in}}%
\pgfpathlineto{\pgfqpoint{-0.048611in}{0.000000in}}%
\pgfusepath{stroke,fill}%
}%
\begin{pgfscope}%
\pgfsys@transformshift{2.952340in}{2.178841in}%
\pgfsys@useobject{currentmarker}{}%
\end{pgfscope}%
\end{pgfscope}%
\begin{pgfscope}%
\definecolor{textcolor}{rgb}{0.000000,0.000000,0.000000}%
\pgfsetstrokecolor{textcolor}%
\pgfsetfillcolor{textcolor}%
\pgftext[x=2.608203in, y=2.130646in, left, base]{\color{textcolor}{\rmfamily\fontsize{10.000000}{12.000000}\selectfont\catcode`\^=\active\def^{\ifmmode\sp\else\^{}\fi}\catcode`\%=\active\def%{\%}$\mathdefault{1.00}$}}%
\end{pgfscope}%
\begin{pgfscope}%
\pgfpathrectangle{\pgfqpoint{2.952340in}{0.524444in}}{\pgfqpoint{1.596161in}{3.308814in}}%
\pgfusepath{clip}%
\pgfsetbuttcap%
\pgfsetroundjoin%
\pgfsetlinewidth{0.501875pt}%
\definecolor{currentstroke}{rgb}{0.501961,0.501961,0.501961}%
\pgfsetstrokecolor{currentstroke}%
\pgfsetstrokeopacity{0.500000}%
\pgfsetdash{{1.850000pt}{0.800000pt}}{0.000000pt}%
\pgfpathmoveto{\pgfqpoint{2.952340in}{2.554845in}}%
\pgfpathlineto{\pgfqpoint{4.548501in}{2.554845in}}%
\pgfusepath{stroke}%
\end{pgfscope}%
\begin{pgfscope}%
\pgfsetbuttcap%
\pgfsetroundjoin%
\definecolor{currentfill}{rgb}{0.000000,0.000000,0.000000}%
\pgfsetfillcolor{currentfill}%
\pgfsetlinewidth{0.803000pt}%
\definecolor{currentstroke}{rgb}{0.000000,0.000000,0.000000}%
\pgfsetstrokecolor{currentstroke}%
\pgfsetdash{}{0pt}%
\pgfsys@defobject{currentmarker}{\pgfqpoint{-0.048611in}{0.000000in}}{\pgfqpoint{-0.000000in}{0.000000in}}{%
\pgfpathmoveto{\pgfqpoint{-0.000000in}{0.000000in}}%
\pgfpathlineto{\pgfqpoint{-0.048611in}{0.000000in}}%
\pgfusepath{stroke,fill}%
}%
\begin{pgfscope}%
\pgfsys@transformshift{2.952340in}{2.554845in}%
\pgfsys@useobject{currentmarker}{}%
\end{pgfscope}%
\end{pgfscope}%
\begin{pgfscope}%
\definecolor{textcolor}{rgb}{0.000000,0.000000,0.000000}%
\pgfsetstrokecolor{textcolor}%
\pgfsetfillcolor{textcolor}%
\pgftext[x=2.608203in, y=2.506650in, left, base]{\color{textcolor}{\rmfamily\fontsize{10.000000}{12.000000}\selectfont\catcode`\^=\active\def^{\ifmmode\sp\else\^{}\fi}\catcode`\%=\active\def%{\%}$\mathdefault{1.25}$}}%
\end{pgfscope}%
\begin{pgfscope}%
\pgfpathrectangle{\pgfqpoint{2.952340in}{0.524444in}}{\pgfqpoint{1.596161in}{3.308814in}}%
\pgfusepath{clip}%
\pgfsetbuttcap%
\pgfsetroundjoin%
\pgfsetlinewidth{0.501875pt}%
\definecolor{currentstroke}{rgb}{0.501961,0.501961,0.501961}%
\pgfsetstrokecolor{currentstroke}%
\pgfsetstrokeopacity{0.500000}%
\pgfsetdash{{1.850000pt}{0.800000pt}}{0.000000pt}%
\pgfpathmoveto{\pgfqpoint{2.952340in}{2.930849in}}%
\pgfpathlineto{\pgfqpoint{4.548501in}{2.930849in}}%
\pgfusepath{stroke}%
\end{pgfscope}%
\begin{pgfscope}%
\pgfsetbuttcap%
\pgfsetroundjoin%
\definecolor{currentfill}{rgb}{0.000000,0.000000,0.000000}%
\pgfsetfillcolor{currentfill}%
\pgfsetlinewidth{0.803000pt}%
\definecolor{currentstroke}{rgb}{0.000000,0.000000,0.000000}%
\pgfsetstrokecolor{currentstroke}%
\pgfsetdash{}{0pt}%
\pgfsys@defobject{currentmarker}{\pgfqpoint{-0.048611in}{0.000000in}}{\pgfqpoint{-0.000000in}{0.000000in}}{%
\pgfpathmoveto{\pgfqpoint{-0.000000in}{0.000000in}}%
\pgfpathlineto{\pgfqpoint{-0.048611in}{0.000000in}}%
\pgfusepath{stroke,fill}%
}%
\begin{pgfscope}%
\pgfsys@transformshift{2.952340in}{2.930849in}%
\pgfsys@useobject{currentmarker}{}%
\end{pgfscope}%
\end{pgfscope}%
\begin{pgfscope}%
\definecolor{textcolor}{rgb}{0.000000,0.000000,0.000000}%
\pgfsetstrokecolor{textcolor}%
\pgfsetfillcolor{textcolor}%
\pgftext[x=2.608203in, y=2.882655in, left, base]{\color{textcolor}{\rmfamily\fontsize{10.000000}{12.000000}\selectfont\catcode`\^=\active\def^{\ifmmode\sp\else\^{}\fi}\catcode`\%=\active\def%{\%}$\mathdefault{1.50}$}}%
\end{pgfscope}%
\begin{pgfscope}%
\pgfpathrectangle{\pgfqpoint{2.952340in}{0.524444in}}{\pgfqpoint{1.596161in}{3.308814in}}%
\pgfusepath{clip}%
\pgfsetbuttcap%
\pgfsetroundjoin%
\pgfsetlinewidth{0.501875pt}%
\definecolor{currentstroke}{rgb}{0.501961,0.501961,0.501961}%
\pgfsetstrokecolor{currentstroke}%
\pgfsetstrokeopacity{0.500000}%
\pgfsetdash{{1.850000pt}{0.800000pt}}{0.000000pt}%
\pgfpathmoveto{\pgfqpoint{2.952340in}{3.306853in}}%
\pgfpathlineto{\pgfqpoint{4.548501in}{3.306853in}}%
\pgfusepath{stroke}%
\end{pgfscope}%
\begin{pgfscope}%
\pgfsetbuttcap%
\pgfsetroundjoin%
\definecolor{currentfill}{rgb}{0.000000,0.000000,0.000000}%
\pgfsetfillcolor{currentfill}%
\pgfsetlinewidth{0.803000pt}%
\definecolor{currentstroke}{rgb}{0.000000,0.000000,0.000000}%
\pgfsetstrokecolor{currentstroke}%
\pgfsetdash{}{0pt}%
\pgfsys@defobject{currentmarker}{\pgfqpoint{-0.048611in}{0.000000in}}{\pgfqpoint{-0.000000in}{0.000000in}}{%
\pgfpathmoveto{\pgfqpoint{-0.000000in}{0.000000in}}%
\pgfpathlineto{\pgfqpoint{-0.048611in}{0.000000in}}%
\pgfusepath{stroke,fill}%
}%
\begin{pgfscope}%
\pgfsys@transformshift{2.952340in}{3.306853in}%
\pgfsys@useobject{currentmarker}{}%
\end{pgfscope}%
\end{pgfscope}%
\begin{pgfscope}%
\definecolor{textcolor}{rgb}{0.000000,0.000000,0.000000}%
\pgfsetstrokecolor{textcolor}%
\pgfsetfillcolor{textcolor}%
\pgftext[x=2.608203in, y=3.258659in, left, base]{\color{textcolor}{\rmfamily\fontsize{10.000000}{12.000000}\selectfont\catcode`\^=\active\def^{\ifmmode\sp\else\^{}\fi}\catcode`\%=\active\def%{\%}$\mathdefault{1.75}$}}%
\end{pgfscope}%
\begin{pgfscope}%
\pgfpathrectangle{\pgfqpoint{2.952340in}{0.524444in}}{\pgfqpoint{1.596161in}{3.308814in}}%
\pgfusepath{clip}%
\pgfsetbuttcap%
\pgfsetroundjoin%
\pgfsetlinewidth{0.501875pt}%
\definecolor{currentstroke}{rgb}{0.501961,0.501961,0.501961}%
\pgfsetstrokecolor{currentstroke}%
\pgfsetstrokeopacity{0.500000}%
\pgfsetdash{{1.850000pt}{0.800000pt}}{0.000000pt}%
\pgfpathmoveto{\pgfqpoint{2.952340in}{3.682857in}}%
\pgfpathlineto{\pgfqpoint{4.548501in}{3.682857in}}%
\pgfusepath{stroke}%
\end{pgfscope}%
\begin{pgfscope}%
\pgfsetbuttcap%
\pgfsetroundjoin%
\definecolor{currentfill}{rgb}{0.000000,0.000000,0.000000}%
\pgfsetfillcolor{currentfill}%
\pgfsetlinewidth{0.803000pt}%
\definecolor{currentstroke}{rgb}{0.000000,0.000000,0.000000}%
\pgfsetstrokecolor{currentstroke}%
\pgfsetdash{}{0pt}%
\pgfsys@defobject{currentmarker}{\pgfqpoint{-0.048611in}{0.000000in}}{\pgfqpoint{-0.000000in}{0.000000in}}{%
\pgfpathmoveto{\pgfqpoint{-0.000000in}{0.000000in}}%
\pgfpathlineto{\pgfqpoint{-0.048611in}{0.000000in}}%
\pgfusepath{stroke,fill}%
}%
\begin{pgfscope}%
\pgfsys@transformshift{2.952340in}{3.682857in}%
\pgfsys@useobject{currentmarker}{}%
\end{pgfscope}%
\end{pgfscope}%
\begin{pgfscope}%
\definecolor{textcolor}{rgb}{0.000000,0.000000,0.000000}%
\pgfsetstrokecolor{textcolor}%
\pgfsetfillcolor{textcolor}%
\pgftext[x=2.608203in, y=3.634663in, left, base]{\color{textcolor}{\rmfamily\fontsize{10.000000}{12.000000}\selectfont\catcode`\^=\active\def^{\ifmmode\sp\else\^{}\fi}\catcode`\%=\active\def%{\%}$\mathdefault{2.00}$}}%
\end{pgfscope}%
\begin{pgfscope}%
\definecolor{textcolor}{rgb}{0.000000,0.000000,0.000000}%
\pgfsetstrokecolor{textcolor}%
\pgfsetfillcolor{textcolor}%
\pgftext[x=2.552648in,y=2.178851in,,bottom,rotate=90.000000]{\color{textcolor}{\rmfamily\fontsize{10.000000}{12.000000}\selectfont\catcode`\^=\active\def^{\ifmmode\sp\else\^{}\fi}\catcode`\%=\active\def%{\%}Valor}}%
\end{pgfscope}%
\begin{pgfscope}%
\pgfpathrectangle{\pgfqpoint{2.952340in}{0.524444in}}{\pgfqpoint{1.596161in}{3.308814in}}%
\pgfusepath{clip}%
\pgfsetrectcap%
\pgfsetroundjoin%
\pgfsetlinewidth{1.505625pt}%
\definecolor{currentstroke}{rgb}{0.121569,0.466667,0.705882}%
\pgfsetstrokecolor{currentstroke}%
\pgfsetdash{}{0pt}%
\pgfpathmoveto{\pgfqpoint{3.024893in}{2.178841in}}%
\pgfpathlineto{\pgfqpoint{3.092835in}{2.614961in}}%
\pgfpathlineto{\pgfqpoint{3.131659in}{2.849655in}}%
\pgfpathlineto{\pgfqpoint{3.160778in}{3.013606in}}%
\pgfpathlineto{\pgfqpoint{3.189896in}{3.164305in}}%
\pgfpathlineto{\pgfqpoint{3.214161in}{3.278030in}}%
\pgfpathlineto{\pgfqpoint{3.233573in}{3.360337in}}%
\pgfpathlineto{\pgfqpoint{3.252985in}{3.434302in}}%
\pgfpathlineto{\pgfqpoint{3.272397in}{3.499401in}}%
\pgfpathlineto{\pgfqpoint{3.291809in}{3.555175in}}%
\pgfpathlineto{\pgfqpoint{3.306368in}{3.590647in}}%
\pgfpathlineto{\pgfqpoint{3.320927in}{3.620510in}}%
\pgfpathlineto{\pgfqpoint{3.335486in}{3.644645in}}%
\pgfpathlineto{\pgfqpoint{3.345193in}{3.657504in}}%
\pgfpathlineto{\pgfqpoint{3.354899in}{3.667752in}}%
\pgfpathlineto{\pgfqpoint{3.364605in}{3.675371in}}%
\pgfpathlineto{\pgfqpoint{3.374311in}{3.680347in}}%
\pgfpathlineto{\pgfqpoint{3.384017in}{3.682671in}}%
\pgfpathlineto{\pgfqpoint{3.393723in}{3.682339in}}%
\pgfpathlineto{\pgfqpoint{3.403429in}{3.679351in}}%
\pgfpathlineto{\pgfqpoint{3.413135in}{3.673714in}}%
\pgfpathlineto{\pgfqpoint{3.422841in}{3.665436in}}%
\pgfpathlineto{\pgfqpoint{3.432547in}{3.654533in}}%
\pgfpathlineto{\pgfqpoint{3.442253in}{3.641024in}}%
\pgfpathlineto{\pgfqpoint{3.456812in}{3.615927in}}%
\pgfpathlineto{\pgfqpoint{3.471371in}{3.585121in}}%
\pgfpathlineto{\pgfqpoint{3.485930in}{3.548727in}}%
\pgfpathlineto{\pgfqpoint{3.500489in}{3.506891in}}%
\pgfpathlineto{\pgfqpoint{3.519902in}{3.442934in}}%
\pgfpathlineto{\pgfqpoint{3.539314in}{3.370050in}}%
\pgfpathlineto{\pgfqpoint{3.558726in}{3.288756in}}%
\pgfpathlineto{\pgfqpoint{3.582991in}{3.176188in}}%
\pgfpathlineto{\pgfqpoint{3.607256in}{3.052621in}}%
\pgfpathlineto{\pgfqpoint{3.636374in}{2.891746in}}%
\pgfpathlineto{\pgfqpoint{3.670345in}{2.689943in}}%
\pgfpathlineto{\pgfqpoint{3.714023in}{2.414901in}}%
\pgfpathlineto{\pgfqpoint{3.840201in}{1.608758in}}%
\pgfpathlineto{\pgfqpoint{3.874173in}{1.410923in}}%
\pgfpathlineto{\pgfqpoint{3.903291in}{1.254398in}}%
\pgfpathlineto{\pgfqpoint{3.927556in}{1.135074in}}%
\pgfpathlineto{\pgfqpoint{3.951821in}{1.027262in}}%
\pgfpathlineto{\pgfqpoint{3.971233in}{0.950104in}}%
\pgfpathlineto{\pgfqpoint{3.990645in}{0.881623in}}%
\pgfpathlineto{\pgfqpoint{4.010057in}{0.822303in}}%
\pgfpathlineto{\pgfqpoint{4.024616in}{0.784078in}}%
\pgfpathlineto{\pgfqpoint{4.039176in}{0.751394in}}%
\pgfpathlineto{\pgfqpoint{4.053735in}{0.724382in}}%
\pgfpathlineto{\pgfqpoint{4.068294in}{0.703148in}}%
\pgfpathlineto{\pgfqpoint{4.078000in}{0.692245in}}%
\pgfpathlineto{\pgfqpoint{4.087706in}{0.683968in}}%
\pgfpathlineto{\pgfqpoint{4.097412in}{0.678330in}}%
\pgfpathlineto{\pgfqpoint{4.107118in}{0.675343in}}%
\pgfpathlineto{\pgfqpoint{4.116824in}{0.675011in}}%
\pgfpathlineto{\pgfqpoint{4.126530in}{0.677335in}}%
\pgfpathlineto{\pgfqpoint{4.136236in}{0.682310in}}%
\pgfpathlineto{\pgfqpoint{4.145942in}{0.689929in}}%
\pgfpathlineto{\pgfqpoint{4.155648in}{0.700177in}}%
\pgfpathlineto{\pgfqpoint{4.165354in}{0.713037in}}%
\pgfpathlineto{\pgfqpoint{4.179913in}{0.737172in}}%
\pgfpathlineto{\pgfqpoint{4.194472in}{0.767035in}}%
\pgfpathlineto{\pgfqpoint{4.209031in}{0.802506in}}%
\pgfpathlineto{\pgfqpoint{4.223591in}{0.843446in}}%
\pgfpathlineto{\pgfqpoint{4.243003in}{0.906255in}}%
\pgfpathlineto{\pgfqpoint{4.262415in}{0.978049in}}%
\pgfpathlineto{\pgfqpoint{4.281827in}{1.058323in}}%
\pgfpathlineto{\pgfqpoint{4.306092in}{1.169720in}}%
\pgfpathlineto{\pgfqpoint{4.330357in}{1.292247in}}%
\pgfpathlineto{\pgfqpoint{4.359475in}{1.452060in}}%
\pgfpathlineto{\pgfqpoint{4.393447in}{1.652904in}}%
\pgfpathlineto{\pgfqpoint{4.437124in}{1.927187in}}%
\pgfpathlineto{\pgfqpoint{4.475948in}{2.178841in}}%
\pgfpathlineto{\pgfqpoint{4.475948in}{2.178841in}}%
\pgfusepath{stroke}%
\end{pgfscope}%
\begin{pgfscope}%
\pgfpathrectangle{\pgfqpoint{2.952340in}{0.524444in}}{\pgfqpoint{1.596161in}{3.308814in}}%
\pgfusepath{clip}%
\pgfsetrectcap%
\pgfsetroundjoin%
\pgfsetlinewidth{1.505625pt}%
\definecolor{currentstroke}{rgb}{1.000000,0.498039,0.054902}%
\pgfsetstrokecolor{currentstroke}%
\pgfsetdash{}{0pt}%
\pgfpathmoveto{\pgfqpoint{3.024893in}{3.682857in}}%
\pgfpathlineto{\pgfqpoint{3.034599in}{3.681529in}}%
\pgfpathlineto{\pgfqpoint{3.044305in}{3.677547in}}%
\pgfpathlineto{\pgfqpoint{3.054011in}{3.670918in}}%
\pgfpathlineto{\pgfqpoint{3.063717in}{3.661654in}}%
\pgfpathlineto{\pgfqpoint{3.073423in}{3.649772in}}%
\pgfpathlineto{\pgfqpoint{3.083129in}{3.635291in}}%
\pgfpathlineto{\pgfqpoint{3.097688in}{3.608757in}}%
\pgfpathlineto{\pgfqpoint{3.112247in}{3.576541in}}%
\pgfpathlineto{\pgfqpoint{3.126806in}{3.538772in}}%
\pgfpathlineto{\pgfqpoint{3.141365in}{3.495601in}}%
\pgfpathlineto{\pgfqpoint{3.160778in}{3.429933in}}%
\pgfpathlineto{\pgfqpoint{3.180190in}{3.355432in}}%
\pgfpathlineto{\pgfqpoint{3.199602in}{3.272622in}}%
\pgfpathlineto{\pgfqpoint{3.223867in}{3.158322in}}%
\pgfpathlineto{\pgfqpoint{3.248132in}{3.033220in}}%
\pgfpathlineto{\pgfqpoint{3.277250in}{2.870786in}}%
\pgfpathlineto{\pgfqpoint{3.311221in}{2.667587in}}%
\pgfpathlineto{\pgfqpoint{3.354899in}{2.391462in}}%
\pgfpathlineto{\pgfqpoint{3.471371in}{1.645509in}}%
\pgfpathlineto{\pgfqpoint{3.505342in}{1.445153in}}%
\pgfpathlineto{\pgfqpoint{3.534461in}{1.285877in}}%
\pgfpathlineto{\pgfqpoint{3.558726in}{1.163875in}}%
\pgfpathlineto{\pgfqpoint{3.582991in}{1.053068in}}%
\pgfpathlineto{\pgfqpoint{3.602403in}{0.973308in}}%
\pgfpathlineto{\pgfqpoint{3.621815in}{0.902061in}}%
\pgfpathlineto{\pgfqpoint{3.641227in}{0.839830in}}%
\pgfpathlineto{\pgfqpoint{3.655786in}{0.799340in}}%
\pgfpathlineto{\pgfqpoint{3.670345in}{0.764330in}}%
\pgfpathlineto{\pgfqpoint{3.684904in}{0.734940in}}%
\pgfpathlineto{\pgfqpoint{3.699464in}{0.711287in}}%
\pgfpathlineto{\pgfqpoint{3.709170in}{0.698753in}}%
\pgfpathlineto{\pgfqpoint{3.718876in}{0.688832in}}%
\pgfpathlineto{\pgfqpoint{3.728582in}{0.681543in}}%
\pgfpathlineto{\pgfqpoint{3.738288in}{0.676899in}}%
\pgfpathlineto{\pgfqpoint{3.747994in}{0.674907in}}%
\pgfpathlineto{\pgfqpoint{3.757700in}{0.675571in}}%
\pgfpathlineto{\pgfqpoint{3.767406in}{0.678890in}}%
\pgfpathlineto{\pgfqpoint{3.777112in}{0.684858in}}%
\pgfpathlineto{\pgfqpoint{3.786818in}{0.693465in}}%
\pgfpathlineto{\pgfqpoint{3.796524in}{0.704695in}}%
\pgfpathlineto{\pgfqpoint{3.806230in}{0.718528in}}%
\pgfpathlineto{\pgfqpoint{3.820789in}{0.744105in}}%
\pgfpathlineto{\pgfqpoint{3.835348in}{0.775382in}}%
\pgfpathlineto{\pgfqpoint{3.849907in}{0.812235in}}%
\pgfpathlineto{\pgfqpoint{3.864467in}{0.854517in}}%
\pgfpathlineto{\pgfqpoint{3.883879in}{0.919046in}}%
\pgfpathlineto{\pgfqpoint{3.903291in}{0.992471in}}%
\pgfpathlineto{\pgfqpoint{3.922703in}{1.074273in}}%
\pgfpathlineto{\pgfqpoint{3.946968in}{1.187421in}}%
\pgfpathlineto{\pgfqpoint{3.971233in}{1.311504in}}%
\pgfpathlineto{\pgfqpoint{4.000351in}{1.472903in}}%
\pgfpathlineto{\pgfqpoint{4.034322in}{1.675176in}}%
\pgfpathlineto{\pgfqpoint{4.078000in}{1.950587in}}%
\pgfpathlineto{\pgfqpoint{4.199325in}{2.726918in}}%
\pgfpathlineto{\pgfqpoint{4.233297in}{2.926283in}}%
\pgfpathlineto{\pgfqpoint{4.262415in}{3.084471in}}%
\pgfpathlineto{\pgfqpoint{4.286680in}{3.205412in}}%
\pgfpathlineto{\pgfqpoint{4.310945in}{3.315030in}}%
\pgfpathlineto{\pgfqpoint{4.330357in}{3.393755in}}%
\pgfpathlineto{\pgfqpoint{4.349769in}{3.463902in}}%
\pgfpathlineto{\pgfqpoint{4.369181in}{3.524974in}}%
\pgfpathlineto{\pgfqpoint{4.383740in}{3.564561in}}%
\pgfpathlineto{\pgfqpoint{4.398300in}{3.598643in}}%
\pgfpathlineto{\pgfqpoint{4.412859in}{3.627084in}}%
\pgfpathlineto{\pgfqpoint{4.427418in}{3.649772in}}%
\pgfpathlineto{\pgfqpoint{4.437124in}{3.661654in}}%
\pgfpathlineto{\pgfqpoint{4.446830in}{3.670918in}}%
\pgfpathlineto{\pgfqpoint{4.456536in}{3.677547in}}%
\pgfpathlineto{\pgfqpoint{4.466242in}{3.681529in}}%
\pgfpathlineto{\pgfqpoint{4.475948in}{3.682857in}}%
\pgfpathlineto{\pgfqpoint{4.475948in}{3.682857in}}%
\pgfusepath{stroke}%
\end{pgfscope}%
\begin{pgfscope}%
\pgfsetrectcap%
\pgfsetmiterjoin%
\pgfsetlinewidth{0.803000pt}%
\definecolor{currentstroke}{rgb}{0.000000,0.000000,0.000000}%
\pgfsetstrokecolor{currentstroke}%
\pgfsetdash{}{0pt}%
\pgfpathmoveto{\pgfqpoint{2.952340in}{0.524444in}}%
\pgfpathlineto{\pgfqpoint{2.952340in}{3.833258in}}%
\pgfusepath{stroke}%
\end{pgfscope}%
\begin{pgfscope}%
\pgfsetrectcap%
\pgfsetmiterjoin%
\pgfsetlinewidth{0.803000pt}%
\definecolor{currentstroke}{rgb}{0.000000,0.000000,0.000000}%
\pgfsetstrokecolor{currentstroke}%
\pgfsetdash{}{0pt}%
\pgfpathmoveto{\pgfqpoint{4.548501in}{0.524444in}}%
\pgfpathlineto{\pgfqpoint{4.548501in}{3.833258in}}%
\pgfusepath{stroke}%
\end{pgfscope}%
\begin{pgfscope}%
\pgfsetrectcap%
\pgfsetmiterjoin%
\pgfsetlinewidth{0.803000pt}%
\definecolor{currentstroke}{rgb}{0.000000,0.000000,0.000000}%
\pgfsetstrokecolor{currentstroke}%
\pgfsetdash{}{0pt}%
\pgfpathmoveto{\pgfqpoint{2.952340in}{0.524444in}}%
\pgfpathlineto{\pgfqpoint{4.548501in}{0.524444in}}%
\pgfusepath{stroke}%
\end{pgfscope}%
\begin{pgfscope}%
\pgfsetrectcap%
\pgfsetmiterjoin%
\pgfsetlinewidth{0.803000pt}%
\definecolor{currentstroke}{rgb}{0.000000,0.000000,0.000000}%
\pgfsetstrokecolor{currentstroke}%
\pgfsetdash{}{0pt}%
\pgfpathmoveto{\pgfqpoint{2.952340in}{3.833258in}}%
\pgfpathlineto{\pgfqpoint{4.548501in}{3.833258in}}%
\pgfusepath{stroke}%
\end{pgfscope}%
\begin{pgfscope}%
\definecolor{textcolor}{rgb}{0.000000,0.000000,0.000000}%
\pgfsetstrokecolor{textcolor}%
\pgfsetfillcolor{textcolor}%
\pgftext[x=3.750420in,y=3.916591in,,base]{\color{textcolor}{\rmfamily\fontsize{12.000000}{14.400000}\selectfont\catcode`\^=\active\def^{\ifmmode\sp\else\^{}\fi}\catcode`\%=\active\def%{\%}Área entre dos curvas}}%
\end{pgfscope}%
\begin{pgfscope}%
\pgfsetbuttcap%
\pgfsetmiterjoin%
\definecolor{currentfill}{rgb}{1.000000,1.000000,1.000000}%
\pgfsetfillcolor{currentfill}%
\pgfsetfillopacity{0.800000}%
\pgfsetlinewidth{1.003750pt}%
\definecolor{currentstroke}{rgb}{0.800000,0.800000,0.800000}%
\pgfsetstrokecolor{currentstroke}%
\pgfsetstrokeopacity{0.800000}%
\pgfsetdash{}{0pt}%
\pgfpathmoveto{\pgfqpoint{3.015584in}{1.740518in}}%
\pgfpathlineto{\pgfqpoint{4.451279in}{1.740518in}}%
\pgfpathquadraticcurveto{\pgfqpoint{4.479056in}{1.740518in}}{\pgfqpoint{4.479056in}{1.768296in}}%
\pgfpathlineto{\pgfqpoint{4.479056in}{2.589406in}}%
\pgfpathquadraticcurveto{\pgfqpoint{4.479056in}{2.617184in}}{\pgfqpoint{4.451279in}{2.617184in}}%
\pgfpathlineto{\pgfqpoint{3.015584in}{2.617184in}}%
\pgfpathquadraticcurveto{\pgfqpoint{2.987806in}{2.617184in}}{\pgfqpoint{2.987806in}{2.589406in}}%
\pgfpathlineto{\pgfqpoint{2.987806in}{1.768296in}}%
\pgfpathquadraticcurveto{\pgfqpoint{2.987806in}{1.740518in}}{\pgfqpoint{3.015584in}{1.740518in}}%
\pgfpathlineto{\pgfqpoint{3.015584in}{1.740518in}}%
\pgfpathclose%
\pgfusepath{stroke,fill}%
\end{pgfscope}%
\begin{pgfscope}%
\pgfsetrectcap%
\pgfsetroundjoin%
\pgfsetlinewidth{1.505625pt}%
\definecolor{currentstroke}{rgb}{0.121569,0.466667,0.705882}%
\pgfsetstrokecolor{currentstroke}%
\pgfsetdash{}{0pt}%
\pgfpathmoveto{\pgfqpoint{3.043362in}{2.513017in}}%
\pgfpathlineto{\pgfqpoint{3.182251in}{2.513017in}}%
\pgfpathlineto{\pgfqpoint{3.321140in}{2.513017in}}%
\pgfusepath{stroke}%
\end{pgfscope}%
\begin{pgfscope}%
\definecolor{textcolor}{rgb}{0.000000,0.000000,0.000000}%
\pgfsetstrokecolor{textcolor}%
\pgfsetfillcolor{textcolor}%
\pgftext[x=3.432251in,y=2.464406in,left,base]{\color{textcolor}{\rmfamily\fontsize{10.000000}{12.000000}\selectfont\catcode`\^=\active\def^{\ifmmode\sp\else\^{}\fi}\catcode`\%=\active\def%{\%}Seno}}%
\end{pgfscope}%
\begin{pgfscope}%
\pgfsetrectcap%
\pgfsetroundjoin%
\pgfsetlinewidth{1.505625pt}%
\definecolor{currentstroke}{rgb}{1.000000,0.498039,0.054902}%
\pgfsetstrokecolor{currentstroke}%
\pgfsetdash{}{0pt}%
\pgfpathmoveto{\pgfqpoint{3.043362in}{2.319406in}}%
\pgfpathlineto{\pgfqpoint{3.182251in}{2.319406in}}%
\pgfpathlineto{\pgfqpoint{3.321140in}{2.319406in}}%
\pgfusepath{stroke}%
\end{pgfscope}%
\begin{pgfscope}%
\definecolor{textcolor}{rgb}{0.000000,0.000000,0.000000}%
\pgfsetstrokecolor{textcolor}%
\pgfsetfillcolor{textcolor}%
\pgftext[x=3.432251in,y=2.270795in,left,base]{\color{textcolor}{\rmfamily\fontsize{10.000000}{12.000000}\selectfont\catcode`\^=\active\def^{\ifmmode\sp\else\^{}\fi}\catcode`\%=\active\def%{\%}Coseno}}%
\end{pgfscope}%
\begin{pgfscope}%
\pgfsetbuttcap%
\pgfsetmiterjoin%
\definecolor{currentfill}{rgb}{0.121569,0.466667,0.705882}%
\pgfsetfillcolor{currentfill}%
\pgfsetfillopacity{0.250000}%
\pgfsetlinewidth{0.000000pt}%
\definecolor{currentstroke}{rgb}{0.000000,0.000000,0.000000}%
\pgfsetstrokecolor{currentstroke}%
\pgfsetstrokeopacity{0.000000}%
\pgfsetdash{}{0pt}%
\pgfpathmoveto{\pgfqpoint{3.043362in}{2.054684in}}%
\pgfpathlineto{\pgfqpoint{3.321140in}{2.054684in}}%
\pgfpathlineto{\pgfqpoint{3.321140in}{2.151906in}}%
\pgfpathlineto{\pgfqpoint{3.043362in}{2.151906in}}%
\pgfpathlineto{\pgfqpoint{3.043362in}{2.054684in}}%
\pgfpathclose%
\pgfusepath{fill}%
\end{pgfscope}%
\begin{pgfscope}%
\definecolor{textcolor}{rgb}{0.000000,0.000000,0.000000}%
\pgfsetstrokecolor{textcolor}%
\pgfsetfillcolor{textcolor}%
\pgftext[x=3.432251in,y=2.054684in,left,base]{\color{textcolor}{\rmfamily\fontsize{10.000000}{12.000000}\selectfont\catcode`\^=\active\def^{\ifmmode\sp\else\^{}\fi}\catcode`\%=\active\def%{\%}Área (y1 >=y2)}}%
\end{pgfscope}%
\begin{pgfscope}%
\pgfsetbuttcap%
\pgfsetmiterjoin%
\definecolor{currentfill}{rgb}{1.000000,0.498039,0.054902}%
\pgfsetfillcolor{currentfill}%
\pgfsetfillopacity{0.250000}%
\pgfsetlinewidth{0.000000pt}%
\definecolor{currentstroke}{rgb}{0.000000,0.000000,0.000000}%
\pgfsetstrokecolor{currentstroke}%
\pgfsetstrokeopacity{0.000000}%
\pgfsetdash{}{0pt}%
\pgfpathmoveto{\pgfqpoint{3.043362in}{1.830796in}}%
\pgfpathlineto{\pgfqpoint{3.321140in}{1.830796in}}%
\pgfpathlineto{\pgfqpoint{3.321140in}{1.928018in}}%
\pgfpathlineto{\pgfqpoint{3.043362in}{1.928018in}}%
\pgfpathlineto{\pgfqpoint{3.043362in}{1.830796in}}%
\pgfpathclose%
\pgfusepath{fill}%
\end{pgfscope}%
\begin{pgfscope}%
\definecolor{textcolor}{rgb}{0.000000,0.000000,0.000000}%
\pgfsetstrokecolor{textcolor}%
\pgfsetfillcolor{textcolor}%
\pgftext[x=3.432251in,y=1.830796in,left,base]{\color{textcolor}{\rmfamily\fontsize{10.000000}{12.000000}\selectfont\catcode`\^=\active\def^{\ifmmode\sp\else\^{}\fi}\catcode`\%=\active\def%{\%}Área (y1 < y2)}}%
\end{pgfscope}%
\begin{pgfscope}%
\definecolor{textcolor}{rgb}{0.000000,0.000000,0.000000}%
\pgfsetstrokecolor{textcolor}%
\pgfsetfillcolor{textcolor}%
\pgftext[x=2.326279in,y=4.712921in,,top]{\color{textcolor}{\rmfamily\fontsize{14.000000}{16.800000}\selectfont\catcode`\^=\active\def^{\ifmmode\sp\else\^{}\fi}\catcode`\%=\active\def%{\%}Gráfico de Áreas}}%
\end{pgfscope}%
\end{pgfpicture}%
\makeatother%
\endgroup%
 
    \caption{Gráfico de Áreas}
    \label{fig:grafico_area_curva}
\end{figure}


\subsection{Gráfico de densidad con histograma (\texttt{.hist(density=True)})}

El histograma es la representación gráfica de la distribución de un conjunto de datos numéricos. Al utilizar el parámetro \texttt{density=True}, normalizamos el gráfico: el eje vertical deja de mostrar la "frecuencia absoluta" y pasa a mostrar la "densidad de probabilidad".

\subsubsection{Conceptos clave del histograma}

\begin{itemize}
\item \textit{Bins} (Contenedores): Es la cantidad de intervalos en los que dividimos los datos. Usar \texttt{bins="auto"} es una excelente práctica, ya que \textit{Matplotlib} utiliza algoritmos estadísticos (como la regla de Sturges o Freedman-Diaconis) para calcular el ancho óptimo según la cantidad de datos.
\item \textit{Normalización:} Es lo que permite comparar el histograma con funciones matemáticas continuas. Sin \texttt{density=True}, las escalas del histograma (frecuencia) y de la curva normal (probabilidad) serían incompatibles.
\end{itemize}

\begin{pycode}{Histograma normalizado}
import numpy as np
import matplotlib.pyplot as plt

np.random.seed(42)
muestras = np.random.normal(loc=0, scale=1, size=2000)\

fig, ax = plt.subplots()

# Histograma como densidad
ax.hist(
    muestras,
    bins="auto",
    density=True, # Con esto lo normalizamos
    edgecolor="black",
    alpha=0.4,
    label="Hist (densidad)"
)

# Curva normal estimada con media y desvío de los datos
mu, sigma = np.mean(muestras), np.std(muestras, ddof=1)
x = np.linspace(mu - 4 * sigma, mu + 4 * sigma, 400)
pdf = (1 / (sigma * np.sqrt(2 * np.pi))) * np.exp(-0.5 * ((x - mu) / sigma) ** 2)

ax.plot(x, pdf, linewidth=2, label=f"N(mu={mu:.2f}, sigma={sigma:.2f})")

ax.set_title("Distribución y densidad (hist con density=True)")
ax.set_xlabel("Valor")
ax.set_ylabel("Densidad")
ax.legend()
plt.show()
\end{pycode}

Vea el histograma normalizado en la figura \ref{fig:grafico_avanzado_hist}

\begin{figure}[ht]
    \centering
    %% Creator: Matplotlib, PGF backend
%%
%% To include the figure in your LaTeX document, write
%%   \input{<filename>.pgf}
%%
%% Make sure the required packages are loaded in your preamble
%%   \usepackage{pgf}
%%
%% Also ensure that all the required font packages are loaded; for instance,
%% the lmodern package is sometimes necessary when using math font.
%%   \usepackage{lmodern}
%%
%% Figures using additional raster images can only be included by \input if
%% they are in the same directory as the main LaTeX file. For loading figures
%% from other directories you can use the `import` package
%%   \usepackage{import}
%%
%% and then include the figures with
%%   \import{<path to file>}{<filename>.pgf}
%%
%% Matplotlib used the following preamble
%%   \def\mathdefault#1{#1}
%%   \everymath=\expandafter{\the\everymath\displaystyle}
%%   \IfFileExists{scrextend.sty}{
%%     \usepackage[fontsize=10.000000pt]{scrextend}
%%   }{
%%     \renewcommand{\normalsize}{\fontsize{10.000000}{12.000000}\selectfont}
%%     \normalsize
%%   }
%%   \usepackage{fontspec}
%%   \setmainfont{CMU Serif}
%%   \makeatletter\@ifpackageloaded{underscore}{}{\usepackage[strings]{underscore}}\makeatother
%%
\begingroup%
\makeatletter%
\begin{pgfpicture}%
\pgfpathrectangle{\pgfpointorigin}{\pgfqpoint{6.045705in}{4.445587in}}%
\pgfusepath{use as bounding box, clip}%
\begin{pgfscope}%
\pgfsetbuttcap%
\pgfsetmiterjoin%
\definecolor{currentfill}{rgb}{1.000000,1.000000,1.000000}%
\pgfsetfillcolor{currentfill}%
\pgfsetlinewidth{0.000000pt}%
\definecolor{currentstroke}{rgb}{1.000000,1.000000,1.000000}%
\pgfsetstrokecolor{currentstroke}%
\pgfsetdash{}{0pt}%
\pgfpathmoveto{\pgfqpoint{0.000000in}{0.000000in}}%
\pgfpathlineto{\pgfqpoint{6.045705in}{0.000000in}}%
\pgfpathlineto{\pgfqpoint{6.045705in}{4.445587in}}%
\pgfpathlineto{\pgfqpoint{0.000000in}{4.445587in}}%
\pgfpathlineto{\pgfqpoint{0.000000in}{0.000000in}}%
\pgfpathclose%
\pgfusepath{fill}%
\end{pgfscope}%
\begin{pgfscope}%
\pgfsetbuttcap%
\pgfsetmiterjoin%
\definecolor{currentfill}{rgb}{1.000000,1.000000,1.000000}%
\pgfsetfillcolor{currentfill}%
\pgfsetlinewidth{0.000000pt}%
\definecolor{currentstroke}{rgb}{0.000000,0.000000,0.000000}%
\pgfsetstrokecolor{currentstroke}%
\pgfsetstrokeopacity{0.000000}%
\pgfsetdash{}{0pt}%
\pgfpathmoveto{\pgfqpoint{0.623025in}{0.499444in}}%
\pgfpathlineto{\pgfqpoint{5.945705in}{0.499444in}}%
\pgfpathlineto{\pgfqpoint{5.945705in}{4.137253in}}%
\pgfpathlineto{\pgfqpoint{0.623025in}{4.137253in}}%
\pgfpathlineto{\pgfqpoint{0.623025in}{0.499444in}}%
\pgfpathclose%
\pgfusepath{fill}%
\end{pgfscope}%
\begin{pgfscope}%
\pgfpathrectangle{\pgfqpoint{0.623025in}{0.499444in}}{\pgfqpoint{5.322679in}{3.637809in}}%
\pgfusepath{clip}%
\pgfsetbuttcap%
\pgfsetmiterjoin%
\definecolor{currentfill}{rgb}{0.121569,0.466667,0.705882}%
\pgfsetfillcolor{currentfill}%
\pgfsetfillopacity{0.400000}%
\pgfsetlinewidth{1.003750pt}%
\definecolor{currentstroke}{rgb}{0.000000,0.000000,0.000000}%
\pgfsetstrokecolor{currentstroke}%
\pgfsetstrokeopacity{0.400000}%
\pgfsetdash{}{0pt}%
\pgfpathmoveto{\pgfqpoint{1.273421in}{0.499444in}}%
\pgfpathlineto{\pgfqpoint{1.397446in}{0.499444in}}%
\pgfpathlineto{\pgfqpoint{1.397446in}{0.519705in}}%
\pgfpathlineto{\pgfqpoint{1.273421in}{0.519705in}}%
\pgfpathlineto{\pgfqpoint{1.273421in}{0.499444in}}%
\pgfpathclose%
\pgfusepath{stroke,fill}%
\end{pgfscope}%
\begin{pgfscope}%
\pgfpathrectangle{\pgfqpoint{0.623025in}{0.499444in}}{\pgfqpoint{5.322679in}{3.637809in}}%
\pgfusepath{clip}%
\pgfsetbuttcap%
\pgfsetmiterjoin%
\definecolor{currentfill}{rgb}{0.121569,0.466667,0.705882}%
\pgfsetfillcolor{currentfill}%
\pgfsetfillopacity{0.400000}%
\pgfsetlinewidth{1.003750pt}%
\definecolor{currentstroke}{rgb}{0.000000,0.000000,0.000000}%
\pgfsetstrokecolor{currentstroke}%
\pgfsetstrokeopacity{0.400000}%
\pgfsetdash{}{0pt}%
\pgfpathmoveto{\pgfqpoint{1.397446in}{0.499444in}}%
\pgfpathlineto{\pgfqpoint{1.521470in}{0.499444in}}%
\pgfpathlineto{\pgfqpoint{1.521470in}{0.600748in}}%
\pgfpathlineto{\pgfqpoint{1.397446in}{0.600748in}}%
\pgfpathlineto{\pgfqpoint{1.397446in}{0.499444in}}%
\pgfpathclose%
\pgfusepath{stroke,fill}%
\end{pgfscope}%
\begin{pgfscope}%
\pgfpathrectangle{\pgfqpoint{0.623025in}{0.499444in}}{\pgfqpoint{5.322679in}{3.637809in}}%
\pgfusepath{clip}%
\pgfsetbuttcap%
\pgfsetmiterjoin%
\definecolor{currentfill}{rgb}{0.121569,0.466667,0.705882}%
\pgfsetfillcolor{currentfill}%
\pgfsetfillopacity{0.400000}%
\pgfsetlinewidth{1.003750pt}%
\definecolor{currentstroke}{rgb}{0.000000,0.000000,0.000000}%
\pgfsetstrokecolor{currentstroke}%
\pgfsetstrokeopacity{0.400000}%
\pgfsetdash{}{0pt}%
\pgfpathmoveto{\pgfqpoint{1.521470in}{0.499444in}}%
\pgfpathlineto{\pgfqpoint{1.645495in}{0.499444in}}%
\pgfpathlineto{\pgfqpoint{1.645495in}{0.560226in}}%
\pgfpathlineto{\pgfqpoint{1.521470in}{0.560226in}}%
\pgfpathlineto{\pgfqpoint{1.521470in}{0.499444in}}%
\pgfpathclose%
\pgfusepath{stroke,fill}%
\end{pgfscope}%
\begin{pgfscope}%
\pgfpathrectangle{\pgfqpoint{0.623025in}{0.499444in}}{\pgfqpoint{5.322679in}{3.637809in}}%
\pgfusepath{clip}%
\pgfsetbuttcap%
\pgfsetmiterjoin%
\definecolor{currentfill}{rgb}{0.121569,0.466667,0.705882}%
\pgfsetfillcolor{currentfill}%
\pgfsetfillopacity{0.400000}%
\pgfsetlinewidth{1.003750pt}%
\definecolor{currentstroke}{rgb}{0.000000,0.000000,0.000000}%
\pgfsetstrokecolor{currentstroke}%
\pgfsetstrokeopacity{0.400000}%
\pgfsetdash{}{0pt}%
\pgfpathmoveto{\pgfqpoint{1.645495in}{0.499444in}}%
\pgfpathlineto{\pgfqpoint{1.769520in}{0.499444in}}%
\pgfpathlineto{\pgfqpoint{1.769520in}{0.621008in}}%
\pgfpathlineto{\pgfqpoint{1.645495in}{0.621008in}}%
\pgfpathlineto{\pgfqpoint{1.645495in}{0.499444in}}%
\pgfpathclose%
\pgfusepath{stroke,fill}%
\end{pgfscope}%
\begin{pgfscope}%
\pgfpathrectangle{\pgfqpoint{0.623025in}{0.499444in}}{\pgfqpoint{5.322679in}{3.637809in}}%
\pgfusepath{clip}%
\pgfsetbuttcap%
\pgfsetmiterjoin%
\definecolor{currentfill}{rgb}{0.121569,0.466667,0.705882}%
\pgfsetfillcolor{currentfill}%
\pgfsetfillopacity{0.400000}%
\pgfsetlinewidth{1.003750pt}%
\definecolor{currentstroke}{rgb}{0.000000,0.000000,0.000000}%
\pgfsetstrokecolor{currentstroke}%
\pgfsetstrokeopacity{0.400000}%
\pgfsetdash{}{0pt}%
\pgfpathmoveto{\pgfqpoint{1.769520in}{0.499444in}}%
\pgfpathlineto{\pgfqpoint{1.893545in}{0.499444in}}%
\pgfpathlineto{\pgfqpoint{1.893545in}{0.580487in}}%
\pgfpathlineto{\pgfqpoint{1.769520in}{0.580487in}}%
\pgfpathlineto{\pgfqpoint{1.769520in}{0.499444in}}%
\pgfpathclose%
\pgfusepath{stroke,fill}%
\end{pgfscope}%
\begin{pgfscope}%
\pgfpathrectangle{\pgfqpoint{0.623025in}{0.499444in}}{\pgfqpoint{5.322679in}{3.637809in}}%
\pgfusepath{clip}%
\pgfsetbuttcap%
\pgfsetmiterjoin%
\definecolor{currentfill}{rgb}{0.121569,0.466667,0.705882}%
\pgfsetfillcolor{currentfill}%
\pgfsetfillopacity{0.400000}%
\pgfsetlinewidth{1.003750pt}%
\definecolor{currentstroke}{rgb}{0.000000,0.000000,0.000000}%
\pgfsetstrokecolor{currentstroke}%
\pgfsetstrokeopacity{0.400000}%
\pgfsetdash{}{0pt}%
\pgfpathmoveto{\pgfqpoint{1.893545in}{0.499444in}}%
\pgfpathlineto{\pgfqpoint{2.017570in}{0.499444in}}%
\pgfpathlineto{\pgfqpoint{2.017570in}{0.803355in}}%
\pgfpathlineto{\pgfqpoint{1.893545in}{0.803355in}}%
\pgfpathlineto{\pgfqpoint{1.893545in}{0.499444in}}%
\pgfpathclose%
\pgfusepath{stroke,fill}%
\end{pgfscope}%
\begin{pgfscope}%
\pgfpathrectangle{\pgfqpoint{0.623025in}{0.499444in}}{\pgfqpoint{5.322679in}{3.637809in}}%
\pgfusepath{clip}%
\pgfsetbuttcap%
\pgfsetmiterjoin%
\definecolor{currentfill}{rgb}{0.121569,0.466667,0.705882}%
\pgfsetfillcolor{currentfill}%
\pgfsetfillopacity{0.400000}%
\pgfsetlinewidth{1.003750pt}%
\definecolor{currentstroke}{rgb}{0.000000,0.000000,0.000000}%
\pgfsetstrokecolor{currentstroke}%
\pgfsetstrokeopacity{0.400000}%
\pgfsetdash{}{0pt}%
\pgfpathmoveto{\pgfqpoint{2.017570in}{0.499444in}}%
\pgfpathlineto{\pgfqpoint{2.141595in}{0.499444in}}%
\pgfpathlineto{\pgfqpoint{2.141595in}{0.904658in}}%
\pgfpathlineto{\pgfqpoint{2.017570in}{0.904658in}}%
\pgfpathlineto{\pgfqpoint{2.017570in}{0.499444in}}%
\pgfpathclose%
\pgfusepath{stroke,fill}%
\end{pgfscope}%
\begin{pgfscope}%
\pgfpathrectangle{\pgfqpoint{0.623025in}{0.499444in}}{\pgfqpoint{5.322679in}{3.637809in}}%
\pgfusepath{clip}%
\pgfsetbuttcap%
\pgfsetmiterjoin%
\definecolor{currentfill}{rgb}{0.121569,0.466667,0.705882}%
\pgfsetfillcolor{currentfill}%
\pgfsetfillopacity{0.400000}%
\pgfsetlinewidth{1.003750pt}%
\definecolor{currentstroke}{rgb}{0.000000,0.000000,0.000000}%
\pgfsetstrokecolor{currentstroke}%
\pgfsetstrokeopacity{0.400000}%
\pgfsetdash{}{0pt}%
\pgfpathmoveto{\pgfqpoint{2.141595in}{0.499444in}}%
\pgfpathlineto{\pgfqpoint{2.265620in}{0.499444in}}%
\pgfpathlineto{\pgfqpoint{2.265620in}{1.087005in}}%
\pgfpathlineto{\pgfqpoint{2.141595in}{1.087005in}}%
\pgfpathlineto{\pgfqpoint{2.141595in}{0.499444in}}%
\pgfpathclose%
\pgfusepath{stroke,fill}%
\end{pgfscope}%
\begin{pgfscope}%
\pgfpathrectangle{\pgfqpoint{0.623025in}{0.499444in}}{\pgfqpoint{5.322679in}{3.637809in}}%
\pgfusepath{clip}%
\pgfsetbuttcap%
\pgfsetmiterjoin%
\definecolor{currentfill}{rgb}{0.121569,0.466667,0.705882}%
\pgfsetfillcolor{currentfill}%
\pgfsetfillopacity{0.400000}%
\pgfsetlinewidth{1.003750pt}%
\definecolor{currentstroke}{rgb}{0.000000,0.000000,0.000000}%
\pgfsetstrokecolor{currentstroke}%
\pgfsetstrokeopacity{0.400000}%
\pgfsetdash{}{0pt}%
\pgfpathmoveto{\pgfqpoint{2.265620in}{0.499444in}}%
\pgfpathlineto{\pgfqpoint{2.389645in}{0.499444in}}%
\pgfpathlineto{\pgfqpoint{2.389645in}{1.593522in}}%
\pgfpathlineto{\pgfqpoint{2.265620in}{1.593522in}}%
\pgfpathlineto{\pgfqpoint{2.265620in}{0.499444in}}%
\pgfpathclose%
\pgfusepath{stroke,fill}%
\end{pgfscope}%
\begin{pgfscope}%
\pgfpathrectangle{\pgfqpoint{0.623025in}{0.499444in}}{\pgfqpoint{5.322679in}{3.637809in}}%
\pgfusepath{clip}%
\pgfsetbuttcap%
\pgfsetmiterjoin%
\definecolor{currentfill}{rgb}{0.121569,0.466667,0.705882}%
\pgfsetfillcolor{currentfill}%
\pgfsetfillopacity{0.400000}%
\pgfsetlinewidth{1.003750pt}%
\definecolor{currentstroke}{rgb}{0.000000,0.000000,0.000000}%
\pgfsetstrokecolor{currentstroke}%
\pgfsetstrokeopacity{0.400000}%
\pgfsetdash{}{0pt}%
\pgfpathmoveto{\pgfqpoint{2.389645in}{0.499444in}}%
\pgfpathlineto{\pgfqpoint{2.513670in}{0.499444in}}%
\pgfpathlineto{\pgfqpoint{2.513670in}{1.816390in}}%
\pgfpathlineto{\pgfqpoint{2.389645in}{1.816390in}}%
\pgfpathlineto{\pgfqpoint{2.389645in}{0.499444in}}%
\pgfpathclose%
\pgfusepath{stroke,fill}%
\end{pgfscope}%
\begin{pgfscope}%
\pgfpathrectangle{\pgfqpoint{0.623025in}{0.499444in}}{\pgfqpoint{5.322679in}{3.637809in}}%
\pgfusepath{clip}%
\pgfsetbuttcap%
\pgfsetmiterjoin%
\definecolor{currentfill}{rgb}{0.121569,0.466667,0.705882}%
\pgfsetfillcolor{currentfill}%
\pgfsetfillopacity{0.400000}%
\pgfsetlinewidth{1.003750pt}%
\definecolor{currentstroke}{rgb}{0.000000,0.000000,0.000000}%
\pgfsetstrokecolor{currentstroke}%
\pgfsetstrokeopacity{0.400000}%
\pgfsetdash{}{0pt}%
\pgfpathmoveto{\pgfqpoint{2.513670in}{0.499444in}}%
\pgfpathlineto{\pgfqpoint{2.637695in}{0.499444in}}%
\pgfpathlineto{\pgfqpoint{2.637695in}{2.039258in}}%
\pgfpathlineto{\pgfqpoint{2.513670in}{2.039258in}}%
\pgfpathlineto{\pgfqpoint{2.513670in}{0.499444in}}%
\pgfpathclose%
\pgfusepath{stroke,fill}%
\end{pgfscope}%
\begin{pgfscope}%
\pgfpathrectangle{\pgfqpoint{0.623025in}{0.499444in}}{\pgfqpoint{5.322679in}{3.637809in}}%
\pgfusepath{clip}%
\pgfsetbuttcap%
\pgfsetmiterjoin%
\definecolor{currentfill}{rgb}{0.121569,0.466667,0.705882}%
\pgfsetfillcolor{currentfill}%
\pgfsetfillopacity{0.400000}%
\pgfsetlinewidth{1.003750pt}%
\definecolor{currentstroke}{rgb}{0.000000,0.000000,0.000000}%
\pgfsetstrokecolor{currentstroke}%
\pgfsetstrokeopacity{0.400000}%
\pgfsetdash{}{0pt}%
\pgfpathmoveto{\pgfqpoint{2.637695in}{0.499444in}}%
\pgfpathlineto{\pgfqpoint{2.761720in}{0.499444in}}%
\pgfpathlineto{\pgfqpoint{2.761720in}{2.566036in}}%
\pgfpathlineto{\pgfqpoint{2.637695in}{2.566036in}}%
\pgfpathlineto{\pgfqpoint{2.637695in}{0.499444in}}%
\pgfpathclose%
\pgfusepath{stroke,fill}%
\end{pgfscope}%
\begin{pgfscope}%
\pgfpathrectangle{\pgfqpoint{0.623025in}{0.499444in}}{\pgfqpoint{5.322679in}{3.637809in}}%
\pgfusepath{clip}%
\pgfsetbuttcap%
\pgfsetmiterjoin%
\definecolor{currentfill}{rgb}{0.121569,0.466667,0.705882}%
\pgfsetfillcolor{currentfill}%
\pgfsetfillopacity{0.400000}%
\pgfsetlinewidth{1.003750pt}%
\definecolor{currentstroke}{rgb}{0.000000,0.000000,0.000000}%
\pgfsetstrokecolor{currentstroke}%
\pgfsetstrokeopacity{0.400000}%
\pgfsetdash{}{0pt}%
\pgfpathmoveto{\pgfqpoint{2.761720in}{0.499444in}}%
\pgfpathlineto{\pgfqpoint{2.885745in}{0.499444in}}%
\pgfpathlineto{\pgfqpoint{2.885745in}{3.133336in}}%
\pgfpathlineto{\pgfqpoint{2.761720in}{3.133336in}}%
\pgfpathlineto{\pgfqpoint{2.761720in}{0.499444in}}%
\pgfpathclose%
\pgfusepath{stroke,fill}%
\end{pgfscope}%
\begin{pgfscope}%
\pgfpathrectangle{\pgfqpoint{0.623025in}{0.499444in}}{\pgfqpoint{5.322679in}{3.637809in}}%
\pgfusepath{clip}%
\pgfsetbuttcap%
\pgfsetmiterjoin%
\definecolor{currentfill}{rgb}{0.121569,0.466667,0.705882}%
\pgfsetfillcolor{currentfill}%
\pgfsetfillopacity{0.400000}%
\pgfsetlinewidth{1.003750pt}%
\definecolor{currentstroke}{rgb}{0.000000,0.000000,0.000000}%
\pgfsetstrokecolor{currentstroke}%
\pgfsetstrokeopacity{0.400000}%
\pgfsetdash{}{0pt}%
\pgfpathmoveto{\pgfqpoint{2.885745in}{0.499444in}}%
\pgfpathlineto{\pgfqpoint{3.009770in}{0.499444in}}%
\pgfpathlineto{\pgfqpoint{3.009770in}{3.457507in}}%
\pgfpathlineto{\pgfqpoint{2.885745in}{3.457507in}}%
\pgfpathlineto{\pgfqpoint{2.885745in}{0.499444in}}%
\pgfpathclose%
\pgfusepath{stroke,fill}%
\end{pgfscope}%
\begin{pgfscope}%
\pgfpathrectangle{\pgfqpoint{0.623025in}{0.499444in}}{\pgfqpoint{5.322679in}{3.637809in}}%
\pgfusepath{clip}%
\pgfsetbuttcap%
\pgfsetmiterjoin%
\definecolor{currentfill}{rgb}{0.121569,0.466667,0.705882}%
\pgfsetfillcolor{currentfill}%
\pgfsetfillopacity{0.400000}%
\pgfsetlinewidth{1.003750pt}%
\definecolor{currentstroke}{rgb}{0.000000,0.000000,0.000000}%
\pgfsetstrokecolor{currentstroke}%
\pgfsetstrokeopacity{0.400000}%
\pgfsetdash{}{0pt}%
\pgfpathmoveto{\pgfqpoint{3.009770in}{0.499444in}}%
\pgfpathlineto{\pgfqpoint{3.133795in}{0.499444in}}%
\pgfpathlineto{\pgfqpoint{3.133795in}{3.498028in}}%
\pgfpathlineto{\pgfqpoint{3.009770in}{3.498028in}}%
\pgfpathlineto{\pgfqpoint{3.009770in}{0.499444in}}%
\pgfpathclose%
\pgfusepath{stroke,fill}%
\end{pgfscope}%
\begin{pgfscope}%
\pgfpathrectangle{\pgfqpoint{0.623025in}{0.499444in}}{\pgfqpoint{5.322679in}{3.637809in}}%
\pgfusepath{clip}%
\pgfsetbuttcap%
\pgfsetmiterjoin%
\definecolor{currentfill}{rgb}{0.121569,0.466667,0.705882}%
\pgfsetfillcolor{currentfill}%
\pgfsetfillopacity{0.400000}%
\pgfsetlinewidth{1.003750pt}%
\definecolor{currentstroke}{rgb}{0.000000,0.000000,0.000000}%
\pgfsetstrokecolor{currentstroke}%
\pgfsetstrokeopacity{0.400000}%
\pgfsetdash{}{0pt}%
\pgfpathmoveto{\pgfqpoint{3.133795in}{0.499444in}}%
\pgfpathlineto{\pgfqpoint{3.257820in}{0.499444in}}%
\pgfpathlineto{\pgfqpoint{3.257820in}{3.801939in}}%
\pgfpathlineto{\pgfqpoint{3.133795in}{3.801939in}}%
\pgfpathlineto{\pgfqpoint{3.133795in}{0.499444in}}%
\pgfpathclose%
\pgfusepath{stroke,fill}%
\end{pgfscope}%
\begin{pgfscope}%
\pgfpathrectangle{\pgfqpoint{0.623025in}{0.499444in}}{\pgfqpoint{5.322679in}{3.637809in}}%
\pgfusepath{clip}%
\pgfsetbuttcap%
\pgfsetmiterjoin%
\definecolor{currentfill}{rgb}{0.121569,0.466667,0.705882}%
\pgfsetfillcolor{currentfill}%
\pgfsetfillopacity{0.400000}%
\pgfsetlinewidth{1.003750pt}%
\definecolor{currentstroke}{rgb}{0.000000,0.000000,0.000000}%
\pgfsetstrokecolor{currentstroke}%
\pgfsetstrokeopacity{0.400000}%
\pgfsetdash{}{0pt}%
\pgfpathmoveto{\pgfqpoint{3.257820in}{0.499444in}}%
\pgfpathlineto{\pgfqpoint{3.381845in}{0.499444in}}%
\pgfpathlineto{\pgfqpoint{3.381845in}{3.964024in}}%
\pgfpathlineto{\pgfqpoint{3.257820in}{3.964024in}}%
\pgfpathlineto{\pgfqpoint{3.257820in}{0.499444in}}%
\pgfpathclose%
\pgfusepath{stroke,fill}%
\end{pgfscope}%
\begin{pgfscope}%
\pgfpathrectangle{\pgfqpoint{0.623025in}{0.499444in}}{\pgfqpoint{5.322679in}{3.637809in}}%
\pgfusepath{clip}%
\pgfsetbuttcap%
\pgfsetmiterjoin%
\definecolor{currentfill}{rgb}{0.121569,0.466667,0.705882}%
\pgfsetfillcolor{currentfill}%
\pgfsetfillopacity{0.400000}%
\pgfsetlinewidth{1.003750pt}%
\definecolor{currentstroke}{rgb}{0.000000,0.000000,0.000000}%
\pgfsetstrokecolor{currentstroke}%
\pgfsetstrokeopacity{0.400000}%
\pgfsetdash{}{0pt}%
\pgfpathmoveto{\pgfqpoint{3.381845in}{0.499444in}}%
\pgfpathlineto{\pgfqpoint{3.505870in}{0.499444in}}%
\pgfpathlineto{\pgfqpoint{3.505870in}{3.741157in}}%
\pgfpathlineto{\pgfqpoint{3.381845in}{3.741157in}}%
\pgfpathlineto{\pgfqpoint{3.381845in}{0.499444in}}%
\pgfpathclose%
\pgfusepath{stroke,fill}%
\end{pgfscope}%
\begin{pgfscope}%
\pgfpathrectangle{\pgfqpoint{0.623025in}{0.499444in}}{\pgfqpoint{5.322679in}{3.637809in}}%
\pgfusepath{clip}%
\pgfsetbuttcap%
\pgfsetmiterjoin%
\definecolor{currentfill}{rgb}{0.121569,0.466667,0.705882}%
\pgfsetfillcolor{currentfill}%
\pgfsetfillopacity{0.400000}%
\pgfsetlinewidth{1.003750pt}%
\definecolor{currentstroke}{rgb}{0.000000,0.000000,0.000000}%
\pgfsetstrokecolor{currentstroke}%
\pgfsetstrokeopacity{0.400000}%
\pgfsetdash{}{0pt}%
\pgfpathmoveto{\pgfqpoint{3.505870in}{0.499444in}}%
\pgfpathlineto{\pgfqpoint{3.629895in}{0.499444in}}%
\pgfpathlineto{\pgfqpoint{3.629895in}{3.437246in}}%
\pgfpathlineto{\pgfqpoint{3.505870in}{3.437246in}}%
\pgfpathlineto{\pgfqpoint{3.505870in}{0.499444in}}%
\pgfpathclose%
\pgfusepath{stroke,fill}%
\end{pgfscope}%
\begin{pgfscope}%
\pgfpathrectangle{\pgfqpoint{0.623025in}{0.499444in}}{\pgfqpoint{5.322679in}{3.637809in}}%
\pgfusepath{clip}%
\pgfsetbuttcap%
\pgfsetmiterjoin%
\definecolor{currentfill}{rgb}{0.121569,0.466667,0.705882}%
\pgfsetfillcolor{currentfill}%
\pgfsetfillopacity{0.400000}%
\pgfsetlinewidth{1.003750pt}%
\definecolor{currentstroke}{rgb}{0.000000,0.000000,0.000000}%
\pgfsetstrokecolor{currentstroke}%
\pgfsetstrokeopacity{0.400000}%
\pgfsetdash{}{0pt}%
\pgfpathmoveto{\pgfqpoint{3.629895in}{0.499444in}}%
\pgfpathlineto{\pgfqpoint{3.753920in}{0.499444in}}%
\pgfpathlineto{\pgfqpoint{3.753920in}{3.052293in}}%
\pgfpathlineto{\pgfqpoint{3.629895in}{3.052293in}}%
\pgfpathlineto{\pgfqpoint{3.629895in}{0.499444in}}%
\pgfpathclose%
\pgfusepath{stroke,fill}%
\end{pgfscope}%
\begin{pgfscope}%
\pgfpathrectangle{\pgfqpoint{0.623025in}{0.499444in}}{\pgfqpoint{5.322679in}{3.637809in}}%
\pgfusepath{clip}%
\pgfsetbuttcap%
\pgfsetmiterjoin%
\definecolor{currentfill}{rgb}{0.121569,0.466667,0.705882}%
\pgfsetfillcolor{currentfill}%
\pgfsetfillopacity{0.400000}%
\pgfsetlinewidth{1.003750pt}%
\definecolor{currentstroke}{rgb}{0.000000,0.000000,0.000000}%
\pgfsetstrokecolor{currentstroke}%
\pgfsetstrokeopacity{0.400000}%
\pgfsetdash{}{0pt}%
\pgfpathmoveto{\pgfqpoint{3.753920in}{0.499444in}}%
\pgfpathlineto{\pgfqpoint{3.877945in}{0.499444in}}%
\pgfpathlineto{\pgfqpoint{3.877945in}{2.647079in}}%
\pgfpathlineto{\pgfqpoint{3.753920in}{2.647079in}}%
\pgfpathlineto{\pgfqpoint{3.753920in}{0.499444in}}%
\pgfpathclose%
\pgfusepath{stroke,fill}%
\end{pgfscope}%
\begin{pgfscope}%
\pgfpathrectangle{\pgfqpoint{0.623025in}{0.499444in}}{\pgfqpoint{5.322679in}{3.637809in}}%
\pgfusepath{clip}%
\pgfsetbuttcap%
\pgfsetmiterjoin%
\definecolor{currentfill}{rgb}{0.121569,0.466667,0.705882}%
\pgfsetfillcolor{currentfill}%
\pgfsetfillopacity{0.400000}%
\pgfsetlinewidth{1.003750pt}%
\definecolor{currentstroke}{rgb}{0.000000,0.000000,0.000000}%
\pgfsetstrokecolor{currentstroke}%
\pgfsetstrokeopacity{0.400000}%
\pgfsetdash{}{0pt}%
\pgfpathmoveto{\pgfqpoint{3.877945in}{0.499444in}}%
\pgfpathlineto{\pgfqpoint{4.001970in}{0.499444in}}%
\pgfpathlineto{\pgfqpoint{4.001970in}{2.424211in}}%
\pgfpathlineto{\pgfqpoint{3.877945in}{2.424211in}}%
\pgfpathlineto{\pgfqpoint{3.877945in}{0.499444in}}%
\pgfpathclose%
\pgfusepath{stroke,fill}%
\end{pgfscope}%
\begin{pgfscope}%
\pgfpathrectangle{\pgfqpoint{0.623025in}{0.499444in}}{\pgfqpoint{5.322679in}{3.637809in}}%
\pgfusepath{clip}%
\pgfsetbuttcap%
\pgfsetmiterjoin%
\definecolor{currentfill}{rgb}{0.121569,0.466667,0.705882}%
\pgfsetfillcolor{currentfill}%
\pgfsetfillopacity{0.400000}%
\pgfsetlinewidth{1.003750pt}%
\definecolor{currentstroke}{rgb}{0.000000,0.000000,0.000000}%
\pgfsetstrokecolor{currentstroke}%
\pgfsetstrokeopacity{0.400000}%
\pgfsetdash{}{0pt}%
\pgfpathmoveto{\pgfqpoint{4.001970in}{0.499444in}}%
\pgfpathlineto{\pgfqpoint{4.125995in}{0.499444in}}%
\pgfpathlineto{\pgfqpoint{4.125995in}{1.634044in}}%
\pgfpathlineto{\pgfqpoint{4.001970in}{1.634044in}}%
\pgfpathlineto{\pgfqpoint{4.001970in}{0.499444in}}%
\pgfpathclose%
\pgfusepath{stroke,fill}%
\end{pgfscope}%
\begin{pgfscope}%
\pgfpathrectangle{\pgfqpoint{0.623025in}{0.499444in}}{\pgfqpoint{5.322679in}{3.637809in}}%
\pgfusepath{clip}%
\pgfsetbuttcap%
\pgfsetmiterjoin%
\definecolor{currentfill}{rgb}{0.121569,0.466667,0.705882}%
\pgfsetfillcolor{currentfill}%
\pgfsetfillopacity{0.400000}%
\pgfsetlinewidth{1.003750pt}%
\definecolor{currentstroke}{rgb}{0.000000,0.000000,0.000000}%
\pgfsetstrokecolor{currentstroke}%
\pgfsetstrokeopacity{0.400000}%
\pgfsetdash{}{0pt}%
\pgfpathmoveto{\pgfqpoint{4.125995in}{0.499444in}}%
\pgfpathlineto{\pgfqpoint{4.250020in}{0.499444in}}%
\pgfpathlineto{\pgfqpoint{4.250020in}{1.634044in}}%
\pgfpathlineto{\pgfqpoint{4.125995in}{1.634044in}}%
\pgfpathlineto{\pgfqpoint{4.125995in}{0.499444in}}%
\pgfpathclose%
\pgfusepath{stroke,fill}%
\end{pgfscope}%
\begin{pgfscope}%
\pgfpathrectangle{\pgfqpoint{0.623025in}{0.499444in}}{\pgfqpoint{5.322679in}{3.637809in}}%
\pgfusepath{clip}%
\pgfsetbuttcap%
\pgfsetmiterjoin%
\definecolor{currentfill}{rgb}{0.121569,0.466667,0.705882}%
\pgfsetfillcolor{currentfill}%
\pgfsetfillopacity{0.400000}%
\pgfsetlinewidth{1.003750pt}%
\definecolor{currentstroke}{rgb}{0.000000,0.000000,0.000000}%
\pgfsetstrokecolor{currentstroke}%
\pgfsetstrokeopacity{0.400000}%
\pgfsetdash{}{0pt}%
\pgfpathmoveto{\pgfqpoint{4.250020in}{0.499444in}}%
\pgfpathlineto{\pgfqpoint{4.374045in}{0.499444in}}%
\pgfpathlineto{\pgfqpoint{4.374045in}{1.309872in}}%
\pgfpathlineto{\pgfqpoint{4.250020in}{1.309872in}}%
\pgfpathlineto{\pgfqpoint{4.250020in}{0.499444in}}%
\pgfpathclose%
\pgfusepath{stroke,fill}%
\end{pgfscope}%
\begin{pgfscope}%
\pgfpathrectangle{\pgfqpoint{0.623025in}{0.499444in}}{\pgfqpoint{5.322679in}{3.637809in}}%
\pgfusepath{clip}%
\pgfsetbuttcap%
\pgfsetmiterjoin%
\definecolor{currentfill}{rgb}{0.121569,0.466667,0.705882}%
\pgfsetfillcolor{currentfill}%
\pgfsetfillopacity{0.400000}%
\pgfsetlinewidth{1.003750pt}%
\definecolor{currentstroke}{rgb}{0.000000,0.000000,0.000000}%
\pgfsetstrokecolor{currentstroke}%
\pgfsetstrokeopacity{0.400000}%
\pgfsetdash{}{0pt}%
\pgfpathmoveto{\pgfqpoint{4.374045in}{0.499444in}}%
\pgfpathlineto{\pgfqpoint{4.498070in}{0.499444in}}%
\pgfpathlineto{\pgfqpoint{4.498070in}{1.046483in}}%
\pgfpathlineto{\pgfqpoint{4.374045in}{1.046483in}}%
\pgfpathlineto{\pgfqpoint{4.374045in}{0.499444in}}%
\pgfpathclose%
\pgfusepath{stroke,fill}%
\end{pgfscope}%
\begin{pgfscope}%
\pgfpathrectangle{\pgfqpoint{0.623025in}{0.499444in}}{\pgfqpoint{5.322679in}{3.637809in}}%
\pgfusepath{clip}%
\pgfsetbuttcap%
\pgfsetmiterjoin%
\definecolor{currentfill}{rgb}{0.121569,0.466667,0.705882}%
\pgfsetfillcolor{currentfill}%
\pgfsetfillopacity{0.400000}%
\pgfsetlinewidth{1.003750pt}%
\definecolor{currentstroke}{rgb}{0.000000,0.000000,0.000000}%
\pgfsetstrokecolor{currentstroke}%
\pgfsetstrokeopacity{0.400000}%
\pgfsetdash{}{0pt}%
\pgfpathmoveto{\pgfqpoint{4.498070in}{0.499444in}}%
\pgfpathlineto{\pgfqpoint{4.622095in}{0.499444in}}%
\pgfpathlineto{\pgfqpoint{4.622095in}{0.985701in}}%
\pgfpathlineto{\pgfqpoint{4.498070in}{0.985701in}}%
\pgfpathlineto{\pgfqpoint{4.498070in}{0.499444in}}%
\pgfpathclose%
\pgfusepath{stroke,fill}%
\end{pgfscope}%
\begin{pgfscope}%
\pgfpathrectangle{\pgfqpoint{0.623025in}{0.499444in}}{\pgfqpoint{5.322679in}{3.637809in}}%
\pgfusepath{clip}%
\pgfsetbuttcap%
\pgfsetmiterjoin%
\definecolor{currentfill}{rgb}{0.121569,0.466667,0.705882}%
\pgfsetfillcolor{currentfill}%
\pgfsetfillopacity{0.400000}%
\pgfsetlinewidth{1.003750pt}%
\definecolor{currentstroke}{rgb}{0.000000,0.000000,0.000000}%
\pgfsetstrokecolor{currentstroke}%
\pgfsetstrokeopacity{0.400000}%
\pgfsetdash{}{0pt}%
\pgfpathmoveto{\pgfqpoint{4.622095in}{0.499444in}}%
\pgfpathlineto{\pgfqpoint{4.746120in}{0.499444in}}%
\pgfpathlineto{\pgfqpoint{4.746120in}{0.661530in}}%
\pgfpathlineto{\pgfqpoint{4.622095in}{0.661530in}}%
\pgfpathlineto{\pgfqpoint{4.622095in}{0.499444in}}%
\pgfpathclose%
\pgfusepath{stroke,fill}%
\end{pgfscope}%
\begin{pgfscope}%
\pgfpathrectangle{\pgfqpoint{0.623025in}{0.499444in}}{\pgfqpoint{5.322679in}{3.637809in}}%
\pgfusepath{clip}%
\pgfsetbuttcap%
\pgfsetmiterjoin%
\definecolor{currentfill}{rgb}{0.121569,0.466667,0.705882}%
\pgfsetfillcolor{currentfill}%
\pgfsetfillopacity{0.400000}%
\pgfsetlinewidth{1.003750pt}%
\definecolor{currentstroke}{rgb}{0.000000,0.000000,0.000000}%
\pgfsetstrokecolor{currentstroke}%
\pgfsetstrokeopacity{0.400000}%
\pgfsetdash{}{0pt}%
\pgfpathmoveto{\pgfqpoint{4.746120in}{0.499444in}}%
\pgfpathlineto{\pgfqpoint{4.870145in}{0.499444in}}%
\pgfpathlineto{\pgfqpoint{4.870145in}{0.762833in}}%
\pgfpathlineto{\pgfqpoint{4.746120in}{0.762833in}}%
\pgfpathlineto{\pgfqpoint{4.746120in}{0.499444in}}%
\pgfpathclose%
\pgfusepath{stroke,fill}%
\end{pgfscope}%
\begin{pgfscope}%
\pgfpathrectangle{\pgfqpoint{0.623025in}{0.499444in}}{\pgfqpoint{5.322679in}{3.637809in}}%
\pgfusepath{clip}%
\pgfsetbuttcap%
\pgfsetmiterjoin%
\definecolor{currentfill}{rgb}{0.121569,0.466667,0.705882}%
\pgfsetfillcolor{currentfill}%
\pgfsetfillopacity{0.400000}%
\pgfsetlinewidth{1.003750pt}%
\definecolor{currentstroke}{rgb}{0.000000,0.000000,0.000000}%
\pgfsetstrokecolor{currentstroke}%
\pgfsetstrokeopacity{0.400000}%
\pgfsetdash{}{0pt}%
\pgfpathmoveto{\pgfqpoint{4.870145in}{0.499444in}}%
\pgfpathlineto{\pgfqpoint{4.994169in}{0.499444in}}%
\pgfpathlineto{\pgfqpoint{4.994169in}{0.539966in}}%
\pgfpathlineto{\pgfqpoint{4.870145in}{0.539966in}}%
\pgfpathlineto{\pgfqpoint{4.870145in}{0.499444in}}%
\pgfpathclose%
\pgfusepath{stroke,fill}%
\end{pgfscope}%
\begin{pgfscope}%
\pgfpathrectangle{\pgfqpoint{0.623025in}{0.499444in}}{\pgfqpoint{5.322679in}{3.637809in}}%
\pgfusepath{clip}%
\pgfsetbuttcap%
\pgfsetmiterjoin%
\definecolor{currentfill}{rgb}{0.121569,0.466667,0.705882}%
\pgfsetfillcolor{currentfill}%
\pgfsetfillopacity{0.400000}%
\pgfsetlinewidth{1.003750pt}%
\definecolor{currentstroke}{rgb}{0.000000,0.000000,0.000000}%
\pgfsetstrokecolor{currentstroke}%
\pgfsetstrokeopacity{0.400000}%
\pgfsetdash{}{0pt}%
\pgfpathmoveto{\pgfqpoint{4.994169in}{0.499444in}}%
\pgfpathlineto{\pgfqpoint{5.118194in}{0.499444in}}%
\pgfpathlineto{\pgfqpoint{5.118194in}{0.499444in}}%
\pgfpathlineto{\pgfqpoint{4.994169in}{0.499444in}}%
\pgfpathlineto{\pgfqpoint{4.994169in}{0.499444in}}%
\pgfpathclose%
\pgfusepath{stroke,fill}%
\end{pgfscope}%
\begin{pgfscope}%
\pgfpathrectangle{\pgfqpoint{0.623025in}{0.499444in}}{\pgfqpoint{5.322679in}{3.637809in}}%
\pgfusepath{clip}%
\pgfsetbuttcap%
\pgfsetmiterjoin%
\definecolor{currentfill}{rgb}{0.121569,0.466667,0.705882}%
\pgfsetfillcolor{currentfill}%
\pgfsetfillopacity{0.400000}%
\pgfsetlinewidth{1.003750pt}%
\definecolor{currentstroke}{rgb}{0.000000,0.000000,0.000000}%
\pgfsetstrokecolor{currentstroke}%
\pgfsetstrokeopacity{0.400000}%
\pgfsetdash{}{0pt}%
\pgfpathmoveto{\pgfqpoint{5.118194in}{0.499444in}}%
\pgfpathlineto{\pgfqpoint{5.242219in}{0.499444in}}%
\pgfpathlineto{\pgfqpoint{5.242219in}{0.560226in}}%
\pgfpathlineto{\pgfqpoint{5.118194in}{0.560226in}}%
\pgfpathlineto{\pgfqpoint{5.118194in}{0.499444in}}%
\pgfpathclose%
\pgfusepath{stroke,fill}%
\end{pgfscope}%
\begin{pgfscope}%
\pgfpathrectangle{\pgfqpoint{0.623025in}{0.499444in}}{\pgfqpoint{5.322679in}{3.637809in}}%
\pgfusepath{clip}%
\pgfsetbuttcap%
\pgfsetmiterjoin%
\definecolor{currentfill}{rgb}{0.121569,0.466667,0.705882}%
\pgfsetfillcolor{currentfill}%
\pgfsetfillopacity{0.400000}%
\pgfsetlinewidth{1.003750pt}%
\definecolor{currentstroke}{rgb}{0.000000,0.000000,0.000000}%
\pgfsetstrokecolor{currentstroke}%
\pgfsetstrokeopacity{0.400000}%
\pgfsetdash{}{0pt}%
\pgfpathmoveto{\pgfqpoint{5.242219in}{0.499444in}}%
\pgfpathlineto{\pgfqpoint{5.366244in}{0.499444in}}%
\pgfpathlineto{\pgfqpoint{5.366244in}{0.499444in}}%
\pgfpathlineto{\pgfqpoint{5.242219in}{0.499444in}}%
\pgfpathlineto{\pgfqpoint{5.242219in}{0.499444in}}%
\pgfpathclose%
\pgfusepath{stroke,fill}%
\end{pgfscope}%
\begin{pgfscope}%
\pgfpathrectangle{\pgfqpoint{0.623025in}{0.499444in}}{\pgfqpoint{5.322679in}{3.637809in}}%
\pgfusepath{clip}%
\pgfsetbuttcap%
\pgfsetmiterjoin%
\definecolor{currentfill}{rgb}{0.121569,0.466667,0.705882}%
\pgfsetfillcolor{currentfill}%
\pgfsetfillopacity{0.400000}%
\pgfsetlinewidth{1.003750pt}%
\definecolor{currentstroke}{rgb}{0.000000,0.000000,0.000000}%
\pgfsetstrokecolor{currentstroke}%
\pgfsetstrokeopacity{0.400000}%
\pgfsetdash{}{0pt}%
\pgfpathmoveto{\pgfqpoint{5.366244in}{0.499444in}}%
\pgfpathlineto{\pgfqpoint{5.490269in}{0.499444in}}%
\pgfpathlineto{\pgfqpoint{5.490269in}{0.499444in}}%
\pgfpathlineto{\pgfqpoint{5.366244in}{0.499444in}}%
\pgfpathlineto{\pgfqpoint{5.366244in}{0.499444in}}%
\pgfpathclose%
\pgfusepath{stroke,fill}%
\end{pgfscope}%
\begin{pgfscope}%
\pgfpathrectangle{\pgfqpoint{0.623025in}{0.499444in}}{\pgfqpoint{5.322679in}{3.637809in}}%
\pgfusepath{clip}%
\pgfsetbuttcap%
\pgfsetmiterjoin%
\definecolor{currentfill}{rgb}{0.121569,0.466667,0.705882}%
\pgfsetfillcolor{currentfill}%
\pgfsetfillopacity{0.400000}%
\pgfsetlinewidth{1.003750pt}%
\definecolor{currentstroke}{rgb}{0.000000,0.000000,0.000000}%
\pgfsetstrokecolor{currentstroke}%
\pgfsetstrokeopacity{0.400000}%
\pgfsetdash{}{0pt}%
\pgfpathmoveto{\pgfqpoint{5.490269in}{0.499444in}}%
\pgfpathlineto{\pgfqpoint{5.614294in}{0.499444in}}%
\pgfpathlineto{\pgfqpoint{5.614294in}{0.519705in}}%
\pgfpathlineto{\pgfqpoint{5.490269in}{0.519705in}}%
\pgfpathlineto{\pgfqpoint{5.490269in}{0.499444in}}%
\pgfpathclose%
\pgfusepath{stroke,fill}%
\end{pgfscope}%
\begin{pgfscope}%
\pgfpathrectangle{\pgfqpoint{0.623025in}{0.499444in}}{\pgfqpoint{5.322679in}{3.637809in}}%
\pgfusepath{clip}%
\pgfsetbuttcap%
\pgfsetroundjoin%
\pgfsetlinewidth{0.501875pt}%
\definecolor{currentstroke}{rgb}{0.501961,0.501961,0.501961}%
\pgfsetstrokecolor{currentstroke}%
\pgfsetstrokeopacity{0.500000}%
\pgfsetdash{{1.850000pt}{0.800000pt}}{0.000000pt}%
\pgfpathmoveto{\pgfqpoint{0.809146in}{0.499444in}}%
\pgfpathlineto{\pgfqpoint{0.809146in}{4.137253in}}%
\pgfusepath{stroke}%
\end{pgfscope}%
\begin{pgfscope}%
\pgfsetbuttcap%
\pgfsetroundjoin%
\definecolor{currentfill}{rgb}{0.000000,0.000000,0.000000}%
\pgfsetfillcolor{currentfill}%
\pgfsetlinewidth{0.803000pt}%
\definecolor{currentstroke}{rgb}{0.000000,0.000000,0.000000}%
\pgfsetstrokecolor{currentstroke}%
\pgfsetdash{}{0pt}%
\pgfsys@defobject{currentmarker}{\pgfqpoint{0.000000in}{-0.048611in}}{\pgfqpoint{0.000000in}{0.000000in}}{%
\pgfpathmoveto{\pgfqpoint{0.000000in}{0.000000in}}%
\pgfpathlineto{\pgfqpoint{0.000000in}{-0.048611in}}%
\pgfusepath{stroke,fill}%
}%
\begin{pgfscope}%
\pgfsys@transformshift{0.809146in}{0.499444in}%
\pgfsys@useobject{currentmarker}{}%
\end{pgfscope}%
\end{pgfscope}%
\begin{pgfscope}%
\definecolor{textcolor}{rgb}{0.000000,0.000000,0.000000}%
\pgfsetstrokecolor{textcolor}%
\pgfsetfillcolor{textcolor}%
\pgftext[x=0.809146in,y=0.402222in,,top]{\color{textcolor}{\rmfamily\fontsize{10.000000}{12.000000}\selectfont\catcode`\^=\active\def^{\ifmmode\sp\else\^{}\fi}\catcode`\%=\active\def%{\%}$\mathdefault{\ensuremath{-}4}$}}%
\end{pgfscope}%
\begin{pgfscope}%
\pgfpathrectangle{\pgfqpoint{0.623025in}{0.499444in}}{\pgfqpoint{5.322679in}{3.637809in}}%
\pgfusepath{clip}%
\pgfsetbuttcap%
\pgfsetroundjoin%
\pgfsetlinewidth{0.501875pt}%
\definecolor{currentstroke}{rgb}{0.501961,0.501961,0.501961}%
\pgfsetstrokecolor{currentstroke}%
\pgfsetstrokeopacity{0.500000}%
\pgfsetdash{{1.850000pt}{0.800000pt}}{0.000000pt}%
\pgfpathmoveto{\pgfqpoint{1.421054in}{0.499444in}}%
\pgfpathlineto{\pgfqpoint{1.421054in}{4.137253in}}%
\pgfusepath{stroke}%
\end{pgfscope}%
\begin{pgfscope}%
\pgfsetbuttcap%
\pgfsetroundjoin%
\definecolor{currentfill}{rgb}{0.000000,0.000000,0.000000}%
\pgfsetfillcolor{currentfill}%
\pgfsetlinewidth{0.803000pt}%
\definecolor{currentstroke}{rgb}{0.000000,0.000000,0.000000}%
\pgfsetstrokecolor{currentstroke}%
\pgfsetdash{}{0pt}%
\pgfsys@defobject{currentmarker}{\pgfqpoint{0.000000in}{-0.048611in}}{\pgfqpoint{0.000000in}{0.000000in}}{%
\pgfpathmoveto{\pgfqpoint{0.000000in}{0.000000in}}%
\pgfpathlineto{\pgfqpoint{0.000000in}{-0.048611in}}%
\pgfusepath{stroke,fill}%
}%
\begin{pgfscope}%
\pgfsys@transformshift{1.421054in}{0.499444in}%
\pgfsys@useobject{currentmarker}{}%
\end{pgfscope}%
\end{pgfscope}%
\begin{pgfscope}%
\definecolor{textcolor}{rgb}{0.000000,0.000000,0.000000}%
\pgfsetstrokecolor{textcolor}%
\pgfsetfillcolor{textcolor}%
\pgftext[x=1.421054in,y=0.402222in,,top]{\color{textcolor}{\rmfamily\fontsize{10.000000}{12.000000}\selectfont\catcode`\^=\active\def^{\ifmmode\sp\else\^{}\fi}\catcode`\%=\active\def%{\%}$\mathdefault{\ensuremath{-}3}$}}%
\end{pgfscope}%
\begin{pgfscope}%
\pgfpathrectangle{\pgfqpoint{0.623025in}{0.499444in}}{\pgfqpoint{5.322679in}{3.637809in}}%
\pgfusepath{clip}%
\pgfsetbuttcap%
\pgfsetroundjoin%
\pgfsetlinewidth{0.501875pt}%
\definecolor{currentstroke}{rgb}{0.501961,0.501961,0.501961}%
\pgfsetstrokecolor{currentstroke}%
\pgfsetstrokeopacity{0.500000}%
\pgfsetdash{{1.850000pt}{0.800000pt}}{0.000000pt}%
\pgfpathmoveto{\pgfqpoint{2.032962in}{0.499444in}}%
\pgfpathlineto{\pgfqpoint{2.032962in}{4.137253in}}%
\pgfusepath{stroke}%
\end{pgfscope}%
\begin{pgfscope}%
\pgfsetbuttcap%
\pgfsetroundjoin%
\definecolor{currentfill}{rgb}{0.000000,0.000000,0.000000}%
\pgfsetfillcolor{currentfill}%
\pgfsetlinewidth{0.803000pt}%
\definecolor{currentstroke}{rgb}{0.000000,0.000000,0.000000}%
\pgfsetstrokecolor{currentstroke}%
\pgfsetdash{}{0pt}%
\pgfsys@defobject{currentmarker}{\pgfqpoint{0.000000in}{-0.048611in}}{\pgfqpoint{0.000000in}{0.000000in}}{%
\pgfpathmoveto{\pgfqpoint{0.000000in}{0.000000in}}%
\pgfpathlineto{\pgfqpoint{0.000000in}{-0.048611in}}%
\pgfusepath{stroke,fill}%
}%
\begin{pgfscope}%
\pgfsys@transformshift{2.032962in}{0.499444in}%
\pgfsys@useobject{currentmarker}{}%
\end{pgfscope}%
\end{pgfscope}%
\begin{pgfscope}%
\definecolor{textcolor}{rgb}{0.000000,0.000000,0.000000}%
\pgfsetstrokecolor{textcolor}%
\pgfsetfillcolor{textcolor}%
\pgftext[x=2.032962in,y=0.402222in,,top]{\color{textcolor}{\rmfamily\fontsize{10.000000}{12.000000}\selectfont\catcode`\^=\active\def^{\ifmmode\sp\else\^{}\fi}\catcode`\%=\active\def%{\%}$\mathdefault{\ensuremath{-}2}$}}%
\end{pgfscope}%
\begin{pgfscope}%
\pgfpathrectangle{\pgfqpoint{0.623025in}{0.499444in}}{\pgfqpoint{5.322679in}{3.637809in}}%
\pgfusepath{clip}%
\pgfsetbuttcap%
\pgfsetroundjoin%
\pgfsetlinewidth{0.501875pt}%
\definecolor{currentstroke}{rgb}{0.501961,0.501961,0.501961}%
\pgfsetstrokecolor{currentstroke}%
\pgfsetstrokeopacity{0.500000}%
\pgfsetdash{{1.850000pt}{0.800000pt}}{0.000000pt}%
\pgfpathmoveto{\pgfqpoint{2.644870in}{0.499444in}}%
\pgfpathlineto{\pgfqpoint{2.644870in}{4.137253in}}%
\pgfusepath{stroke}%
\end{pgfscope}%
\begin{pgfscope}%
\pgfsetbuttcap%
\pgfsetroundjoin%
\definecolor{currentfill}{rgb}{0.000000,0.000000,0.000000}%
\pgfsetfillcolor{currentfill}%
\pgfsetlinewidth{0.803000pt}%
\definecolor{currentstroke}{rgb}{0.000000,0.000000,0.000000}%
\pgfsetstrokecolor{currentstroke}%
\pgfsetdash{}{0pt}%
\pgfsys@defobject{currentmarker}{\pgfqpoint{0.000000in}{-0.048611in}}{\pgfqpoint{0.000000in}{0.000000in}}{%
\pgfpathmoveto{\pgfqpoint{0.000000in}{0.000000in}}%
\pgfpathlineto{\pgfqpoint{0.000000in}{-0.048611in}}%
\pgfusepath{stroke,fill}%
}%
\begin{pgfscope}%
\pgfsys@transformshift{2.644870in}{0.499444in}%
\pgfsys@useobject{currentmarker}{}%
\end{pgfscope}%
\end{pgfscope}%
\begin{pgfscope}%
\definecolor{textcolor}{rgb}{0.000000,0.000000,0.000000}%
\pgfsetstrokecolor{textcolor}%
\pgfsetfillcolor{textcolor}%
\pgftext[x=2.644870in,y=0.402222in,,top]{\color{textcolor}{\rmfamily\fontsize{10.000000}{12.000000}\selectfont\catcode`\^=\active\def^{\ifmmode\sp\else\^{}\fi}\catcode`\%=\active\def%{\%}$\mathdefault{\ensuremath{-}1}$}}%
\end{pgfscope}%
\begin{pgfscope}%
\pgfpathrectangle{\pgfqpoint{0.623025in}{0.499444in}}{\pgfqpoint{5.322679in}{3.637809in}}%
\pgfusepath{clip}%
\pgfsetbuttcap%
\pgfsetroundjoin%
\pgfsetlinewidth{0.501875pt}%
\definecolor{currentstroke}{rgb}{0.501961,0.501961,0.501961}%
\pgfsetstrokecolor{currentstroke}%
\pgfsetstrokeopacity{0.500000}%
\pgfsetdash{{1.850000pt}{0.800000pt}}{0.000000pt}%
\pgfpathmoveto{\pgfqpoint{3.256778in}{0.499444in}}%
\pgfpathlineto{\pgfqpoint{3.256778in}{4.137253in}}%
\pgfusepath{stroke}%
\end{pgfscope}%
\begin{pgfscope}%
\pgfsetbuttcap%
\pgfsetroundjoin%
\definecolor{currentfill}{rgb}{0.000000,0.000000,0.000000}%
\pgfsetfillcolor{currentfill}%
\pgfsetlinewidth{0.803000pt}%
\definecolor{currentstroke}{rgb}{0.000000,0.000000,0.000000}%
\pgfsetstrokecolor{currentstroke}%
\pgfsetdash{}{0pt}%
\pgfsys@defobject{currentmarker}{\pgfqpoint{0.000000in}{-0.048611in}}{\pgfqpoint{0.000000in}{0.000000in}}{%
\pgfpathmoveto{\pgfqpoint{0.000000in}{0.000000in}}%
\pgfpathlineto{\pgfqpoint{0.000000in}{-0.048611in}}%
\pgfusepath{stroke,fill}%
}%
\begin{pgfscope}%
\pgfsys@transformshift{3.256778in}{0.499444in}%
\pgfsys@useobject{currentmarker}{}%
\end{pgfscope}%
\end{pgfscope}%
\begin{pgfscope}%
\definecolor{textcolor}{rgb}{0.000000,0.000000,0.000000}%
\pgfsetstrokecolor{textcolor}%
\pgfsetfillcolor{textcolor}%
\pgftext[x=3.256778in,y=0.402222in,,top]{\color{textcolor}{\rmfamily\fontsize{10.000000}{12.000000}\selectfont\catcode`\^=\active\def^{\ifmmode\sp\else\^{}\fi}\catcode`\%=\active\def%{\%}$\mathdefault{0}$}}%
\end{pgfscope}%
\begin{pgfscope}%
\pgfpathrectangle{\pgfqpoint{0.623025in}{0.499444in}}{\pgfqpoint{5.322679in}{3.637809in}}%
\pgfusepath{clip}%
\pgfsetbuttcap%
\pgfsetroundjoin%
\pgfsetlinewidth{0.501875pt}%
\definecolor{currentstroke}{rgb}{0.501961,0.501961,0.501961}%
\pgfsetstrokecolor{currentstroke}%
\pgfsetstrokeopacity{0.500000}%
\pgfsetdash{{1.850000pt}{0.800000pt}}{0.000000pt}%
\pgfpathmoveto{\pgfqpoint{3.868685in}{0.499444in}}%
\pgfpathlineto{\pgfqpoint{3.868685in}{4.137253in}}%
\pgfusepath{stroke}%
\end{pgfscope}%
\begin{pgfscope}%
\pgfsetbuttcap%
\pgfsetroundjoin%
\definecolor{currentfill}{rgb}{0.000000,0.000000,0.000000}%
\pgfsetfillcolor{currentfill}%
\pgfsetlinewidth{0.803000pt}%
\definecolor{currentstroke}{rgb}{0.000000,0.000000,0.000000}%
\pgfsetstrokecolor{currentstroke}%
\pgfsetdash{}{0pt}%
\pgfsys@defobject{currentmarker}{\pgfqpoint{0.000000in}{-0.048611in}}{\pgfqpoint{0.000000in}{0.000000in}}{%
\pgfpathmoveto{\pgfqpoint{0.000000in}{0.000000in}}%
\pgfpathlineto{\pgfqpoint{0.000000in}{-0.048611in}}%
\pgfusepath{stroke,fill}%
}%
\begin{pgfscope}%
\pgfsys@transformshift{3.868685in}{0.499444in}%
\pgfsys@useobject{currentmarker}{}%
\end{pgfscope}%
\end{pgfscope}%
\begin{pgfscope}%
\definecolor{textcolor}{rgb}{0.000000,0.000000,0.000000}%
\pgfsetstrokecolor{textcolor}%
\pgfsetfillcolor{textcolor}%
\pgftext[x=3.868685in,y=0.402222in,,top]{\color{textcolor}{\rmfamily\fontsize{10.000000}{12.000000}\selectfont\catcode`\^=\active\def^{\ifmmode\sp\else\^{}\fi}\catcode`\%=\active\def%{\%}$\mathdefault{1}$}}%
\end{pgfscope}%
\begin{pgfscope}%
\pgfpathrectangle{\pgfqpoint{0.623025in}{0.499444in}}{\pgfqpoint{5.322679in}{3.637809in}}%
\pgfusepath{clip}%
\pgfsetbuttcap%
\pgfsetroundjoin%
\pgfsetlinewidth{0.501875pt}%
\definecolor{currentstroke}{rgb}{0.501961,0.501961,0.501961}%
\pgfsetstrokecolor{currentstroke}%
\pgfsetstrokeopacity{0.500000}%
\pgfsetdash{{1.850000pt}{0.800000pt}}{0.000000pt}%
\pgfpathmoveto{\pgfqpoint{4.480593in}{0.499444in}}%
\pgfpathlineto{\pgfqpoint{4.480593in}{4.137253in}}%
\pgfusepath{stroke}%
\end{pgfscope}%
\begin{pgfscope}%
\pgfsetbuttcap%
\pgfsetroundjoin%
\definecolor{currentfill}{rgb}{0.000000,0.000000,0.000000}%
\pgfsetfillcolor{currentfill}%
\pgfsetlinewidth{0.803000pt}%
\definecolor{currentstroke}{rgb}{0.000000,0.000000,0.000000}%
\pgfsetstrokecolor{currentstroke}%
\pgfsetdash{}{0pt}%
\pgfsys@defobject{currentmarker}{\pgfqpoint{0.000000in}{-0.048611in}}{\pgfqpoint{0.000000in}{0.000000in}}{%
\pgfpathmoveto{\pgfqpoint{0.000000in}{0.000000in}}%
\pgfpathlineto{\pgfqpoint{0.000000in}{-0.048611in}}%
\pgfusepath{stroke,fill}%
}%
\begin{pgfscope}%
\pgfsys@transformshift{4.480593in}{0.499444in}%
\pgfsys@useobject{currentmarker}{}%
\end{pgfscope}%
\end{pgfscope}%
\begin{pgfscope}%
\definecolor{textcolor}{rgb}{0.000000,0.000000,0.000000}%
\pgfsetstrokecolor{textcolor}%
\pgfsetfillcolor{textcolor}%
\pgftext[x=4.480593in,y=0.402222in,,top]{\color{textcolor}{\rmfamily\fontsize{10.000000}{12.000000}\selectfont\catcode`\^=\active\def^{\ifmmode\sp\else\^{}\fi}\catcode`\%=\active\def%{\%}$\mathdefault{2}$}}%
\end{pgfscope}%
\begin{pgfscope}%
\pgfpathrectangle{\pgfqpoint{0.623025in}{0.499444in}}{\pgfqpoint{5.322679in}{3.637809in}}%
\pgfusepath{clip}%
\pgfsetbuttcap%
\pgfsetroundjoin%
\pgfsetlinewidth{0.501875pt}%
\definecolor{currentstroke}{rgb}{0.501961,0.501961,0.501961}%
\pgfsetstrokecolor{currentstroke}%
\pgfsetstrokeopacity{0.500000}%
\pgfsetdash{{1.850000pt}{0.800000pt}}{0.000000pt}%
\pgfpathmoveto{\pgfqpoint{5.092501in}{0.499444in}}%
\pgfpathlineto{\pgfqpoint{5.092501in}{4.137253in}}%
\pgfusepath{stroke}%
\end{pgfscope}%
\begin{pgfscope}%
\pgfsetbuttcap%
\pgfsetroundjoin%
\definecolor{currentfill}{rgb}{0.000000,0.000000,0.000000}%
\pgfsetfillcolor{currentfill}%
\pgfsetlinewidth{0.803000pt}%
\definecolor{currentstroke}{rgb}{0.000000,0.000000,0.000000}%
\pgfsetstrokecolor{currentstroke}%
\pgfsetdash{}{0pt}%
\pgfsys@defobject{currentmarker}{\pgfqpoint{0.000000in}{-0.048611in}}{\pgfqpoint{0.000000in}{0.000000in}}{%
\pgfpathmoveto{\pgfqpoint{0.000000in}{0.000000in}}%
\pgfpathlineto{\pgfqpoint{0.000000in}{-0.048611in}}%
\pgfusepath{stroke,fill}%
}%
\begin{pgfscope}%
\pgfsys@transformshift{5.092501in}{0.499444in}%
\pgfsys@useobject{currentmarker}{}%
\end{pgfscope}%
\end{pgfscope}%
\begin{pgfscope}%
\definecolor{textcolor}{rgb}{0.000000,0.000000,0.000000}%
\pgfsetstrokecolor{textcolor}%
\pgfsetfillcolor{textcolor}%
\pgftext[x=5.092501in,y=0.402222in,,top]{\color{textcolor}{\rmfamily\fontsize{10.000000}{12.000000}\selectfont\catcode`\^=\active\def^{\ifmmode\sp\else\^{}\fi}\catcode`\%=\active\def%{\%}$\mathdefault{3}$}}%
\end{pgfscope}%
\begin{pgfscope}%
\pgfpathrectangle{\pgfqpoint{0.623025in}{0.499444in}}{\pgfqpoint{5.322679in}{3.637809in}}%
\pgfusepath{clip}%
\pgfsetbuttcap%
\pgfsetroundjoin%
\pgfsetlinewidth{0.501875pt}%
\definecolor{currentstroke}{rgb}{0.501961,0.501961,0.501961}%
\pgfsetstrokecolor{currentstroke}%
\pgfsetstrokeopacity{0.500000}%
\pgfsetdash{{1.850000pt}{0.800000pt}}{0.000000pt}%
\pgfpathmoveto{\pgfqpoint{5.704409in}{0.499444in}}%
\pgfpathlineto{\pgfqpoint{5.704409in}{4.137253in}}%
\pgfusepath{stroke}%
\end{pgfscope}%
\begin{pgfscope}%
\pgfsetbuttcap%
\pgfsetroundjoin%
\definecolor{currentfill}{rgb}{0.000000,0.000000,0.000000}%
\pgfsetfillcolor{currentfill}%
\pgfsetlinewidth{0.803000pt}%
\definecolor{currentstroke}{rgb}{0.000000,0.000000,0.000000}%
\pgfsetstrokecolor{currentstroke}%
\pgfsetdash{}{0pt}%
\pgfsys@defobject{currentmarker}{\pgfqpoint{0.000000in}{-0.048611in}}{\pgfqpoint{0.000000in}{0.000000in}}{%
\pgfpathmoveto{\pgfqpoint{0.000000in}{0.000000in}}%
\pgfpathlineto{\pgfqpoint{0.000000in}{-0.048611in}}%
\pgfusepath{stroke,fill}%
}%
\begin{pgfscope}%
\pgfsys@transformshift{5.704409in}{0.499444in}%
\pgfsys@useobject{currentmarker}{}%
\end{pgfscope}%
\end{pgfscope}%
\begin{pgfscope}%
\definecolor{textcolor}{rgb}{0.000000,0.000000,0.000000}%
\pgfsetstrokecolor{textcolor}%
\pgfsetfillcolor{textcolor}%
\pgftext[x=5.704409in,y=0.402222in,,top]{\color{textcolor}{\rmfamily\fontsize{10.000000}{12.000000}\selectfont\catcode`\^=\active\def^{\ifmmode\sp\else\^{}\fi}\catcode`\%=\active\def%{\%}$\mathdefault{4}$}}%
\end{pgfscope}%
\begin{pgfscope}%
\definecolor{textcolor}{rgb}{0.000000,0.000000,0.000000}%
\pgfsetstrokecolor{textcolor}%
\pgfsetfillcolor{textcolor}%
\pgftext[x=3.284365in,y=0.223333in,,top]{\color{textcolor}{\rmfamily\fontsize{10.000000}{12.000000}\selectfont\catcode`\^=\active\def^{\ifmmode\sp\else\^{}\fi}\catcode`\%=\active\def%{\%}Valor}}%
\end{pgfscope}%
\begin{pgfscope}%
\pgfpathrectangle{\pgfqpoint{0.623025in}{0.499444in}}{\pgfqpoint{5.322679in}{3.637809in}}%
\pgfusepath{clip}%
\pgfsetbuttcap%
\pgfsetroundjoin%
\pgfsetlinewidth{0.501875pt}%
\definecolor{currentstroke}{rgb}{0.501961,0.501961,0.501961}%
\pgfsetstrokecolor{currentstroke}%
\pgfsetstrokeopacity{0.500000}%
\pgfsetdash{{1.850000pt}{0.800000pt}}{0.000000pt}%
\pgfpathmoveto{\pgfqpoint{0.623025in}{0.499444in}}%
\pgfpathlineto{\pgfqpoint{5.945705in}{0.499444in}}%
\pgfusepath{stroke}%
\end{pgfscope}%
\begin{pgfscope}%
\pgfsetbuttcap%
\pgfsetroundjoin%
\definecolor{currentfill}{rgb}{0.000000,0.000000,0.000000}%
\pgfsetfillcolor{currentfill}%
\pgfsetlinewidth{0.803000pt}%
\definecolor{currentstroke}{rgb}{0.000000,0.000000,0.000000}%
\pgfsetstrokecolor{currentstroke}%
\pgfsetdash{}{0pt}%
\pgfsys@defobject{currentmarker}{\pgfqpoint{-0.048611in}{0.000000in}}{\pgfqpoint{-0.000000in}{0.000000in}}{%
\pgfpathmoveto{\pgfqpoint{-0.000000in}{0.000000in}}%
\pgfpathlineto{\pgfqpoint{-0.048611in}{0.000000in}}%
\pgfusepath{stroke,fill}%
}%
\begin{pgfscope}%
\pgfsys@transformshift{0.623025in}{0.499444in}%
\pgfsys@useobject{currentmarker}{}%
\end{pgfscope}%
\end{pgfscope}%
\begin{pgfscope}%
\definecolor{textcolor}{rgb}{0.000000,0.000000,0.000000}%
\pgfsetstrokecolor{textcolor}%
\pgfsetfillcolor{textcolor}%
\pgftext[x=0.278889in, y=0.451250in, left, base]{\color{textcolor}{\rmfamily\fontsize{10.000000}{12.000000}\selectfont\catcode`\^=\active\def^{\ifmmode\sp\else\^{}\fi}\catcode`\%=\active\def%{\%}$\mathdefault{0.00}$}}%
\end{pgfscope}%
\begin{pgfscope}%
\pgfpathrectangle{\pgfqpoint{0.623025in}{0.499444in}}{\pgfqpoint{5.322679in}{3.637809in}}%
\pgfusepath{clip}%
\pgfsetbuttcap%
\pgfsetroundjoin%
\pgfsetlinewidth{0.501875pt}%
\definecolor{currentstroke}{rgb}{0.501961,0.501961,0.501961}%
\pgfsetstrokecolor{currentstroke}%
\pgfsetstrokeopacity{0.500000}%
\pgfsetdash{{1.850000pt}{0.800000pt}}{0.000000pt}%
\pgfpathmoveto{\pgfqpoint{0.623025in}{0.910100in}}%
\pgfpathlineto{\pgfqpoint{5.945705in}{0.910100in}}%
\pgfusepath{stroke}%
\end{pgfscope}%
\begin{pgfscope}%
\pgfsetbuttcap%
\pgfsetroundjoin%
\definecolor{currentfill}{rgb}{0.000000,0.000000,0.000000}%
\pgfsetfillcolor{currentfill}%
\pgfsetlinewidth{0.803000pt}%
\definecolor{currentstroke}{rgb}{0.000000,0.000000,0.000000}%
\pgfsetstrokecolor{currentstroke}%
\pgfsetdash{}{0pt}%
\pgfsys@defobject{currentmarker}{\pgfqpoint{-0.048611in}{0.000000in}}{\pgfqpoint{-0.000000in}{0.000000in}}{%
\pgfpathmoveto{\pgfqpoint{-0.000000in}{0.000000in}}%
\pgfpathlineto{\pgfqpoint{-0.048611in}{0.000000in}}%
\pgfusepath{stroke,fill}%
}%
\begin{pgfscope}%
\pgfsys@transformshift{0.623025in}{0.910100in}%
\pgfsys@useobject{currentmarker}{}%
\end{pgfscope}%
\end{pgfscope}%
\begin{pgfscope}%
\definecolor{textcolor}{rgb}{0.000000,0.000000,0.000000}%
\pgfsetstrokecolor{textcolor}%
\pgfsetfillcolor{textcolor}%
\pgftext[x=0.278889in, y=0.861905in, left, base]{\color{textcolor}{\rmfamily\fontsize{10.000000}{12.000000}\selectfont\catcode`\^=\active\def^{\ifmmode\sp\else\^{}\fi}\catcode`\%=\active\def%{\%}$\mathdefault{0.05}$}}%
\end{pgfscope}%
\begin{pgfscope}%
\pgfpathrectangle{\pgfqpoint{0.623025in}{0.499444in}}{\pgfqpoint{5.322679in}{3.637809in}}%
\pgfusepath{clip}%
\pgfsetbuttcap%
\pgfsetroundjoin%
\pgfsetlinewidth{0.501875pt}%
\definecolor{currentstroke}{rgb}{0.501961,0.501961,0.501961}%
\pgfsetstrokecolor{currentstroke}%
\pgfsetstrokeopacity{0.500000}%
\pgfsetdash{{1.850000pt}{0.800000pt}}{0.000000pt}%
\pgfpathmoveto{\pgfqpoint{0.623025in}{1.320755in}}%
\pgfpathlineto{\pgfqpoint{5.945705in}{1.320755in}}%
\pgfusepath{stroke}%
\end{pgfscope}%
\begin{pgfscope}%
\pgfsetbuttcap%
\pgfsetroundjoin%
\definecolor{currentfill}{rgb}{0.000000,0.000000,0.000000}%
\pgfsetfillcolor{currentfill}%
\pgfsetlinewidth{0.803000pt}%
\definecolor{currentstroke}{rgb}{0.000000,0.000000,0.000000}%
\pgfsetstrokecolor{currentstroke}%
\pgfsetdash{}{0pt}%
\pgfsys@defobject{currentmarker}{\pgfqpoint{-0.048611in}{0.000000in}}{\pgfqpoint{-0.000000in}{0.000000in}}{%
\pgfpathmoveto{\pgfqpoint{-0.000000in}{0.000000in}}%
\pgfpathlineto{\pgfqpoint{-0.048611in}{0.000000in}}%
\pgfusepath{stroke,fill}%
}%
\begin{pgfscope}%
\pgfsys@transformshift{0.623025in}{1.320755in}%
\pgfsys@useobject{currentmarker}{}%
\end{pgfscope}%
\end{pgfscope}%
\begin{pgfscope}%
\definecolor{textcolor}{rgb}{0.000000,0.000000,0.000000}%
\pgfsetstrokecolor{textcolor}%
\pgfsetfillcolor{textcolor}%
\pgftext[x=0.278889in, y=1.272561in, left, base]{\color{textcolor}{\rmfamily\fontsize{10.000000}{12.000000}\selectfont\catcode`\^=\active\def^{\ifmmode\sp\else\^{}\fi}\catcode`\%=\active\def%{\%}$\mathdefault{0.10}$}}%
\end{pgfscope}%
\begin{pgfscope}%
\pgfpathrectangle{\pgfqpoint{0.623025in}{0.499444in}}{\pgfqpoint{5.322679in}{3.637809in}}%
\pgfusepath{clip}%
\pgfsetbuttcap%
\pgfsetroundjoin%
\pgfsetlinewidth{0.501875pt}%
\definecolor{currentstroke}{rgb}{0.501961,0.501961,0.501961}%
\pgfsetstrokecolor{currentstroke}%
\pgfsetstrokeopacity{0.500000}%
\pgfsetdash{{1.850000pt}{0.800000pt}}{0.000000pt}%
\pgfpathmoveto{\pgfqpoint{0.623025in}{1.731411in}}%
\pgfpathlineto{\pgfqpoint{5.945705in}{1.731411in}}%
\pgfusepath{stroke}%
\end{pgfscope}%
\begin{pgfscope}%
\pgfsetbuttcap%
\pgfsetroundjoin%
\definecolor{currentfill}{rgb}{0.000000,0.000000,0.000000}%
\pgfsetfillcolor{currentfill}%
\pgfsetlinewidth{0.803000pt}%
\definecolor{currentstroke}{rgb}{0.000000,0.000000,0.000000}%
\pgfsetstrokecolor{currentstroke}%
\pgfsetdash{}{0pt}%
\pgfsys@defobject{currentmarker}{\pgfqpoint{-0.048611in}{0.000000in}}{\pgfqpoint{-0.000000in}{0.000000in}}{%
\pgfpathmoveto{\pgfqpoint{-0.000000in}{0.000000in}}%
\pgfpathlineto{\pgfqpoint{-0.048611in}{0.000000in}}%
\pgfusepath{stroke,fill}%
}%
\begin{pgfscope}%
\pgfsys@transformshift{0.623025in}{1.731411in}%
\pgfsys@useobject{currentmarker}{}%
\end{pgfscope}%
\end{pgfscope}%
\begin{pgfscope}%
\definecolor{textcolor}{rgb}{0.000000,0.000000,0.000000}%
\pgfsetstrokecolor{textcolor}%
\pgfsetfillcolor{textcolor}%
\pgftext[x=0.278889in, y=1.683216in, left, base]{\color{textcolor}{\rmfamily\fontsize{10.000000}{12.000000}\selectfont\catcode`\^=\active\def^{\ifmmode\sp\else\^{}\fi}\catcode`\%=\active\def%{\%}$\mathdefault{0.15}$}}%
\end{pgfscope}%
\begin{pgfscope}%
\pgfpathrectangle{\pgfqpoint{0.623025in}{0.499444in}}{\pgfqpoint{5.322679in}{3.637809in}}%
\pgfusepath{clip}%
\pgfsetbuttcap%
\pgfsetroundjoin%
\pgfsetlinewidth{0.501875pt}%
\definecolor{currentstroke}{rgb}{0.501961,0.501961,0.501961}%
\pgfsetstrokecolor{currentstroke}%
\pgfsetstrokeopacity{0.500000}%
\pgfsetdash{{1.850000pt}{0.800000pt}}{0.000000pt}%
\pgfpathmoveto{\pgfqpoint{0.623025in}{2.142066in}}%
\pgfpathlineto{\pgfqpoint{5.945705in}{2.142066in}}%
\pgfusepath{stroke}%
\end{pgfscope}%
\begin{pgfscope}%
\pgfsetbuttcap%
\pgfsetroundjoin%
\definecolor{currentfill}{rgb}{0.000000,0.000000,0.000000}%
\pgfsetfillcolor{currentfill}%
\pgfsetlinewidth{0.803000pt}%
\definecolor{currentstroke}{rgb}{0.000000,0.000000,0.000000}%
\pgfsetstrokecolor{currentstroke}%
\pgfsetdash{}{0pt}%
\pgfsys@defobject{currentmarker}{\pgfqpoint{-0.048611in}{0.000000in}}{\pgfqpoint{-0.000000in}{0.000000in}}{%
\pgfpathmoveto{\pgfqpoint{-0.000000in}{0.000000in}}%
\pgfpathlineto{\pgfqpoint{-0.048611in}{0.000000in}}%
\pgfusepath{stroke,fill}%
}%
\begin{pgfscope}%
\pgfsys@transformshift{0.623025in}{2.142066in}%
\pgfsys@useobject{currentmarker}{}%
\end{pgfscope}%
\end{pgfscope}%
\begin{pgfscope}%
\definecolor{textcolor}{rgb}{0.000000,0.000000,0.000000}%
\pgfsetstrokecolor{textcolor}%
\pgfsetfillcolor{textcolor}%
\pgftext[x=0.278889in, y=2.093872in, left, base]{\color{textcolor}{\rmfamily\fontsize{10.000000}{12.000000}\selectfont\catcode`\^=\active\def^{\ifmmode\sp\else\^{}\fi}\catcode`\%=\active\def%{\%}$\mathdefault{0.20}$}}%
\end{pgfscope}%
\begin{pgfscope}%
\pgfpathrectangle{\pgfqpoint{0.623025in}{0.499444in}}{\pgfqpoint{5.322679in}{3.637809in}}%
\pgfusepath{clip}%
\pgfsetbuttcap%
\pgfsetroundjoin%
\pgfsetlinewidth{0.501875pt}%
\definecolor{currentstroke}{rgb}{0.501961,0.501961,0.501961}%
\pgfsetstrokecolor{currentstroke}%
\pgfsetstrokeopacity{0.500000}%
\pgfsetdash{{1.850000pt}{0.800000pt}}{0.000000pt}%
\pgfpathmoveto{\pgfqpoint{0.623025in}{2.552721in}}%
\pgfpathlineto{\pgfqpoint{5.945705in}{2.552721in}}%
\pgfusepath{stroke}%
\end{pgfscope}%
\begin{pgfscope}%
\pgfsetbuttcap%
\pgfsetroundjoin%
\definecolor{currentfill}{rgb}{0.000000,0.000000,0.000000}%
\pgfsetfillcolor{currentfill}%
\pgfsetlinewidth{0.803000pt}%
\definecolor{currentstroke}{rgb}{0.000000,0.000000,0.000000}%
\pgfsetstrokecolor{currentstroke}%
\pgfsetdash{}{0pt}%
\pgfsys@defobject{currentmarker}{\pgfqpoint{-0.048611in}{0.000000in}}{\pgfqpoint{-0.000000in}{0.000000in}}{%
\pgfpathmoveto{\pgfqpoint{-0.000000in}{0.000000in}}%
\pgfpathlineto{\pgfqpoint{-0.048611in}{0.000000in}}%
\pgfusepath{stroke,fill}%
}%
\begin{pgfscope}%
\pgfsys@transformshift{0.623025in}{2.552721in}%
\pgfsys@useobject{currentmarker}{}%
\end{pgfscope}%
\end{pgfscope}%
\begin{pgfscope}%
\definecolor{textcolor}{rgb}{0.000000,0.000000,0.000000}%
\pgfsetstrokecolor{textcolor}%
\pgfsetfillcolor{textcolor}%
\pgftext[x=0.278889in, y=2.504527in, left, base]{\color{textcolor}{\rmfamily\fontsize{10.000000}{12.000000}\selectfont\catcode`\^=\active\def^{\ifmmode\sp\else\^{}\fi}\catcode`\%=\active\def%{\%}$\mathdefault{0.25}$}}%
\end{pgfscope}%
\begin{pgfscope}%
\pgfpathrectangle{\pgfqpoint{0.623025in}{0.499444in}}{\pgfqpoint{5.322679in}{3.637809in}}%
\pgfusepath{clip}%
\pgfsetbuttcap%
\pgfsetroundjoin%
\pgfsetlinewidth{0.501875pt}%
\definecolor{currentstroke}{rgb}{0.501961,0.501961,0.501961}%
\pgfsetstrokecolor{currentstroke}%
\pgfsetstrokeopacity{0.500000}%
\pgfsetdash{{1.850000pt}{0.800000pt}}{0.000000pt}%
\pgfpathmoveto{\pgfqpoint{0.623025in}{2.963377in}}%
\pgfpathlineto{\pgfqpoint{5.945705in}{2.963377in}}%
\pgfusepath{stroke}%
\end{pgfscope}%
\begin{pgfscope}%
\pgfsetbuttcap%
\pgfsetroundjoin%
\definecolor{currentfill}{rgb}{0.000000,0.000000,0.000000}%
\pgfsetfillcolor{currentfill}%
\pgfsetlinewidth{0.803000pt}%
\definecolor{currentstroke}{rgb}{0.000000,0.000000,0.000000}%
\pgfsetstrokecolor{currentstroke}%
\pgfsetdash{}{0pt}%
\pgfsys@defobject{currentmarker}{\pgfqpoint{-0.048611in}{0.000000in}}{\pgfqpoint{-0.000000in}{0.000000in}}{%
\pgfpathmoveto{\pgfqpoint{-0.000000in}{0.000000in}}%
\pgfpathlineto{\pgfqpoint{-0.048611in}{0.000000in}}%
\pgfusepath{stroke,fill}%
}%
\begin{pgfscope}%
\pgfsys@transformshift{0.623025in}{2.963377in}%
\pgfsys@useobject{currentmarker}{}%
\end{pgfscope}%
\end{pgfscope}%
\begin{pgfscope}%
\definecolor{textcolor}{rgb}{0.000000,0.000000,0.000000}%
\pgfsetstrokecolor{textcolor}%
\pgfsetfillcolor{textcolor}%
\pgftext[x=0.278889in, y=2.915182in, left, base]{\color{textcolor}{\rmfamily\fontsize{10.000000}{12.000000}\selectfont\catcode`\^=\active\def^{\ifmmode\sp\else\^{}\fi}\catcode`\%=\active\def%{\%}$\mathdefault{0.30}$}}%
\end{pgfscope}%
\begin{pgfscope}%
\pgfpathrectangle{\pgfqpoint{0.623025in}{0.499444in}}{\pgfqpoint{5.322679in}{3.637809in}}%
\pgfusepath{clip}%
\pgfsetbuttcap%
\pgfsetroundjoin%
\pgfsetlinewidth{0.501875pt}%
\definecolor{currentstroke}{rgb}{0.501961,0.501961,0.501961}%
\pgfsetstrokecolor{currentstroke}%
\pgfsetstrokeopacity{0.500000}%
\pgfsetdash{{1.850000pt}{0.800000pt}}{0.000000pt}%
\pgfpathmoveto{\pgfqpoint{0.623025in}{3.374032in}}%
\pgfpathlineto{\pgfqpoint{5.945705in}{3.374032in}}%
\pgfusepath{stroke}%
\end{pgfscope}%
\begin{pgfscope}%
\pgfsetbuttcap%
\pgfsetroundjoin%
\definecolor{currentfill}{rgb}{0.000000,0.000000,0.000000}%
\pgfsetfillcolor{currentfill}%
\pgfsetlinewidth{0.803000pt}%
\definecolor{currentstroke}{rgb}{0.000000,0.000000,0.000000}%
\pgfsetstrokecolor{currentstroke}%
\pgfsetdash{}{0pt}%
\pgfsys@defobject{currentmarker}{\pgfqpoint{-0.048611in}{0.000000in}}{\pgfqpoint{-0.000000in}{0.000000in}}{%
\pgfpathmoveto{\pgfqpoint{-0.000000in}{0.000000in}}%
\pgfpathlineto{\pgfqpoint{-0.048611in}{0.000000in}}%
\pgfusepath{stroke,fill}%
}%
\begin{pgfscope}%
\pgfsys@transformshift{0.623025in}{3.374032in}%
\pgfsys@useobject{currentmarker}{}%
\end{pgfscope}%
\end{pgfscope}%
\begin{pgfscope}%
\definecolor{textcolor}{rgb}{0.000000,0.000000,0.000000}%
\pgfsetstrokecolor{textcolor}%
\pgfsetfillcolor{textcolor}%
\pgftext[x=0.278889in, y=3.325838in, left, base]{\color{textcolor}{\rmfamily\fontsize{10.000000}{12.000000}\selectfont\catcode`\^=\active\def^{\ifmmode\sp\else\^{}\fi}\catcode`\%=\active\def%{\%}$\mathdefault{0.35}$}}%
\end{pgfscope}%
\begin{pgfscope}%
\pgfpathrectangle{\pgfqpoint{0.623025in}{0.499444in}}{\pgfqpoint{5.322679in}{3.637809in}}%
\pgfusepath{clip}%
\pgfsetbuttcap%
\pgfsetroundjoin%
\pgfsetlinewidth{0.501875pt}%
\definecolor{currentstroke}{rgb}{0.501961,0.501961,0.501961}%
\pgfsetstrokecolor{currentstroke}%
\pgfsetstrokeopacity{0.500000}%
\pgfsetdash{{1.850000pt}{0.800000pt}}{0.000000pt}%
\pgfpathmoveto{\pgfqpoint{0.623025in}{3.784688in}}%
\pgfpathlineto{\pgfqpoint{5.945705in}{3.784688in}}%
\pgfusepath{stroke}%
\end{pgfscope}%
\begin{pgfscope}%
\pgfsetbuttcap%
\pgfsetroundjoin%
\definecolor{currentfill}{rgb}{0.000000,0.000000,0.000000}%
\pgfsetfillcolor{currentfill}%
\pgfsetlinewidth{0.803000pt}%
\definecolor{currentstroke}{rgb}{0.000000,0.000000,0.000000}%
\pgfsetstrokecolor{currentstroke}%
\pgfsetdash{}{0pt}%
\pgfsys@defobject{currentmarker}{\pgfqpoint{-0.048611in}{0.000000in}}{\pgfqpoint{-0.000000in}{0.000000in}}{%
\pgfpathmoveto{\pgfqpoint{-0.000000in}{0.000000in}}%
\pgfpathlineto{\pgfqpoint{-0.048611in}{0.000000in}}%
\pgfusepath{stroke,fill}%
}%
\begin{pgfscope}%
\pgfsys@transformshift{0.623025in}{3.784688in}%
\pgfsys@useobject{currentmarker}{}%
\end{pgfscope}%
\end{pgfscope}%
\begin{pgfscope}%
\definecolor{textcolor}{rgb}{0.000000,0.000000,0.000000}%
\pgfsetstrokecolor{textcolor}%
\pgfsetfillcolor{textcolor}%
\pgftext[x=0.278889in, y=3.736493in, left, base]{\color{textcolor}{\rmfamily\fontsize{10.000000}{12.000000}\selectfont\catcode`\^=\active\def^{\ifmmode\sp\else\^{}\fi}\catcode`\%=\active\def%{\%}$\mathdefault{0.40}$}}%
\end{pgfscope}%
\begin{pgfscope}%
\definecolor{textcolor}{rgb}{0.000000,0.000000,0.000000}%
\pgfsetstrokecolor{textcolor}%
\pgfsetfillcolor{textcolor}%
\pgftext[x=0.223333in,y=2.318349in,,bottom,rotate=90.000000]{\color{textcolor}{\rmfamily\fontsize{10.000000}{12.000000}\selectfont\catcode`\^=\active\def^{\ifmmode\sp\else\^{}\fi}\catcode`\%=\active\def%{\%}Densidad}}%
\end{pgfscope}%
\begin{pgfscope}%
\pgfpathrectangle{\pgfqpoint{0.623025in}{0.499444in}}{\pgfqpoint{5.322679in}{3.637809in}}%
\pgfusepath{clip}%
\pgfsetrectcap%
\pgfsetroundjoin%
\pgfsetlinewidth{2.007500pt}%
\definecolor{currentstroke}{rgb}{1.000000,0.498039,0.054902}%
\pgfsetstrokecolor{currentstroke}%
\pgfsetdash{}{0pt}%
\pgfpathmoveto{\pgfqpoint{0.864965in}{0.500556in}}%
\pgfpathlineto{\pgfqpoint{1.046875in}{0.502983in}}%
\pgfpathlineto{\pgfqpoint{1.168148in}{0.506727in}}%
\pgfpathlineto{\pgfqpoint{1.265167in}{0.512046in}}%
\pgfpathlineto{\pgfqpoint{1.337931in}{0.518139in}}%
\pgfpathlineto{\pgfqpoint{1.398567in}{0.525127in}}%
\pgfpathlineto{\pgfqpoint{1.459204in}{0.534375in}}%
\pgfpathlineto{\pgfqpoint{1.507713in}{0.543796in}}%
\pgfpathlineto{\pgfqpoint{1.556222in}{0.555397in}}%
\pgfpathlineto{\pgfqpoint{1.604732in}{0.569581in}}%
\pgfpathlineto{\pgfqpoint{1.641114in}{0.582181in}}%
\pgfpathlineto{\pgfqpoint{1.677495in}{0.596693in}}%
\pgfpathlineto{\pgfqpoint{1.713877in}{0.613338in}}%
\pgfpathlineto{\pgfqpoint{1.750259in}{0.632350in}}%
\pgfpathlineto{\pgfqpoint{1.786641in}{0.653974in}}%
\pgfpathlineto{\pgfqpoint{1.823023in}{0.678469in}}%
\pgfpathlineto{\pgfqpoint{1.859405in}{0.706097in}}%
\pgfpathlineto{\pgfqpoint{1.895787in}{0.737128in}}%
\pgfpathlineto{\pgfqpoint{1.932169in}{0.771830in}}%
\pgfpathlineto{\pgfqpoint{1.968551in}{0.810472in}}%
\pgfpathlineto{\pgfqpoint{2.004933in}{0.853313in}}%
\pgfpathlineto{\pgfqpoint{2.041315in}{0.900602in}}%
\pgfpathlineto{\pgfqpoint{2.077697in}{0.952567in}}%
\pgfpathlineto{\pgfqpoint{2.114079in}{1.009414in}}%
\pgfpathlineto{\pgfqpoint{2.150461in}{1.071321in}}%
\pgfpathlineto{\pgfqpoint{2.186843in}{1.138428in}}%
\pgfpathlineto{\pgfqpoint{2.223225in}{1.210830in}}%
\pgfpathlineto{\pgfqpoint{2.259607in}{1.288575in}}%
\pgfpathlineto{\pgfqpoint{2.295989in}{1.371656in}}%
\pgfpathlineto{\pgfqpoint{2.332371in}{1.460002in}}%
\pgfpathlineto{\pgfqpoint{2.368753in}{1.553476in}}%
\pgfpathlineto{\pgfqpoint{2.405134in}{1.651869in}}%
\pgfpathlineto{\pgfqpoint{2.453644in}{1.790208in}}%
\pgfpathlineto{\pgfqpoint{2.502153in}{1.935885in}}%
\pgfpathlineto{\pgfqpoint{2.562790in}{2.126530in}}%
\pgfpathlineto{\pgfqpoint{2.635553in}{2.364090in}}%
\pgfpathlineto{\pgfqpoint{2.793209in}{2.883249in}}%
\pgfpathlineto{\pgfqpoint{2.841718in}{3.035493in}}%
\pgfpathlineto{\pgfqpoint{2.890227in}{3.180161in}}%
\pgfpathlineto{\pgfqpoint{2.926609in}{3.282281in}}%
\pgfpathlineto{\pgfqpoint{2.962991in}{3.377858in}}%
\pgfpathlineto{\pgfqpoint{2.999373in}{3.465966in}}%
\pgfpathlineto{\pgfqpoint{3.023628in}{3.520119in}}%
\pgfpathlineto{\pgfqpoint{3.047882in}{3.570320in}}%
\pgfpathlineto{\pgfqpoint{3.072137in}{3.616338in}}%
\pgfpathlineto{\pgfqpoint{3.096392in}{3.657963in}}%
\pgfpathlineto{\pgfqpoint{3.120646in}{3.695001in}}%
\pgfpathlineto{\pgfqpoint{3.144901in}{3.727279in}}%
\pgfpathlineto{\pgfqpoint{3.169155in}{3.754644in}}%
\pgfpathlineto{\pgfqpoint{3.193410in}{3.776967in}}%
\pgfpathlineto{\pgfqpoint{3.217665in}{3.794141in}}%
\pgfpathlineto{\pgfqpoint{3.229792in}{3.800770in}}%
\pgfpathlineto{\pgfqpoint{3.241919in}{3.806083in}}%
\pgfpathlineto{\pgfqpoint{3.254047in}{3.810073in}}%
\pgfpathlineto{\pgfqpoint{3.266174in}{3.812736in}}%
\pgfpathlineto{\pgfqpoint{3.278301in}{3.814068in}}%
\pgfpathlineto{\pgfqpoint{3.290429in}{3.814068in}}%
\pgfpathlineto{\pgfqpoint{3.302556in}{3.812736in}}%
\pgfpathlineto{\pgfqpoint{3.314683in}{3.810073in}}%
\pgfpathlineto{\pgfqpoint{3.326811in}{3.806083in}}%
\pgfpathlineto{\pgfqpoint{3.338938in}{3.800770in}}%
\pgfpathlineto{\pgfqpoint{3.351065in}{3.794141in}}%
\pgfpathlineto{\pgfqpoint{3.375320in}{3.776967in}}%
\pgfpathlineto{\pgfqpoint{3.399574in}{3.754644in}}%
\pgfpathlineto{\pgfqpoint{3.423829in}{3.727279in}}%
\pgfpathlineto{\pgfqpoint{3.448084in}{3.695001in}}%
\pgfpathlineto{\pgfqpoint{3.472338in}{3.657963in}}%
\pgfpathlineto{\pgfqpoint{3.496593in}{3.616338in}}%
\pgfpathlineto{\pgfqpoint{3.520848in}{3.570320in}}%
\pgfpathlineto{\pgfqpoint{3.545102in}{3.520119in}}%
\pgfpathlineto{\pgfqpoint{3.569357in}{3.465966in}}%
\pgfpathlineto{\pgfqpoint{3.605739in}{3.377858in}}%
\pgfpathlineto{\pgfqpoint{3.642121in}{3.282281in}}%
\pgfpathlineto{\pgfqpoint{3.678503in}{3.180161in}}%
\pgfpathlineto{\pgfqpoint{3.727012in}{3.035493in}}%
\pgfpathlineto{\pgfqpoint{3.775521in}{2.883249in}}%
\pgfpathlineto{\pgfqpoint{3.848285in}{2.645805in}}%
\pgfpathlineto{\pgfqpoint{4.005940in}{2.126530in}}%
\pgfpathlineto{\pgfqpoint{4.066577in}{1.935885in}}%
\pgfpathlineto{\pgfqpoint{4.115086in}{1.790208in}}%
\pgfpathlineto{\pgfqpoint{4.163595in}{1.651869in}}%
\pgfpathlineto{\pgfqpoint{4.212105in}{1.521760in}}%
\pgfpathlineto{\pgfqpoint{4.248487in}{1.429975in}}%
\pgfpathlineto{\pgfqpoint{4.284869in}{1.343373in}}%
\pgfpathlineto{\pgfqpoint{4.321251in}{1.262066in}}%
\pgfpathlineto{\pgfqpoint{4.357632in}{1.186103in}}%
\pgfpathlineto{\pgfqpoint{4.394014in}{1.115474in}}%
\pgfpathlineto{\pgfqpoint{4.430396in}{1.050114in}}%
\pgfpathlineto{\pgfqpoint{4.466778in}{0.989911in}}%
\pgfpathlineto{\pgfqpoint{4.503160in}{0.934712in}}%
\pgfpathlineto{\pgfqpoint{4.539542in}{0.884330in}}%
\pgfpathlineto{\pgfqpoint{4.575924in}{0.838551in}}%
\pgfpathlineto{\pgfqpoint{4.612306in}{0.797138in}}%
\pgfpathlineto{\pgfqpoint{4.648688in}{0.759838in}}%
\pgfpathlineto{\pgfqpoint{4.685070in}{0.726390in}}%
\pgfpathlineto{\pgfqpoint{4.721452in}{0.696523in}}%
\pgfpathlineto{\pgfqpoint{4.757834in}{0.669969in}}%
\pgfpathlineto{\pgfqpoint{4.794216in}{0.646460in}}%
\pgfpathlineto{\pgfqpoint{4.830598in}{0.625734in}}%
\pgfpathlineto{\pgfqpoint{4.866980in}{0.607539in}}%
\pgfpathlineto{\pgfqpoint{4.903362in}{0.591630in}}%
\pgfpathlineto{\pgfqpoint{4.939744in}{0.577779in}}%
\pgfpathlineto{\pgfqpoint{4.988253in}{0.562139in}}%
\pgfpathlineto{\pgfqpoint{5.036762in}{0.549300in}}%
\pgfpathlineto{\pgfqpoint{5.085271in}{0.538836in}}%
\pgfpathlineto{\pgfqpoint{5.133781in}{0.530369in}}%
\pgfpathlineto{\pgfqpoint{5.194417in}{0.522091in}}%
\pgfpathlineto{\pgfqpoint{5.267181in}{0.514821in}}%
\pgfpathlineto{\pgfqpoint{5.352072in}{0.509055in}}%
\pgfpathlineto{\pgfqpoint{5.449091in}{0.504928in}}%
\pgfpathlineto{\pgfqpoint{5.570364in}{0.502066in}}%
\pgfpathlineto{\pgfqpoint{5.703765in}{0.500556in}}%
\pgfpathlineto{\pgfqpoint{5.703765in}{0.500556in}}%
\pgfusepath{stroke}%
\end{pgfscope}%
\begin{pgfscope}%
\pgfsetrectcap%
\pgfsetmiterjoin%
\pgfsetlinewidth{0.803000pt}%
\definecolor{currentstroke}{rgb}{0.000000,0.000000,0.000000}%
\pgfsetstrokecolor{currentstroke}%
\pgfsetdash{}{0pt}%
\pgfpathmoveto{\pgfqpoint{0.623025in}{0.499444in}}%
\pgfpathlineto{\pgfqpoint{0.623025in}{4.137253in}}%
\pgfusepath{stroke}%
\end{pgfscope}%
\begin{pgfscope}%
\pgfsetrectcap%
\pgfsetmiterjoin%
\pgfsetlinewidth{0.803000pt}%
\definecolor{currentstroke}{rgb}{0.000000,0.000000,0.000000}%
\pgfsetstrokecolor{currentstroke}%
\pgfsetdash{}{0pt}%
\pgfpathmoveto{\pgfqpoint{5.945705in}{0.499444in}}%
\pgfpathlineto{\pgfqpoint{5.945705in}{4.137253in}}%
\pgfusepath{stroke}%
\end{pgfscope}%
\begin{pgfscope}%
\pgfsetrectcap%
\pgfsetmiterjoin%
\pgfsetlinewidth{0.803000pt}%
\definecolor{currentstroke}{rgb}{0.000000,0.000000,0.000000}%
\pgfsetstrokecolor{currentstroke}%
\pgfsetdash{}{0pt}%
\pgfpathmoveto{\pgfqpoint{0.623025in}{0.499444in}}%
\pgfpathlineto{\pgfqpoint{5.945705in}{0.499444in}}%
\pgfusepath{stroke}%
\end{pgfscope}%
\begin{pgfscope}%
\pgfsetrectcap%
\pgfsetmiterjoin%
\pgfsetlinewidth{0.803000pt}%
\definecolor{currentstroke}{rgb}{0.000000,0.000000,0.000000}%
\pgfsetstrokecolor{currentstroke}%
\pgfsetdash{}{0pt}%
\pgfpathmoveto{\pgfqpoint{0.623025in}{4.137253in}}%
\pgfpathlineto{\pgfqpoint{5.945705in}{4.137253in}}%
\pgfusepath{stroke}%
\end{pgfscope}%
\begin{pgfscope}%
\definecolor{textcolor}{rgb}{0.000000,0.000000,0.000000}%
\pgfsetstrokecolor{textcolor}%
\pgfsetfillcolor{textcolor}%
\pgftext[x=3.284365in,y=4.220587in,,base]{\color{textcolor}{\rmfamily\fontsize{12.000000}{14.400000}\selectfont\catcode`\^=\active\def^{\ifmmode\sp\else\^{}\fi}\catcode`\%=\active\def%{\%}Distribución y densidad (hist con density=True)}}%
\end{pgfscope}%
\end{pgfpicture}%
\makeatother%
\endgroup%
 
    \caption{Histograma normalizado}
    \label{fig:grafico_avanzado_hist}
\end{figure}


\subsection{Diagramas de Error (\textit{.errorbar()})}

La función \texttt{.errorbar()} añade una capa de información sobre la variabilidad o incertidumbre de cada punto graficado.

\subsubsection{Componentes del Gráfico de Error}

\begin{itemize}
\item \texttt{yerr} / \texttt{xerr}: Define la magnitud del error. Puede ser un valor fijo para todos los puntos o un arreglo con un error específico para cada medición.
\item \texttt{fmt} (\textit{format}): Controla la apariencia del marcador y la línea de conexión (ej: \texttt{'o-'} para círculos conectados por una línea sólida).
\item \texttt{capsize}: Agrega pequeñas barras horizontales (tapas) al final de las líneas de error, lo que mejora drásticamente la lectura visual de la dispersión.
\item \texttt{ecolor} / \texttt{elinewidth}: Permite que las barras de error tengan un color y grosor distinto al de la línea principal para evitar que el gráfico se vea sobrecargado.
\end{itemize}

\begin{pycode}{Diagrama de error vertical}
import numpy as np
import matplotlib.pyplot as plt

x = np.arange(1, 6)
y = np.array([2.0, 2.5, 3.7, 3.2, 4.1])

yerr = np.array([0.2, 0.3, 0.25, 0.35, 0.2]) # Error vertical

fig, ax = plt.subplots()
ax.errorbar(
    x,
    y,
    yerr=yerr,
    fmt="o-", # Marcador + línea
    capsize=4, # Lineas horizontales en el extremo de la barra de error
    elinewidth=1.2, # Grosor de la barra de error
    label="Medición con error"
)

ax.set_title("Diagrama de error (errorbar)")
ax.set_xlabel("Muestra")
ax.set_ylabel("Valor medido")
ax.legend()
plt.show()    
\end{pycode}

\begin{pycode}{Diagrama de error vertical y horizontal}
import numpy as np
import matplotlib.pyplot as plt

x = np.linspace(0.5, 2.5, 5)
y = np.array([1.5, 1.8, 2.1, 2.0, 2.3])

# Errores
xerr = 0.1
yerr = [0.05, 0.1, 0.08, 0.12, 0.06]

fig, ax = plt.subplots()
ax.errorbar(
    x,
    y,
    xerr=xerr,
    yerr=yerr,
    fmt="s-",
    capsize=3,
    ecolor="gray",
    label="Error en X e Y"
)

ax.set_title("Barras de error en ambos ejes")
ax.set_xlabel("X")
ax.set_ylabel("Y")
ax.legend()
plt.show()    
\end{pycode}

\begin{nota}
    Si quiere realizar errores asimétricos, esto se realiza utilizando una matriz en lugar de un array simple en \texttt{yerr} o \texttt{xerr}, por ej: \texttt{yerr = [[1, 2, 1], [4, 5, 2]]}.  
\end{nota}

La combinación de \texttt{xerr} y \texttt{yerr} nos da la figura \ref{fig:grafico_diagrama_error}

\begin{figure}[ht]
    \centering
    %% Creator: Matplotlib, PGF backend
%%
%% To include the figure in your LaTeX document, write
%%   \input{<filename>.pgf}
%%
%% Make sure the required packages are loaded in your preamble
%%   \usepackage{pgf}
%%
%% Also ensure that all the required font packages are loaded; for instance,
%% the lmodern package is sometimes necessary when using math font.
%%   \usepackage{lmodern}
%%
%% Figures using additional raster images can only be included by \input if
%% they are in the same directory as the main LaTeX file. For loading figures
%% from other directories you can use the `import` package
%%   \usepackage{import}
%%
%% and then include the figures with
%%   \import{<path to file>}{<filename>.pgf}
%%
%% Matplotlib used the following preamble
%%   \def\mathdefault#1{#1}
%%   \everymath=\expandafter{\the\everymath\displaystyle}
%%   \IfFileExists{scrextend.sty}{
%%     \usepackage[fontsize=10.000000pt]{scrextend}
%%   }{
%%     \renewcommand{\normalsize}{\fontsize{10.000000}{12.000000}\selectfont}
%%     \normalsize
%%   }
%%   \usepackage{fontspec}
%%   \setmainfont{CMU Serif}
%%   \makeatletter\@ifpackageloaded{underscore}{}{\usepackage[strings]{underscore}}\makeatother
%%
\begingroup%
\makeatletter%
\begin{pgfpicture}%
\pgfpathrectangle{\pgfpointorigin}{\pgfqpoint{6.045445in}{4.445478in}}%
\pgfusepath{use as bounding box, clip}%
\begin{pgfscope}%
\pgfsetbuttcap%
\pgfsetmiterjoin%
\definecolor{currentfill}{rgb}{1.000000,1.000000,1.000000}%
\pgfsetfillcolor{currentfill}%
\pgfsetlinewidth{0.000000pt}%
\definecolor{currentstroke}{rgb}{1.000000,1.000000,1.000000}%
\pgfsetstrokecolor{currentstroke}%
\pgfsetdash{}{0pt}%
\pgfpathmoveto{\pgfqpoint{0.000000in}{0.000000in}}%
\pgfpathlineto{\pgfqpoint{6.045445in}{0.000000in}}%
\pgfpathlineto{\pgfqpoint{6.045445in}{4.445478in}}%
\pgfpathlineto{\pgfqpoint{0.000000in}{4.445478in}}%
\pgfpathlineto{\pgfqpoint{0.000000in}{0.000000in}}%
\pgfpathclose%
\pgfusepath{fill}%
\end{pgfscope}%
\begin{pgfscope}%
\pgfsetbuttcap%
\pgfsetmiterjoin%
\definecolor{currentfill}{rgb}{1.000000,1.000000,1.000000}%
\pgfsetfillcolor{currentfill}%
\pgfsetlinewidth{0.000000pt}%
\definecolor{currentstroke}{rgb}{0.000000,0.000000,0.000000}%
\pgfsetstrokecolor{currentstroke}%
\pgfsetstrokeopacity{0.000000}%
\pgfsetdash{}{0pt}%
\pgfpathmoveto{\pgfqpoint{0.553581in}{0.499444in}}%
\pgfpathlineto{\pgfqpoint{5.945445in}{0.499444in}}%
\pgfpathlineto{\pgfqpoint{5.945445in}{4.146478in}}%
\pgfpathlineto{\pgfqpoint{0.553581in}{4.146478in}}%
\pgfpathlineto{\pgfqpoint{0.553581in}{0.499444in}}%
\pgfpathclose%
\pgfusepath{fill}%
\end{pgfscope}%
\begin{pgfscope}%
\pgfpathrectangle{\pgfqpoint{0.553581in}{0.499444in}}{\pgfqpoint{5.391865in}{3.647034in}}%
\pgfusepath{clip}%
\pgfsetbuttcap%
\pgfsetroundjoin%
\pgfsetlinewidth{0.501875pt}%
\definecolor{currentstroke}{rgb}{0.501961,0.501961,0.501961}%
\pgfsetstrokecolor{currentstroke}%
\pgfsetstrokeopacity{0.500000}%
\pgfsetdash{{1.850000pt}{0.800000pt}}{0.000000pt}%
\pgfpathmoveto{\pgfqpoint{1.021470in}{0.499444in}}%
\pgfpathlineto{\pgfqpoint{1.021470in}{4.146478in}}%
\pgfusepath{stroke}%
\end{pgfscope}%
\begin{pgfscope}%
\pgfsetbuttcap%
\pgfsetroundjoin%
\definecolor{currentfill}{rgb}{0.000000,0.000000,0.000000}%
\pgfsetfillcolor{currentfill}%
\pgfsetlinewidth{0.803000pt}%
\definecolor{currentstroke}{rgb}{0.000000,0.000000,0.000000}%
\pgfsetstrokecolor{currentstroke}%
\pgfsetdash{}{0pt}%
\pgfsys@defobject{currentmarker}{\pgfqpoint{0.000000in}{-0.048611in}}{\pgfqpoint{0.000000in}{0.000000in}}{%
\pgfpathmoveto{\pgfqpoint{0.000000in}{0.000000in}}%
\pgfpathlineto{\pgfqpoint{0.000000in}{-0.048611in}}%
\pgfusepath{stroke,fill}%
}%
\begin{pgfscope}%
\pgfsys@transformshift{1.021470in}{0.499444in}%
\pgfsys@useobject{currentmarker}{}%
\end{pgfscope}%
\end{pgfscope}%
\begin{pgfscope}%
\definecolor{textcolor}{rgb}{0.000000,0.000000,0.000000}%
\pgfsetstrokecolor{textcolor}%
\pgfsetfillcolor{textcolor}%
\pgftext[x=1.021470in,y=0.402222in,,top]{\color{textcolor}{\rmfamily\fontsize{10.000000}{12.000000}\selectfont\catcode`\^=\active\def^{\ifmmode\sp\else\^{}\fi}\catcode`\%=\active\def%{\%}$\mathdefault{0.5}$}}%
\end{pgfscope}%
\begin{pgfscope}%
\pgfpathrectangle{\pgfqpoint{0.553581in}{0.499444in}}{\pgfqpoint{5.391865in}{3.647034in}}%
\pgfusepath{clip}%
\pgfsetbuttcap%
\pgfsetroundjoin%
\pgfsetlinewidth{0.501875pt}%
\definecolor{currentstroke}{rgb}{0.501961,0.501961,0.501961}%
\pgfsetstrokecolor{currentstroke}%
\pgfsetstrokeopacity{0.500000}%
\pgfsetdash{{1.850000pt}{0.800000pt}}{0.000000pt}%
\pgfpathmoveto{\pgfqpoint{2.135491in}{0.499444in}}%
\pgfpathlineto{\pgfqpoint{2.135491in}{4.146478in}}%
\pgfusepath{stroke}%
\end{pgfscope}%
\begin{pgfscope}%
\pgfsetbuttcap%
\pgfsetroundjoin%
\definecolor{currentfill}{rgb}{0.000000,0.000000,0.000000}%
\pgfsetfillcolor{currentfill}%
\pgfsetlinewidth{0.803000pt}%
\definecolor{currentstroke}{rgb}{0.000000,0.000000,0.000000}%
\pgfsetstrokecolor{currentstroke}%
\pgfsetdash{}{0pt}%
\pgfsys@defobject{currentmarker}{\pgfqpoint{0.000000in}{-0.048611in}}{\pgfqpoint{0.000000in}{0.000000in}}{%
\pgfpathmoveto{\pgfqpoint{0.000000in}{0.000000in}}%
\pgfpathlineto{\pgfqpoint{0.000000in}{-0.048611in}}%
\pgfusepath{stroke,fill}%
}%
\begin{pgfscope}%
\pgfsys@transformshift{2.135491in}{0.499444in}%
\pgfsys@useobject{currentmarker}{}%
\end{pgfscope}%
\end{pgfscope}%
\begin{pgfscope}%
\definecolor{textcolor}{rgb}{0.000000,0.000000,0.000000}%
\pgfsetstrokecolor{textcolor}%
\pgfsetfillcolor{textcolor}%
\pgftext[x=2.135491in,y=0.402222in,,top]{\color{textcolor}{\rmfamily\fontsize{10.000000}{12.000000}\selectfont\catcode`\^=\active\def^{\ifmmode\sp\else\^{}\fi}\catcode`\%=\active\def%{\%}$\mathdefault{1.0}$}}%
\end{pgfscope}%
\begin{pgfscope}%
\pgfpathrectangle{\pgfqpoint{0.553581in}{0.499444in}}{\pgfqpoint{5.391865in}{3.647034in}}%
\pgfusepath{clip}%
\pgfsetbuttcap%
\pgfsetroundjoin%
\pgfsetlinewidth{0.501875pt}%
\definecolor{currentstroke}{rgb}{0.501961,0.501961,0.501961}%
\pgfsetstrokecolor{currentstroke}%
\pgfsetstrokeopacity{0.500000}%
\pgfsetdash{{1.850000pt}{0.800000pt}}{0.000000pt}%
\pgfpathmoveto{\pgfqpoint{3.249513in}{0.499444in}}%
\pgfpathlineto{\pgfqpoint{3.249513in}{4.146478in}}%
\pgfusepath{stroke}%
\end{pgfscope}%
\begin{pgfscope}%
\pgfsetbuttcap%
\pgfsetroundjoin%
\definecolor{currentfill}{rgb}{0.000000,0.000000,0.000000}%
\pgfsetfillcolor{currentfill}%
\pgfsetlinewidth{0.803000pt}%
\definecolor{currentstroke}{rgb}{0.000000,0.000000,0.000000}%
\pgfsetstrokecolor{currentstroke}%
\pgfsetdash{}{0pt}%
\pgfsys@defobject{currentmarker}{\pgfqpoint{0.000000in}{-0.048611in}}{\pgfqpoint{0.000000in}{0.000000in}}{%
\pgfpathmoveto{\pgfqpoint{0.000000in}{0.000000in}}%
\pgfpathlineto{\pgfqpoint{0.000000in}{-0.048611in}}%
\pgfusepath{stroke,fill}%
}%
\begin{pgfscope}%
\pgfsys@transformshift{3.249513in}{0.499444in}%
\pgfsys@useobject{currentmarker}{}%
\end{pgfscope}%
\end{pgfscope}%
\begin{pgfscope}%
\definecolor{textcolor}{rgb}{0.000000,0.000000,0.000000}%
\pgfsetstrokecolor{textcolor}%
\pgfsetfillcolor{textcolor}%
\pgftext[x=3.249513in,y=0.402222in,,top]{\color{textcolor}{\rmfamily\fontsize{10.000000}{12.000000}\selectfont\catcode`\^=\active\def^{\ifmmode\sp\else\^{}\fi}\catcode`\%=\active\def%{\%}$\mathdefault{1.5}$}}%
\end{pgfscope}%
\begin{pgfscope}%
\pgfpathrectangle{\pgfqpoint{0.553581in}{0.499444in}}{\pgfqpoint{5.391865in}{3.647034in}}%
\pgfusepath{clip}%
\pgfsetbuttcap%
\pgfsetroundjoin%
\pgfsetlinewidth{0.501875pt}%
\definecolor{currentstroke}{rgb}{0.501961,0.501961,0.501961}%
\pgfsetstrokecolor{currentstroke}%
\pgfsetstrokeopacity{0.500000}%
\pgfsetdash{{1.850000pt}{0.800000pt}}{0.000000pt}%
\pgfpathmoveto{\pgfqpoint{4.363535in}{0.499444in}}%
\pgfpathlineto{\pgfqpoint{4.363535in}{4.146478in}}%
\pgfusepath{stroke}%
\end{pgfscope}%
\begin{pgfscope}%
\pgfsetbuttcap%
\pgfsetroundjoin%
\definecolor{currentfill}{rgb}{0.000000,0.000000,0.000000}%
\pgfsetfillcolor{currentfill}%
\pgfsetlinewidth{0.803000pt}%
\definecolor{currentstroke}{rgb}{0.000000,0.000000,0.000000}%
\pgfsetstrokecolor{currentstroke}%
\pgfsetdash{}{0pt}%
\pgfsys@defobject{currentmarker}{\pgfqpoint{0.000000in}{-0.048611in}}{\pgfqpoint{0.000000in}{0.000000in}}{%
\pgfpathmoveto{\pgfqpoint{0.000000in}{0.000000in}}%
\pgfpathlineto{\pgfqpoint{0.000000in}{-0.048611in}}%
\pgfusepath{stroke,fill}%
}%
\begin{pgfscope}%
\pgfsys@transformshift{4.363535in}{0.499444in}%
\pgfsys@useobject{currentmarker}{}%
\end{pgfscope}%
\end{pgfscope}%
\begin{pgfscope}%
\definecolor{textcolor}{rgb}{0.000000,0.000000,0.000000}%
\pgfsetstrokecolor{textcolor}%
\pgfsetfillcolor{textcolor}%
\pgftext[x=4.363535in,y=0.402222in,,top]{\color{textcolor}{\rmfamily\fontsize{10.000000}{12.000000}\selectfont\catcode`\^=\active\def^{\ifmmode\sp\else\^{}\fi}\catcode`\%=\active\def%{\%}$\mathdefault{2.0}$}}%
\end{pgfscope}%
\begin{pgfscope}%
\pgfpathrectangle{\pgfqpoint{0.553581in}{0.499444in}}{\pgfqpoint{5.391865in}{3.647034in}}%
\pgfusepath{clip}%
\pgfsetbuttcap%
\pgfsetroundjoin%
\pgfsetlinewidth{0.501875pt}%
\definecolor{currentstroke}{rgb}{0.501961,0.501961,0.501961}%
\pgfsetstrokecolor{currentstroke}%
\pgfsetstrokeopacity{0.500000}%
\pgfsetdash{{1.850000pt}{0.800000pt}}{0.000000pt}%
\pgfpathmoveto{\pgfqpoint{5.477556in}{0.499444in}}%
\pgfpathlineto{\pgfqpoint{5.477556in}{4.146478in}}%
\pgfusepath{stroke}%
\end{pgfscope}%
\begin{pgfscope}%
\pgfsetbuttcap%
\pgfsetroundjoin%
\definecolor{currentfill}{rgb}{0.000000,0.000000,0.000000}%
\pgfsetfillcolor{currentfill}%
\pgfsetlinewidth{0.803000pt}%
\definecolor{currentstroke}{rgb}{0.000000,0.000000,0.000000}%
\pgfsetstrokecolor{currentstroke}%
\pgfsetdash{}{0pt}%
\pgfsys@defobject{currentmarker}{\pgfqpoint{0.000000in}{-0.048611in}}{\pgfqpoint{0.000000in}{0.000000in}}{%
\pgfpathmoveto{\pgfqpoint{0.000000in}{0.000000in}}%
\pgfpathlineto{\pgfqpoint{0.000000in}{-0.048611in}}%
\pgfusepath{stroke,fill}%
}%
\begin{pgfscope}%
\pgfsys@transformshift{5.477556in}{0.499444in}%
\pgfsys@useobject{currentmarker}{}%
\end{pgfscope}%
\end{pgfscope}%
\begin{pgfscope}%
\definecolor{textcolor}{rgb}{0.000000,0.000000,0.000000}%
\pgfsetstrokecolor{textcolor}%
\pgfsetfillcolor{textcolor}%
\pgftext[x=5.477556in,y=0.402222in,,top]{\color{textcolor}{\rmfamily\fontsize{10.000000}{12.000000}\selectfont\catcode`\^=\active\def^{\ifmmode\sp\else\^{}\fi}\catcode`\%=\active\def%{\%}$\mathdefault{2.5}$}}%
\end{pgfscope}%
\begin{pgfscope}%
\definecolor{textcolor}{rgb}{0.000000,0.000000,0.000000}%
\pgfsetstrokecolor{textcolor}%
\pgfsetfillcolor{textcolor}%
\pgftext[x=3.249513in,y=0.223333in,,top]{\color{textcolor}{\rmfamily\fontsize{10.000000}{12.000000}\selectfont\catcode`\^=\active\def^{\ifmmode\sp\else\^{}\fi}\catcode`\%=\active\def%{\%}X}}%
\end{pgfscope}%
\begin{pgfscope}%
\pgfpathrectangle{\pgfqpoint{0.553581in}{0.499444in}}{\pgfqpoint{5.391865in}{3.647034in}}%
\pgfusepath{clip}%
\pgfsetbuttcap%
\pgfsetroundjoin%
\pgfsetlinewidth{0.501875pt}%
\definecolor{currentstroke}{rgb}{0.501961,0.501961,0.501961}%
\pgfsetstrokecolor{currentstroke}%
\pgfsetstrokeopacity{0.500000}%
\pgfsetdash{{1.850000pt}{0.800000pt}}{0.000000pt}%
\pgfpathmoveto{\pgfqpoint{0.553581in}{1.211727in}}%
\pgfpathlineto{\pgfqpoint{5.945445in}{1.211727in}}%
\pgfusepath{stroke}%
\end{pgfscope}%
\begin{pgfscope}%
\pgfsetbuttcap%
\pgfsetroundjoin%
\definecolor{currentfill}{rgb}{0.000000,0.000000,0.000000}%
\pgfsetfillcolor{currentfill}%
\pgfsetlinewidth{0.803000pt}%
\definecolor{currentstroke}{rgb}{0.000000,0.000000,0.000000}%
\pgfsetstrokecolor{currentstroke}%
\pgfsetdash{}{0pt}%
\pgfsys@defobject{currentmarker}{\pgfqpoint{-0.048611in}{0.000000in}}{\pgfqpoint{-0.000000in}{0.000000in}}{%
\pgfpathmoveto{\pgfqpoint{-0.000000in}{0.000000in}}%
\pgfpathlineto{\pgfqpoint{-0.048611in}{0.000000in}}%
\pgfusepath{stroke,fill}%
}%
\begin{pgfscope}%
\pgfsys@transformshift{0.553581in}{1.211727in}%
\pgfsys@useobject{currentmarker}{}%
\end{pgfscope}%
\end{pgfscope}%
\begin{pgfscope}%
\definecolor{textcolor}{rgb}{0.000000,0.000000,0.000000}%
\pgfsetstrokecolor{textcolor}%
\pgfsetfillcolor{textcolor}%
\pgftext[x=0.278889in, y=1.163533in, left, base]{\color{textcolor}{\rmfamily\fontsize{10.000000}{12.000000}\selectfont\catcode`\^=\active\def^{\ifmmode\sp\else\^{}\fi}\catcode`\%=\active\def%{\%}$\mathdefault{1.6}$}}%
\end{pgfscope}%
\begin{pgfscope}%
\pgfpathrectangle{\pgfqpoint{0.553581in}{0.499444in}}{\pgfqpoint{5.391865in}{3.647034in}}%
\pgfusepath{clip}%
\pgfsetbuttcap%
\pgfsetroundjoin%
\pgfsetlinewidth{0.501875pt}%
\definecolor{currentstroke}{rgb}{0.501961,0.501961,0.501961}%
\pgfsetstrokecolor{currentstroke}%
\pgfsetstrokeopacity{0.500000}%
\pgfsetdash{{1.850000pt}{0.800000pt}}{0.000000pt}%
\pgfpathmoveto{\pgfqpoint{0.553581in}{1.940405in}}%
\pgfpathlineto{\pgfqpoint{5.945445in}{1.940405in}}%
\pgfusepath{stroke}%
\end{pgfscope}%
\begin{pgfscope}%
\pgfsetbuttcap%
\pgfsetroundjoin%
\definecolor{currentfill}{rgb}{0.000000,0.000000,0.000000}%
\pgfsetfillcolor{currentfill}%
\pgfsetlinewidth{0.803000pt}%
\definecolor{currentstroke}{rgb}{0.000000,0.000000,0.000000}%
\pgfsetstrokecolor{currentstroke}%
\pgfsetdash{}{0pt}%
\pgfsys@defobject{currentmarker}{\pgfqpoint{-0.048611in}{0.000000in}}{\pgfqpoint{-0.000000in}{0.000000in}}{%
\pgfpathmoveto{\pgfqpoint{-0.000000in}{0.000000in}}%
\pgfpathlineto{\pgfqpoint{-0.048611in}{0.000000in}}%
\pgfusepath{stroke,fill}%
}%
\begin{pgfscope}%
\pgfsys@transformshift{0.553581in}{1.940405in}%
\pgfsys@useobject{currentmarker}{}%
\end{pgfscope}%
\end{pgfscope}%
\begin{pgfscope}%
\definecolor{textcolor}{rgb}{0.000000,0.000000,0.000000}%
\pgfsetstrokecolor{textcolor}%
\pgfsetfillcolor{textcolor}%
\pgftext[x=0.278889in, y=1.892211in, left, base]{\color{textcolor}{\rmfamily\fontsize{10.000000}{12.000000}\selectfont\catcode`\^=\active\def^{\ifmmode\sp\else\^{}\fi}\catcode`\%=\active\def%{\%}$\mathdefault{1.8}$}}%
\end{pgfscope}%
\begin{pgfscope}%
\pgfpathrectangle{\pgfqpoint{0.553581in}{0.499444in}}{\pgfqpoint{5.391865in}{3.647034in}}%
\pgfusepath{clip}%
\pgfsetbuttcap%
\pgfsetroundjoin%
\pgfsetlinewidth{0.501875pt}%
\definecolor{currentstroke}{rgb}{0.501961,0.501961,0.501961}%
\pgfsetstrokecolor{currentstroke}%
\pgfsetstrokeopacity{0.500000}%
\pgfsetdash{{1.850000pt}{0.800000pt}}{0.000000pt}%
\pgfpathmoveto{\pgfqpoint{0.553581in}{2.669083in}}%
\pgfpathlineto{\pgfqpoint{5.945445in}{2.669083in}}%
\pgfusepath{stroke}%
\end{pgfscope}%
\begin{pgfscope}%
\pgfsetbuttcap%
\pgfsetroundjoin%
\definecolor{currentfill}{rgb}{0.000000,0.000000,0.000000}%
\pgfsetfillcolor{currentfill}%
\pgfsetlinewidth{0.803000pt}%
\definecolor{currentstroke}{rgb}{0.000000,0.000000,0.000000}%
\pgfsetstrokecolor{currentstroke}%
\pgfsetdash{}{0pt}%
\pgfsys@defobject{currentmarker}{\pgfqpoint{-0.048611in}{0.000000in}}{\pgfqpoint{-0.000000in}{0.000000in}}{%
\pgfpathmoveto{\pgfqpoint{-0.000000in}{0.000000in}}%
\pgfpathlineto{\pgfqpoint{-0.048611in}{0.000000in}}%
\pgfusepath{stroke,fill}%
}%
\begin{pgfscope}%
\pgfsys@transformshift{0.553581in}{2.669083in}%
\pgfsys@useobject{currentmarker}{}%
\end{pgfscope}%
\end{pgfscope}%
\begin{pgfscope}%
\definecolor{textcolor}{rgb}{0.000000,0.000000,0.000000}%
\pgfsetstrokecolor{textcolor}%
\pgfsetfillcolor{textcolor}%
\pgftext[x=0.278889in, y=2.620889in, left, base]{\color{textcolor}{\rmfamily\fontsize{10.000000}{12.000000}\selectfont\catcode`\^=\active\def^{\ifmmode\sp\else\^{}\fi}\catcode`\%=\active\def%{\%}$\mathdefault{2.0}$}}%
\end{pgfscope}%
\begin{pgfscope}%
\pgfpathrectangle{\pgfqpoint{0.553581in}{0.499444in}}{\pgfqpoint{5.391865in}{3.647034in}}%
\pgfusepath{clip}%
\pgfsetbuttcap%
\pgfsetroundjoin%
\pgfsetlinewidth{0.501875pt}%
\definecolor{currentstroke}{rgb}{0.501961,0.501961,0.501961}%
\pgfsetstrokecolor{currentstroke}%
\pgfsetstrokeopacity{0.500000}%
\pgfsetdash{{1.850000pt}{0.800000pt}}{0.000000pt}%
\pgfpathmoveto{\pgfqpoint{0.553581in}{3.397762in}}%
\pgfpathlineto{\pgfqpoint{5.945445in}{3.397762in}}%
\pgfusepath{stroke}%
\end{pgfscope}%
\begin{pgfscope}%
\pgfsetbuttcap%
\pgfsetroundjoin%
\definecolor{currentfill}{rgb}{0.000000,0.000000,0.000000}%
\pgfsetfillcolor{currentfill}%
\pgfsetlinewidth{0.803000pt}%
\definecolor{currentstroke}{rgb}{0.000000,0.000000,0.000000}%
\pgfsetstrokecolor{currentstroke}%
\pgfsetdash{}{0pt}%
\pgfsys@defobject{currentmarker}{\pgfqpoint{-0.048611in}{0.000000in}}{\pgfqpoint{-0.000000in}{0.000000in}}{%
\pgfpathmoveto{\pgfqpoint{-0.000000in}{0.000000in}}%
\pgfpathlineto{\pgfqpoint{-0.048611in}{0.000000in}}%
\pgfusepath{stroke,fill}%
}%
\begin{pgfscope}%
\pgfsys@transformshift{0.553581in}{3.397762in}%
\pgfsys@useobject{currentmarker}{}%
\end{pgfscope}%
\end{pgfscope}%
\begin{pgfscope}%
\definecolor{textcolor}{rgb}{0.000000,0.000000,0.000000}%
\pgfsetstrokecolor{textcolor}%
\pgfsetfillcolor{textcolor}%
\pgftext[x=0.278889in, y=3.349567in, left, base]{\color{textcolor}{\rmfamily\fontsize{10.000000}{12.000000}\selectfont\catcode`\^=\active\def^{\ifmmode\sp\else\^{}\fi}\catcode`\%=\active\def%{\%}$\mathdefault{2.2}$}}%
\end{pgfscope}%
\begin{pgfscope}%
\pgfpathrectangle{\pgfqpoint{0.553581in}{0.499444in}}{\pgfqpoint{5.391865in}{3.647034in}}%
\pgfusepath{clip}%
\pgfsetbuttcap%
\pgfsetroundjoin%
\pgfsetlinewidth{0.501875pt}%
\definecolor{currentstroke}{rgb}{0.501961,0.501961,0.501961}%
\pgfsetstrokecolor{currentstroke}%
\pgfsetstrokeopacity{0.500000}%
\pgfsetdash{{1.850000pt}{0.800000pt}}{0.000000pt}%
\pgfpathmoveto{\pgfqpoint{0.553581in}{4.126440in}}%
\pgfpathlineto{\pgfqpoint{5.945445in}{4.126440in}}%
\pgfusepath{stroke}%
\end{pgfscope}%
\begin{pgfscope}%
\pgfsetbuttcap%
\pgfsetroundjoin%
\definecolor{currentfill}{rgb}{0.000000,0.000000,0.000000}%
\pgfsetfillcolor{currentfill}%
\pgfsetlinewidth{0.803000pt}%
\definecolor{currentstroke}{rgb}{0.000000,0.000000,0.000000}%
\pgfsetstrokecolor{currentstroke}%
\pgfsetdash{}{0pt}%
\pgfsys@defobject{currentmarker}{\pgfqpoint{-0.048611in}{0.000000in}}{\pgfqpoint{-0.000000in}{0.000000in}}{%
\pgfpathmoveto{\pgfqpoint{-0.000000in}{0.000000in}}%
\pgfpathlineto{\pgfqpoint{-0.048611in}{0.000000in}}%
\pgfusepath{stroke,fill}%
}%
\begin{pgfscope}%
\pgfsys@transformshift{0.553581in}{4.126440in}%
\pgfsys@useobject{currentmarker}{}%
\end{pgfscope}%
\end{pgfscope}%
\begin{pgfscope}%
\definecolor{textcolor}{rgb}{0.000000,0.000000,0.000000}%
\pgfsetstrokecolor{textcolor}%
\pgfsetfillcolor{textcolor}%
\pgftext[x=0.278889in, y=4.078245in, left, base]{\color{textcolor}{\rmfamily\fontsize{10.000000}{12.000000}\selectfont\catcode`\^=\active\def^{\ifmmode\sp\else\^{}\fi}\catcode`\%=\active\def%{\%}$\mathdefault{2.4}$}}%
\end{pgfscope}%
\begin{pgfscope}%
\definecolor{textcolor}{rgb}{0.000000,0.000000,0.000000}%
\pgfsetstrokecolor{textcolor}%
\pgfsetfillcolor{textcolor}%
\pgftext[x=0.223333in,y=2.322961in,,bottom,rotate=90.000000]{\color{textcolor}{\rmfamily\fontsize{10.000000}{12.000000}\selectfont\catcode`\^=\active\def^{\ifmmode\sp\else\^{}\fi}\catcode`\%=\active\def%{\%}Y}}%
\end{pgfscope}%
\begin{pgfscope}%
\pgfpathrectangle{\pgfqpoint{0.553581in}{0.499444in}}{\pgfqpoint{5.391865in}{3.647034in}}%
\pgfusepath{clip}%
\pgfsetbuttcap%
\pgfsetroundjoin%
\pgfsetlinewidth{1.505625pt}%
\definecolor{currentstroke}{rgb}{0.501961,0.501961,0.501961}%
\pgfsetstrokecolor{currentstroke}%
\pgfsetdash{}{0pt}%
\pgfpathmoveto{\pgfqpoint{0.798665in}{0.847388in}}%
\pgfpathlineto{\pgfqpoint{1.244274in}{0.847388in}}%
\pgfusepath{stroke}%
\end{pgfscope}%
\begin{pgfscope}%
\pgfpathrectangle{\pgfqpoint{0.553581in}{0.499444in}}{\pgfqpoint{5.391865in}{3.647034in}}%
\pgfusepath{clip}%
\pgfsetbuttcap%
\pgfsetroundjoin%
\pgfsetlinewidth{1.505625pt}%
\definecolor{currentstroke}{rgb}{0.501961,0.501961,0.501961}%
\pgfsetstrokecolor{currentstroke}%
\pgfsetdash{}{0pt}%
\pgfpathmoveto{\pgfqpoint{1.912687in}{1.940405in}}%
\pgfpathlineto{\pgfqpoint{2.358296in}{1.940405in}}%
\pgfusepath{stroke}%
\end{pgfscope}%
\begin{pgfscope}%
\pgfpathrectangle{\pgfqpoint{0.553581in}{0.499444in}}{\pgfqpoint{5.391865in}{3.647034in}}%
\pgfusepath{clip}%
\pgfsetbuttcap%
\pgfsetroundjoin%
\pgfsetlinewidth{1.505625pt}%
\definecolor{currentstroke}{rgb}{0.501961,0.501961,0.501961}%
\pgfsetstrokecolor{currentstroke}%
\pgfsetdash{}{0pt}%
\pgfpathmoveto{\pgfqpoint{3.026709in}{3.033423in}}%
\pgfpathlineto{\pgfqpoint{3.472317in}{3.033423in}}%
\pgfusepath{stroke}%
\end{pgfscope}%
\begin{pgfscope}%
\pgfpathrectangle{\pgfqpoint{0.553581in}{0.499444in}}{\pgfqpoint{5.391865in}{3.647034in}}%
\pgfusepath{clip}%
\pgfsetbuttcap%
\pgfsetroundjoin%
\pgfsetlinewidth{1.505625pt}%
\definecolor{currentstroke}{rgb}{0.501961,0.501961,0.501961}%
\pgfsetstrokecolor{currentstroke}%
\pgfsetdash{}{0pt}%
\pgfpathmoveto{\pgfqpoint{4.140730in}{2.669083in}}%
\pgfpathlineto{\pgfqpoint{4.586339in}{2.669083in}}%
\pgfusepath{stroke}%
\end{pgfscope}%
\begin{pgfscope}%
\pgfpathrectangle{\pgfqpoint{0.553581in}{0.499444in}}{\pgfqpoint{5.391865in}{3.647034in}}%
\pgfusepath{clip}%
\pgfsetbuttcap%
\pgfsetroundjoin%
\pgfsetlinewidth{1.505625pt}%
\definecolor{currentstroke}{rgb}{0.501961,0.501961,0.501961}%
\pgfsetstrokecolor{currentstroke}%
\pgfsetdash{}{0pt}%
\pgfpathmoveto{\pgfqpoint{5.254752in}{3.762101in}}%
\pgfpathlineto{\pgfqpoint{5.700360in}{3.762101in}}%
\pgfusepath{stroke}%
\end{pgfscope}%
\begin{pgfscope}%
\pgfpathrectangle{\pgfqpoint{0.553581in}{0.499444in}}{\pgfqpoint{5.391865in}{3.647034in}}%
\pgfusepath{clip}%
\pgfsetbuttcap%
\pgfsetroundjoin%
\pgfsetlinewidth{1.505625pt}%
\definecolor{currentstroke}{rgb}{0.501961,0.501961,0.501961}%
\pgfsetstrokecolor{currentstroke}%
\pgfsetdash{}{0pt}%
\pgfpathmoveto{\pgfqpoint{1.021470in}{0.665218in}}%
\pgfpathlineto{\pgfqpoint{1.021470in}{1.029558in}}%
\pgfusepath{stroke}%
\end{pgfscope}%
\begin{pgfscope}%
\pgfpathrectangle{\pgfqpoint{0.553581in}{0.499444in}}{\pgfqpoint{5.391865in}{3.647034in}}%
\pgfusepath{clip}%
\pgfsetbuttcap%
\pgfsetroundjoin%
\pgfsetlinewidth{1.505625pt}%
\definecolor{currentstroke}{rgb}{0.501961,0.501961,0.501961}%
\pgfsetstrokecolor{currentstroke}%
\pgfsetdash{}{0pt}%
\pgfpathmoveto{\pgfqpoint{2.135491in}{1.576066in}}%
\pgfpathlineto{\pgfqpoint{2.135491in}{2.304744in}}%
\pgfusepath{stroke}%
\end{pgfscope}%
\begin{pgfscope}%
\pgfpathrectangle{\pgfqpoint{0.553581in}{0.499444in}}{\pgfqpoint{5.391865in}{3.647034in}}%
\pgfusepath{clip}%
\pgfsetbuttcap%
\pgfsetroundjoin%
\pgfsetlinewidth{1.505625pt}%
\definecolor{currentstroke}{rgb}{0.501961,0.501961,0.501961}%
\pgfsetstrokecolor{currentstroke}%
\pgfsetdash{}{0pt}%
\pgfpathmoveto{\pgfqpoint{3.249513in}{2.741951in}}%
\pgfpathlineto{\pgfqpoint{3.249513in}{3.324894in}}%
\pgfusepath{stroke}%
\end{pgfscope}%
\begin{pgfscope}%
\pgfpathrectangle{\pgfqpoint{0.553581in}{0.499444in}}{\pgfqpoint{5.391865in}{3.647034in}}%
\pgfusepath{clip}%
\pgfsetbuttcap%
\pgfsetroundjoin%
\pgfsetlinewidth{1.505625pt}%
\definecolor{currentstroke}{rgb}{0.501961,0.501961,0.501961}%
\pgfsetstrokecolor{currentstroke}%
\pgfsetdash{}{0pt}%
\pgfpathmoveto{\pgfqpoint{4.363535in}{2.231877in}}%
\pgfpathlineto{\pgfqpoint{4.363535in}{3.106290in}}%
\pgfusepath{stroke}%
\end{pgfscope}%
\begin{pgfscope}%
\pgfpathrectangle{\pgfqpoint{0.553581in}{0.499444in}}{\pgfqpoint{5.391865in}{3.647034in}}%
\pgfusepath{clip}%
\pgfsetbuttcap%
\pgfsetroundjoin%
\pgfsetlinewidth{1.505625pt}%
\definecolor{currentstroke}{rgb}{0.501961,0.501961,0.501961}%
\pgfsetstrokecolor{currentstroke}%
\pgfsetdash{}{0pt}%
\pgfpathmoveto{\pgfqpoint{5.477556in}{3.543497in}}%
\pgfpathlineto{\pgfqpoint{5.477556in}{3.980704in}}%
\pgfusepath{stroke}%
\end{pgfscope}%
\begin{pgfscope}%
\pgfpathrectangle{\pgfqpoint{0.553581in}{0.499444in}}{\pgfqpoint{5.391865in}{3.647034in}}%
\pgfusepath{clip}%
\pgfsetbuttcap%
\pgfsetroundjoin%
\definecolor{currentfill}{rgb}{0.121569,0.466667,0.705882}%
\pgfsetfillcolor{currentfill}%
\pgfsetlinewidth{1.003750pt}%
\definecolor{currentstroke}{rgb}{0.501961,0.501961,0.501961}%
\pgfsetstrokecolor{currentstroke}%
\pgfsetdash{}{0pt}%
\pgfsys@defobject{currentmarker}{\pgfqpoint{0.000000in}{-0.041667in}}{\pgfqpoint{0.000000in}{0.041667in}}{%
\pgfpathmoveto{\pgfqpoint{0.000000in}{-0.041667in}}%
\pgfpathlineto{\pgfqpoint{0.000000in}{0.041667in}}%
\pgfusepath{stroke,fill}%
}%
\begin{pgfscope}%
\pgfsys@transformshift{0.798665in}{0.847388in}%
\pgfsys@useobject{currentmarker}{}%
\end{pgfscope}%
\begin{pgfscope}%
\pgfsys@transformshift{1.912687in}{1.940405in}%
\pgfsys@useobject{currentmarker}{}%
\end{pgfscope}%
\begin{pgfscope}%
\pgfsys@transformshift{3.026709in}{3.033423in}%
\pgfsys@useobject{currentmarker}{}%
\end{pgfscope}%
\begin{pgfscope}%
\pgfsys@transformshift{4.140730in}{2.669083in}%
\pgfsys@useobject{currentmarker}{}%
\end{pgfscope}%
\begin{pgfscope}%
\pgfsys@transformshift{5.254752in}{3.762101in}%
\pgfsys@useobject{currentmarker}{}%
\end{pgfscope}%
\end{pgfscope}%
\begin{pgfscope}%
\pgfpathrectangle{\pgfqpoint{0.553581in}{0.499444in}}{\pgfqpoint{5.391865in}{3.647034in}}%
\pgfusepath{clip}%
\pgfsetbuttcap%
\pgfsetroundjoin%
\definecolor{currentfill}{rgb}{0.121569,0.466667,0.705882}%
\pgfsetfillcolor{currentfill}%
\pgfsetlinewidth{1.003750pt}%
\definecolor{currentstroke}{rgb}{0.501961,0.501961,0.501961}%
\pgfsetstrokecolor{currentstroke}%
\pgfsetdash{}{0pt}%
\pgfsys@defobject{currentmarker}{\pgfqpoint{0.000000in}{-0.041667in}}{\pgfqpoint{0.000000in}{0.041667in}}{%
\pgfpathmoveto{\pgfqpoint{0.000000in}{-0.041667in}}%
\pgfpathlineto{\pgfqpoint{0.000000in}{0.041667in}}%
\pgfusepath{stroke,fill}%
}%
\begin{pgfscope}%
\pgfsys@transformshift{1.244274in}{0.847388in}%
\pgfsys@useobject{currentmarker}{}%
\end{pgfscope}%
\begin{pgfscope}%
\pgfsys@transformshift{2.358296in}{1.940405in}%
\pgfsys@useobject{currentmarker}{}%
\end{pgfscope}%
\begin{pgfscope}%
\pgfsys@transformshift{3.472317in}{3.033423in}%
\pgfsys@useobject{currentmarker}{}%
\end{pgfscope}%
\begin{pgfscope}%
\pgfsys@transformshift{4.586339in}{2.669083in}%
\pgfsys@useobject{currentmarker}{}%
\end{pgfscope}%
\begin{pgfscope}%
\pgfsys@transformshift{5.700360in}{3.762101in}%
\pgfsys@useobject{currentmarker}{}%
\end{pgfscope}%
\end{pgfscope}%
\begin{pgfscope}%
\pgfpathrectangle{\pgfqpoint{0.553581in}{0.499444in}}{\pgfqpoint{5.391865in}{3.647034in}}%
\pgfusepath{clip}%
\pgfsetbuttcap%
\pgfsetroundjoin%
\definecolor{currentfill}{rgb}{0.121569,0.466667,0.705882}%
\pgfsetfillcolor{currentfill}%
\pgfsetlinewidth{1.003750pt}%
\definecolor{currentstroke}{rgb}{0.501961,0.501961,0.501961}%
\pgfsetstrokecolor{currentstroke}%
\pgfsetdash{}{0pt}%
\pgfsys@defobject{currentmarker}{\pgfqpoint{-0.041667in}{-0.000000in}}{\pgfqpoint{0.041667in}{0.000000in}}{%
\pgfpathmoveto{\pgfqpoint{0.041667in}{-0.000000in}}%
\pgfpathlineto{\pgfqpoint{-0.041667in}{0.000000in}}%
\pgfusepath{stroke,fill}%
}%
\begin{pgfscope}%
\pgfsys@transformshift{1.021470in}{0.665218in}%
\pgfsys@useobject{currentmarker}{}%
\end{pgfscope}%
\begin{pgfscope}%
\pgfsys@transformshift{2.135491in}{1.576066in}%
\pgfsys@useobject{currentmarker}{}%
\end{pgfscope}%
\begin{pgfscope}%
\pgfsys@transformshift{3.249513in}{2.741951in}%
\pgfsys@useobject{currentmarker}{}%
\end{pgfscope}%
\begin{pgfscope}%
\pgfsys@transformshift{4.363535in}{2.231877in}%
\pgfsys@useobject{currentmarker}{}%
\end{pgfscope}%
\begin{pgfscope}%
\pgfsys@transformshift{5.477556in}{3.543497in}%
\pgfsys@useobject{currentmarker}{}%
\end{pgfscope}%
\end{pgfscope}%
\begin{pgfscope}%
\pgfpathrectangle{\pgfqpoint{0.553581in}{0.499444in}}{\pgfqpoint{5.391865in}{3.647034in}}%
\pgfusepath{clip}%
\pgfsetbuttcap%
\pgfsetroundjoin%
\definecolor{currentfill}{rgb}{0.121569,0.466667,0.705882}%
\pgfsetfillcolor{currentfill}%
\pgfsetlinewidth{1.003750pt}%
\definecolor{currentstroke}{rgb}{0.501961,0.501961,0.501961}%
\pgfsetstrokecolor{currentstroke}%
\pgfsetdash{}{0pt}%
\pgfsys@defobject{currentmarker}{\pgfqpoint{-0.041667in}{-0.000000in}}{\pgfqpoint{0.041667in}{0.000000in}}{%
\pgfpathmoveto{\pgfqpoint{0.041667in}{-0.000000in}}%
\pgfpathlineto{\pgfqpoint{-0.041667in}{0.000000in}}%
\pgfusepath{stroke,fill}%
}%
\begin{pgfscope}%
\pgfsys@transformshift{1.021470in}{1.029558in}%
\pgfsys@useobject{currentmarker}{}%
\end{pgfscope}%
\begin{pgfscope}%
\pgfsys@transformshift{2.135491in}{2.304744in}%
\pgfsys@useobject{currentmarker}{}%
\end{pgfscope}%
\begin{pgfscope}%
\pgfsys@transformshift{3.249513in}{3.324894in}%
\pgfsys@useobject{currentmarker}{}%
\end{pgfscope}%
\begin{pgfscope}%
\pgfsys@transformshift{4.363535in}{3.106290in}%
\pgfsys@useobject{currentmarker}{}%
\end{pgfscope}%
\begin{pgfscope}%
\pgfsys@transformshift{5.477556in}{3.980704in}%
\pgfsys@useobject{currentmarker}{}%
\end{pgfscope}%
\end{pgfscope}%
\begin{pgfscope}%
\pgfpathrectangle{\pgfqpoint{0.553581in}{0.499444in}}{\pgfqpoint{5.391865in}{3.647034in}}%
\pgfusepath{clip}%
\pgfsetrectcap%
\pgfsetroundjoin%
\pgfsetlinewidth{1.505625pt}%
\definecolor{currentstroke}{rgb}{0.121569,0.466667,0.705882}%
\pgfsetstrokecolor{currentstroke}%
\pgfsetdash{}{0pt}%
\pgfpathmoveto{\pgfqpoint{1.021470in}{0.847388in}}%
\pgfpathlineto{\pgfqpoint{2.135491in}{1.940405in}}%
\pgfpathlineto{\pgfqpoint{3.249513in}{3.033423in}}%
\pgfpathlineto{\pgfqpoint{4.363535in}{2.669083in}}%
\pgfpathlineto{\pgfqpoint{5.477556in}{3.762101in}}%
\pgfusepath{stroke}%
\end{pgfscope}%
\begin{pgfscope}%
\pgfpathrectangle{\pgfqpoint{0.553581in}{0.499444in}}{\pgfqpoint{5.391865in}{3.647034in}}%
\pgfusepath{clip}%
\pgfsetbuttcap%
\pgfsetmiterjoin%
\definecolor{currentfill}{rgb}{0.121569,0.466667,0.705882}%
\pgfsetfillcolor{currentfill}%
\pgfsetlinewidth{1.003750pt}%
\definecolor{currentstroke}{rgb}{0.121569,0.466667,0.705882}%
\pgfsetstrokecolor{currentstroke}%
\pgfsetdash{}{0pt}%
\pgfsys@defobject{currentmarker}{\pgfqpoint{-0.041667in}{-0.041667in}}{\pgfqpoint{0.041667in}{0.041667in}}{%
\pgfpathmoveto{\pgfqpoint{-0.041667in}{-0.041667in}}%
\pgfpathlineto{\pgfqpoint{0.041667in}{-0.041667in}}%
\pgfpathlineto{\pgfqpoint{0.041667in}{0.041667in}}%
\pgfpathlineto{\pgfqpoint{-0.041667in}{0.041667in}}%
\pgfpathlineto{\pgfqpoint{-0.041667in}{-0.041667in}}%
\pgfpathclose%
\pgfusepath{stroke,fill}%
}%
\begin{pgfscope}%
\pgfsys@transformshift{1.021470in}{0.847388in}%
\pgfsys@useobject{currentmarker}{}%
\end{pgfscope}%
\begin{pgfscope}%
\pgfsys@transformshift{2.135491in}{1.940405in}%
\pgfsys@useobject{currentmarker}{}%
\end{pgfscope}%
\begin{pgfscope}%
\pgfsys@transformshift{3.249513in}{3.033423in}%
\pgfsys@useobject{currentmarker}{}%
\end{pgfscope}%
\begin{pgfscope}%
\pgfsys@transformshift{4.363535in}{2.669083in}%
\pgfsys@useobject{currentmarker}{}%
\end{pgfscope}%
\begin{pgfscope}%
\pgfsys@transformshift{5.477556in}{3.762101in}%
\pgfsys@useobject{currentmarker}{}%
\end{pgfscope}%
\end{pgfscope}%
\begin{pgfscope}%
\pgfsetrectcap%
\pgfsetmiterjoin%
\pgfsetlinewidth{0.803000pt}%
\definecolor{currentstroke}{rgb}{0.000000,0.000000,0.000000}%
\pgfsetstrokecolor{currentstroke}%
\pgfsetdash{}{0pt}%
\pgfpathmoveto{\pgfqpoint{0.553581in}{0.499444in}}%
\pgfpathlineto{\pgfqpoint{0.553581in}{4.146478in}}%
\pgfusepath{stroke}%
\end{pgfscope}%
\begin{pgfscope}%
\pgfsetrectcap%
\pgfsetmiterjoin%
\pgfsetlinewidth{0.803000pt}%
\definecolor{currentstroke}{rgb}{0.000000,0.000000,0.000000}%
\pgfsetstrokecolor{currentstroke}%
\pgfsetdash{}{0pt}%
\pgfpathmoveto{\pgfqpoint{5.945445in}{0.499444in}}%
\pgfpathlineto{\pgfqpoint{5.945445in}{4.146478in}}%
\pgfusepath{stroke}%
\end{pgfscope}%
\begin{pgfscope}%
\pgfsetrectcap%
\pgfsetmiterjoin%
\pgfsetlinewidth{0.803000pt}%
\definecolor{currentstroke}{rgb}{0.000000,0.000000,0.000000}%
\pgfsetstrokecolor{currentstroke}%
\pgfsetdash{}{0pt}%
\pgfpathmoveto{\pgfqpoint{0.553581in}{0.499444in}}%
\pgfpathlineto{\pgfqpoint{5.945445in}{0.499444in}}%
\pgfusepath{stroke}%
\end{pgfscope}%
\begin{pgfscope}%
\pgfsetrectcap%
\pgfsetmiterjoin%
\pgfsetlinewidth{0.803000pt}%
\definecolor{currentstroke}{rgb}{0.000000,0.000000,0.000000}%
\pgfsetstrokecolor{currentstroke}%
\pgfsetdash{}{0pt}%
\pgfpathmoveto{\pgfqpoint{0.553581in}{4.146478in}}%
\pgfpathlineto{\pgfqpoint{5.945445in}{4.146478in}}%
\pgfusepath{stroke}%
\end{pgfscope}%
\begin{pgfscope}%
\definecolor{textcolor}{rgb}{0.000000,0.000000,0.000000}%
\pgfsetstrokecolor{textcolor}%
\pgfsetfillcolor{textcolor}%
\pgftext[x=3.249513in,y=4.229812in,,base]{\color{textcolor}{\rmfamily\fontsize{12.000000}{14.400000}\selectfont\catcode`\^=\active\def^{\ifmmode\sp\else\^{}\fi}\catcode`\%=\active\def%{\%}Barras de error en ambos ejes}}%
\end{pgfscope}%
\begin{pgfscope}%
\pgfsetbuttcap%
\pgfsetmiterjoin%
\definecolor{currentfill}{rgb}{1.000000,1.000000,1.000000}%
\pgfsetfillcolor{currentfill}%
\pgfsetfillopacity{0.800000}%
\pgfsetlinewidth{1.003750pt}%
\definecolor{currentstroke}{rgb}{0.800000,0.800000,0.800000}%
\pgfsetstrokecolor{currentstroke}%
\pgfsetstrokeopacity{0.800000}%
\pgfsetdash{}{0pt}%
\pgfpathmoveto{\pgfqpoint{0.650803in}{3.841756in}}%
\pgfpathlineto{\pgfqpoint{2.015803in}{3.841756in}}%
\pgfpathquadraticcurveto{\pgfqpoint{2.043581in}{3.841756in}}{\pgfqpoint{2.043581in}{3.869534in}}%
\pgfpathlineto{\pgfqpoint{2.043581in}{4.049256in}}%
\pgfpathquadraticcurveto{\pgfqpoint{2.043581in}{4.077034in}}{\pgfqpoint{2.015803in}{4.077034in}}%
\pgfpathlineto{\pgfqpoint{0.650803in}{4.077034in}}%
\pgfpathquadraticcurveto{\pgfqpoint{0.623025in}{4.077034in}}{\pgfqpoint{0.623025in}{4.049256in}}%
\pgfpathlineto{\pgfqpoint{0.623025in}{3.869534in}}%
\pgfpathquadraticcurveto{\pgfqpoint{0.623025in}{3.841756in}}{\pgfqpoint{0.650803in}{3.841756in}}%
\pgfpathlineto{\pgfqpoint{0.650803in}{3.841756in}}%
\pgfpathclose%
\pgfusepath{stroke,fill}%
\end{pgfscope}%
\begin{pgfscope}%
\pgfsetbuttcap%
\pgfsetroundjoin%
\pgfsetlinewidth{1.505625pt}%
\definecolor{currentstroke}{rgb}{0.501961,0.501961,0.501961}%
\pgfsetstrokecolor{currentstroke}%
\pgfsetdash{}{0pt}%
\pgfpathmoveto{\pgfqpoint{0.748025in}{3.972867in}}%
\pgfpathlineto{\pgfqpoint{0.886914in}{3.972867in}}%
\pgfusepath{stroke}%
\end{pgfscope}%
\begin{pgfscope}%
\pgfsetbuttcap%
\pgfsetroundjoin%
\pgfsetlinewidth{1.505625pt}%
\definecolor{currentstroke}{rgb}{0.501961,0.501961,0.501961}%
\pgfsetstrokecolor{currentstroke}%
\pgfsetdash{}{0pt}%
\pgfpathmoveto{\pgfqpoint{0.817470in}{3.903423in}}%
\pgfpathlineto{\pgfqpoint{0.817470in}{4.042312in}}%
\pgfusepath{stroke}%
\end{pgfscope}%
\begin{pgfscope}%
\pgfsetbuttcap%
\pgfsetroundjoin%
\definecolor{currentfill}{rgb}{0.121569,0.466667,0.705882}%
\pgfsetfillcolor{currentfill}%
\pgfsetlinewidth{1.003750pt}%
\definecolor{currentstroke}{rgb}{0.501961,0.501961,0.501961}%
\pgfsetstrokecolor{currentstroke}%
\pgfsetdash{}{0pt}%
\pgfsys@defobject{currentmarker}{\pgfqpoint{0.000000in}{-0.041667in}}{\pgfqpoint{0.000000in}{0.041667in}}{%
\pgfpathmoveto{\pgfqpoint{0.000000in}{-0.041667in}}%
\pgfpathlineto{\pgfqpoint{0.000000in}{0.041667in}}%
\pgfusepath{stroke,fill}%
}%
\begin{pgfscope}%
\pgfsys@transformshift{0.748025in}{3.972867in}%
\pgfsys@useobject{currentmarker}{}%
\end{pgfscope}%
\end{pgfscope}%
\begin{pgfscope}%
\pgfsetbuttcap%
\pgfsetroundjoin%
\definecolor{currentfill}{rgb}{0.121569,0.466667,0.705882}%
\pgfsetfillcolor{currentfill}%
\pgfsetlinewidth{1.003750pt}%
\definecolor{currentstroke}{rgb}{0.501961,0.501961,0.501961}%
\pgfsetstrokecolor{currentstroke}%
\pgfsetdash{}{0pt}%
\pgfsys@defobject{currentmarker}{\pgfqpoint{0.000000in}{-0.041667in}}{\pgfqpoint{0.000000in}{0.041667in}}{%
\pgfpathmoveto{\pgfqpoint{0.000000in}{-0.041667in}}%
\pgfpathlineto{\pgfqpoint{0.000000in}{0.041667in}}%
\pgfusepath{stroke,fill}%
}%
\begin{pgfscope}%
\pgfsys@transformshift{0.886914in}{3.972867in}%
\pgfsys@useobject{currentmarker}{}%
\end{pgfscope}%
\end{pgfscope}%
\begin{pgfscope}%
\pgfsetbuttcap%
\pgfsetroundjoin%
\definecolor{currentfill}{rgb}{0.121569,0.466667,0.705882}%
\pgfsetfillcolor{currentfill}%
\pgfsetlinewidth{1.003750pt}%
\definecolor{currentstroke}{rgb}{0.501961,0.501961,0.501961}%
\pgfsetstrokecolor{currentstroke}%
\pgfsetdash{}{0pt}%
\pgfsys@defobject{currentmarker}{\pgfqpoint{-0.041667in}{-0.000000in}}{\pgfqpoint{0.041667in}{0.000000in}}{%
\pgfpathmoveto{\pgfqpoint{0.041667in}{-0.000000in}}%
\pgfpathlineto{\pgfqpoint{-0.041667in}{0.000000in}}%
\pgfusepath{stroke,fill}%
}%
\begin{pgfscope}%
\pgfsys@transformshift{0.817470in}{3.903423in}%
\pgfsys@useobject{currentmarker}{}%
\end{pgfscope}%
\end{pgfscope}%
\begin{pgfscope}%
\pgfsetbuttcap%
\pgfsetroundjoin%
\definecolor{currentfill}{rgb}{0.121569,0.466667,0.705882}%
\pgfsetfillcolor{currentfill}%
\pgfsetlinewidth{1.003750pt}%
\definecolor{currentstroke}{rgb}{0.501961,0.501961,0.501961}%
\pgfsetstrokecolor{currentstroke}%
\pgfsetdash{}{0pt}%
\pgfsys@defobject{currentmarker}{\pgfqpoint{-0.041667in}{-0.000000in}}{\pgfqpoint{0.041667in}{0.000000in}}{%
\pgfpathmoveto{\pgfqpoint{0.041667in}{-0.000000in}}%
\pgfpathlineto{\pgfqpoint{-0.041667in}{0.000000in}}%
\pgfusepath{stroke,fill}%
}%
\begin{pgfscope}%
\pgfsys@transformshift{0.817470in}{4.042312in}%
\pgfsys@useobject{currentmarker}{}%
\end{pgfscope}%
\end{pgfscope}%
\begin{pgfscope}%
\pgfsetrectcap%
\pgfsetroundjoin%
\pgfsetlinewidth{1.505625pt}%
\definecolor{currentstroke}{rgb}{0.121569,0.466667,0.705882}%
\pgfsetstrokecolor{currentstroke}%
\pgfsetdash{}{0pt}%
\pgfpathmoveto{\pgfqpoint{0.678581in}{3.972867in}}%
\pgfpathlineto{\pgfqpoint{0.956358in}{3.972867in}}%
\pgfusepath{stroke}%
\end{pgfscope}%
\begin{pgfscope}%
\pgfsetbuttcap%
\pgfsetmiterjoin%
\definecolor{currentfill}{rgb}{0.121569,0.466667,0.705882}%
\pgfsetfillcolor{currentfill}%
\pgfsetlinewidth{1.003750pt}%
\definecolor{currentstroke}{rgb}{0.121569,0.466667,0.705882}%
\pgfsetstrokecolor{currentstroke}%
\pgfsetdash{}{0pt}%
\pgfsys@defobject{currentmarker}{\pgfqpoint{-0.041667in}{-0.041667in}}{\pgfqpoint{0.041667in}{0.041667in}}{%
\pgfpathmoveto{\pgfqpoint{-0.041667in}{-0.041667in}}%
\pgfpathlineto{\pgfqpoint{0.041667in}{-0.041667in}}%
\pgfpathlineto{\pgfqpoint{0.041667in}{0.041667in}}%
\pgfpathlineto{\pgfqpoint{-0.041667in}{0.041667in}}%
\pgfpathlineto{\pgfqpoint{-0.041667in}{-0.041667in}}%
\pgfpathclose%
\pgfusepath{stroke,fill}%
}%
\begin{pgfscope}%
\pgfsys@transformshift{0.817470in}{3.972867in}%
\pgfsys@useobject{currentmarker}{}%
\end{pgfscope}%
\end{pgfscope}%
\begin{pgfscope}%
\definecolor{textcolor}{rgb}{0.000000,0.000000,0.000000}%
\pgfsetstrokecolor{textcolor}%
\pgfsetfillcolor{textcolor}%
\pgftext[x=1.067470in,y=3.924256in,left,base]{\color{textcolor}{\rmfamily\fontsize{10.000000}{12.000000}\selectfont\catcode`\^=\active\def^{\ifmmode\sp\else\^{}\fi}\catcode`\%=\active\def%{\%}Error en X e Y}}%
\end{pgfscope}%
\end{pgfpicture}%
\makeatother%
\endgroup%
 
    \caption{Gráfico de ejes compartidos}
    \label{fig:grafico_diagrama_error}
\end{figure}



\subsection{Diseño con \texttt{GridSpec}}

GridSpec es un módulo de \textit{Matplotlib} que permite al usuario realizar la arquitectura de la figura. Si usted hiciera \texttt{[plt.GridSpec(3, 3)]}, lo que esta haciendo es trazar líneas imaginarias sobre su figura para dividirla en 9 áreas iguales (3 filas y 3 columnas).

\subsubsection{Asignación de espacios}

La diferencia de \texttt{.GridSpect} con la utilización de objetos de \textit{Matplotlib} esta en la forma en la que se pueden \textit{asignar los espacios}. Usted puede decirle a cada gráfico (\texttt{ax}) en qué coordenadas va, para esop utilizamos el \textit{Slicing}. El grafico puede tomar más de un área.

Si Usted quiere que su grafico ocupe toda la parte de arriba, es decir las 3 columnas de la primera fila, utiliza \texttt{[0, :]}. Si quiere que el gráfico ocupe toda una columna, por ejemplo la ultima columna, utiliza \texttt{[:, 2]}. Si no termino de entender, aquí le dejo otras posibles posiciones:
\begin{itemize}
    \item \texttt{[0, 0]}: Solo la celda de la esquina superior izquierda.
    \item \texttt{[1, :]}: Toda la segunda fila completa.
    \item \texttt{[:, -1]}: Toda la última columna completa.
    \item \texttt{[1:, :2]}: Desde la segunda fila hasta el final, y desde la primera columna hasta la segunda.
\end{itemize}

La diferencia con \texttt{subplots} es que todas las grillas son flexibles, por lo tanto puede enfatizar gráficos asignandole más espacios.

Para utilizar \texttt{GridSpec} se suele utilizar la siguiente estructura:
\begin{verbatim}
gs = GridSpec(nrows=?, ncols=?, figure=?)
\end{verbatim}

Para controlar las proporciones, es decir determinar el tamaño de las columnas o filas, puede utilizar dentro de \texttt{GridSpec} los parametros de \texttt{width\_ratios} y \texttt{height\_ratios}. 
\begin{ejemplo}
    Si creara una grilla con \texttt{GridSpec} de \texttt{2 x 2} y quisiera que la primera columna sea tres veces más ancha que la segunda columna, utilizara \texttt{width\_ratios=[3,1]}. Si quiere que la segunda fila sea el doble de alta que la primera fila utilizaria \texttt{height\_ratios=[1,2]}
\end{ejemplo}

\begin{pycode}{Estructura GridSpec}
import numpy as np
import matplotlib.pyplot as plt
from matplotlib.gridspec import GridSpec

# Datos
t = np.linspace(0, 6 * np.pi, 600)

# Creamos la figura
fig = plt.figure(figsize=(10, 6), layout='constrained')

# Creamos un GridSpec de 2 filas y 3 columnas con dimensiones personalizadas
gs = GridSpec(
    nrows=2,
    ncols=3,
    figure=fig,
    width_ratios=[1, 2, 2],
    height_ratios=[1, 1]
)

# Primer Gráfico:
ax_left = fig.add_subplot(gs[:, 0]) # Ocupa todas las filas de la primera columna.
ax_left.plot(t, np.sin(t), linewidth=1)
ax_left.set_title("Panel izquierdo (2 filas)")

# Segundo Gráfico
ax_c_top = fig.add_subplot(gs[0, 1]) # Primera fila segunda columna
ax_c_top.plot(t, np.cos(t), color="orange")
ax_c_top.set_title("Centro arriba")

# Tercer Gráfico
ax_c_bot = fig.add_subplot(gs[1, :]) # Segunda fila columna 1 y 2
ax_c_bot.plot(t, np.sin(2 * t), color="green")
ax_c_bot.set_title("Centro abajo")

# Cuarto Gráfico
ax_right = fig.add_subplot(gs[0, 2]) # Primera fila, 3era columna
ax_right.plot(t, np.cos(2 * t), color="red")
ax_right.set_title("Derecha arriba")

# Título general
fig.suptitle("GridSpec con anchos asimétricos y spans de filas")
plt.show()
\end{pycode}

El código deja la figura \ref{fig:grafico_con_GridSpec}.

\begin{nota}
    Tambien se puede utilizar las funciones de \texttt{sharex} y \texttt{sharey} dentro de grid spec.
\end{nota}

\begin{figure}[ht]
    \centering
    %% Creator: Matplotlib, PGF backend
%%
%% To include the figure in your LaTeX document, write
%%   \input{<filename>.pgf}
%%
%% Make sure the required packages are loaded in your preamble
%%   \usepackage{pgf}
%%
%% Also ensure that all the required font packages are loaded; for instance,
%% the lmodern package is sometimes necessary when using math font.
%%   \usepackage{lmodern}
%%
%% Figures using additional raster images can only be included by \input if
%% they are in the same directory as the main LaTeX file. For loading figures
%% from other directories you can use the `import` package
%%   \usepackage{import}
%%
%% and then include the figures with
%%   \import{<path to file>}{<filename>.pgf}
%%
%% Matplotlib used the following preamble
%%   \def\mathdefault#1{#1}
%%   \everymath=\expandafter{\the\everymath\displaystyle}
%%   \IfFileExists{scrextend.sty}{
%%     \usepackage[fontsize=10.000000pt]{scrextend}
%%   }{
%%     \renewcommand{\normalsize}{\fontsize{10.000000}{12.000000}\selectfont}
%%     \normalsize
%%   }
%%   \usepackage{fontspec}
%%   \setmainfont{CMU Serif}
%%   \makeatletter\@ifpackageloaded{underscore}{}{\usepackage[strings]{underscore}}\makeatother
%%
\begingroup%
\makeatletter%
\begin{pgfpicture}%
\pgfpathrectangle{\pgfpointorigin}{\pgfqpoint{5.692717in}{4.822560in}}%
\pgfusepath{use as bounding box, clip}%
\begin{pgfscope}%
\pgfsetbuttcap%
\pgfsetmiterjoin%
\definecolor{currentfill}{rgb}{1.000000,1.000000,1.000000}%
\pgfsetfillcolor{currentfill}%
\pgfsetlinewidth{0.000000pt}%
\definecolor{currentstroke}{rgb}{1.000000,1.000000,1.000000}%
\pgfsetstrokecolor{currentstroke}%
\pgfsetdash{}{0pt}%
\pgfpathmoveto{\pgfqpoint{0.000000in}{0.000000in}}%
\pgfpathlineto{\pgfqpoint{5.692717in}{0.000000in}}%
\pgfpathlineto{\pgfqpoint{5.692717in}{4.822560in}}%
\pgfpathlineto{\pgfqpoint{0.000000in}{4.822560in}}%
\pgfpathlineto{\pgfqpoint{0.000000in}{0.000000in}}%
\pgfpathclose%
\pgfusepath{fill}%
\end{pgfscope}%
\begin{pgfscope}%
\pgfsetbuttcap%
\pgfsetmiterjoin%
\definecolor{currentfill}{rgb}{1.000000,1.000000,1.000000}%
\pgfsetfillcolor{currentfill}%
\pgfsetlinewidth{0.000000pt}%
\definecolor{currentstroke}{rgb}{0.000000,0.000000,0.000000}%
\pgfsetstrokecolor{currentstroke}%
\pgfsetstrokeopacity{0.000000}%
\pgfsetdash{}{0pt}%
\pgfpathmoveto{\pgfqpoint{0.599108in}{0.320555in}}%
\pgfpathlineto{\pgfqpoint{1.334059in}{0.320555in}}%
\pgfpathlineto{\pgfqpoint{1.334059in}{3.886945in}}%
\pgfpathlineto{\pgfqpoint{0.599108in}{3.886945in}}%
\pgfpathlineto{\pgfqpoint{0.599108in}{0.320555in}}%
\pgfpathclose%
\pgfusepath{fill}%
\end{pgfscope}%
\begin{pgfscope}%
\pgfpathrectangle{\pgfqpoint{0.599108in}{0.320555in}}{\pgfqpoint{0.734951in}{3.566389in}}%
\pgfusepath{clip}%
\pgfsetbuttcap%
\pgfsetroundjoin%
\pgfsetlinewidth{0.501875pt}%
\definecolor{currentstroke}{rgb}{0.501961,0.501961,0.501961}%
\pgfsetstrokecolor{currentstroke}%
\pgfsetstrokeopacity{0.500000}%
\pgfsetdash{{1.850000pt}{0.800000pt}}{0.000000pt}%
\pgfpathmoveto{\pgfqpoint{0.632515in}{0.320555in}}%
\pgfpathlineto{\pgfqpoint{0.632515in}{3.886945in}}%
\pgfusepath{stroke}%
\end{pgfscope}%
\begin{pgfscope}%
\pgfsetbuttcap%
\pgfsetroundjoin%
\definecolor{currentfill}{rgb}{0.000000,0.000000,0.000000}%
\pgfsetfillcolor{currentfill}%
\pgfsetlinewidth{0.803000pt}%
\definecolor{currentstroke}{rgb}{0.000000,0.000000,0.000000}%
\pgfsetstrokecolor{currentstroke}%
\pgfsetdash{}{0pt}%
\pgfsys@defobject{currentmarker}{\pgfqpoint{0.000000in}{-0.048611in}}{\pgfqpoint{0.000000in}{0.000000in}}{%
\pgfpathmoveto{\pgfqpoint{0.000000in}{0.000000in}}%
\pgfpathlineto{\pgfqpoint{0.000000in}{-0.048611in}}%
\pgfusepath{stroke,fill}%
}%
\begin{pgfscope}%
\pgfsys@transformshift{0.632515in}{0.320555in}%
\pgfsys@useobject{currentmarker}{}%
\end{pgfscope}%
\end{pgfscope}%
\begin{pgfscope}%
\definecolor{textcolor}{rgb}{0.000000,0.000000,0.000000}%
\pgfsetstrokecolor{textcolor}%
\pgfsetfillcolor{textcolor}%
\pgftext[x=0.632515in,y=0.223333in,,top]{\color{textcolor}{\rmfamily\fontsize{10.000000}{12.000000}\selectfont\catcode`\^=\active\def^{\ifmmode\sp\else\^{}\fi}\catcode`\%=\active\def%{\%}$\mathdefault{0}$}}%
\end{pgfscope}%
\begin{pgfscope}%
\pgfpathrectangle{\pgfqpoint{0.599108in}{0.320555in}}{\pgfqpoint{0.734951in}{3.566389in}}%
\pgfusepath{clip}%
\pgfsetbuttcap%
\pgfsetroundjoin%
\pgfsetlinewidth{0.501875pt}%
\definecolor{currentstroke}{rgb}{0.501961,0.501961,0.501961}%
\pgfsetstrokecolor{currentstroke}%
\pgfsetstrokeopacity{0.500000}%
\pgfsetdash{{1.850000pt}{0.800000pt}}{0.000000pt}%
\pgfpathmoveto{\pgfqpoint{0.986973in}{0.320555in}}%
\pgfpathlineto{\pgfqpoint{0.986973in}{3.886945in}}%
\pgfusepath{stroke}%
\end{pgfscope}%
\begin{pgfscope}%
\pgfsetbuttcap%
\pgfsetroundjoin%
\definecolor{currentfill}{rgb}{0.000000,0.000000,0.000000}%
\pgfsetfillcolor{currentfill}%
\pgfsetlinewidth{0.803000pt}%
\definecolor{currentstroke}{rgb}{0.000000,0.000000,0.000000}%
\pgfsetstrokecolor{currentstroke}%
\pgfsetdash{}{0pt}%
\pgfsys@defobject{currentmarker}{\pgfqpoint{0.000000in}{-0.048611in}}{\pgfqpoint{0.000000in}{0.000000in}}{%
\pgfpathmoveto{\pgfqpoint{0.000000in}{0.000000in}}%
\pgfpathlineto{\pgfqpoint{0.000000in}{-0.048611in}}%
\pgfusepath{stroke,fill}%
}%
\begin{pgfscope}%
\pgfsys@transformshift{0.986973in}{0.320555in}%
\pgfsys@useobject{currentmarker}{}%
\end{pgfscope}%
\end{pgfscope}%
\begin{pgfscope}%
\definecolor{textcolor}{rgb}{0.000000,0.000000,0.000000}%
\pgfsetstrokecolor{textcolor}%
\pgfsetfillcolor{textcolor}%
\pgftext[x=0.986973in,y=0.223333in,,top]{\color{textcolor}{\rmfamily\fontsize{10.000000}{12.000000}\selectfont\catcode`\^=\active\def^{\ifmmode\sp\else\^{}\fi}\catcode`\%=\active\def%{\%}$\mathdefault{10}$}}%
\end{pgfscope}%
\begin{pgfscope}%
\pgfpathrectangle{\pgfqpoint{0.599108in}{0.320555in}}{\pgfqpoint{0.734951in}{3.566389in}}%
\pgfusepath{clip}%
\pgfsetbuttcap%
\pgfsetroundjoin%
\pgfsetlinewidth{0.501875pt}%
\definecolor{currentstroke}{rgb}{0.501961,0.501961,0.501961}%
\pgfsetstrokecolor{currentstroke}%
\pgfsetstrokeopacity{0.500000}%
\pgfsetdash{{1.850000pt}{0.800000pt}}{0.000000pt}%
\pgfpathmoveto{\pgfqpoint{0.599108in}{0.482658in}}%
\pgfpathlineto{\pgfqpoint{1.334059in}{0.482658in}}%
\pgfusepath{stroke}%
\end{pgfscope}%
\begin{pgfscope}%
\pgfsetbuttcap%
\pgfsetroundjoin%
\definecolor{currentfill}{rgb}{0.000000,0.000000,0.000000}%
\pgfsetfillcolor{currentfill}%
\pgfsetlinewidth{0.803000pt}%
\definecolor{currentstroke}{rgb}{0.000000,0.000000,0.000000}%
\pgfsetstrokecolor{currentstroke}%
\pgfsetdash{}{0pt}%
\pgfsys@defobject{currentmarker}{\pgfqpoint{-0.048611in}{0.000000in}}{\pgfqpoint{-0.000000in}{0.000000in}}{%
\pgfpathmoveto{\pgfqpoint{-0.000000in}{0.000000in}}%
\pgfpathlineto{\pgfqpoint{-0.048611in}{0.000000in}}%
\pgfusepath{stroke,fill}%
}%
\begin{pgfscope}%
\pgfsys@transformshift{0.599108in}{0.482658in}%
\pgfsys@useobject{currentmarker}{}%
\end{pgfscope}%
\end{pgfscope}%
\begin{pgfscope}%
\definecolor{textcolor}{rgb}{0.000000,0.000000,0.000000}%
\pgfsetstrokecolor{textcolor}%
\pgfsetfillcolor{textcolor}%
\pgftext[x=0.146947in, y=0.434464in, left, base]{\color{textcolor}{\rmfamily\fontsize{10.000000}{12.000000}\selectfont\catcode`\^=\active\def^{\ifmmode\sp\else\^{}\fi}\catcode`\%=\active\def%{\%}$\mathdefault{\ensuremath{-}1.00}$}}%
\end{pgfscope}%
\begin{pgfscope}%
\pgfpathrectangle{\pgfqpoint{0.599108in}{0.320555in}}{\pgfqpoint{0.734951in}{3.566389in}}%
\pgfusepath{clip}%
\pgfsetbuttcap%
\pgfsetroundjoin%
\pgfsetlinewidth{0.501875pt}%
\definecolor{currentstroke}{rgb}{0.501961,0.501961,0.501961}%
\pgfsetstrokecolor{currentstroke}%
\pgfsetstrokeopacity{0.500000}%
\pgfsetdash{{1.850000pt}{0.800000pt}}{0.000000pt}%
\pgfpathmoveto{\pgfqpoint{0.599108in}{0.887931in}}%
\pgfpathlineto{\pgfqpoint{1.334059in}{0.887931in}}%
\pgfusepath{stroke}%
\end{pgfscope}%
\begin{pgfscope}%
\pgfsetbuttcap%
\pgfsetroundjoin%
\definecolor{currentfill}{rgb}{0.000000,0.000000,0.000000}%
\pgfsetfillcolor{currentfill}%
\pgfsetlinewidth{0.803000pt}%
\definecolor{currentstroke}{rgb}{0.000000,0.000000,0.000000}%
\pgfsetstrokecolor{currentstroke}%
\pgfsetdash{}{0pt}%
\pgfsys@defobject{currentmarker}{\pgfqpoint{-0.048611in}{0.000000in}}{\pgfqpoint{-0.000000in}{0.000000in}}{%
\pgfpathmoveto{\pgfqpoint{-0.000000in}{0.000000in}}%
\pgfpathlineto{\pgfqpoint{-0.048611in}{0.000000in}}%
\pgfusepath{stroke,fill}%
}%
\begin{pgfscope}%
\pgfsys@transformshift{0.599108in}{0.887931in}%
\pgfsys@useobject{currentmarker}{}%
\end{pgfscope}%
\end{pgfscope}%
\begin{pgfscope}%
\definecolor{textcolor}{rgb}{0.000000,0.000000,0.000000}%
\pgfsetstrokecolor{textcolor}%
\pgfsetfillcolor{textcolor}%
\pgftext[x=0.146947in, y=0.839737in, left, base]{\color{textcolor}{\rmfamily\fontsize{10.000000}{12.000000}\selectfont\catcode`\^=\active\def^{\ifmmode\sp\else\^{}\fi}\catcode`\%=\active\def%{\%}$\mathdefault{\ensuremath{-}0.75}$}}%
\end{pgfscope}%
\begin{pgfscope}%
\pgfpathrectangle{\pgfqpoint{0.599108in}{0.320555in}}{\pgfqpoint{0.734951in}{3.566389in}}%
\pgfusepath{clip}%
\pgfsetbuttcap%
\pgfsetroundjoin%
\pgfsetlinewidth{0.501875pt}%
\definecolor{currentstroke}{rgb}{0.501961,0.501961,0.501961}%
\pgfsetstrokecolor{currentstroke}%
\pgfsetstrokeopacity{0.500000}%
\pgfsetdash{{1.850000pt}{0.800000pt}}{0.000000pt}%
\pgfpathmoveto{\pgfqpoint{0.599108in}{1.293204in}}%
\pgfpathlineto{\pgfqpoint{1.334059in}{1.293204in}}%
\pgfusepath{stroke}%
\end{pgfscope}%
\begin{pgfscope}%
\pgfsetbuttcap%
\pgfsetroundjoin%
\definecolor{currentfill}{rgb}{0.000000,0.000000,0.000000}%
\pgfsetfillcolor{currentfill}%
\pgfsetlinewidth{0.803000pt}%
\definecolor{currentstroke}{rgb}{0.000000,0.000000,0.000000}%
\pgfsetstrokecolor{currentstroke}%
\pgfsetdash{}{0pt}%
\pgfsys@defobject{currentmarker}{\pgfqpoint{-0.048611in}{0.000000in}}{\pgfqpoint{-0.000000in}{0.000000in}}{%
\pgfpathmoveto{\pgfqpoint{-0.000000in}{0.000000in}}%
\pgfpathlineto{\pgfqpoint{-0.048611in}{0.000000in}}%
\pgfusepath{stroke,fill}%
}%
\begin{pgfscope}%
\pgfsys@transformshift{0.599108in}{1.293204in}%
\pgfsys@useobject{currentmarker}{}%
\end{pgfscope}%
\end{pgfscope}%
\begin{pgfscope}%
\definecolor{textcolor}{rgb}{0.000000,0.000000,0.000000}%
\pgfsetstrokecolor{textcolor}%
\pgfsetfillcolor{textcolor}%
\pgftext[x=0.146947in, y=1.245010in, left, base]{\color{textcolor}{\rmfamily\fontsize{10.000000}{12.000000}\selectfont\catcode`\^=\active\def^{\ifmmode\sp\else\^{}\fi}\catcode`\%=\active\def%{\%}$\mathdefault{\ensuremath{-}0.50}$}}%
\end{pgfscope}%
\begin{pgfscope}%
\pgfpathrectangle{\pgfqpoint{0.599108in}{0.320555in}}{\pgfqpoint{0.734951in}{3.566389in}}%
\pgfusepath{clip}%
\pgfsetbuttcap%
\pgfsetroundjoin%
\pgfsetlinewidth{0.501875pt}%
\definecolor{currentstroke}{rgb}{0.501961,0.501961,0.501961}%
\pgfsetstrokecolor{currentstroke}%
\pgfsetstrokeopacity{0.500000}%
\pgfsetdash{{1.850000pt}{0.800000pt}}{0.000000pt}%
\pgfpathmoveto{\pgfqpoint{0.599108in}{1.698477in}}%
\pgfpathlineto{\pgfqpoint{1.334059in}{1.698477in}}%
\pgfusepath{stroke}%
\end{pgfscope}%
\begin{pgfscope}%
\pgfsetbuttcap%
\pgfsetroundjoin%
\definecolor{currentfill}{rgb}{0.000000,0.000000,0.000000}%
\pgfsetfillcolor{currentfill}%
\pgfsetlinewidth{0.803000pt}%
\definecolor{currentstroke}{rgb}{0.000000,0.000000,0.000000}%
\pgfsetstrokecolor{currentstroke}%
\pgfsetdash{}{0pt}%
\pgfsys@defobject{currentmarker}{\pgfqpoint{-0.048611in}{0.000000in}}{\pgfqpoint{-0.000000in}{0.000000in}}{%
\pgfpathmoveto{\pgfqpoint{-0.000000in}{0.000000in}}%
\pgfpathlineto{\pgfqpoint{-0.048611in}{0.000000in}}%
\pgfusepath{stroke,fill}%
}%
\begin{pgfscope}%
\pgfsys@transformshift{0.599108in}{1.698477in}%
\pgfsys@useobject{currentmarker}{}%
\end{pgfscope}%
\end{pgfscope}%
\begin{pgfscope}%
\definecolor{textcolor}{rgb}{0.000000,0.000000,0.000000}%
\pgfsetstrokecolor{textcolor}%
\pgfsetfillcolor{textcolor}%
\pgftext[x=0.146947in, y=1.650283in, left, base]{\color{textcolor}{\rmfamily\fontsize{10.000000}{12.000000}\selectfont\catcode`\^=\active\def^{\ifmmode\sp\else\^{}\fi}\catcode`\%=\active\def%{\%}$\mathdefault{\ensuremath{-}0.25}$}}%
\end{pgfscope}%
\begin{pgfscope}%
\pgfpathrectangle{\pgfqpoint{0.599108in}{0.320555in}}{\pgfqpoint{0.734951in}{3.566389in}}%
\pgfusepath{clip}%
\pgfsetbuttcap%
\pgfsetroundjoin%
\pgfsetlinewidth{0.501875pt}%
\definecolor{currentstroke}{rgb}{0.501961,0.501961,0.501961}%
\pgfsetstrokecolor{currentstroke}%
\pgfsetstrokeopacity{0.500000}%
\pgfsetdash{{1.850000pt}{0.800000pt}}{0.000000pt}%
\pgfpathmoveto{\pgfqpoint{0.599108in}{2.103750in}}%
\pgfpathlineto{\pgfqpoint{1.334059in}{2.103750in}}%
\pgfusepath{stroke}%
\end{pgfscope}%
\begin{pgfscope}%
\pgfsetbuttcap%
\pgfsetroundjoin%
\definecolor{currentfill}{rgb}{0.000000,0.000000,0.000000}%
\pgfsetfillcolor{currentfill}%
\pgfsetlinewidth{0.803000pt}%
\definecolor{currentstroke}{rgb}{0.000000,0.000000,0.000000}%
\pgfsetstrokecolor{currentstroke}%
\pgfsetdash{}{0pt}%
\pgfsys@defobject{currentmarker}{\pgfqpoint{-0.048611in}{0.000000in}}{\pgfqpoint{-0.000000in}{0.000000in}}{%
\pgfpathmoveto{\pgfqpoint{-0.000000in}{0.000000in}}%
\pgfpathlineto{\pgfqpoint{-0.048611in}{0.000000in}}%
\pgfusepath{stroke,fill}%
}%
\begin{pgfscope}%
\pgfsys@transformshift{0.599108in}{2.103750in}%
\pgfsys@useobject{currentmarker}{}%
\end{pgfscope}%
\end{pgfscope}%
\begin{pgfscope}%
\definecolor{textcolor}{rgb}{0.000000,0.000000,0.000000}%
\pgfsetstrokecolor{textcolor}%
\pgfsetfillcolor{textcolor}%
\pgftext[x=0.254972in, y=2.055556in, left, base]{\color{textcolor}{\rmfamily\fontsize{10.000000}{12.000000}\selectfont\catcode`\^=\active\def^{\ifmmode\sp\else\^{}\fi}\catcode`\%=\active\def%{\%}$\mathdefault{0.00}$}}%
\end{pgfscope}%
\begin{pgfscope}%
\pgfpathrectangle{\pgfqpoint{0.599108in}{0.320555in}}{\pgfqpoint{0.734951in}{3.566389in}}%
\pgfusepath{clip}%
\pgfsetbuttcap%
\pgfsetroundjoin%
\pgfsetlinewidth{0.501875pt}%
\definecolor{currentstroke}{rgb}{0.501961,0.501961,0.501961}%
\pgfsetstrokecolor{currentstroke}%
\pgfsetstrokeopacity{0.500000}%
\pgfsetdash{{1.850000pt}{0.800000pt}}{0.000000pt}%
\pgfpathmoveto{\pgfqpoint{0.599108in}{2.509023in}}%
\pgfpathlineto{\pgfqpoint{1.334059in}{2.509023in}}%
\pgfusepath{stroke}%
\end{pgfscope}%
\begin{pgfscope}%
\pgfsetbuttcap%
\pgfsetroundjoin%
\definecolor{currentfill}{rgb}{0.000000,0.000000,0.000000}%
\pgfsetfillcolor{currentfill}%
\pgfsetlinewidth{0.803000pt}%
\definecolor{currentstroke}{rgb}{0.000000,0.000000,0.000000}%
\pgfsetstrokecolor{currentstroke}%
\pgfsetdash{}{0pt}%
\pgfsys@defobject{currentmarker}{\pgfqpoint{-0.048611in}{0.000000in}}{\pgfqpoint{-0.000000in}{0.000000in}}{%
\pgfpathmoveto{\pgfqpoint{-0.000000in}{0.000000in}}%
\pgfpathlineto{\pgfqpoint{-0.048611in}{0.000000in}}%
\pgfusepath{stroke,fill}%
}%
\begin{pgfscope}%
\pgfsys@transformshift{0.599108in}{2.509023in}%
\pgfsys@useobject{currentmarker}{}%
\end{pgfscope}%
\end{pgfscope}%
\begin{pgfscope}%
\definecolor{textcolor}{rgb}{0.000000,0.000000,0.000000}%
\pgfsetstrokecolor{textcolor}%
\pgfsetfillcolor{textcolor}%
\pgftext[x=0.254972in, y=2.460829in, left, base]{\color{textcolor}{\rmfamily\fontsize{10.000000}{12.000000}\selectfont\catcode`\^=\active\def^{\ifmmode\sp\else\^{}\fi}\catcode`\%=\active\def%{\%}$\mathdefault{0.25}$}}%
\end{pgfscope}%
\begin{pgfscope}%
\pgfpathrectangle{\pgfqpoint{0.599108in}{0.320555in}}{\pgfqpoint{0.734951in}{3.566389in}}%
\pgfusepath{clip}%
\pgfsetbuttcap%
\pgfsetroundjoin%
\pgfsetlinewidth{0.501875pt}%
\definecolor{currentstroke}{rgb}{0.501961,0.501961,0.501961}%
\pgfsetstrokecolor{currentstroke}%
\pgfsetstrokeopacity{0.500000}%
\pgfsetdash{{1.850000pt}{0.800000pt}}{0.000000pt}%
\pgfpathmoveto{\pgfqpoint{0.599108in}{2.914296in}}%
\pgfpathlineto{\pgfqpoint{1.334059in}{2.914296in}}%
\pgfusepath{stroke}%
\end{pgfscope}%
\begin{pgfscope}%
\pgfsetbuttcap%
\pgfsetroundjoin%
\definecolor{currentfill}{rgb}{0.000000,0.000000,0.000000}%
\pgfsetfillcolor{currentfill}%
\pgfsetlinewidth{0.803000pt}%
\definecolor{currentstroke}{rgb}{0.000000,0.000000,0.000000}%
\pgfsetstrokecolor{currentstroke}%
\pgfsetdash{}{0pt}%
\pgfsys@defobject{currentmarker}{\pgfqpoint{-0.048611in}{0.000000in}}{\pgfqpoint{-0.000000in}{0.000000in}}{%
\pgfpathmoveto{\pgfqpoint{-0.000000in}{0.000000in}}%
\pgfpathlineto{\pgfqpoint{-0.048611in}{0.000000in}}%
\pgfusepath{stroke,fill}%
}%
\begin{pgfscope}%
\pgfsys@transformshift{0.599108in}{2.914296in}%
\pgfsys@useobject{currentmarker}{}%
\end{pgfscope}%
\end{pgfscope}%
\begin{pgfscope}%
\definecolor{textcolor}{rgb}{0.000000,0.000000,0.000000}%
\pgfsetstrokecolor{textcolor}%
\pgfsetfillcolor{textcolor}%
\pgftext[x=0.254972in, y=2.866102in, left, base]{\color{textcolor}{\rmfamily\fontsize{10.000000}{12.000000}\selectfont\catcode`\^=\active\def^{\ifmmode\sp\else\^{}\fi}\catcode`\%=\active\def%{\%}$\mathdefault{0.50}$}}%
\end{pgfscope}%
\begin{pgfscope}%
\pgfpathrectangle{\pgfqpoint{0.599108in}{0.320555in}}{\pgfqpoint{0.734951in}{3.566389in}}%
\pgfusepath{clip}%
\pgfsetbuttcap%
\pgfsetroundjoin%
\pgfsetlinewidth{0.501875pt}%
\definecolor{currentstroke}{rgb}{0.501961,0.501961,0.501961}%
\pgfsetstrokecolor{currentstroke}%
\pgfsetstrokeopacity{0.500000}%
\pgfsetdash{{1.850000pt}{0.800000pt}}{0.000000pt}%
\pgfpathmoveto{\pgfqpoint{0.599108in}{3.319569in}}%
\pgfpathlineto{\pgfqpoint{1.334059in}{3.319569in}}%
\pgfusepath{stroke}%
\end{pgfscope}%
\begin{pgfscope}%
\pgfsetbuttcap%
\pgfsetroundjoin%
\definecolor{currentfill}{rgb}{0.000000,0.000000,0.000000}%
\pgfsetfillcolor{currentfill}%
\pgfsetlinewidth{0.803000pt}%
\definecolor{currentstroke}{rgb}{0.000000,0.000000,0.000000}%
\pgfsetstrokecolor{currentstroke}%
\pgfsetdash{}{0pt}%
\pgfsys@defobject{currentmarker}{\pgfqpoint{-0.048611in}{0.000000in}}{\pgfqpoint{-0.000000in}{0.000000in}}{%
\pgfpathmoveto{\pgfqpoint{-0.000000in}{0.000000in}}%
\pgfpathlineto{\pgfqpoint{-0.048611in}{0.000000in}}%
\pgfusepath{stroke,fill}%
}%
\begin{pgfscope}%
\pgfsys@transformshift{0.599108in}{3.319569in}%
\pgfsys@useobject{currentmarker}{}%
\end{pgfscope}%
\end{pgfscope}%
\begin{pgfscope}%
\definecolor{textcolor}{rgb}{0.000000,0.000000,0.000000}%
\pgfsetstrokecolor{textcolor}%
\pgfsetfillcolor{textcolor}%
\pgftext[x=0.254972in, y=3.271374in, left, base]{\color{textcolor}{\rmfamily\fontsize{10.000000}{12.000000}\selectfont\catcode`\^=\active\def^{\ifmmode\sp\else\^{}\fi}\catcode`\%=\active\def%{\%}$\mathdefault{0.75}$}}%
\end{pgfscope}%
\begin{pgfscope}%
\pgfpathrectangle{\pgfqpoint{0.599108in}{0.320555in}}{\pgfqpoint{0.734951in}{3.566389in}}%
\pgfusepath{clip}%
\pgfsetbuttcap%
\pgfsetroundjoin%
\pgfsetlinewidth{0.501875pt}%
\definecolor{currentstroke}{rgb}{0.501961,0.501961,0.501961}%
\pgfsetstrokecolor{currentstroke}%
\pgfsetstrokeopacity{0.500000}%
\pgfsetdash{{1.850000pt}{0.800000pt}}{0.000000pt}%
\pgfpathmoveto{\pgfqpoint{0.599108in}{3.724842in}}%
\pgfpathlineto{\pgfqpoint{1.334059in}{3.724842in}}%
\pgfusepath{stroke}%
\end{pgfscope}%
\begin{pgfscope}%
\pgfsetbuttcap%
\pgfsetroundjoin%
\definecolor{currentfill}{rgb}{0.000000,0.000000,0.000000}%
\pgfsetfillcolor{currentfill}%
\pgfsetlinewidth{0.803000pt}%
\definecolor{currentstroke}{rgb}{0.000000,0.000000,0.000000}%
\pgfsetstrokecolor{currentstroke}%
\pgfsetdash{}{0pt}%
\pgfsys@defobject{currentmarker}{\pgfqpoint{-0.048611in}{0.000000in}}{\pgfqpoint{-0.000000in}{0.000000in}}{%
\pgfpathmoveto{\pgfqpoint{-0.000000in}{0.000000in}}%
\pgfpathlineto{\pgfqpoint{-0.048611in}{0.000000in}}%
\pgfusepath{stroke,fill}%
}%
\begin{pgfscope}%
\pgfsys@transformshift{0.599108in}{3.724842in}%
\pgfsys@useobject{currentmarker}{}%
\end{pgfscope}%
\end{pgfscope}%
\begin{pgfscope}%
\definecolor{textcolor}{rgb}{0.000000,0.000000,0.000000}%
\pgfsetstrokecolor{textcolor}%
\pgfsetfillcolor{textcolor}%
\pgftext[x=0.254972in, y=3.676647in, left, base]{\color{textcolor}{\rmfamily\fontsize{10.000000}{12.000000}\selectfont\catcode`\^=\active\def^{\ifmmode\sp\else\^{}\fi}\catcode`\%=\active\def%{\%}$\mathdefault{1.00}$}}%
\end{pgfscope}%
\begin{pgfscope}%
\pgfpathrectangle{\pgfqpoint{0.599108in}{0.320555in}}{\pgfqpoint{0.734951in}{3.566389in}}%
\pgfusepath{clip}%
\pgfsetrectcap%
\pgfsetroundjoin%
\pgfsetlinewidth{1.003750pt}%
\definecolor{currentstroke}{rgb}{0.121569,0.466667,0.705882}%
\pgfsetstrokecolor{currentstroke}%
\pgfsetdash{}{0pt}%
\pgfpathmoveto{\pgfqpoint{0.632515in}{2.103750in}}%
\pgfpathlineto{\pgfqpoint{0.652593in}{2.973666in}}%
\pgfpathlineto{\pgfqpoint{0.662631in}{3.321266in}}%
\pgfpathlineto{\pgfqpoint{0.670439in}{3.525714in}}%
\pgfpathlineto{\pgfqpoint{0.677132in}{3.646548in}}%
\pgfpathlineto{\pgfqpoint{0.681593in}{3.696825in}}%
\pgfpathlineto{\pgfqpoint{0.684940in}{3.718019in}}%
\pgfpathlineto{\pgfqpoint{0.687171in}{3.724167in}}%
\pgfpathlineto{\pgfqpoint{0.688286in}{3.724836in}}%
\pgfpathlineto{\pgfqpoint{0.689401in}{3.723900in}}%
\pgfpathlineto{\pgfqpoint{0.691632in}{3.717217in}}%
\pgfpathlineto{\pgfqpoint{0.694979in}{3.695229in}}%
\pgfpathlineto{\pgfqpoint{0.699440in}{3.643916in}}%
\pgfpathlineto{\pgfqpoint{0.705017in}{3.545636in}}%
\pgfpathlineto{\pgfqpoint{0.711710in}{3.380955in}}%
\pgfpathlineto{\pgfqpoint{0.719518in}{3.131954in}}%
\pgfpathlineto{\pgfqpoint{0.730672in}{2.693551in}}%
\pgfpathlineto{\pgfqpoint{0.754096in}{1.642588in}}%
\pgfpathlineto{\pgfqpoint{0.767481in}{1.102063in}}%
\pgfpathlineto{\pgfqpoint{0.776404in}{0.816144in}}%
\pgfpathlineto{\pgfqpoint{0.784212in}{0.632057in}}%
\pgfpathlineto{\pgfqpoint{0.789789in}{0.543724in}}%
\pgfpathlineto{\pgfqpoint{0.794251in}{0.500735in}}%
\pgfpathlineto{\pgfqpoint{0.797597in}{0.485116in}}%
\pgfpathlineto{\pgfqpoint{0.799828in}{0.482709in}}%
\pgfpathlineto{\pgfqpoint{0.802059in}{0.486720in}}%
\pgfpathlineto{\pgfqpoint{0.804290in}{0.497135in}}%
\pgfpathlineto{\pgfqpoint{0.807636in}{0.524665in}}%
\pgfpathlineto{\pgfqpoint{0.812098in}{0.583187in}}%
\pgfpathlineto{\pgfqpoint{0.817675in}{0.690030in}}%
\pgfpathlineto{\pgfqpoint{0.824367in}{0.864050in}}%
\pgfpathlineto{\pgfqpoint{0.833291in}{1.163320in}}%
\pgfpathlineto{\pgfqpoint{0.845560in}{1.667097in}}%
\pgfpathlineto{\pgfqpoint{0.880138in}{3.151546in}}%
\pgfpathlineto{\pgfqpoint{0.889062in}{3.426639in}}%
\pgfpathlineto{\pgfqpoint{0.895754in}{3.578989in}}%
\pgfpathlineto{\pgfqpoint{0.901331in}{3.666066in}}%
\pgfpathlineto{\pgfqpoint{0.905793in}{3.708010in}}%
\pgfpathlineto{\pgfqpoint{0.909139in}{3.722830in}}%
\pgfpathlineto{\pgfqpoint{0.911370in}{3.724703in}}%
\pgfpathlineto{\pgfqpoint{0.913601in}{3.720156in}}%
\pgfpathlineto{\pgfqpoint{0.916947in}{3.701349in}}%
\pgfpathlineto{\pgfqpoint{0.921409in}{3.654187in}}%
\pgfpathlineto{\pgfqpoint{0.926986in}{3.560861in}}%
\pgfpathlineto{\pgfqpoint{0.933678in}{3.401617in}}%
\pgfpathlineto{\pgfqpoint{0.941486in}{3.158019in}}%
\pgfpathlineto{\pgfqpoint{0.951525in}{2.771899in}}%
\pgfpathlineto{\pgfqpoint{0.968257in}{2.027259in}}%
\pgfpathlineto{\pgfqpoint{0.987219in}{1.212420in}}%
\pgfpathlineto{\pgfqpoint{0.997257in}{0.869545in}}%
\pgfpathlineto{\pgfqpoint{1.005065in}{0.669715in}}%
\pgfpathlineto{\pgfqpoint{1.011758in}{0.553313in}}%
\pgfpathlineto{\pgfqpoint{1.016220in}{0.506151in}}%
\pgfpathlineto{\pgfqpoint{1.019566in}{0.487344in}}%
\pgfpathlineto{\pgfqpoint{1.021797in}{0.482798in}}%
\pgfpathlineto{\pgfqpoint{1.022912in}{0.482932in}}%
\pgfpathlineto{\pgfqpoint{1.025143in}{0.488012in}}%
\pgfpathlineto{\pgfqpoint{1.028489in}{0.507616in}}%
\pgfpathlineto{\pgfqpoint{1.032951in}{0.555817in}}%
\pgfpathlineto{\pgfqpoint{1.038528in}{0.650386in}}%
\pgfpathlineto{\pgfqpoint{1.045221in}{0.810996in}}%
\pgfpathlineto{\pgfqpoint{1.053028in}{1.055954in}}%
\pgfpathlineto{\pgfqpoint{1.063067in}{1.443357in}}%
\pgfpathlineto{\pgfqpoint{1.080914in}{2.239626in}}%
\pgfpathlineto{\pgfqpoint{1.098761in}{3.002170in}}%
\pgfpathlineto{\pgfqpoint{1.108800in}{3.343451in}}%
\pgfpathlineto{\pgfqpoint{1.116607in}{3.541731in}}%
\pgfpathlineto{\pgfqpoint{1.123300in}{3.656649in}}%
\pgfpathlineto{\pgfqpoint{1.127762in}{3.702769in}}%
\pgfpathlineto{\pgfqpoint{1.131108in}{3.720780in}}%
\pgfpathlineto{\pgfqpoint{1.133339in}{3.724792in}}%
\pgfpathlineto{\pgfqpoint{1.134454in}{3.724390in}}%
\pgfpathlineto{\pgfqpoint{1.136685in}{3.718776in}}%
\pgfpathlineto{\pgfqpoint{1.140031in}{3.698377in}}%
\pgfpathlineto{\pgfqpoint{1.144493in}{3.649137in}}%
\pgfpathlineto{\pgfqpoint{1.150070in}{3.553328in}}%
\pgfpathlineto{\pgfqpoint{1.156763in}{3.391357in}}%
\pgfpathlineto{\pgfqpoint{1.164571in}{3.145044in}}%
\pgfpathlineto{\pgfqpoint{1.174609in}{2.756369in}}%
\pgfpathlineto{\pgfqpoint{1.192456in}{1.959404in}}%
\pgfpathlineto{\pgfqpoint{1.209187in}{1.241020in}}%
\pgfpathlineto{\pgfqpoint{1.219226in}{0.891865in}}%
\pgfpathlineto{\pgfqpoint{1.227034in}{0.685889in}}%
\pgfpathlineto{\pgfqpoint{1.233727in}{0.563584in}}%
\pgfpathlineto{\pgfqpoint{1.238188in}{0.512271in}}%
\pgfpathlineto{\pgfqpoint{1.241535in}{0.490283in}}%
\pgfpathlineto{\pgfqpoint{1.243765in}{0.483600in}}%
\pgfpathlineto{\pgfqpoint{1.244881in}{0.482664in}}%
\pgfpathlineto{\pgfqpoint{1.245996in}{0.483333in}}%
\pgfpathlineto{\pgfqpoint{1.248227in}{0.489482in}}%
\pgfpathlineto{\pgfqpoint{1.251573in}{0.510676in}}%
\pgfpathlineto{\pgfqpoint{1.256035in}{0.560953in}}%
\pgfpathlineto{\pgfqpoint{1.261612in}{0.657998in}}%
\pgfpathlineto{\pgfqpoint{1.268305in}{0.821327in}}%
\pgfpathlineto{\pgfqpoint{1.276113in}{1.068987in}}%
\pgfpathlineto{\pgfqpoint{1.287267in}{1.506038in}}%
\pgfpathlineto{\pgfqpoint{1.300652in}{2.103750in}}%
\pgfpathlineto{\pgfqpoint{1.300652in}{2.103750in}}%
\pgfusepath{stroke}%
\end{pgfscope}%
\begin{pgfscope}%
\pgfsetrectcap%
\pgfsetmiterjoin%
\pgfsetlinewidth{0.803000pt}%
\definecolor{currentstroke}{rgb}{0.000000,0.000000,0.000000}%
\pgfsetstrokecolor{currentstroke}%
\pgfsetdash{}{0pt}%
\pgfpathmoveto{\pgfqpoint{0.599108in}{0.320555in}}%
\pgfpathlineto{\pgfqpoint{0.599108in}{3.886945in}}%
\pgfusepath{stroke}%
\end{pgfscope}%
\begin{pgfscope}%
\pgfsetrectcap%
\pgfsetmiterjoin%
\pgfsetlinewidth{0.803000pt}%
\definecolor{currentstroke}{rgb}{0.000000,0.000000,0.000000}%
\pgfsetstrokecolor{currentstroke}%
\pgfsetdash{}{0pt}%
\pgfpathmoveto{\pgfqpoint{1.334059in}{0.320555in}}%
\pgfpathlineto{\pgfqpoint{1.334059in}{3.886945in}}%
\pgfusepath{stroke}%
\end{pgfscope}%
\begin{pgfscope}%
\pgfsetrectcap%
\pgfsetmiterjoin%
\pgfsetlinewidth{0.803000pt}%
\definecolor{currentstroke}{rgb}{0.000000,0.000000,0.000000}%
\pgfsetstrokecolor{currentstroke}%
\pgfsetdash{}{0pt}%
\pgfpathmoveto{\pgfqpoint{0.599108in}{0.320555in}}%
\pgfpathlineto{\pgfqpoint{1.334059in}{0.320555in}}%
\pgfusepath{stroke}%
\end{pgfscope}%
\begin{pgfscope}%
\pgfsetrectcap%
\pgfsetmiterjoin%
\pgfsetlinewidth{0.803000pt}%
\definecolor{currentstroke}{rgb}{0.000000,0.000000,0.000000}%
\pgfsetstrokecolor{currentstroke}%
\pgfsetdash{}{0pt}%
\pgfpathmoveto{\pgfqpoint{0.599108in}{3.886945in}}%
\pgfpathlineto{\pgfqpoint{1.334059in}{3.886945in}}%
\pgfusepath{stroke}%
\end{pgfscope}%
\begin{pgfscope}%
\definecolor{textcolor}{rgb}{0.000000,0.000000,0.000000}%
\pgfsetstrokecolor{textcolor}%
\pgfsetfillcolor{textcolor}%
\pgftext[x=0.966583in,y=3.970278in,,base]{\color{textcolor}{\rmfamily\fontsize{12.000000}{14.400000}\selectfont\catcode`\^=\active\def^{\ifmmode\sp\else\^{}\fi}\catcode`\%=\active\def%{\%}Panel izquierdo (2 filas)}}%
\end{pgfscope}%
\begin{pgfscope}%
\pgfsetbuttcap%
\pgfsetmiterjoin%
\definecolor{currentfill}{rgb}{1.000000,1.000000,1.000000}%
\pgfsetfillcolor{currentfill}%
\pgfsetlinewidth{0.000000pt}%
\definecolor{currentstroke}{rgb}{0.000000,0.000000,0.000000}%
\pgfsetstrokecolor{currentstroke}%
\pgfsetstrokeopacity{0.000000}%
\pgfsetdash{}{0pt}%
\pgfpathmoveto{\pgfqpoint{1.993487in}{2.452069in}}%
\pgfpathlineto{\pgfqpoint{3.463388in}{2.452069in}}%
\pgfpathlineto{\pgfqpoint{3.463388in}{3.886945in}}%
\pgfpathlineto{\pgfqpoint{1.993487in}{3.886945in}}%
\pgfpathlineto{\pgfqpoint{1.993487in}{2.452069in}}%
\pgfpathclose%
\pgfusepath{fill}%
\end{pgfscope}%
\begin{pgfscope}%
\pgfpathrectangle{\pgfqpoint{1.993487in}{2.452069in}}{\pgfqpoint{1.469901in}{1.434876in}}%
\pgfusepath{clip}%
\pgfsetbuttcap%
\pgfsetroundjoin%
\pgfsetlinewidth{0.501875pt}%
\definecolor{currentstroke}{rgb}{0.501961,0.501961,0.501961}%
\pgfsetstrokecolor{currentstroke}%
\pgfsetstrokeopacity{0.500000}%
\pgfsetdash{{1.850000pt}{0.800000pt}}{0.000000pt}%
\pgfpathmoveto{\pgfqpoint{2.060300in}{2.452069in}}%
\pgfpathlineto{\pgfqpoint{2.060300in}{3.886945in}}%
\pgfusepath{stroke}%
\end{pgfscope}%
\begin{pgfscope}%
\pgfsetbuttcap%
\pgfsetroundjoin%
\definecolor{currentfill}{rgb}{0.000000,0.000000,0.000000}%
\pgfsetfillcolor{currentfill}%
\pgfsetlinewidth{0.803000pt}%
\definecolor{currentstroke}{rgb}{0.000000,0.000000,0.000000}%
\pgfsetstrokecolor{currentstroke}%
\pgfsetdash{}{0pt}%
\pgfsys@defobject{currentmarker}{\pgfqpoint{0.000000in}{-0.048611in}}{\pgfqpoint{0.000000in}{0.000000in}}{%
\pgfpathmoveto{\pgfqpoint{0.000000in}{0.000000in}}%
\pgfpathlineto{\pgfqpoint{0.000000in}{-0.048611in}}%
\pgfusepath{stroke,fill}%
}%
\begin{pgfscope}%
\pgfsys@transformshift{2.060300in}{2.452069in}%
\pgfsys@useobject{currentmarker}{}%
\end{pgfscope}%
\end{pgfscope}%
\begin{pgfscope}%
\definecolor{textcolor}{rgb}{0.000000,0.000000,0.000000}%
\pgfsetstrokecolor{textcolor}%
\pgfsetfillcolor{textcolor}%
\pgftext[x=2.060300in,y=2.354846in,,top]{\color{textcolor}{\rmfamily\fontsize{10.000000}{12.000000}\selectfont\catcode`\^=\active\def^{\ifmmode\sp\else\^{}\fi}\catcode`\%=\active\def%{\%}$\mathdefault{0}$}}%
\end{pgfscope}%
\begin{pgfscope}%
\pgfpathrectangle{\pgfqpoint{1.993487in}{2.452069in}}{\pgfqpoint{1.469901in}{1.434876in}}%
\pgfusepath{clip}%
\pgfsetbuttcap%
\pgfsetroundjoin%
\pgfsetlinewidth{0.501875pt}%
\definecolor{currentstroke}{rgb}{0.501961,0.501961,0.501961}%
\pgfsetstrokecolor{currentstroke}%
\pgfsetstrokeopacity{0.500000}%
\pgfsetdash{{1.850000pt}{0.800000pt}}{0.000000pt}%
\pgfpathmoveto{\pgfqpoint{2.769216in}{2.452069in}}%
\pgfpathlineto{\pgfqpoint{2.769216in}{3.886945in}}%
\pgfusepath{stroke}%
\end{pgfscope}%
\begin{pgfscope}%
\pgfsetbuttcap%
\pgfsetroundjoin%
\definecolor{currentfill}{rgb}{0.000000,0.000000,0.000000}%
\pgfsetfillcolor{currentfill}%
\pgfsetlinewidth{0.803000pt}%
\definecolor{currentstroke}{rgb}{0.000000,0.000000,0.000000}%
\pgfsetstrokecolor{currentstroke}%
\pgfsetdash{}{0pt}%
\pgfsys@defobject{currentmarker}{\pgfqpoint{0.000000in}{-0.048611in}}{\pgfqpoint{0.000000in}{0.000000in}}{%
\pgfpathmoveto{\pgfqpoint{0.000000in}{0.000000in}}%
\pgfpathlineto{\pgfqpoint{0.000000in}{-0.048611in}}%
\pgfusepath{stroke,fill}%
}%
\begin{pgfscope}%
\pgfsys@transformshift{2.769216in}{2.452069in}%
\pgfsys@useobject{currentmarker}{}%
\end{pgfscope}%
\end{pgfscope}%
\begin{pgfscope}%
\definecolor{textcolor}{rgb}{0.000000,0.000000,0.000000}%
\pgfsetstrokecolor{textcolor}%
\pgfsetfillcolor{textcolor}%
\pgftext[x=2.769216in,y=2.354846in,,top]{\color{textcolor}{\rmfamily\fontsize{10.000000}{12.000000}\selectfont\catcode`\^=\active\def^{\ifmmode\sp\else\^{}\fi}\catcode`\%=\active\def%{\%}$\mathdefault{10}$}}%
\end{pgfscope}%
\begin{pgfscope}%
\pgfpathrectangle{\pgfqpoint{1.993487in}{2.452069in}}{\pgfqpoint{1.469901in}{1.434876in}}%
\pgfusepath{clip}%
\pgfsetbuttcap%
\pgfsetroundjoin%
\pgfsetlinewidth{0.501875pt}%
\definecolor{currentstroke}{rgb}{0.501961,0.501961,0.501961}%
\pgfsetstrokecolor{currentstroke}%
\pgfsetstrokeopacity{0.500000}%
\pgfsetdash{{1.850000pt}{0.800000pt}}{0.000000pt}%
\pgfpathmoveto{\pgfqpoint{1.993487in}{2.517281in}}%
\pgfpathlineto{\pgfqpoint{3.463388in}{2.517281in}}%
\pgfusepath{stroke}%
\end{pgfscope}%
\begin{pgfscope}%
\pgfsetbuttcap%
\pgfsetroundjoin%
\definecolor{currentfill}{rgb}{0.000000,0.000000,0.000000}%
\pgfsetfillcolor{currentfill}%
\pgfsetlinewidth{0.803000pt}%
\definecolor{currentstroke}{rgb}{0.000000,0.000000,0.000000}%
\pgfsetstrokecolor{currentstroke}%
\pgfsetdash{}{0pt}%
\pgfsys@defobject{currentmarker}{\pgfqpoint{-0.048611in}{0.000000in}}{\pgfqpoint{-0.000000in}{0.000000in}}{%
\pgfpathmoveto{\pgfqpoint{-0.000000in}{0.000000in}}%
\pgfpathlineto{\pgfqpoint{-0.048611in}{0.000000in}}%
\pgfusepath{stroke,fill}%
}%
\begin{pgfscope}%
\pgfsys@transformshift{1.993487in}{2.517281in}%
\pgfsys@useobject{currentmarker}{}%
\end{pgfscope}%
\end{pgfscope}%
\begin{pgfscope}%
\definecolor{textcolor}{rgb}{0.000000,0.000000,0.000000}%
\pgfsetstrokecolor{textcolor}%
\pgfsetfillcolor{textcolor}%
\pgftext[x=1.610770in, y=2.469087in, left, base]{\color{textcolor}{\rmfamily\fontsize{10.000000}{12.000000}\selectfont\catcode`\^=\active\def^{\ifmmode\sp\else\^{}\fi}\catcode`\%=\active\def%{\%}$\mathdefault{\ensuremath{-}1.0}$}}%
\end{pgfscope}%
\begin{pgfscope}%
\pgfpathrectangle{\pgfqpoint{1.993487in}{2.452069in}}{\pgfqpoint{1.469901in}{1.434876in}}%
\pgfusepath{clip}%
\pgfsetbuttcap%
\pgfsetroundjoin%
\pgfsetlinewidth{0.501875pt}%
\definecolor{currentstroke}{rgb}{0.501961,0.501961,0.501961}%
\pgfsetstrokecolor{currentstroke}%
\pgfsetstrokeopacity{0.500000}%
\pgfsetdash{{1.850000pt}{0.800000pt}}{0.000000pt}%
\pgfpathmoveto{\pgfqpoint{1.993487in}{2.843392in}}%
\pgfpathlineto{\pgfqpoint{3.463388in}{2.843392in}}%
\pgfusepath{stroke}%
\end{pgfscope}%
\begin{pgfscope}%
\pgfsetbuttcap%
\pgfsetroundjoin%
\definecolor{currentfill}{rgb}{0.000000,0.000000,0.000000}%
\pgfsetfillcolor{currentfill}%
\pgfsetlinewidth{0.803000pt}%
\definecolor{currentstroke}{rgb}{0.000000,0.000000,0.000000}%
\pgfsetstrokecolor{currentstroke}%
\pgfsetdash{}{0pt}%
\pgfsys@defobject{currentmarker}{\pgfqpoint{-0.048611in}{0.000000in}}{\pgfqpoint{-0.000000in}{0.000000in}}{%
\pgfpathmoveto{\pgfqpoint{-0.000000in}{0.000000in}}%
\pgfpathlineto{\pgfqpoint{-0.048611in}{0.000000in}}%
\pgfusepath{stroke,fill}%
}%
\begin{pgfscope}%
\pgfsys@transformshift{1.993487in}{2.843392in}%
\pgfsys@useobject{currentmarker}{}%
\end{pgfscope}%
\end{pgfscope}%
\begin{pgfscope}%
\definecolor{textcolor}{rgb}{0.000000,0.000000,0.000000}%
\pgfsetstrokecolor{textcolor}%
\pgfsetfillcolor{textcolor}%
\pgftext[x=1.610770in, y=2.795197in, left, base]{\color{textcolor}{\rmfamily\fontsize{10.000000}{12.000000}\selectfont\catcode`\^=\active\def^{\ifmmode\sp\else\^{}\fi}\catcode`\%=\active\def%{\%}$\mathdefault{\ensuremath{-}0.5}$}}%
\end{pgfscope}%
\begin{pgfscope}%
\pgfpathrectangle{\pgfqpoint{1.993487in}{2.452069in}}{\pgfqpoint{1.469901in}{1.434876in}}%
\pgfusepath{clip}%
\pgfsetbuttcap%
\pgfsetroundjoin%
\pgfsetlinewidth{0.501875pt}%
\definecolor{currentstroke}{rgb}{0.501961,0.501961,0.501961}%
\pgfsetstrokecolor{currentstroke}%
\pgfsetstrokeopacity{0.500000}%
\pgfsetdash{{1.850000pt}{0.800000pt}}{0.000000pt}%
\pgfpathmoveto{\pgfqpoint{1.993487in}{3.169502in}}%
\pgfpathlineto{\pgfqpoint{3.463388in}{3.169502in}}%
\pgfusepath{stroke}%
\end{pgfscope}%
\begin{pgfscope}%
\pgfsetbuttcap%
\pgfsetroundjoin%
\definecolor{currentfill}{rgb}{0.000000,0.000000,0.000000}%
\pgfsetfillcolor{currentfill}%
\pgfsetlinewidth{0.803000pt}%
\definecolor{currentstroke}{rgb}{0.000000,0.000000,0.000000}%
\pgfsetstrokecolor{currentstroke}%
\pgfsetdash{}{0pt}%
\pgfsys@defobject{currentmarker}{\pgfqpoint{-0.048611in}{0.000000in}}{\pgfqpoint{-0.000000in}{0.000000in}}{%
\pgfpathmoveto{\pgfqpoint{-0.000000in}{0.000000in}}%
\pgfpathlineto{\pgfqpoint{-0.048611in}{0.000000in}}%
\pgfusepath{stroke,fill}%
}%
\begin{pgfscope}%
\pgfsys@transformshift{1.993487in}{3.169502in}%
\pgfsys@useobject{currentmarker}{}%
\end{pgfscope}%
\end{pgfscope}%
\begin{pgfscope}%
\definecolor{textcolor}{rgb}{0.000000,0.000000,0.000000}%
\pgfsetstrokecolor{textcolor}%
\pgfsetfillcolor{textcolor}%
\pgftext[x=1.718795in, y=3.121308in, left, base]{\color{textcolor}{\rmfamily\fontsize{10.000000}{12.000000}\selectfont\catcode`\^=\active\def^{\ifmmode\sp\else\^{}\fi}\catcode`\%=\active\def%{\%}$\mathdefault{0.0}$}}%
\end{pgfscope}%
\begin{pgfscope}%
\pgfpathrectangle{\pgfqpoint{1.993487in}{2.452069in}}{\pgfqpoint{1.469901in}{1.434876in}}%
\pgfusepath{clip}%
\pgfsetbuttcap%
\pgfsetroundjoin%
\pgfsetlinewidth{0.501875pt}%
\definecolor{currentstroke}{rgb}{0.501961,0.501961,0.501961}%
\pgfsetstrokecolor{currentstroke}%
\pgfsetstrokeopacity{0.500000}%
\pgfsetdash{{1.850000pt}{0.800000pt}}{0.000000pt}%
\pgfpathmoveto{\pgfqpoint{1.993487in}{3.495613in}}%
\pgfpathlineto{\pgfqpoint{3.463388in}{3.495613in}}%
\pgfusepath{stroke}%
\end{pgfscope}%
\begin{pgfscope}%
\pgfsetbuttcap%
\pgfsetroundjoin%
\definecolor{currentfill}{rgb}{0.000000,0.000000,0.000000}%
\pgfsetfillcolor{currentfill}%
\pgfsetlinewidth{0.803000pt}%
\definecolor{currentstroke}{rgb}{0.000000,0.000000,0.000000}%
\pgfsetstrokecolor{currentstroke}%
\pgfsetdash{}{0pt}%
\pgfsys@defobject{currentmarker}{\pgfqpoint{-0.048611in}{0.000000in}}{\pgfqpoint{-0.000000in}{0.000000in}}{%
\pgfpathmoveto{\pgfqpoint{-0.000000in}{0.000000in}}%
\pgfpathlineto{\pgfqpoint{-0.048611in}{0.000000in}}%
\pgfusepath{stroke,fill}%
}%
\begin{pgfscope}%
\pgfsys@transformshift{1.993487in}{3.495613in}%
\pgfsys@useobject{currentmarker}{}%
\end{pgfscope}%
\end{pgfscope}%
\begin{pgfscope}%
\definecolor{textcolor}{rgb}{0.000000,0.000000,0.000000}%
\pgfsetstrokecolor{textcolor}%
\pgfsetfillcolor{textcolor}%
\pgftext[x=1.718795in, y=3.447418in, left, base]{\color{textcolor}{\rmfamily\fontsize{10.000000}{12.000000}\selectfont\catcode`\^=\active\def^{\ifmmode\sp\else\^{}\fi}\catcode`\%=\active\def%{\%}$\mathdefault{0.5}$}}%
\end{pgfscope}%
\begin{pgfscope}%
\pgfpathrectangle{\pgfqpoint{1.993487in}{2.452069in}}{\pgfqpoint{1.469901in}{1.434876in}}%
\pgfusepath{clip}%
\pgfsetbuttcap%
\pgfsetroundjoin%
\pgfsetlinewidth{0.501875pt}%
\definecolor{currentstroke}{rgb}{0.501961,0.501961,0.501961}%
\pgfsetstrokecolor{currentstroke}%
\pgfsetstrokeopacity{0.500000}%
\pgfsetdash{{1.850000pt}{0.800000pt}}{0.000000pt}%
\pgfpathmoveto{\pgfqpoint{1.993487in}{3.821723in}}%
\pgfpathlineto{\pgfqpoint{3.463388in}{3.821723in}}%
\pgfusepath{stroke}%
\end{pgfscope}%
\begin{pgfscope}%
\pgfsetbuttcap%
\pgfsetroundjoin%
\definecolor{currentfill}{rgb}{0.000000,0.000000,0.000000}%
\pgfsetfillcolor{currentfill}%
\pgfsetlinewidth{0.803000pt}%
\definecolor{currentstroke}{rgb}{0.000000,0.000000,0.000000}%
\pgfsetstrokecolor{currentstroke}%
\pgfsetdash{}{0pt}%
\pgfsys@defobject{currentmarker}{\pgfqpoint{-0.048611in}{0.000000in}}{\pgfqpoint{-0.000000in}{0.000000in}}{%
\pgfpathmoveto{\pgfqpoint{-0.000000in}{0.000000in}}%
\pgfpathlineto{\pgfqpoint{-0.048611in}{0.000000in}}%
\pgfusepath{stroke,fill}%
}%
\begin{pgfscope}%
\pgfsys@transformshift{1.993487in}{3.821723in}%
\pgfsys@useobject{currentmarker}{}%
\end{pgfscope}%
\end{pgfscope}%
\begin{pgfscope}%
\definecolor{textcolor}{rgb}{0.000000,0.000000,0.000000}%
\pgfsetstrokecolor{textcolor}%
\pgfsetfillcolor{textcolor}%
\pgftext[x=1.718795in, y=3.773529in, left, base]{\color{textcolor}{\rmfamily\fontsize{10.000000}{12.000000}\selectfont\catcode`\^=\active\def^{\ifmmode\sp\else\^{}\fi}\catcode`\%=\active\def%{\%}$\mathdefault{1.0}$}}%
\end{pgfscope}%
\begin{pgfscope}%
\pgfpathrectangle{\pgfqpoint{1.993487in}{2.452069in}}{\pgfqpoint{1.469901in}{1.434876in}}%
\pgfusepath{clip}%
\pgfsetrectcap%
\pgfsetroundjoin%
\pgfsetlinewidth{1.505625pt}%
\definecolor{currentstroke}{rgb}{1.000000,0.647059,0.000000}%
\pgfsetstrokecolor{currentstroke}%
\pgfsetdash{}{0pt}%
\pgfpathmoveto{\pgfqpoint{2.060300in}{3.821723in}}%
\pgfpathlineto{\pgfqpoint{2.064762in}{3.820432in}}%
\pgfpathlineto{\pgfqpoint{2.069224in}{3.816563in}}%
\pgfpathlineto{\pgfqpoint{2.073686in}{3.810132in}}%
\pgfpathlineto{\pgfqpoint{2.080378in}{3.795740in}}%
\pgfpathlineto{\pgfqpoint{2.087071in}{3.775771in}}%
\pgfpathlineto{\pgfqpoint{2.095994in}{3.740784in}}%
\pgfpathlineto{\pgfqpoint{2.104917in}{3.696758in}}%
\pgfpathlineto{\pgfqpoint{2.116071in}{3.630087in}}%
\pgfpathlineto{\pgfqpoint{2.129457in}{3.535227in}}%
\pgfpathlineto{\pgfqpoint{2.145072in}{3.408392in}}%
\pgfpathlineto{\pgfqpoint{2.169612in}{3.188314in}}%
\pgfpathlineto{\pgfqpoint{2.203074in}{2.889821in}}%
\pgfpathlineto{\pgfqpoint{2.218690in}{2.767835in}}%
\pgfpathlineto{\pgfqpoint{2.232075in}{2.678526in}}%
\pgfpathlineto{\pgfqpoint{2.243229in}{2.617316in}}%
\pgfpathlineto{\pgfqpoint{2.252153in}{2.578109in}}%
\pgfpathlineto{\pgfqpoint{2.261076in}{2.548259in}}%
\pgfpathlineto{\pgfqpoint{2.267769in}{2.532302in}}%
\pgfpathlineto{\pgfqpoint{2.274461in}{2.522021in}}%
\pgfpathlineto{\pgfqpoint{2.278923in}{2.518366in}}%
\pgfpathlineto{\pgfqpoint{2.283385in}{2.517290in}}%
\pgfpathlineto{\pgfqpoint{2.287846in}{2.518797in}}%
\pgfpathlineto{\pgfqpoint{2.292308in}{2.522880in}}%
\pgfpathlineto{\pgfqpoint{2.296770in}{2.529523in}}%
\pgfpathlineto{\pgfqpoint{2.303462in}{2.544229in}}%
\pgfpathlineto{\pgfqpoint{2.310155in}{2.564503in}}%
\pgfpathlineto{\pgfqpoint{2.319078in}{2.599879in}}%
\pgfpathlineto{\pgfqpoint{2.328001in}{2.644268in}}%
\pgfpathlineto{\pgfqpoint{2.339156in}{2.711346in}}%
\pgfpathlineto{\pgfqpoint{2.352541in}{2.806615in}}%
\pgfpathlineto{\pgfqpoint{2.370387in}{2.953050in}}%
\pgfpathlineto{\pgfqpoint{2.399388in}{3.215643in}}%
\pgfpathlineto{\pgfqpoint{2.426158in}{3.452270in}}%
\pgfpathlineto{\pgfqpoint{2.441774in}{3.573859in}}%
\pgfpathlineto{\pgfqpoint{2.455159in}{3.662723in}}%
\pgfpathlineto{\pgfqpoint{2.466314in}{3.723501in}}%
\pgfpathlineto{\pgfqpoint{2.475237in}{3.762330in}}%
\pgfpathlineto{\pgfqpoint{2.484160in}{3.791779in}}%
\pgfpathlineto{\pgfqpoint{2.490853in}{3.807423in}}%
\pgfpathlineto{\pgfqpoint{2.497545in}{3.817386in}}%
\pgfpathlineto{\pgfqpoint{2.502007in}{3.820826in}}%
\pgfpathlineto{\pgfqpoint{2.506469in}{3.821687in}}%
\pgfpathlineto{\pgfqpoint{2.510930in}{3.819966in}}%
\pgfpathlineto{\pgfqpoint{2.515392in}{3.815669in}}%
\pgfpathlineto{\pgfqpoint{2.519854in}{3.808813in}}%
\pgfpathlineto{\pgfqpoint{2.526546in}{3.793794in}}%
\pgfpathlineto{\pgfqpoint{2.533239in}{3.773215in}}%
\pgfpathlineto{\pgfqpoint{2.542162in}{3.737452in}}%
\pgfpathlineto{\pgfqpoint{2.551085in}{3.692702in}}%
\pgfpathlineto{\pgfqpoint{2.562240in}{3.625218in}}%
\pgfpathlineto{\pgfqpoint{2.575625in}{3.529543in}}%
\pgfpathlineto{\pgfqpoint{2.593471in}{3.382725in}}%
\pgfpathlineto{\pgfqpoint{2.622472in}{3.119950in}}%
\pgfpathlineto{\pgfqpoint{2.649242in}{2.883656in}}%
\pgfpathlineto{\pgfqpoint{2.664858in}{2.762467in}}%
\pgfpathlineto{\pgfqpoint{2.678243in}{2.674050in}}%
\pgfpathlineto{\pgfqpoint{2.689398in}{2.613705in}}%
\pgfpathlineto{\pgfqpoint{2.698321in}{2.575256in}}%
\pgfpathlineto{\pgfqpoint{2.707244in}{2.546210in}}%
\pgfpathlineto{\pgfqpoint{2.713937in}{2.530878in}}%
\pgfpathlineto{\pgfqpoint{2.720629in}{2.521233in}}%
\pgfpathlineto{\pgfqpoint{2.725091in}{2.518008in}}%
\pgfpathlineto{\pgfqpoint{2.729553in}{2.517362in}}%
\pgfpathlineto{\pgfqpoint{2.734014in}{2.519299in}}%
\pgfpathlineto{\pgfqpoint{2.738476in}{2.523810in}}%
\pgfpathlineto{\pgfqpoint{2.745169in}{2.535362in}}%
\pgfpathlineto{\pgfqpoint{2.751861in}{2.552562in}}%
\pgfpathlineto{\pgfqpoint{2.760785in}{2.584008in}}%
\pgfpathlineto{\pgfqpoint{2.769708in}{2.624719in}}%
\pgfpathlineto{\pgfqpoint{2.780862in}{2.687641in}}%
\pgfpathlineto{\pgfqpoint{2.794247in}{2.778703in}}%
\pgfpathlineto{\pgfqpoint{2.809863in}{2.902242in}}%
\pgfpathlineto{\pgfqpoint{2.832171in}{3.099512in}}%
\pgfpathlineto{\pgfqpoint{2.872327in}{3.458420in}}%
\pgfpathlineto{\pgfqpoint{2.887942in}{3.579205in}}%
\pgfpathlineto{\pgfqpoint{2.901327in}{3.667173in}}%
\pgfpathlineto{\pgfqpoint{2.912482in}{3.727081in}}%
\pgfpathlineto{\pgfqpoint{2.921405in}{3.765150in}}%
\pgfpathlineto{\pgfqpoint{2.930328in}{3.793794in}}%
\pgfpathlineto{\pgfqpoint{2.937021in}{3.808813in}}%
\pgfpathlineto{\pgfqpoint{2.943713in}{3.818138in}}%
\pgfpathlineto{\pgfqpoint{2.948175in}{3.821149in}}%
\pgfpathlineto{\pgfqpoint{2.952637in}{3.821580in}}%
\pgfpathlineto{\pgfqpoint{2.957099in}{3.819428in}}%
\pgfpathlineto{\pgfqpoint{2.961560in}{3.814703in}}%
\pgfpathlineto{\pgfqpoint{2.968253in}{3.802834in}}%
\pgfpathlineto{\pgfqpoint{2.974945in}{3.785324in}}%
\pgfpathlineto{\pgfqpoint{2.983869in}{3.753481in}}%
\pgfpathlineto{\pgfqpoint{2.992792in}{3.712397in}}%
\pgfpathlineto{\pgfqpoint{3.003946in}{3.649052in}}%
\pgfpathlineto{\pgfqpoint{3.017331in}{3.557557in}}%
\pgfpathlineto{\pgfqpoint{3.032947in}{3.433638in}}%
\pgfpathlineto{\pgfqpoint{3.055256in}{3.236090in}}%
\pgfpathlineto{\pgfqpoint{3.095411in}{2.877522in}}%
\pgfpathlineto{\pgfqpoint{3.111027in}{2.757144in}}%
\pgfpathlineto{\pgfqpoint{3.124412in}{2.669628in}}%
\pgfpathlineto{\pgfqpoint{3.135566in}{2.610156in}}%
\pgfpathlineto{\pgfqpoint{3.144489in}{2.572469in}}%
\pgfpathlineto{\pgfqpoint{3.153413in}{2.544229in}}%
\pgfpathlineto{\pgfqpoint{3.160105in}{2.529523in}}%
\pgfpathlineto{\pgfqpoint{3.166798in}{2.520517in}}%
\pgfpathlineto{\pgfqpoint{3.171259in}{2.517721in}}%
\pgfpathlineto{\pgfqpoint{3.175721in}{2.517506in}}%
\pgfpathlineto{\pgfqpoint{3.180183in}{2.519872in}}%
\pgfpathlineto{\pgfqpoint{3.184644in}{2.524811in}}%
\pgfpathlineto{\pgfqpoint{3.191337in}{2.536997in}}%
\pgfpathlineto{\pgfqpoint{3.198029in}{2.554816in}}%
\pgfpathlineto{\pgfqpoint{3.206953in}{2.587055in}}%
\pgfpathlineto{\pgfqpoint{3.215876in}{2.628510in}}%
\pgfpathlineto{\pgfqpoint{3.227030in}{2.692278in}}%
\pgfpathlineto{\pgfqpoint{3.240415in}{2.784202in}}%
\pgfpathlineto{\pgfqpoint{3.256031in}{2.908498in}}%
\pgfpathlineto{\pgfqpoint{3.278340in}{3.106318in}}%
\pgfpathlineto{\pgfqpoint{3.318495in}{3.464537in}}%
\pgfpathlineto{\pgfqpoint{3.334111in}{3.584505in}}%
\pgfpathlineto{\pgfqpoint{3.347496in}{3.671567in}}%
\pgfpathlineto{\pgfqpoint{3.358650in}{3.730600in}}%
\pgfpathlineto{\pgfqpoint{3.367573in}{3.767904in}}%
\pgfpathlineto{\pgfqpoint{3.376497in}{3.795740in}}%
\pgfpathlineto{\pgfqpoint{3.383189in}{3.810132in}}%
\pgfpathlineto{\pgfqpoint{3.389882in}{3.818819in}}%
\pgfpathlineto{\pgfqpoint{3.394343in}{3.821400in}}%
\pgfpathlineto{\pgfqpoint{3.396574in}{3.821723in}}%
\pgfpathlineto{\pgfqpoint{3.396574in}{3.821723in}}%
\pgfusepath{stroke}%
\end{pgfscope}%
\begin{pgfscope}%
\pgfsetrectcap%
\pgfsetmiterjoin%
\pgfsetlinewidth{0.803000pt}%
\definecolor{currentstroke}{rgb}{0.000000,0.000000,0.000000}%
\pgfsetstrokecolor{currentstroke}%
\pgfsetdash{}{0pt}%
\pgfpathmoveto{\pgfqpoint{1.993487in}{2.452069in}}%
\pgfpathlineto{\pgfqpoint{1.993487in}{3.886945in}}%
\pgfusepath{stroke}%
\end{pgfscope}%
\begin{pgfscope}%
\pgfsetrectcap%
\pgfsetmiterjoin%
\pgfsetlinewidth{0.803000pt}%
\definecolor{currentstroke}{rgb}{0.000000,0.000000,0.000000}%
\pgfsetstrokecolor{currentstroke}%
\pgfsetdash{}{0pt}%
\pgfpathmoveto{\pgfqpoint{3.463388in}{2.452069in}}%
\pgfpathlineto{\pgfqpoint{3.463388in}{3.886945in}}%
\pgfusepath{stroke}%
\end{pgfscope}%
\begin{pgfscope}%
\pgfsetrectcap%
\pgfsetmiterjoin%
\pgfsetlinewidth{0.803000pt}%
\definecolor{currentstroke}{rgb}{0.000000,0.000000,0.000000}%
\pgfsetstrokecolor{currentstroke}%
\pgfsetdash{}{0pt}%
\pgfpathmoveto{\pgfqpoint{1.993487in}{2.452069in}}%
\pgfpathlineto{\pgfqpoint{3.463388in}{2.452069in}}%
\pgfusepath{stroke}%
\end{pgfscope}%
\begin{pgfscope}%
\pgfsetrectcap%
\pgfsetmiterjoin%
\pgfsetlinewidth{0.803000pt}%
\definecolor{currentstroke}{rgb}{0.000000,0.000000,0.000000}%
\pgfsetstrokecolor{currentstroke}%
\pgfsetdash{}{0pt}%
\pgfpathmoveto{\pgfqpoint{1.993487in}{3.886945in}}%
\pgfpathlineto{\pgfqpoint{3.463388in}{3.886945in}}%
\pgfusepath{stroke}%
\end{pgfscope}%
\begin{pgfscope}%
\definecolor{textcolor}{rgb}{0.000000,0.000000,0.000000}%
\pgfsetstrokecolor{textcolor}%
\pgfsetfillcolor{textcolor}%
\pgftext[x=2.728437in,y=3.970278in,,base]{\color{textcolor}{\rmfamily\fontsize{12.000000}{14.400000}\selectfont\catcode`\^=\active\def^{\ifmmode\sp\else\^{}\fi}\catcode`\%=\active\def%{\%}Centro arriba}}%
\end{pgfscope}%
\begin{pgfscope}%
\pgfsetbuttcap%
\pgfsetmiterjoin%
\definecolor{currentfill}{rgb}{1.000000,1.000000,1.000000}%
\pgfsetfillcolor{currentfill}%
\pgfsetlinewidth{0.000000pt}%
\definecolor{currentstroke}{rgb}{0.000000,0.000000,0.000000}%
\pgfsetstrokecolor{currentstroke}%
\pgfsetstrokeopacity{0.000000}%
\pgfsetdash{}{0pt}%
\pgfpathmoveto{\pgfqpoint{1.993487in}{0.320555in}}%
\pgfpathlineto{\pgfqpoint{5.592717in}{0.320555in}}%
\pgfpathlineto{\pgfqpoint{5.592717in}{1.755432in}}%
\pgfpathlineto{\pgfqpoint{1.993487in}{1.755432in}}%
\pgfpathlineto{\pgfqpoint{1.993487in}{0.320555in}}%
\pgfpathclose%
\pgfusepath{fill}%
\end{pgfscope}%
\begin{pgfscope}%
\pgfpathrectangle{\pgfqpoint{1.993487in}{0.320555in}}{\pgfqpoint{3.599230in}{1.434876in}}%
\pgfusepath{clip}%
\pgfsetbuttcap%
\pgfsetroundjoin%
\pgfsetlinewidth{0.501875pt}%
\definecolor{currentstroke}{rgb}{0.501961,0.501961,0.501961}%
\pgfsetstrokecolor{currentstroke}%
\pgfsetstrokeopacity{0.500000}%
\pgfsetdash{{1.850000pt}{0.800000pt}}{0.000000pt}%
\pgfpathmoveto{\pgfqpoint{2.157088in}{0.320555in}}%
\pgfpathlineto{\pgfqpoint{2.157088in}{1.755432in}}%
\pgfusepath{stroke}%
\end{pgfscope}%
\begin{pgfscope}%
\pgfsetbuttcap%
\pgfsetroundjoin%
\definecolor{currentfill}{rgb}{0.000000,0.000000,0.000000}%
\pgfsetfillcolor{currentfill}%
\pgfsetlinewidth{0.803000pt}%
\definecolor{currentstroke}{rgb}{0.000000,0.000000,0.000000}%
\pgfsetstrokecolor{currentstroke}%
\pgfsetdash{}{0pt}%
\pgfsys@defobject{currentmarker}{\pgfqpoint{0.000000in}{-0.048611in}}{\pgfqpoint{0.000000in}{0.000000in}}{%
\pgfpathmoveto{\pgfqpoint{0.000000in}{0.000000in}}%
\pgfpathlineto{\pgfqpoint{0.000000in}{-0.048611in}}%
\pgfusepath{stroke,fill}%
}%
\begin{pgfscope}%
\pgfsys@transformshift{2.157088in}{0.320555in}%
\pgfsys@useobject{currentmarker}{}%
\end{pgfscope}%
\end{pgfscope}%
\begin{pgfscope}%
\definecolor{textcolor}{rgb}{0.000000,0.000000,0.000000}%
\pgfsetstrokecolor{textcolor}%
\pgfsetfillcolor{textcolor}%
\pgftext[x=2.157088in,y=0.223333in,,top]{\color{textcolor}{\rmfamily\fontsize{10.000000}{12.000000}\selectfont\catcode`\^=\active\def^{\ifmmode\sp\else\^{}\fi}\catcode`\%=\active\def%{\%}$\mathdefault{0}$}}%
\end{pgfscope}%
\begin{pgfscope}%
\pgfpathrectangle{\pgfqpoint{1.993487in}{0.320555in}}{\pgfqpoint{3.599230in}{1.434876in}}%
\pgfusepath{clip}%
\pgfsetbuttcap%
\pgfsetroundjoin%
\pgfsetlinewidth{0.501875pt}%
\definecolor{currentstroke}{rgb}{0.501961,0.501961,0.501961}%
\pgfsetstrokecolor{currentstroke}%
\pgfsetstrokeopacity{0.500000}%
\pgfsetdash{{1.850000pt}{0.800000pt}}{0.000000pt}%
\pgfpathmoveto{\pgfqpoint{3.025020in}{0.320555in}}%
\pgfpathlineto{\pgfqpoint{3.025020in}{1.755432in}}%
\pgfusepath{stroke}%
\end{pgfscope}%
\begin{pgfscope}%
\pgfsetbuttcap%
\pgfsetroundjoin%
\definecolor{currentfill}{rgb}{0.000000,0.000000,0.000000}%
\pgfsetfillcolor{currentfill}%
\pgfsetlinewidth{0.803000pt}%
\definecolor{currentstroke}{rgb}{0.000000,0.000000,0.000000}%
\pgfsetstrokecolor{currentstroke}%
\pgfsetdash{}{0pt}%
\pgfsys@defobject{currentmarker}{\pgfqpoint{0.000000in}{-0.048611in}}{\pgfqpoint{0.000000in}{0.000000in}}{%
\pgfpathmoveto{\pgfqpoint{0.000000in}{0.000000in}}%
\pgfpathlineto{\pgfqpoint{0.000000in}{-0.048611in}}%
\pgfusepath{stroke,fill}%
}%
\begin{pgfscope}%
\pgfsys@transformshift{3.025020in}{0.320555in}%
\pgfsys@useobject{currentmarker}{}%
\end{pgfscope}%
\end{pgfscope}%
\begin{pgfscope}%
\definecolor{textcolor}{rgb}{0.000000,0.000000,0.000000}%
\pgfsetstrokecolor{textcolor}%
\pgfsetfillcolor{textcolor}%
\pgftext[x=3.025020in,y=0.223333in,,top]{\color{textcolor}{\rmfamily\fontsize{10.000000}{12.000000}\selectfont\catcode`\^=\active\def^{\ifmmode\sp\else\^{}\fi}\catcode`\%=\active\def%{\%}$\mathdefault{5}$}}%
\end{pgfscope}%
\begin{pgfscope}%
\pgfpathrectangle{\pgfqpoint{1.993487in}{0.320555in}}{\pgfqpoint{3.599230in}{1.434876in}}%
\pgfusepath{clip}%
\pgfsetbuttcap%
\pgfsetroundjoin%
\pgfsetlinewidth{0.501875pt}%
\definecolor{currentstroke}{rgb}{0.501961,0.501961,0.501961}%
\pgfsetstrokecolor{currentstroke}%
\pgfsetstrokeopacity{0.500000}%
\pgfsetdash{{1.850000pt}{0.800000pt}}{0.000000pt}%
\pgfpathmoveto{\pgfqpoint{3.892953in}{0.320555in}}%
\pgfpathlineto{\pgfqpoint{3.892953in}{1.755432in}}%
\pgfusepath{stroke}%
\end{pgfscope}%
\begin{pgfscope}%
\pgfsetbuttcap%
\pgfsetroundjoin%
\definecolor{currentfill}{rgb}{0.000000,0.000000,0.000000}%
\pgfsetfillcolor{currentfill}%
\pgfsetlinewidth{0.803000pt}%
\definecolor{currentstroke}{rgb}{0.000000,0.000000,0.000000}%
\pgfsetstrokecolor{currentstroke}%
\pgfsetdash{}{0pt}%
\pgfsys@defobject{currentmarker}{\pgfqpoint{0.000000in}{-0.048611in}}{\pgfqpoint{0.000000in}{0.000000in}}{%
\pgfpathmoveto{\pgfqpoint{0.000000in}{0.000000in}}%
\pgfpathlineto{\pgfqpoint{0.000000in}{-0.048611in}}%
\pgfusepath{stroke,fill}%
}%
\begin{pgfscope}%
\pgfsys@transformshift{3.892953in}{0.320555in}%
\pgfsys@useobject{currentmarker}{}%
\end{pgfscope}%
\end{pgfscope}%
\begin{pgfscope}%
\definecolor{textcolor}{rgb}{0.000000,0.000000,0.000000}%
\pgfsetstrokecolor{textcolor}%
\pgfsetfillcolor{textcolor}%
\pgftext[x=3.892953in,y=0.223333in,,top]{\color{textcolor}{\rmfamily\fontsize{10.000000}{12.000000}\selectfont\catcode`\^=\active\def^{\ifmmode\sp\else\^{}\fi}\catcode`\%=\active\def%{\%}$\mathdefault{10}$}}%
\end{pgfscope}%
\begin{pgfscope}%
\pgfpathrectangle{\pgfqpoint{1.993487in}{0.320555in}}{\pgfqpoint{3.599230in}{1.434876in}}%
\pgfusepath{clip}%
\pgfsetbuttcap%
\pgfsetroundjoin%
\pgfsetlinewidth{0.501875pt}%
\definecolor{currentstroke}{rgb}{0.501961,0.501961,0.501961}%
\pgfsetstrokecolor{currentstroke}%
\pgfsetstrokeopacity{0.500000}%
\pgfsetdash{{1.850000pt}{0.800000pt}}{0.000000pt}%
\pgfpathmoveto{\pgfqpoint{4.760885in}{0.320555in}}%
\pgfpathlineto{\pgfqpoint{4.760885in}{1.755432in}}%
\pgfusepath{stroke}%
\end{pgfscope}%
\begin{pgfscope}%
\pgfsetbuttcap%
\pgfsetroundjoin%
\definecolor{currentfill}{rgb}{0.000000,0.000000,0.000000}%
\pgfsetfillcolor{currentfill}%
\pgfsetlinewidth{0.803000pt}%
\definecolor{currentstroke}{rgb}{0.000000,0.000000,0.000000}%
\pgfsetstrokecolor{currentstroke}%
\pgfsetdash{}{0pt}%
\pgfsys@defobject{currentmarker}{\pgfqpoint{0.000000in}{-0.048611in}}{\pgfqpoint{0.000000in}{0.000000in}}{%
\pgfpathmoveto{\pgfqpoint{0.000000in}{0.000000in}}%
\pgfpathlineto{\pgfqpoint{0.000000in}{-0.048611in}}%
\pgfusepath{stroke,fill}%
}%
\begin{pgfscope}%
\pgfsys@transformshift{4.760885in}{0.320555in}%
\pgfsys@useobject{currentmarker}{}%
\end{pgfscope}%
\end{pgfscope}%
\begin{pgfscope}%
\definecolor{textcolor}{rgb}{0.000000,0.000000,0.000000}%
\pgfsetstrokecolor{textcolor}%
\pgfsetfillcolor{textcolor}%
\pgftext[x=4.760885in,y=0.223333in,,top]{\color{textcolor}{\rmfamily\fontsize{10.000000}{12.000000}\selectfont\catcode`\^=\active\def^{\ifmmode\sp\else\^{}\fi}\catcode`\%=\active\def%{\%}$\mathdefault{15}$}}%
\end{pgfscope}%
\begin{pgfscope}%
\pgfpathrectangle{\pgfqpoint{1.993487in}{0.320555in}}{\pgfqpoint{3.599230in}{1.434876in}}%
\pgfusepath{clip}%
\pgfsetbuttcap%
\pgfsetroundjoin%
\pgfsetlinewidth{0.501875pt}%
\definecolor{currentstroke}{rgb}{0.501961,0.501961,0.501961}%
\pgfsetstrokecolor{currentstroke}%
\pgfsetstrokeopacity{0.500000}%
\pgfsetdash{{1.850000pt}{0.800000pt}}{0.000000pt}%
\pgfpathmoveto{\pgfqpoint{1.993487in}{0.385775in}}%
\pgfpathlineto{\pgfqpoint{5.592717in}{0.385775in}}%
\pgfusepath{stroke}%
\end{pgfscope}%
\begin{pgfscope}%
\pgfsetbuttcap%
\pgfsetroundjoin%
\definecolor{currentfill}{rgb}{0.000000,0.000000,0.000000}%
\pgfsetfillcolor{currentfill}%
\pgfsetlinewidth{0.803000pt}%
\definecolor{currentstroke}{rgb}{0.000000,0.000000,0.000000}%
\pgfsetstrokecolor{currentstroke}%
\pgfsetdash{}{0pt}%
\pgfsys@defobject{currentmarker}{\pgfqpoint{-0.048611in}{0.000000in}}{\pgfqpoint{-0.000000in}{0.000000in}}{%
\pgfpathmoveto{\pgfqpoint{-0.000000in}{0.000000in}}%
\pgfpathlineto{\pgfqpoint{-0.048611in}{0.000000in}}%
\pgfusepath{stroke,fill}%
}%
\begin{pgfscope}%
\pgfsys@transformshift{1.993487in}{0.385775in}%
\pgfsys@useobject{currentmarker}{}%
\end{pgfscope}%
\end{pgfscope}%
\begin{pgfscope}%
\definecolor{textcolor}{rgb}{0.000000,0.000000,0.000000}%
\pgfsetstrokecolor{textcolor}%
\pgfsetfillcolor{textcolor}%
\pgftext[x=1.610770in, y=0.337580in, left, base]{\color{textcolor}{\rmfamily\fontsize{10.000000}{12.000000}\selectfont\catcode`\^=\active\def^{\ifmmode\sp\else\^{}\fi}\catcode`\%=\active\def%{\%}$\mathdefault{\ensuremath{-}1.0}$}}%
\end{pgfscope}%
\begin{pgfscope}%
\pgfpathrectangle{\pgfqpoint{1.993487in}{0.320555in}}{\pgfqpoint{3.599230in}{1.434876in}}%
\pgfusepath{clip}%
\pgfsetbuttcap%
\pgfsetroundjoin%
\pgfsetlinewidth{0.501875pt}%
\definecolor{currentstroke}{rgb}{0.501961,0.501961,0.501961}%
\pgfsetstrokecolor{currentstroke}%
\pgfsetstrokeopacity{0.500000}%
\pgfsetdash{{1.850000pt}{0.800000pt}}{0.000000pt}%
\pgfpathmoveto{\pgfqpoint{1.993487in}{0.711884in}}%
\pgfpathlineto{\pgfqpoint{5.592717in}{0.711884in}}%
\pgfusepath{stroke}%
\end{pgfscope}%
\begin{pgfscope}%
\pgfsetbuttcap%
\pgfsetroundjoin%
\definecolor{currentfill}{rgb}{0.000000,0.000000,0.000000}%
\pgfsetfillcolor{currentfill}%
\pgfsetlinewidth{0.803000pt}%
\definecolor{currentstroke}{rgb}{0.000000,0.000000,0.000000}%
\pgfsetstrokecolor{currentstroke}%
\pgfsetdash{}{0pt}%
\pgfsys@defobject{currentmarker}{\pgfqpoint{-0.048611in}{0.000000in}}{\pgfqpoint{-0.000000in}{0.000000in}}{%
\pgfpathmoveto{\pgfqpoint{-0.000000in}{0.000000in}}%
\pgfpathlineto{\pgfqpoint{-0.048611in}{0.000000in}}%
\pgfusepath{stroke,fill}%
}%
\begin{pgfscope}%
\pgfsys@transformshift{1.993487in}{0.711884in}%
\pgfsys@useobject{currentmarker}{}%
\end{pgfscope}%
\end{pgfscope}%
\begin{pgfscope}%
\definecolor{textcolor}{rgb}{0.000000,0.000000,0.000000}%
\pgfsetstrokecolor{textcolor}%
\pgfsetfillcolor{textcolor}%
\pgftext[x=1.610770in, y=0.663690in, left, base]{\color{textcolor}{\rmfamily\fontsize{10.000000}{12.000000}\selectfont\catcode`\^=\active\def^{\ifmmode\sp\else\^{}\fi}\catcode`\%=\active\def%{\%}$\mathdefault{\ensuremath{-}0.5}$}}%
\end{pgfscope}%
\begin{pgfscope}%
\pgfpathrectangle{\pgfqpoint{1.993487in}{0.320555in}}{\pgfqpoint{3.599230in}{1.434876in}}%
\pgfusepath{clip}%
\pgfsetbuttcap%
\pgfsetroundjoin%
\pgfsetlinewidth{0.501875pt}%
\definecolor{currentstroke}{rgb}{0.501961,0.501961,0.501961}%
\pgfsetstrokecolor{currentstroke}%
\pgfsetstrokeopacity{0.500000}%
\pgfsetdash{{1.850000pt}{0.800000pt}}{0.000000pt}%
\pgfpathmoveto{\pgfqpoint{1.993487in}{1.037994in}}%
\pgfpathlineto{\pgfqpoint{5.592717in}{1.037994in}}%
\pgfusepath{stroke}%
\end{pgfscope}%
\begin{pgfscope}%
\pgfsetbuttcap%
\pgfsetroundjoin%
\definecolor{currentfill}{rgb}{0.000000,0.000000,0.000000}%
\pgfsetfillcolor{currentfill}%
\pgfsetlinewidth{0.803000pt}%
\definecolor{currentstroke}{rgb}{0.000000,0.000000,0.000000}%
\pgfsetstrokecolor{currentstroke}%
\pgfsetdash{}{0pt}%
\pgfsys@defobject{currentmarker}{\pgfqpoint{-0.048611in}{0.000000in}}{\pgfqpoint{-0.000000in}{0.000000in}}{%
\pgfpathmoveto{\pgfqpoint{-0.000000in}{0.000000in}}%
\pgfpathlineto{\pgfqpoint{-0.048611in}{0.000000in}}%
\pgfusepath{stroke,fill}%
}%
\begin{pgfscope}%
\pgfsys@transformshift{1.993487in}{1.037994in}%
\pgfsys@useobject{currentmarker}{}%
\end{pgfscope}%
\end{pgfscope}%
\begin{pgfscope}%
\definecolor{textcolor}{rgb}{0.000000,0.000000,0.000000}%
\pgfsetstrokecolor{textcolor}%
\pgfsetfillcolor{textcolor}%
\pgftext[x=1.718795in, y=0.989799in, left, base]{\color{textcolor}{\rmfamily\fontsize{10.000000}{12.000000}\selectfont\catcode`\^=\active\def^{\ifmmode\sp\else\^{}\fi}\catcode`\%=\active\def%{\%}$\mathdefault{0.0}$}}%
\end{pgfscope}%
\begin{pgfscope}%
\pgfpathrectangle{\pgfqpoint{1.993487in}{0.320555in}}{\pgfqpoint{3.599230in}{1.434876in}}%
\pgfusepath{clip}%
\pgfsetbuttcap%
\pgfsetroundjoin%
\pgfsetlinewidth{0.501875pt}%
\definecolor{currentstroke}{rgb}{0.501961,0.501961,0.501961}%
\pgfsetstrokecolor{currentstroke}%
\pgfsetstrokeopacity{0.500000}%
\pgfsetdash{{1.850000pt}{0.800000pt}}{0.000000pt}%
\pgfpathmoveto{\pgfqpoint{1.993487in}{1.364103in}}%
\pgfpathlineto{\pgfqpoint{5.592717in}{1.364103in}}%
\pgfusepath{stroke}%
\end{pgfscope}%
\begin{pgfscope}%
\pgfsetbuttcap%
\pgfsetroundjoin%
\definecolor{currentfill}{rgb}{0.000000,0.000000,0.000000}%
\pgfsetfillcolor{currentfill}%
\pgfsetlinewidth{0.803000pt}%
\definecolor{currentstroke}{rgb}{0.000000,0.000000,0.000000}%
\pgfsetstrokecolor{currentstroke}%
\pgfsetdash{}{0pt}%
\pgfsys@defobject{currentmarker}{\pgfqpoint{-0.048611in}{0.000000in}}{\pgfqpoint{-0.000000in}{0.000000in}}{%
\pgfpathmoveto{\pgfqpoint{-0.000000in}{0.000000in}}%
\pgfpathlineto{\pgfqpoint{-0.048611in}{0.000000in}}%
\pgfusepath{stroke,fill}%
}%
\begin{pgfscope}%
\pgfsys@transformshift{1.993487in}{1.364103in}%
\pgfsys@useobject{currentmarker}{}%
\end{pgfscope}%
\end{pgfscope}%
\begin{pgfscope}%
\definecolor{textcolor}{rgb}{0.000000,0.000000,0.000000}%
\pgfsetstrokecolor{textcolor}%
\pgfsetfillcolor{textcolor}%
\pgftext[x=1.718795in, y=1.315908in, left, base]{\color{textcolor}{\rmfamily\fontsize{10.000000}{12.000000}\selectfont\catcode`\^=\active\def^{\ifmmode\sp\else\^{}\fi}\catcode`\%=\active\def%{\%}$\mathdefault{0.5}$}}%
\end{pgfscope}%
\begin{pgfscope}%
\pgfpathrectangle{\pgfqpoint{1.993487in}{0.320555in}}{\pgfqpoint{3.599230in}{1.434876in}}%
\pgfusepath{clip}%
\pgfsetbuttcap%
\pgfsetroundjoin%
\pgfsetlinewidth{0.501875pt}%
\definecolor{currentstroke}{rgb}{0.501961,0.501961,0.501961}%
\pgfsetstrokecolor{currentstroke}%
\pgfsetstrokeopacity{0.500000}%
\pgfsetdash{{1.850000pt}{0.800000pt}}{0.000000pt}%
\pgfpathmoveto{\pgfqpoint{1.993487in}{1.690212in}}%
\pgfpathlineto{\pgfqpoint{5.592717in}{1.690212in}}%
\pgfusepath{stroke}%
\end{pgfscope}%
\begin{pgfscope}%
\pgfsetbuttcap%
\pgfsetroundjoin%
\definecolor{currentfill}{rgb}{0.000000,0.000000,0.000000}%
\pgfsetfillcolor{currentfill}%
\pgfsetlinewidth{0.803000pt}%
\definecolor{currentstroke}{rgb}{0.000000,0.000000,0.000000}%
\pgfsetstrokecolor{currentstroke}%
\pgfsetdash{}{0pt}%
\pgfsys@defobject{currentmarker}{\pgfqpoint{-0.048611in}{0.000000in}}{\pgfqpoint{-0.000000in}{0.000000in}}{%
\pgfpathmoveto{\pgfqpoint{-0.000000in}{0.000000in}}%
\pgfpathlineto{\pgfqpoint{-0.048611in}{0.000000in}}%
\pgfusepath{stroke,fill}%
}%
\begin{pgfscope}%
\pgfsys@transformshift{1.993487in}{1.690212in}%
\pgfsys@useobject{currentmarker}{}%
\end{pgfscope}%
\end{pgfscope}%
\begin{pgfscope}%
\definecolor{textcolor}{rgb}{0.000000,0.000000,0.000000}%
\pgfsetstrokecolor{textcolor}%
\pgfsetfillcolor{textcolor}%
\pgftext[x=1.718795in, y=1.642018in, left, base]{\color{textcolor}{\rmfamily\fontsize{10.000000}{12.000000}\selectfont\catcode`\^=\active\def^{\ifmmode\sp\else\^{}\fi}\catcode`\%=\active\def%{\%}$\mathdefault{1.0}$}}%
\end{pgfscope}%
\begin{pgfscope}%
\pgfpathrectangle{\pgfqpoint{1.993487in}{0.320555in}}{\pgfqpoint{3.599230in}{1.434876in}}%
\pgfusepath{clip}%
\pgfsetrectcap%
\pgfsetroundjoin%
\pgfsetlinewidth{1.505625pt}%
\definecolor{currentstroke}{rgb}{0.000000,0.501961,0.000000}%
\pgfsetstrokecolor{currentstroke}%
\pgfsetdash{}{0pt}%
\pgfpathmoveto{\pgfqpoint{2.157088in}{1.037994in}}%
\pgfpathlineto{\pgfqpoint{2.189863in}{1.278473in}}%
\pgfpathlineto{\pgfqpoint{2.211713in}{1.421911in}}%
\pgfpathlineto{\pgfqpoint{2.228100in}{1.514049in}}%
\pgfpathlineto{\pgfqpoint{2.239025in}{1.566252in}}%
\pgfpathlineto{\pgfqpoint{2.249950in}{1.610097in}}%
\pgfpathlineto{\pgfqpoint{2.260875in}{1.644889in}}%
\pgfpathlineto{\pgfqpoint{2.271800in}{1.670077in}}%
\pgfpathlineto{\pgfqpoint{2.277263in}{1.678940in}}%
\pgfpathlineto{\pgfqpoint{2.282725in}{1.685265in}}%
\pgfpathlineto{\pgfqpoint{2.288188in}{1.689026in}}%
\pgfpathlineto{\pgfqpoint{2.293650in}{1.690210in}}%
\pgfpathlineto{\pgfqpoint{2.299113in}{1.688811in}}%
\pgfpathlineto{\pgfqpoint{2.304575in}{1.684835in}}%
\pgfpathlineto{\pgfqpoint{2.310038in}{1.678298in}}%
\pgfpathlineto{\pgfqpoint{2.315500in}{1.669225in}}%
\pgfpathlineto{\pgfqpoint{2.326425in}{1.643627in}}%
\pgfpathlineto{\pgfqpoint{2.337350in}{1.608446in}}%
\pgfpathlineto{\pgfqpoint{2.348275in}{1.564239in}}%
\pgfpathlineto{\pgfqpoint{2.359200in}{1.511704in}}%
\pgfpathlineto{\pgfqpoint{2.375587in}{1.419141in}}%
\pgfpathlineto{\pgfqpoint{2.391975in}{1.313030in}}%
\pgfpathlineto{\pgfqpoint{2.413825in}{1.157047in}}%
\pgfpathlineto{\pgfqpoint{2.468450in}{0.756769in}}%
\pgfpathlineto{\pgfqpoint{2.484837in}{0.651316in}}%
\pgfpathlineto{\pgfqpoint{2.501225in}{0.559607in}}%
\pgfpathlineto{\pgfqpoint{2.512150in}{0.507736in}}%
\pgfpathlineto{\pgfqpoint{2.523075in}{0.464256in}}%
\pgfpathlineto{\pgfqpoint{2.533999in}{0.429854in}}%
\pgfpathlineto{\pgfqpoint{2.544924in}{0.405075in}}%
\pgfpathlineto{\pgfqpoint{2.550387in}{0.396423in}}%
\pgfpathlineto{\pgfqpoint{2.555849in}{0.390311in}}%
\pgfpathlineto{\pgfqpoint{2.561312in}{0.386764in}}%
\pgfpathlineto{\pgfqpoint{2.566774in}{0.385795in}}%
\pgfpathlineto{\pgfqpoint{2.572237in}{0.387409in}}%
\pgfpathlineto{\pgfqpoint{2.577699in}{0.391599in}}%
\pgfpathlineto{\pgfqpoint{2.583162in}{0.398349in}}%
\pgfpathlineto{\pgfqpoint{2.588624in}{0.407631in}}%
\pgfpathlineto{\pgfqpoint{2.599549in}{0.433638in}}%
\pgfpathlineto{\pgfqpoint{2.610474in}{0.469207in}}%
\pgfpathlineto{\pgfqpoint{2.621399in}{0.513777in}}%
\pgfpathlineto{\pgfqpoint{2.632324in}{0.566641in}}%
\pgfpathlineto{\pgfqpoint{2.648712in}{0.659627in}}%
\pgfpathlineto{\pgfqpoint{2.665099in}{0.766062in}}%
\pgfpathlineto{\pgfqpoint{2.692412in}{0.962905in}}%
\pgfpathlineto{\pgfqpoint{2.736111in}{1.284819in}}%
\pgfpathlineto{\pgfqpoint{2.757961in}{1.427421in}}%
\pgfpathlineto{\pgfqpoint{2.774349in}{1.518699in}}%
\pgfpathlineto{\pgfqpoint{2.785274in}{1.570236in}}%
\pgfpathlineto{\pgfqpoint{2.796199in}{1.613350in}}%
\pgfpathlineto{\pgfqpoint{2.807124in}{1.647361in}}%
\pgfpathlineto{\pgfqpoint{2.818049in}{1.671729in}}%
\pgfpathlineto{\pgfqpoint{2.823511in}{1.680171in}}%
\pgfpathlineto{\pgfqpoint{2.828974in}{1.686070in}}%
\pgfpathlineto{\pgfqpoint{2.834436in}{1.689403in}}%
\pgfpathlineto{\pgfqpoint{2.839899in}{1.690156in}}%
\pgfpathlineto{\pgfqpoint{2.845361in}{1.688327in}}%
\pgfpathlineto{\pgfqpoint{2.850824in}{1.683923in}}%
\pgfpathlineto{\pgfqpoint{2.856286in}{1.676961in}}%
\pgfpathlineto{\pgfqpoint{2.861748in}{1.667469in}}%
\pgfpathlineto{\pgfqpoint{2.872673in}{1.641055in}}%
\pgfpathlineto{\pgfqpoint{2.883598in}{1.605098in}}%
\pgfpathlineto{\pgfqpoint{2.894523in}{1.560168in}}%
\pgfpathlineto{\pgfqpoint{2.905448in}{1.506976in}}%
\pgfpathlineto{\pgfqpoint{2.921836in}{1.413568in}}%
\pgfpathlineto{\pgfqpoint{2.938223in}{1.306812in}}%
\pgfpathlineto{\pgfqpoint{2.965536in}{1.109683in}}%
\pgfpathlineto{\pgfqpoint{3.009236in}{0.788005in}}%
\pgfpathlineto{\pgfqpoint{3.031085in}{0.645828in}}%
\pgfpathlineto{\pgfqpoint{3.047473in}{0.554983in}}%
\pgfpathlineto{\pgfqpoint{3.058398in}{0.503782in}}%
\pgfpathlineto{\pgfqpoint{3.069323in}{0.461034in}}%
\pgfpathlineto{\pgfqpoint{3.080248in}{0.427415in}}%
\pgfpathlineto{\pgfqpoint{3.091173in}{0.403458in}}%
\pgfpathlineto{\pgfqpoint{3.096635in}{0.395227in}}%
\pgfpathlineto{\pgfqpoint{3.102098in}{0.389541in}}%
\pgfpathlineto{\pgfqpoint{3.107560in}{0.386423in}}%
\pgfpathlineto{\pgfqpoint{3.113023in}{0.385885in}}%
\pgfpathlineto{\pgfqpoint{3.118485in}{0.387929in}}%
\pgfpathlineto{\pgfqpoint{3.123948in}{0.392547in}}%
\pgfpathlineto{\pgfqpoint{3.129410in}{0.399721in}}%
\pgfpathlineto{\pgfqpoint{3.134873in}{0.409422in}}%
\pgfpathlineto{\pgfqpoint{3.145798in}{0.436243in}}%
\pgfpathlineto{\pgfqpoint{3.156723in}{0.472586in}}%
\pgfpathlineto{\pgfqpoint{3.167648in}{0.517876in}}%
\pgfpathlineto{\pgfqpoint{3.178573in}{0.571395in}}%
\pgfpathlineto{\pgfqpoint{3.194960in}{0.665221in}}%
\pgfpathlineto{\pgfqpoint{3.211347in}{0.772296in}}%
\pgfpathlineto{\pgfqpoint{3.238660in}{0.969705in}}%
\pgfpathlineto{\pgfqpoint{3.282360in}{1.291138in}}%
\pgfpathlineto{\pgfqpoint{3.304210in}{1.432887in}}%
\pgfpathlineto{\pgfqpoint{3.320597in}{1.523296in}}%
\pgfpathlineto{\pgfqpoint{3.331522in}{1.574160in}}%
\pgfpathlineto{\pgfqpoint{3.342447in}{1.616541in}}%
\pgfpathlineto{\pgfqpoint{3.353372in}{1.649766in}}%
\pgfpathlineto{\pgfqpoint{3.364297in}{1.673312in}}%
\pgfpathlineto{\pgfqpoint{3.369759in}{1.681332in}}%
\pgfpathlineto{\pgfqpoint{3.375222in}{1.686804in}}%
\pgfpathlineto{\pgfqpoint{3.380684in}{1.689708in}}%
\pgfpathlineto{\pgfqpoint{3.386147in}{1.690031in}}%
\pgfpathlineto{\pgfqpoint{3.391609in}{1.687772in}}%
\pgfpathlineto{\pgfqpoint{3.397072in}{1.682940in}}%
\pgfpathlineto{\pgfqpoint{3.402534in}{1.675554in}}%
\pgfpathlineto{\pgfqpoint{3.407997in}{1.665644in}}%
\pgfpathlineto{\pgfqpoint{3.418922in}{1.638416in}}%
\pgfpathlineto{\pgfqpoint{3.429847in}{1.601688in}}%
\pgfpathlineto{\pgfqpoint{3.440772in}{1.556040in}}%
\pgfpathlineto{\pgfqpoint{3.451697in}{1.502195in}}%
\pgfpathlineto{\pgfqpoint{3.468084in}{1.407954in}}%
\pgfpathlineto{\pgfqpoint{3.484472in}{1.300564in}}%
\pgfpathlineto{\pgfqpoint{3.511784in}{1.102880in}}%
\pgfpathlineto{\pgfqpoint{3.555484in}{0.781700in}}%
\pgfpathlineto{\pgfqpoint{3.577334in}{0.640383in}}%
\pgfpathlineto{\pgfqpoint{3.593721in}{0.550412in}}%
\pgfpathlineto{\pgfqpoint{3.604646in}{0.499886in}}%
\pgfpathlineto{\pgfqpoint{3.615571in}{0.457875in}}%
\pgfpathlineto{\pgfqpoint{3.626496in}{0.425043in}}%
\pgfpathlineto{\pgfqpoint{3.637421in}{0.401911in}}%
\pgfpathlineto{\pgfqpoint{3.642884in}{0.394102in}}%
\pgfpathlineto{\pgfqpoint{3.648346in}{0.388843in}}%
\pgfpathlineto{\pgfqpoint{3.653809in}{0.386154in}}%
\pgfpathlineto{\pgfqpoint{3.659271in}{0.386046in}}%
\pgfpathlineto{\pgfqpoint{3.664734in}{0.388520in}}%
\pgfpathlineto{\pgfqpoint{3.670196in}{0.393566in}}%
\pgfpathlineto{\pgfqpoint{3.675659in}{0.401163in}}%
\pgfpathlineto{\pgfqpoint{3.681121in}{0.411282in}}%
\pgfpathlineto{\pgfqpoint{3.692046in}{0.438915in}}%
\pgfpathlineto{\pgfqpoint{3.702971in}{0.476028in}}%
\pgfpathlineto{\pgfqpoint{3.713896in}{0.522032in}}%
\pgfpathlineto{\pgfqpoint{3.724821in}{0.576201in}}%
\pgfpathlineto{\pgfqpoint{3.741208in}{0.670855in}}%
\pgfpathlineto{\pgfqpoint{3.757596in}{0.778558in}}%
\pgfpathlineto{\pgfqpoint{3.784908in}{0.976512in}}%
\pgfpathlineto{\pgfqpoint{3.828608in}{1.297429in}}%
\pgfpathlineto{\pgfqpoint{3.850458in}{1.438310in}}%
\pgfpathlineto{\pgfqpoint{3.866845in}{1.527840in}}%
\pgfpathlineto{\pgfqpoint{3.877770in}{1.578026in}}%
\pgfpathlineto{\pgfqpoint{3.888695in}{1.619667in}}%
\pgfpathlineto{\pgfqpoint{3.899620in}{1.652104in}}%
\pgfpathlineto{\pgfqpoint{3.910545in}{1.674824in}}%
\pgfpathlineto{\pgfqpoint{3.916008in}{1.682421in}}%
\pgfpathlineto{\pgfqpoint{3.921470in}{1.687467in}}%
\pgfpathlineto{\pgfqpoint{3.926933in}{1.689941in}}%
\pgfpathlineto{\pgfqpoint{3.932395in}{1.689833in}}%
\pgfpathlineto{\pgfqpoint{3.937858in}{1.687145in}}%
\pgfpathlineto{\pgfqpoint{3.943320in}{1.681885in}}%
\pgfpathlineto{\pgfqpoint{3.948783in}{1.674077in}}%
\pgfpathlineto{\pgfqpoint{3.954245in}{1.663749in}}%
\pgfpathlineto{\pgfqpoint{3.965170in}{1.635711in}}%
\pgfpathlineto{\pgfqpoint{3.976095in}{1.598216in}}%
\pgfpathlineto{\pgfqpoint{3.987020in}{1.551855in}}%
\pgfpathlineto{\pgfqpoint{3.997945in}{1.497364in}}%
\pgfpathlineto{\pgfqpoint{4.014332in}{1.402300in}}%
\pgfpathlineto{\pgfqpoint{4.030720in}{1.294287in}}%
\pgfpathlineto{\pgfqpoint{4.058032in}{1.096069in}}%
\pgfpathlineto{\pgfqpoint{4.101732in}{0.775423in}}%
\pgfpathlineto{\pgfqpoint{4.118120in}{0.668033in}}%
\pgfpathlineto{\pgfqpoint{4.134507in}{0.573792in}}%
\pgfpathlineto{\pgfqpoint{4.145432in}{0.519947in}}%
\pgfpathlineto{\pgfqpoint{4.156357in}{0.474299in}}%
\pgfpathlineto{\pgfqpoint{4.167282in}{0.437571in}}%
\pgfpathlineto{\pgfqpoint{4.178207in}{0.410344in}}%
\pgfpathlineto{\pgfqpoint{4.183669in}{0.400433in}}%
\pgfpathlineto{\pgfqpoint{4.189132in}{0.393047in}}%
\pgfpathlineto{\pgfqpoint{4.194594in}{0.388215in}}%
\pgfpathlineto{\pgfqpoint{4.200057in}{0.385956in}}%
\pgfpathlineto{\pgfqpoint{4.205519in}{0.386279in}}%
\pgfpathlineto{\pgfqpoint{4.210982in}{0.389183in}}%
\pgfpathlineto{\pgfqpoint{4.216444in}{0.394655in}}%
\pgfpathlineto{\pgfqpoint{4.221907in}{0.402675in}}%
\pgfpathlineto{\pgfqpoint{4.227369in}{0.413211in}}%
\pgfpathlineto{\pgfqpoint{4.238294in}{0.441653in}}%
\pgfpathlineto{\pgfqpoint{4.249219in}{0.479531in}}%
\pgfpathlineto{\pgfqpoint{4.260144in}{0.526245in}}%
\pgfpathlineto{\pgfqpoint{4.271069in}{0.581058in}}%
\pgfpathlineto{\pgfqpoint{4.287457in}{0.676530in}}%
\pgfpathlineto{\pgfqpoint{4.303844in}{0.784849in}}%
\pgfpathlineto{\pgfqpoint{4.331157in}{0.983326in}}%
\pgfpathlineto{\pgfqpoint{4.374856in}{1.303691in}}%
\pgfpathlineto{\pgfqpoint{4.391244in}{1.410767in}}%
\pgfpathlineto{\pgfqpoint{4.407631in}{1.504592in}}%
\pgfpathlineto{\pgfqpoint{4.418556in}{1.558111in}}%
\pgfpathlineto{\pgfqpoint{4.429481in}{1.603401in}}%
\pgfpathlineto{\pgfqpoint{4.440406in}{1.639744in}}%
\pgfpathlineto{\pgfqpoint{4.451331in}{1.666565in}}%
\pgfpathlineto{\pgfqpoint{4.456794in}{1.676266in}}%
\pgfpathlineto{\pgfqpoint{4.462256in}{1.683440in}}%
\pgfpathlineto{\pgfqpoint{4.467719in}{1.688058in}}%
\pgfpathlineto{\pgfqpoint{4.473181in}{1.690102in}}%
\pgfpathlineto{\pgfqpoint{4.478644in}{1.689564in}}%
\pgfpathlineto{\pgfqpoint{4.484106in}{1.686446in}}%
\pgfpathlineto{\pgfqpoint{4.489569in}{1.680760in}}%
\pgfpathlineto{\pgfqpoint{4.495031in}{1.672529in}}%
\pgfpathlineto{\pgfqpoint{4.500494in}{1.661786in}}%
\pgfpathlineto{\pgfqpoint{4.511418in}{1.632941in}}%
\pgfpathlineto{\pgfqpoint{4.522343in}{1.594682in}}%
\pgfpathlineto{\pgfqpoint{4.533268in}{1.547614in}}%
\pgfpathlineto{\pgfqpoint{4.544193in}{1.492482in}}%
\pgfpathlineto{\pgfqpoint{4.560581in}{1.396605in}}%
\pgfpathlineto{\pgfqpoint{4.582431in}{1.249598in}}%
\pgfpathlineto{\pgfqpoint{4.615206in}{1.007219in}}%
\pgfpathlineto{\pgfqpoint{4.647981in}{0.769175in}}%
\pgfpathlineto{\pgfqpoint{4.664368in}{0.662419in}}%
\pgfpathlineto{\pgfqpoint{4.680755in}{0.569012in}}%
\pgfpathlineto{\pgfqpoint{4.691680in}{0.515819in}}%
\pgfpathlineto{\pgfqpoint{4.702605in}{0.470889in}}%
\pgfpathlineto{\pgfqpoint{4.713530in}{0.434932in}}%
\pgfpathlineto{\pgfqpoint{4.724455in}{0.408518in}}%
\pgfpathlineto{\pgfqpoint{4.729918in}{0.399026in}}%
\pgfpathlineto{\pgfqpoint{4.735380in}{0.392064in}}%
\pgfpathlineto{\pgfqpoint{4.740843in}{0.387660in}}%
\pgfpathlineto{\pgfqpoint{4.746305in}{0.385831in}}%
\pgfpathlineto{\pgfqpoint{4.751768in}{0.386584in}}%
\pgfpathlineto{\pgfqpoint{4.757230in}{0.389917in}}%
\pgfpathlineto{\pgfqpoint{4.762693in}{0.395816in}}%
\pgfpathlineto{\pgfqpoint{4.768155in}{0.404258in}}%
\pgfpathlineto{\pgfqpoint{4.779080in}{0.428626in}}%
\pgfpathlineto{\pgfqpoint{4.790005in}{0.462637in}}%
\pgfpathlineto{\pgfqpoint{4.800930in}{0.505752in}}%
\pgfpathlineto{\pgfqpoint{4.811855in}{0.557288in}}%
\pgfpathlineto{\pgfqpoint{4.828243in}{0.648566in}}%
\pgfpathlineto{\pgfqpoint{4.844630in}{0.753686in}}%
\pgfpathlineto{\pgfqpoint{4.866480in}{0.908865in}}%
\pgfpathlineto{\pgfqpoint{4.926567in}{1.346672in}}%
\pgfpathlineto{\pgfqpoint{4.942955in}{1.449024in}}%
\pgfpathlineto{\pgfqpoint{4.959342in}{1.536766in}}%
\pgfpathlineto{\pgfqpoint{4.970267in}{1.585579in}}%
\pgfpathlineto{\pgfqpoint{4.981192in}{1.625728in}}%
\pgfpathlineto{\pgfqpoint{4.992117in}{1.656577in}}%
\pgfpathlineto{\pgfqpoint{5.003042in}{1.677638in}}%
\pgfpathlineto{\pgfqpoint{5.008504in}{1.684388in}}%
\pgfpathlineto{\pgfqpoint{5.013967in}{1.688578in}}%
\pgfpathlineto{\pgfqpoint{5.019429in}{1.690192in}}%
\pgfpathlineto{\pgfqpoint{5.024892in}{1.689224in}}%
\pgfpathlineto{\pgfqpoint{5.030354in}{1.685676in}}%
\pgfpathlineto{\pgfqpoint{5.035817in}{1.679564in}}%
\pgfpathlineto{\pgfqpoint{5.041279in}{1.670912in}}%
\pgfpathlineto{\pgfqpoint{5.052204in}{1.646133in}}%
\pgfpathlineto{\pgfqpoint{5.063129in}{1.611731in}}%
\pgfpathlineto{\pgfqpoint{5.074054in}{1.568251in}}%
\pgfpathlineto{\pgfqpoint{5.084979in}{1.516381in}}%
\pgfpathlineto{\pgfqpoint{5.101367in}{1.424671in}}%
\pgfpathlineto{\pgfqpoint{5.117754in}{1.319218in}}%
\pgfpathlineto{\pgfqpoint{5.139604in}{1.163767in}}%
\pgfpathlineto{\pgfqpoint{5.199691in}{0.726306in}}%
\pgfpathlineto{\pgfqpoint{5.216079in}{0.624313in}}%
\pgfpathlineto{\pgfqpoint{5.232466in}{0.537024in}}%
\pgfpathlineto{\pgfqpoint{5.243391in}{0.488557in}}%
\pgfpathlineto{\pgfqpoint{5.254316in}{0.448784in}}%
\pgfpathlineto{\pgfqpoint{5.265241in}{0.418334in}}%
\pgfpathlineto{\pgfqpoint{5.276166in}{0.397689in}}%
\pgfpathlineto{\pgfqpoint{5.281629in}{0.391152in}}%
\pgfpathlineto{\pgfqpoint{5.287091in}{0.387176in}}%
\pgfpathlineto{\pgfqpoint{5.292554in}{0.385777in}}%
\pgfpathlineto{\pgfqpoint{5.298016in}{0.386961in}}%
\pgfpathlineto{\pgfqpoint{5.303479in}{0.390722in}}%
\pgfpathlineto{\pgfqpoint{5.308941in}{0.397047in}}%
\pgfpathlineto{\pgfqpoint{5.314404in}{0.405910in}}%
\pgfpathlineto{\pgfqpoint{5.325328in}{0.431099in}}%
\pgfpathlineto{\pgfqpoint{5.336253in}{0.465891in}}%
\pgfpathlineto{\pgfqpoint{5.347178in}{0.509735in}}%
\pgfpathlineto{\pgfqpoint{5.358103in}{0.561938in}}%
\pgfpathlineto{\pgfqpoint{5.374491in}{0.654076in}}%
\pgfpathlineto{\pgfqpoint{5.390878in}{0.759859in}}%
\pgfpathlineto{\pgfqpoint{5.412728in}{0.915578in}}%
\pgfpathlineto{\pgfqpoint{5.429116in}{1.037994in}}%
\pgfpathlineto{\pgfqpoint{5.429116in}{1.037994in}}%
\pgfusepath{stroke}%
\end{pgfscope}%
\begin{pgfscope}%
\pgfsetrectcap%
\pgfsetmiterjoin%
\pgfsetlinewidth{0.803000pt}%
\definecolor{currentstroke}{rgb}{0.000000,0.000000,0.000000}%
\pgfsetstrokecolor{currentstroke}%
\pgfsetdash{}{0pt}%
\pgfpathmoveto{\pgfqpoint{1.993487in}{0.320555in}}%
\pgfpathlineto{\pgfqpoint{1.993487in}{1.755432in}}%
\pgfusepath{stroke}%
\end{pgfscope}%
\begin{pgfscope}%
\pgfsetrectcap%
\pgfsetmiterjoin%
\pgfsetlinewidth{0.803000pt}%
\definecolor{currentstroke}{rgb}{0.000000,0.000000,0.000000}%
\pgfsetstrokecolor{currentstroke}%
\pgfsetdash{}{0pt}%
\pgfpathmoveto{\pgfqpoint{5.592717in}{0.320555in}}%
\pgfpathlineto{\pgfqpoint{5.592717in}{1.755432in}}%
\pgfusepath{stroke}%
\end{pgfscope}%
\begin{pgfscope}%
\pgfsetrectcap%
\pgfsetmiterjoin%
\pgfsetlinewidth{0.803000pt}%
\definecolor{currentstroke}{rgb}{0.000000,0.000000,0.000000}%
\pgfsetstrokecolor{currentstroke}%
\pgfsetdash{}{0pt}%
\pgfpathmoveto{\pgfqpoint{1.993487in}{0.320555in}}%
\pgfpathlineto{\pgfqpoint{5.592717in}{0.320555in}}%
\pgfusepath{stroke}%
\end{pgfscope}%
\begin{pgfscope}%
\pgfsetrectcap%
\pgfsetmiterjoin%
\pgfsetlinewidth{0.803000pt}%
\definecolor{currentstroke}{rgb}{0.000000,0.000000,0.000000}%
\pgfsetstrokecolor{currentstroke}%
\pgfsetdash{}{0pt}%
\pgfpathmoveto{\pgfqpoint{1.993487in}{1.755432in}}%
\pgfpathlineto{\pgfqpoint{5.592717in}{1.755432in}}%
\pgfusepath{stroke}%
\end{pgfscope}%
\begin{pgfscope}%
\definecolor{textcolor}{rgb}{0.000000,0.000000,0.000000}%
\pgfsetstrokecolor{textcolor}%
\pgfsetfillcolor{textcolor}%
\pgftext[x=3.793102in,y=1.838765in,,base]{\color{textcolor}{\rmfamily\fontsize{12.000000}{14.400000}\selectfont\catcode`\^=\active\def^{\ifmmode\sp\else\^{}\fi}\catcode`\%=\active\def%{\%}Centro abajo}}%
\end{pgfscope}%
\begin{pgfscope}%
\pgfsetbuttcap%
\pgfsetmiterjoin%
\definecolor{currentfill}{rgb}{1.000000,1.000000,1.000000}%
\pgfsetfillcolor{currentfill}%
\pgfsetlinewidth{0.000000pt}%
\definecolor{currentstroke}{rgb}{0.000000,0.000000,0.000000}%
\pgfsetstrokecolor{currentstroke}%
\pgfsetstrokeopacity{0.000000}%
\pgfsetdash{}{0pt}%
\pgfpathmoveto{\pgfqpoint{4.122816in}{2.452069in}}%
\pgfpathlineto{\pgfqpoint{5.592717in}{2.452069in}}%
\pgfpathlineto{\pgfqpoint{5.592717in}{3.886945in}}%
\pgfpathlineto{\pgfqpoint{4.122816in}{3.886945in}}%
\pgfpathlineto{\pgfqpoint{4.122816in}{2.452069in}}%
\pgfpathclose%
\pgfusepath{fill}%
\end{pgfscope}%
\begin{pgfscope}%
\pgfpathrectangle{\pgfqpoint{4.122816in}{2.452069in}}{\pgfqpoint{1.469901in}{1.434876in}}%
\pgfusepath{clip}%
\pgfsetbuttcap%
\pgfsetroundjoin%
\pgfsetlinewidth{0.501875pt}%
\definecolor{currentstroke}{rgb}{0.501961,0.501961,0.501961}%
\pgfsetstrokecolor{currentstroke}%
\pgfsetstrokeopacity{0.500000}%
\pgfsetdash{{1.850000pt}{0.800000pt}}{0.000000pt}%
\pgfpathmoveto{\pgfqpoint{4.189630in}{2.452069in}}%
\pgfpathlineto{\pgfqpoint{4.189630in}{3.886945in}}%
\pgfusepath{stroke}%
\end{pgfscope}%
\begin{pgfscope}%
\pgfsetbuttcap%
\pgfsetroundjoin%
\definecolor{currentfill}{rgb}{0.000000,0.000000,0.000000}%
\pgfsetfillcolor{currentfill}%
\pgfsetlinewidth{0.803000pt}%
\definecolor{currentstroke}{rgb}{0.000000,0.000000,0.000000}%
\pgfsetstrokecolor{currentstroke}%
\pgfsetdash{}{0pt}%
\pgfsys@defobject{currentmarker}{\pgfqpoint{0.000000in}{-0.048611in}}{\pgfqpoint{0.000000in}{0.000000in}}{%
\pgfpathmoveto{\pgfqpoint{0.000000in}{0.000000in}}%
\pgfpathlineto{\pgfqpoint{0.000000in}{-0.048611in}}%
\pgfusepath{stroke,fill}%
}%
\begin{pgfscope}%
\pgfsys@transformshift{4.189630in}{2.452069in}%
\pgfsys@useobject{currentmarker}{}%
\end{pgfscope}%
\end{pgfscope}%
\begin{pgfscope}%
\definecolor{textcolor}{rgb}{0.000000,0.000000,0.000000}%
\pgfsetstrokecolor{textcolor}%
\pgfsetfillcolor{textcolor}%
\pgftext[x=4.189630in,y=2.354846in,,top]{\color{textcolor}{\rmfamily\fontsize{10.000000}{12.000000}\selectfont\catcode`\^=\active\def^{\ifmmode\sp\else\^{}\fi}\catcode`\%=\active\def%{\%}$\mathdefault{0}$}}%
\end{pgfscope}%
\begin{pgfscope}%
\pgfpathrectangle{\pgfqpoint{4.122816in}{2.452069in}}{\pgfqpoint{1.469901in}{1.434876in}}%
\pgfusepath{clip}%
\pgfsetbuttcap%
\pgfsetroundjoin%
\pgfsetlinewidth{0.501875pt}%
\definecolor{currentstroke}{rgb}{0.501961,0.501961,0.501961}%
\pgfsetstrokecolor{currentstroke}%
\pgfsetstrokeopacity{0.500000}%
\pgfsetdash{{1.850000pt}{0.800000pt}}{0.000000pt}%
\pgfpathmoveto{\pgfqpoint{4.898545in}{2.452069in}}%
\pgfpathlineto{\pgfqpoint{4.898545in}{3.886945in}}%
\pgfusepath{stroke}%
\end{pgfscope}%
\begin{pgfscope}%
\pgfsetbuttcap%
\pgfsetroundjoin%
\definecolor{currentfill}{rgb}{0.000000,0.000000,0.000000}%
\pgfsetfillcolor{currentfill}%
\pgfsetlinewidth{0.803000pt}%
\definecolor{currentstroke}{rgb}{0.000000,0.000000,0.000000}%
\pgfsetstrokecolor{currentstroke}%
\pgfsetdash{}{0pt}%
\pgfsys@defobject{currentmarker}{\pgfqpoint{0.000000in}{-0.048611in}}{\pgfqpoint{0.000000in}{0.000000in}}{%
\pgfpathmoveto{\pgfqpoint{0.000000in}{0.000000in}}%
\pgfpathlineto{\pgfqpoint{0.000000in}{-0.048611in}}%
\pgfusepath{stroke,fill}%
}%
\begin{pgfscope}%
\pgfsys@transformshift{4.898545in}{2.452069in}%
\pgfsys@useobject{currentmarker}{}%
\end{pgfscope}%
\end{pgfscope}%
\begin{pgfscope}%
\definecolor{textcolor}{rgb}{0.000000,0.000000,0.000000}%
\pgfsetstrokecolor{textcolor}%
\pgfsetfillcolor{textcolor}%
\pgftext[x=4.898545in,y=2.354846in,,top]{\color{textcolor}{\rmfamily\fontsize{10.000000}{12.000000}\selectfont\catcode`\^=\active\def^{\ifmmode\sp\else\^{}\fi}\catcode`\%=\active\def%{\%}$\mathdefault{10}$}}%
\end{pgfscope}%
\begin{pgfscope}%
\pgfpathrectangle{\pgfqpoint{4.122816in}{2.452069in}}{\pgfqpoint{1.469901in}{1.434876in}}%
\pgfusepath{clip}%
\pgfsetbuttcap%
\pgfsetroundjoin%
\pgfsetlinewidth{0.501875pt}%
\definecolor{currentstroke}{rgb}{0.501961,0.501961,0.501961}%
\pgfsetstrokecolor{currentstroke}%
\pgfsetstrokeopacity{0.500000}%
\pgfsetdash{{1.850000pt}{0.800000pt}}{0.000000pt}%
\pgfpathmoveto{\pgfqpoint{4.122816in}{2.517281in}}%
\pgfpathlineto{\pgfqpoint{5.592717in}{2.517281in}}%
\pgfusepath{stroke}%
\end{pgfscope}%
\begin{pgfscope}%
\pgfsetbuttcap%
\pgfsetroundjoin%
\definecolor{currentfill}{rgb}{0.000000,0.000000,0.000000}%
\pgfsetfillcolor{currentfill}%
\pgfsetlinewidth{0.803000pt}%
\definecolor{currentstroke}{rgb}{0.000000,0.000000,0.000000}%
\pgfsetstrokecolor{currentstroke}%
\pgfsetdash{}{0pt}%
\pgfsys@defobject{currentmarker}{\pgfqpoint{-0.048611in}{0.000000in}}{\pgfqpoint{-0.000000in}{0.000000in}}{%
\pgfpathmoveto{\pgfqpoint{-0.000000in}{0.000000in}}%
\pgfpathlineto{\pgfqpoint{-0.048611in}{0.000000in}}%
\pgfusepath{stroke,fill}%
}%
\begin{pgfscope}%
\pgfsys@transformshift{4.122816in}{2.517281in}%
\pgfsys@useobject{currentmarker}{}%
\end{pgfscope}%
\end{pgfscope}%
\begin{pgfscope}%
\definecolor{textcolor}{rgb}{0.000000,0.000000,0.000000}%
\pgfsetstrokecolor{textcolor}%
\pgfsetfillcolor{textcolor}%
\pgftext[x=3.740099in, y=2.469087in, left, base]{\color{textcolor}{\rmfamily\fontsize{10.000000}{12.000000}\selectfont\catcode`\^=\active\def^{\ifmmode\sp\else\^{}\fi}\catcode`\%=\active\def%{\%}$\mathdefault{\ensuremath{-}1.0}$}}%
\end{pgfscope}%
\begin{pgfscope}%
\pgfpathrectangle{\pgfqpoint{4.122816in}{2.452069in}}{\pgfqpoint{1.469901in}{1.434876in}}%
\pgfusepath{clip}%
\pgfsetbuttcap%
\pgfsetroundjoin%
\pgfsetlinewidth{0.501875pt}%
\definecolor{currentstroke}{rgb}{0.501961,0.501961,0.501961}%
\pgfsetstrokecolor{currentstroke}%
\pgfsetstrokeopacity{0.500000}%
\pgfsetdash{{1.850000pt}{0.800000pt}}{0.000000pt}%
\pgfpathmoveto{\pgfqpoint{4.122816in}{2.843392in}}%
\pgfpathlineto{\pgfqpoint{5.592717in}{2.843392in}}%
\pgfusepath{stroke}%
\end{pgfscope}%
\begin{pgfscope}%
\pgfsetbuttcap%
\pgfsetroundjoin%
\definecolor{currentfill}{rgb}{0.000000,0.000000,0.000000}%
\pgfsetfillcolor{currentfill}%
\pgfsetlinewidth{0.803000pt}%
\definecolor{currentstroke}{rgb}{0.000000,0.000000,0.000000}%
\pgfsetstrokecolor{currentstroke}%
\pgfsetdash{}{0pt}%
\pgfsys@defobject{currentmarker}{\pgfqpoint{-0.048611in}{0.000000in}}{\pgfqpoint{-0.000000in}{0.000000in}}{%
\pgfpathmoveto{\pgfqpoint{-0.000000in}{0.000000in}}%
\pgfpathlineto{\pgfqpoint{-0.048611in}{0.000000in}}%
\pgfusepath{stroke,fill}%
}%
\begin{pgfscope}%
\pgfsys@transformshift{4.122816in}{2.843392in}%
\pgfsys@useobject{currentmarker}{}%
\end{pgfscope}%
\end{pgfscope}%
\begin{pgfscope}%
\definecolor{textcolor}{rgb}{0.000000,0.000000,0.000000}%
\pgfsetstrokecolor{textcolor}%
\pgfsetfillcolor{textcolor}%
\pgftext[x=3.740099in, y=2.795197in, left, base]{\color{textcolor}{\rmfamily\fontsize{10.000000}{12.000000}\selectfont\catcode`\^=\active\def^{\ifmmode\sp\else\^{}\fi}\catcode`\%=\active\def%{\%}$\mathdefault{\ensuremath{-}0.5}$}}%
\end{pgfscope}%
\begin{pgfscope}%
\pgfpathrectangle{\pgfqpoint{4.122816in}{2.452069in}}{\pgfqpoint{1.469901in}{1.434876in}}%
\pgfusepath{clip}%
\pgfsetbuttcap%
\pgfsetroundjoin%
\pgfsetlinewidth{0.501875pt}%
\definecolor{currentstroke}{rgb}{0.501961,0.501961,0.501961}%
\pgfsetstrokecolor{currentstroke}%
\pgfsetstrokeopacity{0.500000}%
\pgfsetdash{{1.850000pt}{0.800000pt}}{0.000000pt}%
\pgfpathmoveto{\pgfqpoint{4.122816in}{3.169502in}}%
\pgfpathlineto{\pgfqpoint{5.592717in}{3.169502in}}%
\pgfusepath{stroke}%
\end{pgfscope}%
\begin{pgfscope}%
\pgfsetbuttcap%
\pgfsetroundjoin%
\definecolor{currentfill}{rgb}{0.000000,0.000000,0.000000}%
\pgfsetfillcolor{currentfill}%
\pgfsetlinewidth{0.803000pt}%
\definecolor{currentstroke}{rgb}{0.000000,0.000000,0.000000}%
\pgfsetstrokecolor{currentstroke}%
\pgfsetdash{}{0pt}%
\pgfsys@defobject{currentmarker}{\pgfqpoint{-0.048611in}{0.000000in}}{\pgfqpoint{-0.000000in}{0.000000in}}{%
\pgfpathmoveto{\pgfqpoint{-0.000000in}{0.000000in}}%
\pgfpathlineto{\pgfqpoint{-0.048611in}{0.000000in}}%
\pgfusepath{stroke,fill}%
}%
\begin{pgfscope}%
\pgfsys@transformshift{4.122816in}{3.169502in}%
\pgfsys@useobject{currentmarker}{}%
\end{pgfscope}%
\end{pgfscope}%
\begin{pgfscope}%
\definecolor{textcolor}{rgb}{0.000000,0.000000,0.000000}%
\pgfsetstrokecolor{textcolor}%
\pgfsetfillcolor{textcolor}%
\pgftext[x=3.848124in, y=3.121308in, left, base]{\color{textcolor}{\rmfamily\fontsize{10.000000}{12.000000}\selectfont\catcode`\^=\active\def^{\ifmmode\sp\else\^{}\fi}\catcode`\%=\active\def%{\%}$\mathdefault{0.0}$}}%
\end{pgfscope}%
\begin{pgfscope}%
\pgfpathrectangle{\pgfqpoint{4.122816in}{2.452069in}}{\pgfqpoint{1.469901in}{1.434876in}}%
\pgfusepath{clip}%
\pgfsetbuttcap%
\pgfsetroundjoin%
\pgfsetlinewidth{0.501875pt}%
\definecolor{currentstroke}{rgb}{0.501961,0.501961,0.501961}%
\pgfsetstrokecolor{currentstroke}%
\pgfsetstrokeopacity{0.500000}%
\pgfsetdash{{1.850000pt}{0.800000pt}}{0.000000pt}%
\pgfpathmoveto{\pgfqpoint{4.122816in}{3.495613in}}%
\pgfpathlineto{\pgfqpoint{5.592717in}{3.495613in}}%
\pgfusepath{stroke}%
\end{pgfscope}%
\begin{pgfscope}%
\pgfsetbuttcap%
\pgfsetroundjoin%
\definecolor{currentfill}{rgb}{0.000000,0.000000,0.000000}%
\pgfsetfillcolor{currentfill}%
\pgfsetlinewidth{0.803000pt}%
\definecolor{currentstroke}{rgb}{0.000000,0.000000,0.000000}%
\pgfsetstrokecolor{currentstroke}%
\pgfsetdash{}{0pt}%
\pgfsys@defobject{currentmarker}{\pgfqpoint{-0.048611in}{0.000000in}}{\pgfqpoint{-0.000000in}{0.000000in}}{%
\pgfpathmoveto{\pgfqpoint{-0.000000in}{0.000000in}}%
\pgfpathlineto{\pgfqpoint{-0.048611in}{0.000000in}}%
\pgfusepath{stroke,fill}%
}%
\begin{pgfscope}%
\pgfsys@transformshift{4.122816in}{3.495613in}%
\pgfsys@useobject{currentmarker}{}%
\end{pgfscope}%
\end{pgfscope}%
\begin{pgfscope}%
\definecolor{textcolor}{rgb}{0.000000,0.000000,0.000000}%
\pgfsetstrokecolor{textcolor}%
\pgfsetfillcolor{textcolor}%
\pgftext[x=3.848124in, y=3.447418in, left, base]{\color{textcolor}{\rmfamily\fontsize{10.000000}{12.000000}\selectfont\catcode`\^=\active\def^{\ifmmode\sp\else\^{}\fi}\catcode`\%=\active\def%{\%}$\mathdefault{0.5}$}}%
\end{pgfscope}%
\begin{pgfscope}%
\pgfpathrectangle{\pgfqpoint{4.122816in}{2.452069in}}{\pgfqpoint{1.469901in}{1.434876in}}%
\pgfusepath{clip}%
\pgfsetbuttcap%
\pgfsetroundjoin%
\pgfsetlinewidth{0.501875pt}%
\definecolor{currentstroke}{rgb}{0.501961,0.501961,0.501961}%
\pgfsetstrokecolor{currentstroke}%
\pgfsetstrokeopacity{0.500000}%
\pgfsetdash{{1.850000pt}{0.800000pt}}{0.000000pt}%
\pgfpathmoveto{\pgfqpoint{4.122816in}{3.821723in}}%
\pgfpathlineto{\pgfqpoint{5.592717in}{3.821723in}}%
\pgfusepath{stroke}%
\end{pgfscope}%
\begin{pgfscope}%
\pgfsetbuttcap%
\pgfsetroundjoin%
\definecolor{currentfill}{rgb}{0.000000,0.000000,0.000000}%
\pgfsetfillcolor{currentfill}%
\pgfsetlinewidth{0.803000pt}%
\definecolor{currentstroke}{rgb}{0.000000,0.000000,0.000000}%
\pgfsetstrokecolor{currentstroke}%
\pgfsetdash{}{0pt}%
\pgfsys@defobject{currentmarker}{\pgfqpoint{-0.048611in}{0.000000in}}{\pgfqpoint{-0.000000in}{0.000000in}}{%
\pgfpathmoveto{\pgfqpoint{-0.000000in}{0.000000in}}%
\pgfpathlineto{\pgfqpoint{-0.048611in}{0.000000in}}%
\pgfusepath{stroke,fill}%
}%
\begin{pgfscope}%
\pgfsys@transformshift{4.122816in}{3.821723in}%
\pgfsys@useobject{currentmarker}{}%
\end{pgfscope}%
\end{pgfscope}%
\begin{pgfscope}%
\definecolor{textcolor}{rgb}{0.000000,0.000000,0.000000}%
\pgfsetstrokecolor{textcolor}%
\pgfsetfillcolor{textcolor}%
\pgftext[x=3.848124in, y=3.773529in, left, base]{\color{textcolor}{\rmfamily\fontsize{10.000000}{12.000000}\selectfont\catcode`\^=\active\def^{\ifmmode\sp\else\^{}\fi}\catcode`\%=\active\def%{\%}$\mathdefault{1.0}$}}%
\end{pgfscope}%
\begin{pgfscope}%
\pgfpathrectangle{\pgfqpoint{4.122816in}{2.452069in}}{\pgfqpoint{1.469901in}{1.434876in}}%
\pgfusepath{clip}%
\pgfsetrectcap%
\pgfsetroundjoin%
\pgfsetlinewidth{1.505625pt}%
\definecolor{currentstroke}{rgb}{1.000000,0.000000,0.000000}%
\pgfsetstrokecolor{currentstroke}%
\pgfsetdash{}{0pt}%
\pgfpathmoveto{\pgfqpoint{4.189630in}{3.821723in}}%
\pgfpathlineto{\pgfqpoint{4.191860in}{3.820432in}}%
\pgfpathlineto{\pgfqpoint{4.194091in}{3.816563in}}%
\pgfpathlineto{\pgfqpoint{4.198553in}{3.801164in}}%
\pgfpathlineto{\pgfqpoint{4.203015in}{3.775771in}}%
\pgfpathlineto{\pgfqpoint{4.209707in}{3.719860in}}%
\pgfpathlineto{\pgfqpoint{4.216400in}{3.644389in}}%
\pgfpathlineto{\pgfqpoint{4.225323in}{3.518055in}}%
\pgfpathlineto{\pgfqpoint{4.238708in}{3.290238in}}%
\pgfpathlineto{\pgfqpoint{4.265478in}{2.818063in}}%
\pgfpathlineto{\pgfqpoint{4.276632in}{2.665261in}}%
\pgfpathlineto{\pgfqpoint{4.283325in}{2.596578in}}%
\pgfpathlineto{\pgfqpoint{4.290017in}{2.548259in}}%
\pgfpathlineto{\pgfqpoint{4.294479in}{2.528239in}}%
\pgfpathlineto{\pgfqpoint{4.298941in}{2.518366in}}%
\pgfpathlineto{\pgfqpoint{4.301172in}{2.517290in}}%
\pgfpathlineto{\pgfqpoint{4.303403in}{2.518797in}}%
\pgfpathlineto{\pgfqpoint{4.305633in}{2.522880in}}%
\pgfpathlineto{\pgfqpoint{4.310095in}{2.538701in}}%
\pgfpathlineto{\pgfqpoint{4.314557in}{2.564503in}}%
\pgfpathlineto{\pgfqpoint{4.321249in}{2.620987in}}%
\pgfpathlineto{\pgfqpoint{4.327942in}{2.696967in}}%
\pgfpathlineto{\pgfqpoint{4.336865in}{2.823846in}}%
\pgfpathlineto{\pgfqpoint{4.350250in}{3.052130in}}%
\pgfpathlineto{\pgfqpoint{4.377020in}{3.523818in}}%
\pgfpathlineto{\pgfqpoint{4.388174in}{3.675906in}}%
\pgfpathlineto{\pgfqpoint{4.394867in}{3.744053in}}%
\pgfpathlineto{\pgfqpoint{4.401560in}{3.791779in}}%
\pgfpathlineto{\pgfqpoint{4.406021in}{3.811381in}}%
\pgfpathlineto{\pgfqpoint{4.410483in}{3.820826in}}%
\pgfpathlineto{\pgfqpoint{4.412714in}{3.821687in}}%
\pgfpathlineto{\pgfqpoint{4.414945in}{3.819966in}}%
\pgfpathlineto{\pgfqpoint{4.417175in}{3.815669in}}%
\pgfpathlineto{\pgfqpoint{4.421637in}{3.799426in}}%
\pgfpathlineto{\pgfqpoint{4.426099in}{3.773215in}}%
\pgfpathlineto{\pgfqpoint{4.432791in}{3.716159in}}%
\pgfpathlineto{\pgfqpoint{4.439484in}{3.639673in}}%
\pgfpathlineto{\pgfqpoint{4.448407in}{3.512253in}}%
\pgfpathlineto{\pgfqpoint{4.461792in}{3.283508in}}%
\pgfpathlineto{\pgfqpoint{4.488562in}{2.812319in}}%
\pgfpathlineto{\pgfqpoint{4.499717in}{2.660949in}}%
\pgfpathlineto{\pgfqpoint{4.506409in}{2.593340in}}%
\pgfpathlineto{\pgfqpoint{4.513102in}{2.546210in}}%
\pgfpathlineto{\pgfqpoint{4.517563in}{2.527026in}}%
\pgfpathlineto{\pgfqpoint{4.522025in}{2.518008in}}%
\pgfpathlineto{\pgfqpoint{4.524256in}{2.517362in}}%
\pgfpathlineto{\pgfqpoint{4.526487in}{2.519299in}}%
\pgfpathlineto{\pgfqpoint{4.528717in}{2.523810in}}%
\pgfpathlineto{\pgfqpoint{4.533179in}{2.540474in}}%
\pgfpathlineto{\pgfqpoint{4.537641in}{2.567092in}}%
\pgfpathlineto{\pgfqpoint{4.544333in}{2.624719in}}%
\pgfpathlineto{\pgfqpoint{4.551026in}{2.701709in}}%
\pgfpathlineto{\pgfqpoint{4.559949in}{2.829666in}}%
\pgfpathlineto{\pgfqpoint{4.573334in}{3.058866in}}%
\pgfpathlineto{\pgfqpoint{4.600104in}{3.529543in}}%
\pgfpathlineto{\pgfqpoint{4.609028in}{3.653662in}}%
\pgfpathlineto{\pgfqpoint{4.617951in}{3.747260in}}%
\pgfpathlineto{\pgfqpoint{4.624644in}{3.793794in}}%
\pgfpathlineto{\pgfqpoint{4.629105in}{3.812559in}}%
\pgfpathlineto{\pgfqpoint{4.633567in}{3.821149in}}%
\pgfpathlineto{\pgfqpoint{4.635798in}{3.821580in}}%
\pgfpathlineto{\pgfqpoint{4.638029in}{3.819428in}}%
\pgfpathlineto{\pgfqpoint{4.642490in}{3.807423in}}%
\pgfpathlineto{\pgfqpoint{4.646952in}{3.785324in}}%
\pgfpathlineto{\pgfqpoint{4.653645in}{3.734057in}}%
\pgfpathlineto{\pgfqpoint{4.660337in}{3.662723in}}%
\pgfpathlineto{\pgfqpoint{4.669260in}{3.540872in}}%
\pgfpathlineto{\pgfqpoint{4.682645in}{3.317016in}}%
\pgfpathlineto{\pgfqpoint{4.713877in}{2.773247in}}%
\pgfpathlineto{\pgfqpoint{4.722801in}{2.656694in}}%
\pgfpathlineto{\pgfqpoint{4.729493in}{2.590166in}}%
\pgfpathlineto{\pgfqpoint{4.736186in}{2.544229in}}%
\pgfpathlineto{\pgfqpoint{4.740647in}{2.525883in}}%
\pgfpathlineto{\pgfqpoint{4.745109in}{2.517721in}}%
\pgfpathlineto{\pgfqpoint{4.747340in}{2.517506in}}%
\pgfpathlineto{\pgfqpoint{4.749571in}{2.519872in}}%
\pgfpathlineto{\pgfqpoint{4.754032in}{2.532302in}}%
\pgfpathlineto{\pgfqpoint{4.758494in}{2.554816in}}%
\pgfpathlineto{\pgfqpoint{4.765187in}{2.606668in}}%
\pgfpathlineto{\pgfqpoint{4.771879in}{2.678526in}}%
\pgfpathlineto{\pgfqpoint{4.780802in}{2.800950in}}%
\pgfpathlineto{\pgfqpoint{4.794188in}{3.025322in}}%
\pgfpathlineto{\pgfqpoint{4.823188in}{3.535227in}}%
\pgfpathlineto{\pgfqpoint{4.832112in}{3.658220in}}%
\pgfpathlineto{\pgfqpoint{4.841035in}{3.750402in}}%
\pgfpathlineto{\pgfqpoint{4.847728in}{3.795740in}}%
\pgfpathlineto{\pgfqpoint{4.852189in}{3.813667in}}%
\pgfpathlineto{\pgfqpoint{4.856651in}{3.821400in}}%
\pgfpathlineto{\pgfqpoint{4.858882in}{3.821400in}}%
\pgfpathlineto{\pgfqpoint{4.861113in}{3.818819in}}%
\pgfpathlineto{\pgfqpoint{4.865574in}{3.805963in}}%
\pgfpathlineto{\pgfqpoint{4.870036in}{3.783037in}}%
\pgfpathlineto{\pgfqpoint{4.876729in}{3.730600in}}%
\pgfpathlineto{\pgfqpoint{4.883421in}{3.658220in}}%
\pgfpathlineto{\pgfqpoint{4.892345in}{3.535227in}}%
\pgfpathlineto{\pgfqpoint{4.905730in}{3.310344in}}%
\pgfpathlineto{\pgfqpoint{4.934731in}{2.800950in}}%
\pgfpathlineto{\pgfqpoint{4.943654in}{2.678526in}}%
\pgfpathlineto{\pgfqpoint{4.952577in}{2.587055in}}%
\pgfpathlineto{\pgfqpoint{4.959270in}{2.542317in}}%
\pgfpathlineto{\pgfqpoint{4.963731in}{2.524811in}}%
\pgfpathlineto{\pgfqpoint{4.968193in}{2.517506in}}%
\pgfpathlineto{\pgfqpoint{4.970424in}{2.517721in}}%
\pgfpathlineto{\pgfqpoint{4.972655in}{2.520517in}}%
\pgfpathlineto{\pgfqpoint{4.977116in}{2.533797in}}%
\pgfpathlineto{\pgfqpoint{4.981578in}{2.557137in}}%
\pgfpathlineto{\pgfqpoint{4.988271in}{2.610156in}}%
\pgfpathlineto{\pgfqpoint{4.994963in}{2.683057in}}%
\pgfpathlineto{\pgfqpoint{5.003887in}{2.806615in}}%
\pgfpathlineto{\pgfqpoint{5.017272in}{3.032002in}}%
\pgfpathlineto{\pgfqpoint{5.046273in}{3.540872in}}%
\pgfpathlineto{\pgfqpoint{5.055196in}{3.662723in}}%
\pgfpathlineto{\pgfqpoint{5.064119in}{3.753481in}}%
\pgfpathlineto{\pgfqpoint{5.070812in}{3.797617in}}%
\pgfpathlineto{\pgfqpoint{5.075273in}{3.814703in}}%
\pgfpathlineto{\pgfqpoint{5.079735in}{3.821580in}}%
\pgfpathlineto{\pgfqpoint{5.081966in}{3.821149in}}%
\pgfpathlineto{\pgfqpoint{5.084197in}{3.818138in}}%
\pgfpathlineto{\pgfqpoint{5.088659in}{3.804434in}}%
\pgfpathlineto{\pgfqpoint{5.093120in}{3.780682in}}%
\pgfpathlineto{\pgfqpoint{5.099813in}{3.727081in}}%
\pgfpathlineto{\pgfqpoint{5.106505in}{3.653662in}}%
\pgfpathlineto{\pgfqpoint{5.115429in}{3.529543in}}%
\pgfpathlineto{\pgfqpoint{5.128814in}{3.303656in}}%
\pgfpathlineto{\pgfqpoint{5.157815in}{2.795326in}}%
\pgfpathlineto{\pgfqpoint{5.166738in}{2.674050in}}%
\pgfpathlineto{\pgfqpoint{5.175661in}{2.584008in}}%
\pgfpathlineto{\pgfqpoint{5.182354in}{2.540474in}}%
\pgfpathlineto{\pgfqpoint{5.186816in}{2.523810in}}%
\pgfpathlineto{\pgfqpoint{5.189046in}{2.519299in}}%
\pgfpathlineto{\pgfqpoint{5.191277in}{2.517362in}}%
\pgfpathlineto{\pgfqpoint{5.193508in}{2.518008in}}%
\pgfpathlineto{\pgfqpoint{5.195739in}{2.521233in}}%
\pgfpathlineto{\pgfqpoint{5.200201in}{2.535362in}}%
\pgfpathlineto{\pgfqpoint{5.204662in}{2.559525in}}%
\pgfpathlineto{\pgfqpoint{5.211355in}{2.613705in}}%
\pgfpathlineto{\pgfqpoint{5.218047in}{2.687641in}}%
\pgfpathlineto{\pgfqpoint{5.226971in}{2.812319in}}%
\pgfpathlineto{\pgfqpoint{5.240356in}{3.038698in}}%
\pgfpathlineto{\pgfqpoint{5.269357in}{3.546475in}}%
\pgfpathlineto{\pgfqpoint{5.278280in}{3.667173in}}%
\pgfpathlineto{\pgfqpoint{5.284973in}{3.737452in}}%
\pgfpathlineto{\pgfqpoint{5.291665in}{3.787544in}}%
\pgfpathlineto{\pgfqpoint{5.296127in}{3.808813in}}%
\pgfpathlineto{\pgfqpoint{5.300588in}{3.819966in}}%
\pgfpathlineto{\pgfqpoint{5.302819in}{3.821687in}}%
\pgfpathlineto{\pgfqpoint{5.305050in}{3.820826in}}%
\pgfpathlineto{\pgfqpoint{5.307281in}{3.817386in}}%
\pgfpathlineto{\pgfqpoint{5.311743in}{3.802834in}}%
\pgfpathlineto{\pgfqpoint{5.316204in}{3.778260in}}%
\pgfpathlineto{\pgfqpoint{5.322897in}{3.723501in}}%
\pgfpathlineto{\pgfqpoint{5.329589in}{3.649052in}}%
\pgfpathlineto{\pgfqpoint{5.338513in}{3.523818in}}%
\pgfpathlineto{\pgfqpoint{5.351898in}{3.296954in}}%
\pgfpathlineto{\pgfqpoint{5.380899in}{2.789743in}}%
\pgfpathlineto{\pgfqpoint{5.389822in}{2.669628in}}%
\pgfpathlineto{\pgfqpoint{5.396515in}{2.599879in}}%
\pgfpathlineto{\pgfqpoint{5.403207in}{2.550376in}}%
\pgfpathlineto{\pgfqpoint{5.407669in}{2.529523in}}%
\pgfpathlineto{\pgfqpoint{5.412130in}{2.518797in}}%
\pgfpathlineto{\pgfqpoint{5.414361in}{2.517290in}}%
\pgfpathlineto{\pgfqpoint{5.416592in}{2.518366in}}%
\pgfpathlineto{\pgfqpoint{5.418823in}{2.522021in}}%
\pgfpathlineto{\pgfqpoint{5.423285in}{2.536997in}}%
\pgfpathlineto{\pgfqpoint{5.427746in}{2.561981in}}%
\pgfpathlineto{\pgfqpoint{5.434439in}{2.617316in}}%
\pgfpathlineto{\pgfqpoint{5.441131in}{2.692278in}}%
\pgfpathlineto{\pgfqpoint{5.450055in}{2.818063in}}%
\pgfpathlineto{\pgfqpoint{5.463440in}{3.045407in}}%
\pgfpathlineto{\pgfqpoint{5.492441in}{3.552037in}}%
\pgfpathlineto{\pgfqpoint{5.501364in}{3.671567in}}%
\pgfpathlineto{\pgfqpoint{5.508057in}{3.740784in}}%
\pgfpathlineto{\pgfqpoint{5.514749in}{3.789695in}}%
\pgfpathlineto{\pgfqpoint{5.519211in}{3.810132in}}%
\pgfpathlineto{\pgfqpoint{5.523673in}{3.820432in}}%
\pgfpathlineto{\pgfqpoint{5.525903in}{3.821723in}}%
\pgfpathlineto{\pgfqpoint{5.525903in}{3.821723in}}%
\pgfusepath{stroke}%
\end{pgfscope}%
\begin{pgfscope}%
\pgfsetrectcap%
\pgfsetmiterjoin%
\pgfsetlinewidth{0.803000pt}%
\definecolor{currentstroke}{rgb}{0.000000,0.000000,0.000000}%
\pgfsetstrokecolor{currentstroke}%
\pgfsetdash{}{0pt}%
\pgfpathmoveto{\pgfqpoint{4.122816in}{2.452069in}}%
\pgfpathlineto{\pgfqpoint{4.122816in}{3.886945in}}%
\pgfusepath{stroke}%
\end{pgfscope}%
\begin{pgfscope}%
\pgfsetrectcap%
\pgfsetmiterjoin%
\pgfsetlinewidth{0.803000pt}%
\definecolor{currentstroke}{rgb}{0.000000,0.000000,0.000000}%
\pgfsetstrokecolor{currentstroke}%
\pgfsetdash{}{0pt}%
\pgfpathmoveto{\pgfqpoint{5.592717in}{2.452069in}}%
\pgfpathlineto{\pgfqpoint{5.592717in}{3.886945in}}%
\pgfusepath{stroke}%
\end{pgfscope}%
\begin{pgfscope}%
\pgfsetrectcap%
\pgfsetmiterjoin%
\pgfsetlinewidth{0.803000pt}%
\definecolor{currentstroke}{rgb}{0.000000,0.000000,0.000000}%
\pgfsetstrokecolor{currentstroke}%
\pgfsetdash{}{0pt}%
\pgfpathmoveto{\pgfqpoint{4.122816in}{2.452069in}}%
\pgfpathlineto{\pgfqpoint{5.592717in}{2.452069in}}%
\pgfusepath{stroke}%
\end{pgfscope}%
\begin{pgfscope}%
\pgfsetrectcap%
\pgfsetmiterjoin%
\pgfsetlinewidth{0.803000pt}%
\definecolor{currentstroke}{rgb}{0.000000,0.000000,0.000000}%
\pgfsetstrokecolor{currentstroke}%
\pgfsetdash{}{0pt}%
\pgfpathmoveto{\pgfqpoint{4.122816in}{3.886945in}}%
\pgfpathlineto{\pgfqpoint{5.592717in}{3.886945in}}%
\pgfusepath{stroke}%
\end{pgfscope}%
\begin{pgfscope}%
\definecolor{textcolor}{rgb}{0.000000,0.000000,0.000000}%
\pgfsetstrokecolor{textcolor}%
\pgfsetfillcolor{textcolor}%
\pgftext[x=4.857767in,y=3.970278in,,base]{\color{textcolor}{\rmfamily\fontsize{12.000000}{14.400000}\selectfont\catcode`\^=\active\def^{\ifmmode\sp\else\^{}\fi}\catcode`\%=\active\def%{\%}Derecha arriba}}%
\end{pgfscope}%
\begin{pgfscope}%
\definecolor{textcolor}{rgb}{0.000000,0.000000,0.000000}%
\pgfsetstrokecolor{textcolor}%
\pgfsetfillcolor{textcolor}%
\pgftext[x=2.870495in,y=4.722560in,,top]{\color{textcolor}{\rmfamily\fontsize{14.000000}{16.800000}\selectfont\catcode`\^=\active\def^{\ifmmode\sp\else\^{}\fi}\catcode`\%=\active\def%{\%}GridSpec con anchos asimétricos y spans de filas}}%
\end{pgfscope}%
\end{pgfpicture}%
\makeatother%
\endgroup%
 
    \caption{Gráfico Utilizando GridSpec}
    \label{fig:grafico_con_GridSpec}
\end{figure}